% Framework source amendment v3.068: scoped geometry/operator additions, physical occurrence unchanged.
% Paper source Paper5C-InteractionResponse-v2_181.tex: v3.040 S3 literature mathematical/conditional comparison; physical EM packet unresolved.
% Framework v3.027: exact 600-cell/Hurwitz return channels, Cartan rank, native clutch, Clifford constraints and conditional physical gates.
% Paper-series v2.382: reintegrated vetted mathematics from prior repair batches 5--8 into the latest 31-paper branch.
% Paper-series v2.381: complete algebraic exceptional torsion atlas and all-cocycle twist-stability; original claims retained.
% Paper-series v2.380: unified bibliography, exact exceptional double/length/gauge mathematics in I-B; previous theorem labels retained.
% Paper-series v2.379: full equivariant Dirac/Hodge/eta/torsion atlas in Paper I-B; paper-root references synchronized.
% Paper-series v2.378: 2026 exceptional literature consolidation and reference synchronization. The mathematical additions are located in Paper I-B.
% Paper-series v2.377: reference-only rebind to current Paper I-B v1.1; original mathematics retained except new Paper I-B comparison chapter.
% Framework v2.371: conditional rooted Petersen same-action residual test; no new vertex.
% Framework v2.370: derived exceptional-root/Hurwitz comparison in I-A/I-C; other papers reference-only synchronized.
% Framework v2.368: preserve prior II-B convention/unital fixes and verified II-D corrections; current references synchronized.
% Framework v2.365: scoped Hurwitz/600-cell hadron-source integration; physical APBI/decay frontier remains open.
% Framework v2.358: current 30-root source reference synchronization; selectively updated III-E, III-F, VI-B, VI-C.
% Framework v2.356: bibliography/source-registry synchronization; mathematical body unchanged.
% Framework v2.348: Paper V audit, terminal observable scope, canonical Gram eigenvalues and exact weak-memory variance.
% Framework v2.338: exact Petersen 120-carrier, oriented two-source and source-ownership results. No hadron mass prediction.
% Framework v2.321: Paper V Petersen rooted 3+12 induction, exact sector compression, conjugation orbit refinement and source-incidence firewall. Original source theorem bodies preserved.
% Framework v2.275: exact positive dense-twin word cubature, finite repeatability and eighth-moment classification, nested sampling frames, joint action and POVM-jet exports; preserves v2.274.
% Framework v2.274: two-scalar positive convolution certificates, first-leaked exceptional coset detection, branch-symmetric twin mixing, and finite endpoint moments; all v2.273 branches preserved.
% Framework v2.273: restored missing Problem 3 theorems from v2.268; exact exceptional frames, twin geometry, convolution repeatability and projector/occurrence firewalls; merged on v2.272 without overwriting independent branches.
% Framework v2.272: merged v2.270 covariant native-readout mathematics with v2.271 Problem 4 finite-action Hopf hypercharge/holonomy corrections; no universal bulk-gauge uniqueness claim.
% Independent branches: v2.270 readout/guard exports and v2.270 K16 action-consumed charge closure; all prior scoped assumptions retained.
% Framework v2.267: exact Paper II-A converse/observable certificates, III-F action tests, V-B right orbit and V-C native gauge scope; prior v2.266 preserved.
% Framework v2.266: exact Paper II-C outcome-count/readout converse results and scoped downstream imports; 30-root references synchronized.
% Framework v2.265: II-C coordinate-stage constitution and exact first-symbol host inheritance; prior mathematical domains retained.
% Framework v2.264: II-B reverse and action-bridge integration; preserves prior mathematical domains; active 30-root references synchronized.
% Framework v2.263: Paper I-D adds bridge-controlled root-carry parity and source-relative old-cap cancellation; other root mathematical bodies retained; active 30-root references synchronized.
% Framework v2.262: merges Paper II-F three-gap closure from v2.260 into v2.261, preserving newer II-B/II-D repairs; 30-root references synchronized.
% Framework v2.261: merged Paper II-B GP2-level/geometric repairs; preserves v2.260 II-D repairs; active references synchronized.
% Framework v2.260: Paper II-D signed RS, finite-symbol typing, and conditional curvature/Jordan repairs; references synchronized.
% Framework v2.259: conditional native SU(2) atlas and explicit later discharge; v2.258 mathematical bodies retained elsewhere.
% Framework v2.258: preserves parallel II-A/dynamics update and adds complete readout/frequency closures.
% Framework v2.257: exact minimal complete position readout and arbitrary-preparation frequency law.
% Framework v2.256: exact native inertia and sharp bounded ordered-source closure; preserves concurrent v2.255 work.
% Framework v2.255: preserves parallel detector integration and adds exact finite readout closures.
% Framework v2.254: shared bibliography/reference sync; substantive Pasterski detector-Ward additions only IV-B/IV-C/VII-C.
% Framework v2.253: typed raw coordinate defect in I-A and synchronized VI-A import; active references synchronized.
% Framework v2.252: optional shell-frame source-return calculation in III-G; active references synchronized.
% Framework v2.251: forward common-action execution, joint source jets and scoped coupled control.
% Framework v2.251: exact full-heat variations, fermion source jets and explicit physical completion frontier.
% Framework v2.251: certified interacting source jets, pole stability and scoped completion proofs.
% Framework v2.251: interacting finite-source evaluation and coupled local vacuum classification.
% Framework v2.251: exact finite source execution and native causal common-parent bridge.
% Framework v2.251: constructive parent and scoped source-incidence audit; see release report.
% Framework v2.251: common-action two-jet, corrected source/response closure and synchronized references.
% Framework v2.244: scoped rooted source-order Schur diagnostics and synchronized paper references.
% Framework v2.242: signed rooted phase lift, complete transverse return and exact family cubic.
% Framework v2.241: proved rooted conormal geometry, conditional cap realization and scoped action tests.
% Framework v2.240: observable-minimal action and response certificates integrated; inherited v2.239 geometry retained.
% Framework v2.234: restructured main physics owner, technical proofs in V-D.
\documentclass[12pt,a4paper]{article}

\usepackage{mathtools}
\usepackage{microtype}
\PassOptionsToPackage{hidelinks}{hyperref}

% Canonical series preamble.
%% Macros-v2_201: downstream full-carrier geometry update, five current root bindings, no mathematical macro change.
%% Macros-v2_200: Paper III-F v1.280 current-source binding for framework v3.074; mathematical macro definitions unchanged.
%% Macros-v2_198: current project v3.072 III-D and bibliography binding; mathematical macros unchanged.
%% Current project v3.070 source-reference bindings; previous macro bodies unchanged.
% Current II-D reference root for PapersMinimal v3.063; mathematical macros otherwise unchanged.
% Macros-v2_188: re-based current paper identities and A--G family; no change to geometry.
% Macros-v2_184: active Paper III-E v1.200/III-F v1.267 binding for framework v3.031; mathematical macros unchanged.
% Macros-v2_183: updated shared cross-reference bindings for Paper1B/III-E/III-F v3.030; mathematical macros unchanged.
% Macros-v2_182: active Paper II-F/V-A cross-references for framework v3.029; definitions unchanged.
% Paper-series v2.381: new I-B torsion/modular results, bibliography and current source roots.
% Paper-series v2.380: new exceptional bibliography and synchronized paper references.
% Paper-series v2.379 current 31-paper macro/aux registry.
% Paper-series v2.378: current 31-paper AUX/bibliography registry updated for Paper I-B v1.2 literature additions.
% Framework v2.371: current-source reference registry refreshed; old theorem labels retained.
% Framework v2.370: current 30-paper source/AUX root rebinding; native mathematical macros unchanged.
% Framework v2.368: synchronized merged 30-root registry; mathematical notation unchanged.
% Framework v2.358: active 30-paper source registry synchronized; existing mathematical macros unchanged.
% Framework v2.357: selective new source-owned certificates; active reference roots updated for five revised papers.
% Framework v2.356: active 30-root registry synchronized to the 30 papers in this release.
% Framework v2.275: exact positive dense-twin word cubature, finite repeatability and eighth-moment classification, nested sampling frames, joint action and POVM-jet exports; preserves v2.274.
% Framework v2.274: filename and root-source synchronization.
% Framework v2.273: restored missing Problem 3 theorems from v2.268; exact exceptional frames, twin geometry, convolution repeatability and projector/occurrence firewalls; merged on v2.272 without overwriting independent branches.
% Framework v2.272: merged v2.270 covariant native-readout mathematics with v2.271 Problem 4 finite-action Hopf hypercharge/holonomy corrections; no universal bulk-gauge uniqueness claim.
% Independent branches: v2.270 readout/guard exports and v2.270 K16 action-consumed charge closure; all prior scoped assumptions retained.
% Framework v2.269: synchronized 30 active AUX roots; mathematics and bibliography preserved from Macros-v2_161.
% Framework v2.266: Paper II-C converse-derived outcome resources; II-E, III-C, IV-A/B, V-A and VI-B downstream mathematical closures; 30-root registry synchronized.
% Framework v2.265: II-C coordinate-stage constitution and exact first-symbol host inheritance; prior mathematical domains retained.
% Framework v2.264: paper II-B reverse/audit and IV-A observable-minimal exchange; synchronized external references to 30 active roots.
% Framework v2.263: Paper I-C adds bridge-controlled root-carry parity and source-relative old-cap cancellation; other root mathematical bodies retained; active 30-root references synchronized.
% Framework v2.262: integrates Paper II-F radius, Casimir and readout results; active 30-root references synchronized.
% Framework v2.261: merged Paper II-B GP2-level/geometric repairs; preserves v2.260 II-D repairs; active references synchronized.
% Framework v2.260: Paper II-D signed RS, finite-symbol typing, and conditional curvature/Jordan repairs; references synchronized.
% Framework v2.259: conditional native SU(2) atlas and explicit later discharge; v2.258 mathematical bodies retained elsewhere.
% Framework v2.258: preserves parallel II-A/dynamics update and adds complete readout/frequency closures.
% Framework v2.257: exact minimal complete position readout and arbitrary-preparation frequency law.
% Framework v2.256: exact native inertia and sharp bounded ordered-source closure; preserves concurrent v2.255 work.
% Framework v2.255: preserves parallel detector integration and adds exact finite readout closures.
% Framework v2.254: Pasterski bibliography and synchronized 30-paper external-reference registry.
% Framework v2.253: typed raw coordinate defect in I-A and synchronized VI-A import; active references synchronized.
% Framework v2.252: synchronized source registry; optional III-G shell-frame calculation.
% Framework v2.251: forward common-action execution, joint source jets and scoped coupled control.
% Framework v2.251: exact full-heat variations, fermion source jets and explicit physical completion frontier.
% Framework v2.251: certified interacting source jets, pole stability and scoped completion proofs.
% Framework v2.251: synchronized active source registry; mathematical notation unchanged.
% Framework v2.251: synchronized active source registry; mathematical notation unchanged.
% Framework v2.251: synchronized active source registry; mathematical notation unchanged.
% Framework v2.251: synchronized source registry; no mathematical notation changes.
% Framework v2.244: synchronized 30-paper AUX registry and rooted phase-source action-order tests.
% Framework v2.242: signed rooted phase lift, complete transverse return and exact family cubic.
% Framework v2.241: proved rooted conormal geometry, conditional cap realization and scoped action tests.
% Framework v2.239: all 30 active source roots synchronized; no mathematical macro changes.
% Framework v2.238: Paper IV A/B/C split and 30-root reference registry.
% Framework v2.237: Paper VII-D technical split and 28-root reference registry.
% Framework v2.236: Paper VI-C technical-volume reference added; all current root bindings advanced.
% Framework v2.235: 26-root AUX and bibliography bindings synchronized; mathematical macros unchanged.
% BEGIN INLINE SOURCE: Paper00-TeXBuildCompatibility-v1_0.tex
% Rendering-only repairs for inherited long appendices and bitmap text fonts.
% Integer counters through 26 retain the original labels; later values use AA,
% AB, ... . This file changes no action coefficient, state or theorem premise.
\RequirePackage{alphalph}
\makeatletter
\let\BHPrimitiveAlph\@Alph
\newalphalph{\BHExtendedAlph}{\BHPrimitiveAlph}{26}
\let\@Alph\BHExtendedAlph
\AtBeginDocument{%
  \@ifpackageloaded{microtype}{\microtypesetup{expansion=false}}{}%
}
\makeatother

% END INLINE SOURCE: Paper00-TeXBuildCompatibility-v1_0.tex

% Framework v2.164: current reference/bibliography profile; mathematical definitions retained.
% Framework v2.162: finite metric, functional Hilbert and scoped Ward integration.
% Framework v2.161: inherits Macros-v2_57.tex exactly except the current external-reference binding.
% Macros v2.57 / framework v2.156: current reference registry only; mathematical definitions unchanged.
% Framework v2.143 / Macros v2.55: current external-reference bindings; mathematical notation unchanged.
% Framework v2.110 / Macros v2.54: current external-reference registry and bibliography; mathematical notation inherited.
% Framework v2.67 / Macros-v2_53: adds Born-action-cell and joint-EH readout symbols; absolute unit realization is metrological; MATCH_{J/scal}^{EH} remains the sole G24 normalization bridge.
% Framework v2.35 / Macros-v2_52: grouped paper numbering, active AUX registry, and Born-v6_20 bibliography binding synchronized; notation unchanged.
% Framework v2.26: active-source AUX and bibliography bindings synchronized; notation unchanged.
% Framework v2.09: reviewed early-reciprocity integration; scoped physical imports.
% Framework v2.08: Born--Euler mathematics and conditional complete-action potential criterion.
% Framework v2.07: active AUX registry includes Paper I-B v1.2; shared bibliography Born-v6_18.
% Macros-v2_46 / framework v2.04: mathematical atlas and downstream realization correction.
% Macros-v2_46 / framework v2.03: scoped structure integration and current dependency bindings.
% Macros-v2_46 / framework v2.01: shared source bindings; prior revision comments are historical.
% v2.42 / framework v1.99: Synchronized RG/AS owner and bibliography bindings.
% v2.41 / framework v1.99: active manuscript AUX bindings and geometric-closure bibliography synchronized.
% Release-v1.80 centralized source synchronization for scalar-circle/comparison-band integration.
% Release v1.79: Hurwitz hierarchy and kinetic-normalized action response integration; physical mass matching remains conditional.
% Release v1.78: approved M8 closure-batch integration; physical and analytic domains retained.
% Current v1.78 registry: all active TeX roots are synchronized below.
% Earlier release comments are historical scope records, not current status.
% Release-v1.77 synchronization: native same-owner Noether closure, generator--flow/history typing, analytic-certification firewalls, scalaron/K16, LF/QCD, quantum-Ward, and RAR boundary-owner transport are registered; shared mathematical notation is unchanged.
% Release-v1.76 synchronization: registry points to the v1.76 active paper roots.
% Release-v1.76 synchronization: registers the local A+Phi+g Noether/constraint closure and dynamics-frontier propagation; shared mathematical notation is unchanged.
% Release-v1.75 synchronization: registry points to the v1.75 active paper roots.
% Release-v1.75 synchronization: integrates the joint-Dirac Einstein bulk real form and the local odd-scalar heat-coefficient action on explicit analytic first-jet/interior domains; v1.74 Yang--Mills closure is retained. Fermionic reality, pure-metric curvature-squared/exact-nonlocal terms, noncanonical boundaries, physical vacuum selection, and actual solution-history tokens remain open.
% Release-v1.74 synchronization: registry points to the v1.74 active paper roots and retains the v1.73 theorem labels.
% Release-v1.74 synchronization: canonical Macros-v2_34.tex; generated local nonlinear Yang--Mills action, constrained connection-field phase space, source-free Gauss momentum map, and Legendre closure are integrated on their proved domains. The remaining non-gauge ACT_cl sectors and actual solution-history selection remain open.
% Release-v1.71 synchronization: closes the former JOINT_x/JOINT_5
% connection-occurrence gap by canonical structure-group extension of the
% nonflat Born connection, faithful associated representation, and the local
% cotangent--Clifford operator constructor.  Actual fluctuation members,
% actions, vacua, solutions, calibrations, and analytic QFT domains remain
% separately typed.  The active-AUX registry is advanced once.
% Release-v1.71 canonical macro and active-AUX registry.
% ======================================================================
% Macros-v2_39.tex -- canonical publication macro file
% Release v1.71 retains the v1.70 null-ray/QCD boundaries and additionally
% removes the duplicate structural connection actuality premise.  It does not
% claim global bundle topology, select a gauge/Higgs member, or promote the
% finite construction to an interacting continuum theory.
% Release v1.68 constructs the finite light-front total-fraction operators and
% parton number/momentum operator measures, proves their simplex and Hausdorff-limit
% theorems, and separates positive cut measures from scheme-renormalized
% light-ray distributions.  It also separates finite real Schur poles, stable
% norm-resolvent/Riesz pole-residue limits, and continued resonance
% pole/residue limits.  The Paper-VII-C scattering package is claim-indexed:
% local QFT and spectral claims import no S-matrix or infrared premise, while
% asymptotic amplitudes, LSZ, S-matrix poles, completeness, and cross sections
% import only their named massive or massless/long-range clauses.
% Release v1.67 factorizes BND48_H into a same-action Hessian/Schur
% descendant and still-open typed boundary-trace, critical-lift, and
% refinement-limit conditions.  It proves the finite-collar Dirichlet-to-
% Neumann law sqrt(y)tanh(L sqrt(y)), its nonuniform critical limit, and the
% finite-return ultra-infrared turnover.  The Paper-I-A Hardy constant mode is
% an exact geometric near-hit, not yet the physical metric boundary mode.
% The current local f(R) Hessian remains massive and therefore still refutes
% a RAR prediction on the minimal action branch.
% Release v1.66 derives the transverse depth-six normalizer bridge from the
% already generated C6 Reynolds image and Paper-I-A binary-dihedral normalizer.
% The bridge is the unique least 2D3-stable scalar-containing realization of
% the faithful quotient 2D3/C6=C2, has trace defect six and hence K=48, but is
% transverse to the nested active carrier.  Exact K=44 and K=52 contraction,
% semigroup, and energy-midpoint identities are registered as structural
% readouts rather than projector events.  The active E48 carrier gap remains
% unchanged.  A same-action boundary-symbol/Calderon intertwiner is isolated
% as the necessary and still unproved owner bridge to the static DtN sort;
% representation equivalence alone neither derives RAR nor creates a K=48
% action junction, pole, occupation, or cosmological transition.
% Release v1.65 derives the native E_48-centered carrier mass gap from the
% GP1--GP3-generated active spectrum, distinguishes its carrier-free interval
% from a particle located at E_48, and withdraws the unsupported promotion to
% a positive continuum Yang--Mills mass gap.  The latter remains a valid
% conditional transfer theorem on an enlarged continuum branch, but it is not
% an active framework prediction or an unfinished native-gap proof.
% Release v1.64 synchronizes Derivation Batches 14--22.  It registers the
% GP1--GP3 heat-profile orbit, same-Hessian electroweak response and pole
% ownership, K16/K48 action-response corrections, perturbation state and
% non-Gaussianity typing, finite-vacuum/gap/QCD limit repairs, factorized
% position-area ownership, the actual static metric response, and the ordered
% Wilson/string-tension limits.  Strict continuum convergence, positive
% physical gap, surviving higher cumulants, and positive tension remain open.
% Release v1.61 closes TIMEWIDTH16 as an independent direct-production datum.
% JUNC_{2|7} continues to own the finite intrinsic microscopic profile/width,
% while the actual same-action background/production history determines its
% physical-time pullback, support, and duration.  The current Paper-I-A--VII-C AUX
% registry is advanced once; shared mathematical notation is unchanged.
% Release v1.60 derives the final-octet primitive-remnant depletion, canonical
% Hopf-current representative, exact Hopf descent to dOmega/(4pi), and
% closed-FLRW a^{-3} neutral-current shape.  It retypes K=48 as the onset
% of the final one-unit geometric transfer and records cDM/RAR as distinct
% physical projections of that common neutral parent while preserving
% occurrence, Hilbert-owner, stress-matching, and static-kernel premises.
% Release v1.59 transports the v1.58 exceptional-shell/K16 closure through
% Papers III-F and VI-B--VII-C.  It replaces Paper VI-B's obsolete mixed-Hessian
% output typing by the third-variation pair vertex and Bogoliubov occupation
% support, derives the post-seed Higgs-pair feedback vertex and its exact
% alpha--delta kernel factors, propagates the unique K16/K48 action-junction
% split, identifies the QCD delta count with the upstream 2O shell atom only
% as a common-scalar numerical compatibility, and narrows THERM16 to dynamical
% approach/beta selection because the finite Gibbs/KMS target is already
% unique on a supplied simple observable algebra once its Hamiltonian and
% inverse temperature are supplied; direct-sum central weights remain typed.
% Release v1.58 registers the exceptional nearest-shell closure, the derived
% transverse K=16 action event, the unique interior action junction, the
% type-correct ordered pair vertex, the thin-layer squeeze law, and the
% radiation-like/matter-like kinematic split.  REG_0, transfer magnitude,
% physical layer width, and thermalization remain explicitly conditional.
% Release v1.57 registers the active-depth-six obstruction, the native
% K=48 carrier desert, the typed 7|9 owner reduction, and the dynamical
% reheating handoff across Papers III-B--VI-B.
% Release v1.56 registers the least-support n_2 inflation dressing of the
% persistent n_7 Laurent response, the analytic global-convexity and positive
% endpoint-measure proofs, the derived inflation core and K=16 coefficient
% freeze, and the no-particle dark-sector split into static RAR plus the
% conditional homogeneous effective-curvature response at reflected K=48.
% There is no n_6 action, scalaron-DM mass, or second Einstein--Hilbert term.
% Release v1.55 registers the protected-octet resolved/residual reflection
% q_sp+Z_W=8, its j<->8-j and K<->64-K scale covariance, and the derived
% rank-one-exit quadratic grade C_1^2/U=C_2 in Papers III-C, III-E, and VIII.
% The reflection is a typed trace/scale theorem, not a time reversal or a
% physical event-occurrence theorem; reheating transfer remains conditional.
% Release v1.54 registers the exact positive-part Laurent indices, the
% active-domain/terminal-endpoint action split, the no-second-EH firewall, and
% the q_DM=6/K_DM=48 scalaron readout typing across Papers III-D/E/F, VI,
% VIII, and IX.
% Release v1.53 registers the common nonnegative Laurent obstacle projection
% for n_1^+ and n_7^+, discharges the three former exit/freeze assumptions,
% and advances the Paper-III-D/E/F, VIII, and IX consumers.  Shared notation
% and bibliography ownership are unchanged.
% Release v1.51 registers the PQTI derivation batch and its updated
% Paper-III-A/C/D/E/F, VIII, and IX consumers.  Shared notation and
% bibliography ownership are unchanged.
% Release v1.50 advances the Paper-III-B AUX target to v1.40 after the
% rigorous audit repair; shared notation and bibliography ownership are unchanged.
% Release v1.49 advances the Paper-II-E/F AUX target to v3.14 after the
% minimal audit repair; shared notation and bibliography ownership are unchanged.
% Release v1.46 synchronizes Paper II-A's minimal frozen constructor-origin
% theorem and its five dependent AUX targets.  Shared notation is unchanged.
% Release v1.45 updates the Paper-II-B AUX target to v2.23 after the
% GP3-rigidity and appendix-placement repair; shared notation is unchanged.
% Release v1.44 advances the centralized Paper-I-A--VII-C AUX registry after the
% Paper-I-A v5.16 protected-completion scope repair.  Every current
% paper imports this one file; no paper-level opt-in or \externaldocument list
% is permitted.  The bibliography remains centralized here through Born-v6_13.
% Shared bibliography, theorem infrastructure, colors, notation alignment,
% and current Paper-I-A--VII-C external-reference registration.
% This file centralizes common imports, bibliography sorting, colors, theorem
% infrastructure, and the notation shared by the complete series.
% v2.9 (release v1.43) made the Paper-I-A--VII-C AUX registry automatic and
% centralized the bibliography and all external-reference ownership.
% v2.8 release-registry refresh (v1.33) registers the rooted generated-parent
% cascade and the POL/REG/free-law/11+4/Born-group closure revisions.  No
% shared mathematical notation or Paper-II-A fundamental principle is changed.
% v2.8 release-registry refresh (v1.32) registers the universal generated-
% continuation revisions in Papers II-A, III-A, III-C, IV, and VIII.  No
% shared mathematical notation or Paper-II-A fundamental principle is changed.
% v2.8 release-registry refresh (v1.30) registers the CRP4/projective-state
% source revisions and their direct registry-only consumers.  No shared
% mathematical notation or Paper-II-A axiom is changed.
% v2.7 release-registry refresh (v1.27) updates only the current external
% Paper-I-A--VII-C basenames after the placement/assumption synchronization; no
% shared mathematical notation or Paper-II-A axiom is changed.
% v2.7 synchronizes the geometric defect-repair revision: Paper-II-A v3.42,
% Paper-III-C v1.26, Paper-III-D v1.27, Paper-III-E v1.32, and Paper-III-F
% v1.34.  It registers the common trace-defect/spectral-filtration theorem,
% its owner-typed response pullback, and the reduced GRED and TERM interfaces.
% No shared mathematical notation is changed.
% v2.6 synchronizes Paper-III-E v1.31's cone-type repair, minimal-response
% classification, and explicit classical action-role boundary, together with
% Paper-III-F v1.32's downstream response-atlas synchronization.  No shared
% mathematical notation is changed.
% v2.5 synchronizes the Paper-III-D v1.26 faithful-trace-duality proof,
% exact ordered-owner joint locus, and typed alternative-presentation
% bookkeeping repair, together with its Paper-III-E and closure-frontier
% references.  No shared mathematical notation is changed.
% v2.4 synchronizes the Paper-III-C v1.25 orbit/process/owner repair, its
% downstream Papers III-B--F, V, VI, VIII, and IX revisions, and Born-v6_13.
% v2.3 synchronizes the Paper-II-A v3.41 witness/support/principal inheritance
% layer, the Paper-III-A v1.24 joint-realizability/readout firewall, all dependent
% host-rate revisions, and the current Paper-I-A--VII-B source registry.
% v2.2 synchronizes the generated active-realization, determinant-depth,
% dimension-weighted hierarchy, and typed cosmological-event interfaces.
% v2.1 synchronizes the exceptional-support interface revisions: Paper I-A
% v5.9, Paper II-A v3.30, Paper II-E/F v3.7, and Paper IV v4.16.
% v2.0 introduces the stage-safe G0 occurrence notation, synchronizes the
% active Paper-I-A--IV source registry (current rigorous Paper-III closure:
% III-A v1.17, III-B v1.22, III-C--E v1.17, III-F v1.14, IV v4.14), and makes
% external-reference self-skip depend on logical paper role rather than a
% versioned job name.
% ======================================================================

% ======================================================================
% macros.tex — Shared macros for the Born Geometry series
% ======================================================================

% ======================================================================
% 1. PACKAGES, BIBLIOGRAPHY, COLORS, AND UNIT DEFINITIONS
% ======================================================================
% Core packages shared across the complete Paper-I-A--VII-C series.  Loading them here makes
% the macro file a genuine series preamble while remaining compatible with existing
% paper preambles that load the same packages before \input.
\usepackage{amsmath}
\usepackage{amssymb}
\usepackage{amsthm}
\usepackage{mathrsfs}
\usepackage{array}
\usepackage{booktabs}
\usepackage{xcolor}
\makeatletter
\@ifpackageloaded{geometry}{}{\usepackage{geometry}\geometry{margin=1in}}
\makeatother
\usepackage{mdframed}

% Shared support packages required by the macros below and by most papers.
\usepackage{ragged2e}
\usepackage{longtable}
\usepackage{tabularx}
\usepackage{enumitem}
\usepackage[most]{tcolorbox}

% Series-wide numeric bibliography convention using the portable BibTeX path.
% All Papers I-A--VII-C use the ordinary \cite interface.
\usepackage[numbers,sort&compress]{natbib}
\providecommand{\printbibliography}{%
  \bibliographystyle{plainnat}%
  \bibliography{Born-v6_60}%
}

% Shared colors.  \providecolor avoids clobbering a paper-specific override.
\providecolor{postulatecolor}{RGB}{230,242,255}
\providecolor{compare}{RGB}{230,245,230}
\providecolor{definitioncolor}{RGB}{230,242,255}

\providecolor{openriskcolor}{RGB}{255,235,235}
\providecolor{bridgecolor}{RGB}{240,240,255}
\providecolor{gold}{RGB}{255,215,0}

% ======================================================================
% 2. CUSTOM ENVIRONMENTS
% ======================================================================
% Common table column types used in all papers.  Guarded so that a paper may
% define its own versions before loading this file.
\makeatletter
\@ifundefined{NC@find@L}{%
  \newcolumntype{L}[1]{>{\RaggedRight\arraybackslash}p{#1}}%
}{}
\@ifundefined{NC@find@C}{%
  \newcolumntype{C}[1]{>{\Centering\arraybackslash}p{#1}}%
}{}
\@ifundefined{NC@find@Y}{%
  \newcolumntype{Y}{>{\RaggedRight\arraybackslash}X}%
}{}
\makeatother

% ======================================================================
% 3. THEOREM STYLES AND ENVIRONMENTS
% ======================================================================
% Custom underlined theorem style (useful for remarks or notes)
\newtheoremstyle{underlined}
  {3pt}   % Space above
  {3pt}   % Space below
  {}      % Body font (empty means normal/roman text, typical for remarks)
  {}      % Indent amount
  {}      % Theorem head font (can be \itshape, \bfseries, etc.)
  {.}     % Punctuation after theorem head
  { }     % Space after theorem head
  {\underline{\thmname{#1}\thmnumber{ #2}}\thmnote{ (#3)}} % Theorem head formatting

% Standard plain theorem style environments.  Each environment is guarded so
% the shared macro file can be loaded after a paper-specific theorem setup.
\theoremstyle{plain}
\makeatletter
\@ifundefined{theorem}{\newtheorem{theorem}{Theorem}[section]}{}
\@ifundefined{lemma}{\newtheorem{lemma}{Lemma}[section]}{}
\@ifundefined{proposition}{\newtheorem{proposition}{Proposition}[section]}{}
\@ifundefined{corollary}{\newtheorem{corollary}{Corollary}[section]}{}
\makeatother

% Definition style environments (typically non-italicized body)
\theoremstyle{definition}
\makeatletter
\@ifundefined{definition}{\newtheorem{definition}{Definition}[section]}{}
\@ifundefined{postulate}{\newtheorem{postulate}{Postulate}[section]}{}
\@ifundefined{closure}{\newtheorem{closure}{Closure}[section]}{}
\makeatother

% Applying the custom underlined style
\theoremstyle{underlined}
\makeatletter
\@ifundefined{remark}{\newtheorem{remark}{Remark}[section]}{}
\@ifundefined{convention}{\newtheorem{convention}{Convention}[section]}{}
\@ifundefined{principle}{\newtheorem{principle}{Principle}[section]}{}
\makeatother

% Framed proof environment.
% Every \begin{proof} ... \end{proof} is automatically framed/shaded.
% This avoids mdframed/trivlist nesting, which can drop the top border after
% theorem-like environments.
\makeatletter
\renewenvironment{proof}[1][\proofname]{%
  \par\addvspace{6pt}%
  \noindent
  \begin{tcolorbox}[%
    enhanced,
    breakable,
    colback=black!4,
    colframe=black!35,
    boxrule=0.35pt,
    arc=1pt,
    left=6pt,
    right=6pt,
    top=4pt,
    bottom=4pt,
    boxsep=0pt,
    before skip=0pt,
    after skip=6pt,
    parbox=false
  ]%
  \pushQED{\qed}%
  \normalfont\itshape #1\@addpunct{.}\enspace
  \normalfont\ignorespaces
}{%
  \unskip\nobreak\hfill\popQED
  \end{tcolorbox}%
  \par\addvspace{2pt}%
}
\makeatother

\mdfdefinestyle{postulatebox}{
  backgroundcolor=yellow!10,
  linecolor=black,
  linewidth=0.4pt,
  skipabove=\baselineskip,
  skipbelow=\baselineskip,
  nobreak=true
}

% ======================================================================
% PROOF-READABILITY HELPERS
% ======================================================================

\newcommand{\ProofPlan}[1]{%
  \par\smallskip
  \noindent\textbf{Proof plan. }#1%
  \par\smallskip
}

\newlist{proofsteps}{enumerate}{2}
\setlist[proofsteps,1]{
  label=\textbf{(\alph*)},
  ref=\alph*,
  leftmargin=2.2em,
  itemsep=0.45\baselineskip,
  topsep=0.35\baselineskip
}
\setlist[proofsteps,2]{
  label=\textbf{(\roman*)},
  ref=\roman*,
  leftmargin=2.0em,
  itemsep=0.25\baselineskip,
  topsep=0.2\baselineskip
}


% ======================================================================
% 4. NUMBER SYSTEMS AND ALPHABETS
% ======================================================================
\providecommand{\RR}{\mathbb{R}} % Real numbers
\providecommand{\PP}{\mathbb{P}} % Projective space / Probability
\providecommand{\CC}{\mathbb{C}} % Complex numbers
\providecommand{\ZZ}{\mathbb{Z}} % Integers
\providecommand{\NN}{\mathbb{N}} % Natural numbers
\providecommand{\HH}{\mathbb{H}} % Quaternions / Hyperbolic space
\providecommand{\LFRayGeom}{\mathsf{LF\!\mbox{-}RAY\!\mbox{-}GEOM}}
\providecommand{\LFRep}{\mathsf{LF\!\mbox{-}REP}}
\providecommand{\LFState}{\mathsf{LF\!\mbox{-}STATE}}
\providecommand{\LFLim}{\mathsf{LF\!\mbox{-}LIM}}
\providecommand{\LFScat}{\mathsf{LF\!\mbox{-}SCAT}}
% Number systems / projective spaces
\providecommand{\CP}{\mathbb{CP}}
\providecommand{\RP}{\mathbb{RP}}

% ======================================================================
% 5. GROUPS AND LIE THEORY
% ======================================================================
\providecommand{\SO}{\ensuremath{\mathrm{SO}}}     % Special Orthogonal group
\providecommand{\SU}{\ensuremath{\mathrm{SU}}}     % Special Unitary group
\providecommand{\Spin}{\ensuremath{\mathrm{Spin}}} % Spin group
\providecommand{\CS}{\ensuremath{\mathrm{CS}}}     % Chern-Simons (capitalized)
\providecommand{\cs}{\ensuremath{\mathrm{cs}}}     % Chern-Simons (lowercase)

% ======================================================================
% 6. COMMON MANIFOLDS AND TOPOLOGY
% ======================================================================
\providecommand{\Sone}{S^1}                 % Standard circle

% ======================================================================
% 7. MATH OPERATORS (Declared via \DeclareMathOperator)
% ======================================================================
\DeclareMathOperator{\End}{End}   % Endomorphisms
\DeclareMathOperator{\Aut}{Aut}   % Automorphisms
\DeclareMathOperator{\spec}{spec} % Lowercase spectrum
\DeclareMathOperator{\tr}{tr}     % Lowercase trace (if you already have \Tr)
\DeclareMathOperator{\supp}{supp} % Support
\DeclareMathOperator{\TrA}{Tr_{A}} % Partial trace over subsystem A
\DeclareMathOperator{\TrB}{Tr_{B}} % Partial trace over subsystem B
% Lie/algebra operators
\DeclareMathOperator{\Ad}{Ad}
\DeclareMathOperator{\Sym}{Sym}
\DeclareMathOperator{\su}{su}
% ======================================================================
% 8. STANDARD OPERATOR-LIKE SYMBOLS (Safe definitions via \providecommand)
% ======================================================================
\providecommand{\vol}{\operatorname{vol}}   % Volume form
\providecommand{\Vol}{\operatorname{Vol}}   % Total volume
\providecommand{\Ric}{\operatorname{Ric}}   % Ricci tensor
\providecommand{\Rm}{\operatorname{Rm}}     % Riemann tensor
\providecommand{\Tr}{\operatorname{Tr}}     % Uppercase trace
\providecommand{\Spec}{\operatorname{Spec}} % Uppercase spectrum
\providecommand{\diag}{\operatorname{diag}} % Diagonal matrix
\providecommand{\rank}{\operatorname{rank}} % Rank
\providecommand{\scal}{\mathcal R} % Scalar curvature
\providecommand{\Cl}{\operatorname{Cl}}     % Clifford algebra
\providecommand{\diam}{\operatorname{diam}} % Diameter

% ======================================================================
% 9. FUNCTIONS, TRIGONOMETRY, AND CALCULUS
% ======================================================================

% Codifferential naming conventions
% (Avoid \ddag: commonly defined by physics.sty and symbol packages)
\providecommand{\codiff}{\mathrm{d}^{\dagger}} % Codifferential operator

% Quantum mechanics bra-ket notation (needed for basis tables)
\providecommand{\ket}[1]{\lvert #1 \rangle}
\providecommand{\bra}[1]{\langle #1 \rvert}
\providecommand{\braket}[2]{\langle #1 \mid #2 \rangle}

% ======================================================================
% 10. PHYSICAL SCALES AND CONSTANTS (Paper 1c / Born Geometry)
% ======================================================================
\providecommand{\Rp}{\ensuremath{R_{\mathrm P}}}     % matched common EH curvature scale (not defined from observed G)
\providecommand{\hBorn}{\ensuremath{h_{\mathrm B}}} % physical representative 2\pi\hbar of one primitive Born scalar action cell
\providecommand{\UEHjoint}{\ensuremath{U_{J}}} % joint-Dirac EH curvature-coordinate scale 8 kappa_J
\providecommand{\Gread}{\ensuremath{G_{\rm scal}}} % scalar-action coordinate c_op^3/(h_B U); physical matching explicit
\providecommand{\GreadJoint}{\ensuremath{G_{J,\mathrm{raw}}}} % raw ordinary-trace coordinate = G_scal/96
\providecommand{\Rds}{R_{\mathrm{dS}}}    % de Sitter Curvature
\providecommand{\Mpl}{M_{\mathrm{Pl}}}    % Planck mass
\providecommand{\lp}{\ell_{\mathrm{Pl}}}  % Planck length
\providecommand{\Pmax}{P_{\mathrm{max}}}  % Maximum momentum
\providecommand{\Mplbar}{\bar{M}_{\mathrm{Pl}}} % Reduced Planck mass
\providecommand{\eff}{\mathrm{eff}}       % Subscript for effective quantities
\providecommand{\Cosmo}{\ensuremath{\Lambda}}                % Cosmological constant (use \Cosmo in source)
\providecommand{\LCDM}{\ensuremath{\Cosmo\mathrm{CDM}}}  % \LambdaCDM (typesets with \Cosmo)
\providecommand{\Lcut}{\ensuremath{\Lambda_{\mathrm{cut}}}}  % Cutoff scale (brace when used as script)
\providecommand{\Rcs}{\ensuremath{R_{\mathrm{cs}}}}      % Configuration-space radius
\providecommand{\CosmoEff}{\ensuremath{\Cosmo_{\eff}}}   % Effective cosmological constant

% Disambiguation: \Rmom is the momentum-space (S^3) radius. 
% Reserve plain R for the 4D curvature scalar in f(R).
\providecommand{\Rmom}{\ensuremath{R_{\mathrm{mom}}}} 
\providecommand{\Punit}{P_{0}} % Momentum unit set by \Rmom (= \hbar/\Rmom)
\providecommand{\Rvec}{\ensuremath{R_{\mathrm{vec}}}} % invariant vector-space radius in the reduced round sector
% Specific numeric values and expressions
\providecommand{\Vlogcap}{L_{\mathrm{Hopf}}^{(1)}}
\providecommand{\Hopf}{\operatorname{Hopf}}
% Scalar phase (S^1) sector + universal cover notation (Paper I-A)
\providecommand{\tper}{T_\star}      % Time period/scale
\providecommand{\Eper}{E_\star}      % Energy period/scale
\providecommand{\ttilde}{\tilde t}   % Universal cover time
\providecommand{\Etilde}{\tilde E}   % Universal cover energy

% ======================================================================
% 10A. SERIES-WIDE LOGICAL / NOTATIONAL DISAMBIGUATIONS
% ======================================================================
% Use these symbols to avoid collisions among: (i) the Paper--III
% no-insertion/geometric-closure class, (ii) Paper--I-A trimer diagonals,
% (iii) Paper--V-A/V-C family-space diagonals, (iv) scalar-index running,
% and (v) protected shell-clock labels.
\providecommand{\Rprim}{\ensuremath{\mathsf R_{\mathrm{prim}}^{\Sigma}}} % primitive finite structural readouts
\providecommand{\Rloop}{\ensuremath{\mathsf R_{\circ}^{\Sigma}}} % closed-loop readout role
\providecommand{\Rbdry}{\ensuremath{\mathsf R_{\partial}^{\Sigma}}} % boundary-conditioned readout role
\providecommand{\Rgen}{\ensuremath{\mathsf R_{\mathrm{fin}}^{\Sigma}}} % finite generated closure of primitive roles
\providecommand{\PhysOcc}{\operatorname{PhysOcc}}
\providecommand{\Actual}{\operatorname{Actual}}
\providecommand{\GeomLic}{\operatorname{GeomLic}}
\providecommand{\InvDist}{\operatorname{InvDist}}
\providecommand{\Dcl}{\ensuremath{\mathsf{D}_1}}      % unmarked positive Dirac-even closed spectral-readout form class
\providecommand{\Dheat}{\ensuremath{\mathsf D_{\mathrm{heat}}}} % one-parameter heat-semigroup subclass
\providecommand{\DtriI}{\ensuremath{D^{\mathrm{tri}}_1}} % Paper--I-A cyclic-trimer diagonal
\providecommand{\DtriII}{\ensuremath{D^{\mathrm{tri}}_2}}
\providecommand{\DfamI}{\ensuremath{D^{\mathrm{fam}}_1}} % Family-space diagonal selector
\providecommand{\DfamII}{\ensuremath{D^{\mathrm{fam}}_2}}
\providecommand{\alpharun}{\ensuremath{\alpha_{\mathrm{run}}}} % scalar spectral-index running
\providecommand{\Ksh}{\ensuremath{K_{\mathrm{sh}}}} % protected shell-clock coordinate
\providecommand{\ddef}{\ensuremath{\delta_{\mathrm{def}}}} % universal defect scalar after Paper--III realization

% ======================================================================
% 10B. SERIES-WIDE LOGICAL STATUS TAGS
% ======================================================================
% Upright small-caps wrapper for audit/status labels.  It prevents LaTeX from
% trying to use unavailable small-caps italic font shapes inside theorem bodies
% such as TU/lmr/m/scit.
\providecommand{\AuditLabel}[1]{}
\renewcommand{\AuditLabel}[1]{{\normalfont\scshape #1}}

% These tags are bookkeeping labels for consistency audits. They do not
% create new mathematical hypotheses; they record how a datum is being used
% in the proof chain.
\providecommand{\StatusPrimitive}{}
\renewcommand{\StatusPrimitive}{\AuditLabel{Primitive}}
\providecommand{\StatusDerived}{}
\renewcommand{\StatusDerived}{\AuditLabel{Derived}}
\providecommand{\StatusImported}{}
\renewcommand{\StatusImported}{\AuditLabel{Imported}}
\providecommand{\StatusConditional}{}
\renewcommand{\StatusConditional}{\AuditLabel{Conditional}}
\providecommand{\StatusBridge}{}
\renewcommand{\StatusBridge}{\AuditLabel{Bridge-derived}}
\providecommand{\StatusComparison}{}
\renewcommand{\StatusComparison}{\AuditLabel{Comparison-only}}

% ======================================================================
% 10C. SERIES-WIDE CLAIM-WORDING AUDIT LABELS
% ======================================================================
% These are textual audit labels used to keep modal language consistent
% across the series.  They introduce no new mathematical assumptions.
\providecommand{\ClaimForced}{}
\renewcommand{\ClaimForced}{\AuditLabel{forced}}
\providecommand{\ClaimUnique}{}
\renewcommand{\ClaimUnique}{\AuditLabel{unique}}
\providecommand{\ClaimCanonical}{}
\renewcommand{\ClaimCanonical}{\AuditLabel{canonical}}
\providecommand{\ClaimMinimal}{}
\renewcommand{\ClaimMinimal}{\AuditLabel{minimal}}
\providecommand{\ClaimClosed}{}
\renewcommand{\ClaimClosed}{\AuditLabel{closed}}
\providecommand{\ClaimPrediction}{}
\renewcommand{\ClaimPrediction}{\AuditLabel{prediction}}

% ======================================================================
% 10D. SERIES-WIDE THEOREM-CHAIN AUDIT LABELS
% ======================================================================
% These labels mark the direction of load-bearing theorem chains. They are
% audit annotations only; they introduce no new mathematical assumptions.
\providecommand{\ChainInput}{}
\renewcommand{\ChainInput}{\AuditLabel{Chain input}}
\providecommand{\ChainOutput}{}
\renewcommand{\ChainOutput}{\AuditLabel{Chain output}}
\providecommand{\ChainGuard}{}
\renewcommand{\ChainGuard}{\AuditLabel{Chain guard}}
\providecommand{\ChainNoBackflow}{}
\renewcommand{\ChainNoBackflow}{\AuditLabel{No backflow}}

% ======================================================================
% 10E. SERIES-WIDE LOCAL THEOREM-AUDIT LABELS
% ======================================================================
% These labels mark a statement-level dependency audit around individual
% theorem/definition nodes. They are bookkeeping labels only.
\providecommand{\LocalAudit}{}
\renewcommand{\LocalAudit}{\AuditLabel{Local audit}}
\providecommand{\UsesOnly}{}
\renewcommand{\UsesOnly}{\AuditLabel{Uses only}}
\providecommand{\Exports}{}
\renewcommand{\Exports}{\AuditLabel{Exports}}
\providecommand{\DoesNotUse}{}
\renewcommand{\DoesNotUse}{\AuditLabel{Does not use}}

% ======================================================================
% 10F. SERIES-WIDE NUMERICAL-ROW AUDIT LABELS
% ======================================================================
% These labels classify individual rows in numerical, benchmark, status,
% and comparison tables.  They are audit annotations only and do not change
% any theorem, value, formula, or comparison number.
\providecommand{\RowAudit}{}
\renewcommand{\RowAudit}{\AuditLabel{Numerical-row audit}}
\providecommand{\RowTheorem}{}
\renewcommand{\RowTheorem}{\AuditLabel{Theorem-fixed row}}
\providecommand{\RowBridge}{}
\renewcommand{\RowBridge}{\AuditLabel{Bridge-fixed row}}
\providecommand{\RowRepresentative}{}
\renewcommand{\RowRepresentative}{\AuditLabel{Representative row}}
\providecommand{\RowComparison}{}
\renewcommand{\RowComparison}{\AuditLabel{Comparison row}}
\providecommand{\RowCalibrated}{}
\renewcommand{\RowCalibrated}{\AuditLabel{Calibrated row}}
\providecommand{\RowEstimate}{}
\renewcommand{\RowEstimate}{\AuditLabel{Estimate row}}
\providecommand{\RowOpen}{}
\renewcommand{\RowOpen}{\AuditLabel{Open row}}
\providecommand{\RowNoSelector}{}
\renewcommand{\RowNoSelector}{\AuditLabel{No selector}}

% ======================================================================
% 10G. TABLE HELPERS
% ======================================================================
% Table helpers used by Paper V-C.  The unused readout-label preset has been
% removed; only active table helpers remain.
\providecommand{\gpsection}[1]{%
  \addlinespace[0.55em]
  \multicolumn{3}{@{}l}{\textsc{#1}}\\[-0.25em]
  \midrule
}
\providecommand{\mathcell}[1]{%
  \(\begin{gathered}#1\end{gathered}\)%
}

\newcommand{\DerBy}[1]{%
  \begin{tabular}[t]{@{}l@{}}#1\end{tabular}%
}

\newcommand{\TableNotes}[1]{%
  \par\medskip
  \begin{minipage}{0.94\linewidth}
  \footnotesize\emph{Notes.} #1
  \end{minipage}\par
}
% ======================================================================
% 10H. PAPER VI-A / QCD NOTATION
% ======================================================================
% Shared symbols used by Paper VI-A and downstream QCD/Yang--Mills handoffs.
\providecommand{\Pqcd}{\mathfrak P_{\rm VIII}^{\rm QCD}}
\providecommand{\Iqcd}{\mathfrak I_{\rm QCD}^{(7)}}
\providecommand{\muQCD}{\mu_{\rm QCD}^{\rm int}}
\providecommand{\LambdaQCD}{\Lambda_{\rm QCD}}
\providecommand{\LambdaBH}{\Lambda_{\rm BH}^{\rm QCD}}
\providecommand{\HQCD}{\mathcal H_{\rm QCD}^{\rm BH}}
\providecommand{\HQCDfin}{H_{\rm QCD}^{\rm fin}}
\providecommand{\Hraw}{\mathcal H_{\rm raw}}
\providecommand{\HphysQCD}{\mathcal H_{\rm phys}^{\rm QCD}}
\providecommand{\Vcol}{V_C}
\providecommand{\Ps}{P_{\rm sing}}
\providecommand{\PiOct}{\Pi_8}
\providecommand{\MSbar}{\overline{\rm MS}}
\providecommand{\YM}{\mathrm{YM}}
\providecommand{\Ktri}{K_{\triangle}}
\providecommand{\KYM}{K_{\YM}^{\rm BH}}
\providecommand{\aBH}{a_{\rm BH}}
\providecommand{\epsBH}{\varepsilon_{\rm BH}}
\newcommand{\RowConstructive}{\textnormal{\textsc{Constructive}}}
\newcommand{\RowObligation}{\textnormal{\textsc{Bridge obligation}}}
\newcommand{\RowTarget}{\textnormal{\textsc{Target row}}}
\newcommand{\RowValidation}{\textnormal{\textsc{Validation row}}}
\providecommand{\slashed}[1]{\not{#1}}

% ======================================================================
% 10I. PAPER VI-B / CONTINUOUS-QFT NOTATION
% ======================================================================
% Shared symbols used by Paper VI-B and any downstream continuous-QFT or
% Yang--Mills handoff paper.  All definitions are guarded so individual
% papers may override them locally before loading this file if needed.
\providecommand{\A}{\mathfrak A}
\providecommand{\Hphys}{\mathcal H_{\rm phys}}
\providecommand{\QCD}{\mathrm{QCD}}
\providecommand{\BH}{\mathrm{BH}}
\providecommand{\loc}{\mathrm{loc}}

\providecommand{\RBH}{\mathcal R_{\rm BH}}
\providecommand{\DG}{\mathfrak D_G}
\providecommand{\PW}{\mathcal{P\!W}} % Peter--Weyl spectral-symbol subspace
\providecommand{\Toe}{\operatorname{Toe}} % Toeplitz compression/representative
\providecommand{\Hol}{\operatorname{Hol}} % Holonomy

% ======================================================================
% 11. LATE-LOADED PACKAGES (Must be at the end of the preamble)
% ======================================================================
% Load late packages only if a paper has not already loaded them, avoiding
% option clashes in paper-specific preambles.
\makeatletter
\@ifpackageloaded{hyperref}{}{\usepackage{hyperref}}
\makeatother
% --- Series-wide disambiguations (Pass 3) ---
\providecommand{\diracdelta}{\delta_{\rm D}} % Dirac delta distribution (delta reserved for hierarchy)
\providecommand{\Hg}{H_{\mathrm g}} % Hodge harmonic projector (keep H for Hubble)
\providecommand{\Gr}{G_{\mathrm r}} % Green operator (keep G for Newton)
\providecommand{\FD}{\mathrm{D}} % functional/variational derivative symbol (delta reserved)
\providecommand{\Id}{\mathbf{1}} % identity/Kronecker tensor (avoid \delta_{ij})


% BEGIN INLINE SOURCE: Paper00-ActiveExternalReferences-v2_130.tex
% Framework v2.191: current reference binding; existing definitions retained.
% Framework v2.190: current22-root registry with native clock-evolution proofs.
% Framework v2.189: exact current22-root external registry, including retained v2.188 hardening owners.
% Framework v2.187: current registry synchronization only; source release v2.186 preserved.
% Framework v2.186: authoritative current 22-paper registry.
% Framework v2.182: all current 22 logical paper AUX bindings.
% Framework v2.170: current I-B v1.43 external AUX binding.
% Framework v2.167: exact current 22-root registry after CP--BL2 integration.
% Framework v2.166: exact current 22-root compact-release AUX registry.
% Framework v2.165: authoritative current root bindings.
% Framework v2.164: exact current 22 manuscript AUX bindings.
% Framework v2.162: finite metric, functional Hilbert and scoped Ward integration.
% Framework v2.161: current 22 manuscript AUX bindings; source only.
% Framework v2.156 / registry v2.113: synchronized current 22 manuscript roots; source-only registry.
% Framework v2.143 / registry v2.111: active manuscript AUX basenames; no generated AUX or PDF included.
% ======================================================================
% Central external-reference imports, numbers only, no PDF jumping
% ======================================================================

% xr-hyper is loaded centrally so that papers may retain their existing
% \usepackage{xr-hyper} lines without redefining \externaldocument after this
% registry is installed.  If plain xr was loaded earlier, retain it.
\makeatletter
\@ifpackageloaded{xr-hyper}{}{%
  \@ifpackageloaded{xr}{}{\usepackage{xr-hyper}}%
}
\makeatother

% Make external-document registration idempotent.  The registry identifies a
% logical paper by its Paper<number><optional letter> prefix.  Prefix-qualified
% imports remain independent.  Current paper roots contain no local
% \externaldocument calls; this guard also makes accidental duplicates harmless.
\ExplSyntaxOn
\prop_new:N \g_bh_external_refs_prop
\cs_new_eq:NN \BHOriginalExternalDocument \externaldocument

\cs_new_protected:Npn \bh_external_document:nnnn #1#2#3#4
 {
  \tl_set:Nn \l_tmpa_tl {#3}
  \regex_replace_once:nnN
   { \A (?: xref/ )? (Paper[0-9]+[A-G]?) .* \Z }
   { \1 }
   \l_tmpa_tl
  \tl_set:Nx \l_tmpb_tl { #1 | \tl_use:N \l_tmpa_tl }
  \prop_if_in:NVTF \g_bh_external_refs_prop \l_tmpb_tl
   {
    \PackageInfo{BornHopf}
     {Skipping~duplicate~external~refs~for~#3}
   }
   {
    \prop_gput:NVn \g_bh_external_refs_prop \l_tmpb_tl {#3}
    \IfNoValueTF{#4}
     { \BHOriginalExternalDocument[#1][#2]{#3} }
     { \BHOriginalExternalDocument[#1][#2]{#3}[#4] }
   }
 }

\RenewDocumentCommand{\externaldocument}{O{} O{} m o}
 { \bh_external_document:nnnn {#1}{#2}{#3}{#4} }

\NewDocumentCommand{\BHExternalRefs}{m}
 {
  % Do not import the aux file of the logical paper currently being compiled,
  % even when the job and registry use different version suffixes.
  \tl_set:Nx \l_tmpa_tl { \c_sys_jobname_str }
  \regex_replace_once:nnN
   { \A (Paper[0-9]+[A-G]?) .* \Z }
   { \1 }
   \l_tmpa_tl
  \tl_set:Nn \l_tmpb_tl { #1 }
  \regex_replace_once:nnN
   { \A (Paper[0-9]+[A-G]?) .* \Z }
   { \1 }
   \l_tmpb_tl
  \str_if_eq:eeTF { \tl_use:N \l_tmpa_tl } { \tl_use:N \l_tmpb_tl }
   {
    \PackageInfo{BornHopf}
     {Skipping~external~refs~for~#1~because~this~is~the~current~document}
   }
   {
    \IfFileExists{xref/#1.aux}
     { \externaldocument[][nocite]{xref/#1}[#1.pdf] }
     {
      \IfFileExists{#1.aux}
       { \externaldocument[][nocite]{#1} }
       {
        \PackageInfo{BornHopf}
         {External~aux~file~#1.aux~not~found;~no~labels~imported}
       }
     }
   }
 }
\ExplSyntaxOff

% External references are automatic.  This file is the sole owner of both the
% bibliography and the current Paper-I-A--VII-D AUX registry (30 roots including I-C).  Missing auxiliaries
% are informational on the first pass and are populated by the source-only
% two-pass series build.  Basenames deliberately omit temporary download/copy
% suffixes such as ``(1)'' or ``-1''.
  % Paper I-A: round S3 compendium.
  \BHExternalRefs{Paper1A-Geometry-v5_181}

  % Paper I-B: complete exceptional geometry and finite spectral structures.
  \BHExternalRefs{Paper1B-ExceptionalGeometry-v1_26}
  % Paper I-D: technical companion of I-A through I-C.
  \BHExternalRefs{Paper1D-TechnicalMaterial-v1_69}
  % Paper I-C: reciprocal M8 compendium.
  \BHExternalRefs{Paper1C-M8Geometry-v1_146}

\ifdefined\BHMathOnly
  % Papers I-A and I-B import only the mathematical atlas labels.
\else
  % Paper II: five physical-mathematical main papers plus one technical proof volume.
  \BHExternalRefs{Paper2A-GeometricConstitution-v3_227}
  \BHExternalRefs{Paper2B-MomentumCarrier-v2_147}
  \BHExternalRefs{Paper2C-ReciprocalLocalization-v2_176}
  \BHExternalRefs{Paper2D-SymmetrySpectra-v2_145}
  % The former II-F export labels remain in II-E; current II-F is the technical volume.
  \BHExternalRefs{Paper2E-KinematicExports-v3_140}
  % Paper II-F: complete technical proof and calculation owner for II-A--II-E.
  \BHExternalRefs{Paper2F-TechnicalMaterial-v1_71}

  % Paper III: Dynamics and Intrinsic Geometric Constitution (six components).
  \BHExternalRefs{Paper3A-GeneratedDynamics-v1_181}
  \BHExternalRefs{Paper3B-SpectralAction-v1_192}
  \BHExternalRefs{Paper3C-ProtectedHierarchy-v1_196}
  \BHExternalRefs{Paper3D-RunningResponse-v1_168}
  \BHExternalRefs{Paper3E-GravityAction-v1_204}
  \BHExternalRefs{Paper3F-CompleteDynamics-v1_280}
  % Paper III-G: complete preserved technical proof owner.
  \BHExternalRefs{Paper3G-TechnicalMaterial-v1_57}

  % Paper IV-A: generated finite quantum mechanics.
  \BHExternalRefs{Paper4A-FiniteQuantumMechanics-v4_217}
  % Paper IV-B: complete finite quantum processes and physical response.
  \BHExternalRefs{Paper4B-QuantumResponse-v1_55}
  % Paper IV-C: extended proofs and finite technical calculations.
  \BHExternalRefs{Paper4C-TechnicalMaterial-v1_50}

  % Paper V-A: The Standard Model.
  \BHExternalRefs{Paper5A-FiniteMatterGeometry-v3_219}

  % Paper V-B: Finite Spectrum and Flavor Closure.
  \BHExternalRefs{Paper5B-TerminalSpectrum-v5_202}

  % Paper V-C: Finite Interaction Boundaries of the Standard Model.
  \BHExternalRefs{Paper5C-InteractionResponse-v2_185}
  % Paper V-D: preserved technical proof and calculation owner.
  \BHExternalRefs{Paper5D-TechnicalMaterial-v1_58}

  % Paper VI-A: Gravity and Closed Cosmology.
  \BHExternalRefs{Paper6A-GravityCosmology-v2_268}

  % Paper VI-B: Geometric Dark Sector and Cosmological Readouts.
  \BHExternalRefs{Paper6B-CosmologyReadouts-v1_215}

  % Paper VI-C: complete technical proof and comparison owner.
  \BHExternalRefs{Paper6C-TechnicalMaterial-v1_53}

  % Paper VII-A: Finite QCD Kernel and Nonperturbative Closure.
\BHExternalRefs{Paper7A-QCDKernel-v1_184}

  % Paper VII-B: Hadronic, Thermal, and Nuclear QCD Readouts.
  \BHExternalRefs{Paper7B-QCDReadouts-v1_168}

  % Paper VII-C: Continuous Quantum Field Theory.
  \BHExternalRefs{Paper7C-ContinuousQFT-v2_208}

  % Paper VII-D: complete technical material for VII-A through VII-C.
  \BHExternalRefs{Paper7D-TechnicalMaterial-v1_60}

\fi

% END INLINE SOURCE: Paper00-ActiveExternalReferences-v2_130.tex



\title{\textbf{Born--Hopf Geometry in a Closed Universe V-C: Finite Interactions and Observable Response}\\[0.45em]
\large Joint action, coherent currents, boundary receivers, and measured response}

\author{Peter Kugelmann\thanks{\texttt{pkugelmann@proton.me}}}
\date{Version 2.185 --- 20 November 2026 \\ Framework v3.075}
\newcommand{\PaperSevenRevision}{v2.106 gauge variation and source normalization}
\newcommand{\Hkin}{\mathsf{H}_{\rm kin}}
\newcommand{\Hqft}{\mathsf{H}_{\rm QFT}}
\newcommand{\Hem}{\mathsf{H}_{\rm EM}}
\newcommand{\Hvac}{\mathsf{H}_{\rm vac}}
\newcommand{\Hcp}{\mathsf{H}_{\rm CP}}
\newcommand{\Hnu}{\mathsf{H}_{\nu}}
\newcommand{\Hgauge}{\mathsf{H}_{\rm gauge}}
\newcommand{\Aem}{\mathcal A_{\rm EM}}
\newcommand{\YN}{Y_\nu}
\newcommand{\Gfive}{\mathsf G_5^{\rm fin}}
\newcommand{\Rrep}{\mathsf R_5^{\rm rep}}
\newcommand{\Rconn}{\mathsf J_5^{\rm conn,der}}
\newcommand{\Rssqft}{\mathsf R_5^{\rm SSQFT}}
\newcommand{\JOINTfiveocc}{\mathsf{JOINT}_{5}^{\rm der}}
\newcommand{\CALgauge}{\mathsf{CAL}_{\rm gauge}}
\newcommand{\REGsig}{\mathsf{REG}_{\sigma}}
\newcommand{\REGzero}{\mathsf{REG}_{0}}
\newcommand{\Ustr}{\mathcal U_6^{\rm str}}
\newcommand{\Urole}{\mathcal U_6^{\rm role}}
\newcommand{\Uratio}{\mathcal U_6^{\rm ratio}}
\newcommand{\Uabs}{\mathcal U_6^{\rm abs}}
\newcommand{\Uroleratio}{\mathcal U_6^{\rm role,ratio}}
\newcommand{\Uroleabs}{\mathcal U_6^{\rm role,abs}}
\newcommand{\Istr}{\mathfrak I_{5C}^{\rm str}}
\newcommand{\Irole}{\mathfrak I_{5C}^{\rm role}}
\newcommand{\Brawsix}{\mathsf B_{6}^{\rm hist}}
\newcommand{\Arawsix}{\mathfrak A_{6,{\rm fin}}^{\rm hist}}
\newcommand{\Acompsix}{\mathfrak A_{6,{\rm fin}}^{\rm comp}}
\newcommand{\CMPexp}{\mathsf{CMP}_{\rm exp}}
\newcommand{\Hresopt}{\mathsf H_{\rm res}^{\rm opt}}
\newcommand{\Hhopfopt}{\mathsf H_{\rm Hopf}^{\rm opt}}

\begin{document}
% Publication line-breaking flexibility for technical prose (no math scaling).
\emergencystretch=2em
\maketitle
\begin{abstract}
Paper V-C assembles the generated finite interaction vertices and
their common-action electromagnetic, weak, strong, scalar and boundary
responses. Observable current and density returns are distinguished from
physical production--propagation--detection maps; a generated CP invariant
is not by itself a baryogenesis yield. The precise gauge and source variation,
weak-current returns, thermal boundary, and existing negative comparison
tests retain their explicit hypotheses. Paper V-D preserves auxiliary
calculations without supplying a new occurrence mechanism.
\end{abstract}
\tableofcontents
\clearpage

\section{Claim status, imports and the actual response target}
\label{p5c:sec:organized-scope}
\subsection{Upstream package and dependency boundary}
\label{p7v2:sec:imports}

Paper~V-C uses the framework principles in their established meanings:
GP1 is total boundedness of the constituted geometry and its admissible
operations, GP2 is geometric primacy, and GP3 is minimality.  No additional
meaning is assigned to any of the three principles in this paper.

The upstream input is split by theorem domain before any physical role name is
used.  The structural and named-role domains are
\begin{equation}
\label{p7v2:eq:upstream-package}
\boxed{
\Ustr=\mathfrak F_{\leq4}+\Gfive,
\qquad
\Urole=\Ustr,
\qquad
\Rrep:=\varnothing.
}
\end{equation}
On \(\Ustr\), Paper~V-C imports the structural output
\begin{equation}
\label{p7v2:eq:structural-import}
\Istr
=
\left(
 \mathcal C_{\rm fin},\,
 \mathcal A_{\rm act}^{(4,5,7)},\,
 \mathcal H_N^{\rm odd},\,
 \mathcal R_{\rm carrier}^{\rm str}
\right),
\end{equation}
where the active support chain is
\((2T,C_6,\mathbb R\mathbf1)\) at \(K=32,40,56\), with terminal scalar Haar
projector, and \(\mathcal R_{\rm carrier}^{\rm str}\) contains the tensor
carrier, structural projectors, color incidence, Abelian candidate plane and
kernel, and the generated even relative-chirality ancestry.  Named physical
hypercharge and the hypercharge-typed matter/Yukawa roles, and hence the
physical notation \(G_{\rm SM}\) and \(\mathcal R_{\rm matter}\), are licensed only on
\begin{equation}
\label{p7v2:eq:physical-role-import}
\Irole=(\Istr,G_{\rm SM},\mathcal R_{\rm matter})
\qquad\text{over }\Urole.
\end{equation}
This is a domain statement, not an additional numerical selector.

The determinant-depth character is structural, but its interpretation as the
physical energy law is separately typed.  Define
\begin{equation}
\label{p7v2:eq:energy-import-domains}
\boxed{\begin{aligned}
\Uratio&=\Ustr+(\REGsig\in\mathfrak A_{\rm REG}^{+})
 +(\mathrm{ANC}_E),\\
\Uabs&=\Uratio+(\mathrm{MOMENERGY}_{0}^{\rm der}),\\
\Uroleratio&=\Urole+(\REGsig\in\mathfrak A_{\rm REG}^{+})
 +(\mathrm{ANC}_E),\\
\Uroleabs&=\Uroleratio+(\mathrm{MOMENERGY}_{0}^{\rm der}).
\end{aligned}}
\end{equation}
The tag \(\mathrm{MOMENERGY}_{0}^{\rm der}\) denotes the theorem-level import
\(E_0=\hbar c_{\rm op}\sqrt U=\sqrt{18/\pi}\,\hbar c_{\rm op}/R_{\rm mom}\); it is
not a calibration premise.
Thus \(E(K)/E(K_0)\) is a physical energy ratio only on the corresponding
ratio domain, while an absolute native shell value additionally imports the derived G24
anchor \(E_0=\sqrt{18/\pi}\,\hbar c_{\rm op}/R_{\rm mom}\); numerical unit conversion
is metrological.  Dimension counting and numerical calibration do
not supply \((\mathrm{ANC}_E)\).

The numerical spectrum/flavor coordinates used in conditional Paper--V-C comparison readouts
belong instead to a frozen Paper--V-B raw comparison branch.  The flavor entries in this tuple are now the native current maps themselves; conventional CKM/PMNS matrices are optional coordinate sections of these maps,
\begin{equation}
\label{p7v2:eq:upstream-raw-branch}
\Brawsix
=
\left(
 v,\lambda_H,\,
 D_e,D_u,D_d,D_\nu,\,
 [C_{ud}],[C_{e\nu}],\,
 \mathcal M_N^{\rm aff}
\right),
\end{equation}
whose finite/raw enumeration atlas is
\begin{equation}
\label{p7v2:eq:paper6-selector-atlas}
\boxed{
\Arawsix
=
\mathfrak A_{\rm SLOT}
\times\mathfrak A_{\rm TMF}
\times\mathfrak A_{\rm CLK}
\times\mathfrak A_{\mathrm{CMP,exp}}
\times\mathfrak A_H
\times\mathfrak A_F.
}
\end{equation}
Here the comparison member is exactly \(\CMPexp\): the particle exponent
function modulo the proved word-lift redundancy.  Paper~V-B defines the actual
finite/raw candidate class by joint compatibility.  Paper~V-C therefore
imports only
\begin{equation}
\boxed{\Acompsix\subseteq\Arawsix}
\end{equation}
of tuples satisfying Paper~V-B's typing, covariance, positivity,
self-adjointness, and cross-sector conditions.  A formal Cartesian tuple
outside \(\Acompsix\) is not a physical upstream branch.  The fixed-shell
scheme is the separate family
\begin{equation}
\label{p7v2:eq:reg-family-import}
\mathfrak A_{\rm REG}
=
\{\REGsig:\sigma_{\rm fin}\in\mathbb R\},
\qquad
\delta_\sigma
=
\frac1{6\sqrt2}+\sigma_{\rm fin},
\end{equation}

% BEGIN M8 v1.78 VII_IMPORT
\begin{remark}[Anchored REG and finite interaction scope]
\label{p7:rem:m8-batch-interaction-inputs}
The active Hurwitz/Hopf branch inherits \(\mathrm{REG}_0\)
from Theorem~\ref{p3c:thm:occurring-hurwitz-reg-origin};
the scheme family in the generic formulas is the unanchored
domain, not a second choice on that branch.
Proposition~\ref{p4:prop:paired-berezin-determinant-line}
and Theorem~\ref{p4:thm:finite-fermion-clock-real-form}
apply to their generated finite field/dual owner.
They do not select \(\bar\theta\), an unpaired determinant
phase, a nonzero Yukawa/Majorana member or a compatible
chiral regulator.  The conditional quadratic stabilizer
in V leaves those physical interfaces explicit.
Vertices and their Noether identities derive from the same
claim-complete action, with its sources and boundary roles.
\end{remark}
% END M8 v1.78 VII_IMPORT


with \(\REGzero\) the source member \(\sigma_{\rm fin}=0\).  It is the seventh
historical Paper--V-B owner family, but it is not silently folded into the
finite/raw bridge or into the structural theorem domain.  A printed source
number that uses \(1/(6\sqrt2)\) is a \(\REGzero\) specialization, not a
generic \(\mathrm{REG}\) theorem.

In particular, after a raw branch is frozen the heavy affine operator may be
written
\begin{equation}
\label{p7v2:eq:heavy-affine-import}
 M_{N_1}=E_N,\qquad
 M_{N_2}=E_N-\frac{55}{56}\Delta_M,\qquad
 M_{N_3}=E_N-\frac{57}{56}\Delta_M,
\end{equation}
with threshold number \(x_N=1\).  The quotient-odd carrier and the affine
coefficient grammar are structural; named dimensionful masses
\(E_N,\Delta_M\), comparison flavor-current coordinates, charged root, neutral operator, scalar
normalization, and other entries of
\eqref{p7v2:eq:upstream-raw-branch} require \(\Uroleabs\) and one
\(b_6\in\Acompsix\).  A thermal identification
\(E_N=k_BT_{\rm reh}\) is a later matching relation between independently
obtained quantities and is not part of either upstream constitution.

Unless a theorem explicitly states otherwise, every Paper--V-C named
spectrum/flavor formula consumes \(\Urole+b_6\) with
\(b_6\in\Acompsix\); every absolute numerical mass additionally consumes
\(\Uroleabs\).  Interpreting the scalar coordinate \(v\) as a physical
nonzero Higgs expectation also consumes \(\Hvac\).  This convention preserves
exact conditional mathematics without promoting a raw Cartesian tuple,
scheme member, ancestry port, or calibration into a unique prefix.

The dependency order used throughout is
\begin{equation}
\label{p7v2:eq:dependency-order}
 \begin{gathered}
 \text{Papers I-A--V-B}\\[-0.15em]
 \big\downarrow\\[-0.15em]
 \text{Paper V-C finite interaction boundary}\\[-0.15em]
 \big\downarrow\\[-0.15em]
 \begin{array}{c}
   \text{Papers VII-A--VII-B: QCD/hadronic execution},\\
   \text{Paper VI-A: cosmological background},\\
   \text{Paper VII-C: local continuum reconstruction}.
 \end{array}
 \end{gathered}
\end{equation}
No result of Papers VI-A, VII-A, VII-B, or VII-C is used to establish a Paper--V-C
finite-core claim.  Later synthesis may combine their outputs with the boundary
defined here, but it may not feed those outputs backward into \(\Istr\),
\(\Irole\), or any unresolved member of \(\Acompsix\).  Later data may discharge
an earlier family only when constituted independently and genuinely
discriminating; the earlier chronological status remains unresolved at its own
stage.

\subsection{Five claim statuses}
\label{p7v2:sec:status-contract}

\begin{definition}[Claim-status contract]
\label{p7v2:def:claim-status}
Every load-bearing statement in this paper has one of the following statuses.
\begin{longtable}{@{}L{0.16\textwidth}L{0.29\textwidth}L{0.47\textwidth}@{}}
\caption[Paper--V-C claim-status contract]{Claim-status contract used in the
revised Paper~V-C.}\label{p7v2:tab:status-contract}\\
\toprule
Status & Meaning & Evidential requirement\\
\midrule
\endfirsthead
\toprule
Status & Meaning & Evidential requirement\\
\midrule
\endhead
\(\mathsf T\): theorem
& Finite or mathematical consequence
& Proved from explicitly named upstream theorems, definitions, and ordinary
mathematics; no unstated physical bridge is permitted.\\
\(\mathsf Q\): QFT interface
& Conditional continuum result
& Follows after a stated local-QFT action, normalization, field content,
renormalization scheme, or kinetic equation is supplied.\\
\(\mathsf A\): ansatz
& Finite model choice
& A precise admissible formula whose uniqueness or physical identification has
not been proved.\\
\(\mathsf B\): benchmark
& Arithmetic or sensitivity readout
& Reproducible after named ansatzes or interface inputs are imposed; it is not a
prediction.\\
\(\mathsf O\): open obligation
& Missing operator or dynamics
& Requires a construction, theorem, or evolution not supplied by the present
framework chain.\\
\bottomrule
\end{longtable}
\end{definition}

The words ``derived'', ``fixed'', ``closed'', and ``prediction'' are reserved
for status \(\mathsf T\), or for a conditional conclusion whose hypotheses are
displayed in the same statement.  Minimality can exclude an inserted object
only after the admissible class and objective are fixed.  It does not turn an
undetermined parameter into a distinguished endpoint, uniform weight, or
extremum.

\subsection{The minimal interface data}
\label{p7v2:sec:minimal-interface-data}

\begin{definition}[Interface hypotheses]
\label{p7v2:def:interface-hypotheses}
The following symbols name the minimal Paper--V-C-owned physical data when
their corresponding target is requested:
\begin{enumerate}[label=(\roman*),leftmargin=2.5em]
  \item \(\Hgauge\): the typed gauge-normalization package.  On the
  premise-minimal source-native common-heat branch,
  \[
    \Hgauge=\Rconn,
    \qquad \Rconn=(\JOINTfiveocc),
  \]
  because Papers~III-F and V-A derive the least self-adjoint heat carrier,
  ordinary trace, common heat coefficient, and absolute gauge matching
  readout.  On the broader parity-even kinetic-class branch one instead has
  \[
    \Hgauge=\Rconn+\CALgauge.
  \]
  Thus \(\CALgauge\) is a branch-specific calibration of the larger kinetic
  class, not an independent datum on the source-native heat branch;
  \item \(\Hqft\): \(\Urole+\Hgauge\) together with the analytic Lorentzian/local-QFT
  theorem domain \(\Rssqft\), specialized to the renormalizable
  Standard--Model QFT in a named subtraction and matching scheme.  The symbol
  \(\Rssqft\) denotes hypotheses, not a selector occurrence;
  \item \(\Hem\): the dependent electromagnetic endpoint certificate.  On the
  preferred physical route it consists of the same-effective-action quadratic
  response, the generated $U(1)_Q$ restriction, and a proved unique
  $q^2\to0$ infrared limit/matching.  A finite boundary-operator/trace/branch
  tuple may be used only as a representation of that same response after the
  identification is proved; it is not an independent physical owner;
  \item \(\Hvac\): a claim-scoped scalar stationary/effective-action
  certificate, consumed only by claims which explicitly require a stationary
  scalar equilibrium, global energetic preference, equilibrium scalar Hessian,
  or actual relaxation/history.  It is not an owner required to constitute the
  native electroweak stabilizer or gauge-mass kernel;
  \item \(\Hcp\): the claim-scoped regulated strong-phase/action/measure
  transport certificate compatible with the generated weak-family CP.
  An exact nonanomalous CP/parity comparison branch is one sufficient
  route where explicitly imposed; it is not a mandatory CP symmetry of
  the complete weak sector and does not reopen the fixed native finite
  strong character;
  \item \(\Hnu\): one member of the high-energy neutral lift fiber over a
  compatible Paper--V-B light/heavy spectrum member.  The physical \(Y_\nu\)
  realization is the fiber datum; a Casas--Ibarra \(R\) matrix is a coordinate
  on that member, not an additional selector beside it;
  \item \(\Hkin\): the actual observable algebra, state and complete
  interaction/cosmological evolution consumed by the declared kinetic or
  thermal target.  On the native branch, collision/washout and spectator
  maps, when used, are properties or certified reductions of that same
  evolution.  On an explicitly declared standard kinetic comparison they
  remain supplied interface data.
\end{enumerate}
These are not additional geometric principles.  They are the mathematical
objects required to define the corresponding requested physical calculations.

Two optional extensions are deliberately excluded from this minimal list.
The symbol \(\Hresopt\) denotes an imposed regulator prescription such as
exact saturation when physical masses and widths have not independently fixed
it.  Once \(M_i,M_j,\Gamma_j\) are physical, the standard regulator is a
derived readout.  The symbol \(\Hhopfopt\) denotes a typed heavy-pair Hopf
operator and expectation functional for a Born-specific source model.  No
typed Hopf-to-\(\varepsilon_{23}\) map is used in the standard
\(\Hnu+\Hkin\) baryogenesis target, so this branch is optional rather than an
irreducible closure premise.
\end{definition}

\begin{definition}[Selector-status partition]
\label{p7v2:def:selector-status-partition}
Paper~V-C distinguishes five logical classes.
\begin{enumerate}[label=(\roman*),leftmargin=2.5em]
 \item \emph{Derived/co-realized structural core}: \(\Istr\), the active anchors,
 color representation data, finite defect counts, and exact algebraic
 identities.
 \item \emph{Imported domains and conditional data}: structural hypercharge
 and the named representation roles are already generated upstream.  Numerical
 claims separately consume one \(b_6\in\Acompsix\), a positive \(\REGsig\)
 member, and energy ancestry/calibration when energy readouts are claimed.  Local gauge
 branches use the derived certificate \(\Rconn=(\JOINTfiveocc)\).
 \item \emph{Conditional Paper--V-C physical interfaces}: for their
 corresponding comparison targets, the three interface classes
 \[
  \boxed{\mathrm{CP},\ \mathrm{NU},\ \mathrm{KIN}}.
 \]
 They are represented by \(\Hcp,\Hnu,\Hkin\); this list is not three
 independent native mechanisms or three mandatory remaining gaps.  The scalar-stationary \(\Hvac\) branch is listed separately because it is consumed only by claims that explicitly require scalar equilibrium/history.  The electromagnetic quantity \(\Hem\) is instead a derived-dependent endpoint certificate whose preferred owner is the same physical action and occurring solder, without an upstream stable-vacuum prerequisite.  On the
 source-native common-heat branch the gauge coefficient is derived with
 \(\Rconn\); only the broader kinetic-class branch adds the branch-specific
 \(\CALgauge\) calibration.  This is an owner-family ledger, not an
 independent-coordinate count, and genuinely dimensionful calibrations are
 not counted as structural assumptions.
 \item \emph{Optional ansatz extensions}: \(\Hresopt\) and \(\Hhopfopt\).
 They are retained only when the corresponding exactly saturated or
 Born-specific benchmark is studied and are not premises of the standard
 physical target.
 \item \emph{Analytic theorem domains and derived imports}: in particular
 \(\Rssqft\) is an analytic domain and \(\Rconn\) is a derived import; neither
 is a finite realization selector or Paper--V-C owner.  Calibrations and
 downstream interfaces include absolute scale choices,
 later nonperturbative QCD/hadronic data, continuum vertices, and later
 matching relations.
\end{enumerate}
Thus the exact finite interaction-boundary core has zero structural selectors,
whereas the conditional comparison/interface package has three named classes;
the actual native source/action/state prefixes are reused on their proved
domains and only consumed unproved arrows remain.  The scalar-stationary
branch is claim-scoped.  The imported derived joint gauge datum is not
a sixth Paper--V-C-owned family.  Product-valued packages are counted once at this level, not once per
scalar coordinate; their independent physical freedom is determined only
after joint compatibility, exact redundancies, and optional-branch
inapplicability are imposed.
\end{definition}

\paragraph{Complete native owner before comparison.}
\label{p5c:rem:v2187-complete-owner-before-comparison}
The source-native action, current and pair constructors retain their complete
coefficients and contacts before a scalar or carrier-only readout.  The
current cyclic charge target uses the inherited state and the direct complete
return of Corollary~\ref{p5c:cor:v280-native-baryogenesis-frontier}; its
actual field restriction and expectation still require evaluation.  Native
electroweak geometry is the source stabilizer and positive gauge Gram of
Theorem~\ref{p7v2:thm:v258-native-ew-solder-reduction}, without an added
scalar-vacuum premise.  The carrier-only examples in
Theorem~\ref{p7v2:thm:type-separation} are comparison diagnostics in the
appendix; they do not license competing native actions, states or kinetic
laws after those generated owners are retained.

\label{p7v2:part:gauge-em}

\subsection{weighted-flavor and vacuum firewall}
\label{p7:sec:post-v125-flavor-vacuum}

The actual flavor-member boundary and the optional scalar-vacuum problem are separately typed.

A metric-compatible/unitary family connection cannot generate a
three-generation hierarchy: $Y_aY_a^*=I$ (or $|y_a|^2I$ after one scalar edge
weight).  The weighted quiver--Dirac correspondence is retained as the generated owner/readout geometry, with
nontrivial positive polar parts and sector-relative unitary data.  For the
quark sector, the presentation-free low-energy object is the simultaneous
conjugacy class of
\[
 H_u=Y_uY_u^*,\qquad H_d=Y_dY_d^*.
\]
A separate high-energy neutral lift remains necessary on the type-I/Majorana
branch.

Solving this flavor problem does not by itself solve the scalar vacuum.  On
the allowed Dirac-neutrino branch the exact positive common-heat contribution
is strictly decreasing in \(r=H^\dagger H\) whenever the squared Yukawa
operator is nonzero, and its quartic truncation places any formal stationary
point outside a uniformly small heat-expansion domain.  The full sign and
remainder proof is retained in
Appendix~\ref{p7:app:exact-heat-vacuum-diagnostic}; it is a negative diagnostic,
not part of the positive vacuum construction.

Accordingly the source-native ``finite-Dirac plus exact heat'' route is not a
general EWSB theorem.  The live scalar target is one generated renormalized or
otherwise physical effective action with a negative symmetric Hessian, a
stable finite nonzero gauge orbit, and global energetic preference.  Once it
generates $v_*$, the former scalar coordinate is only the readout
$\kappa_*=v_*/E_{32}$.



% BEGIN EM-SAME-SOURCE-BOUNDARY-RESPONSE v2.71
\label{p7v2:part:scope}

\paragraph{Technical proof owner and claim-status guard.}
Detailed auxiliary matrix, representation, finite-spectral and comparison
calculations formerly embedded in this manuscript now belong to
Paper~V-D. The original named theorem/proof labels remain
cross-referenced, and no optional finite comparison is promoted
to physical source occurrence, interacting mass pole or measured response.

\section{Joint interaction action and constituted vertices}
\label{p5c:sec:organized-vertices}
\subsection{Mixed CP--pair source incidence and generated coefficient tests}
\label{p5c:sec:v2167-cpbl2-source-incidence}
The intrinsic even Hopf--Hurwitz comparison of Paper~I-A supplies the
common geometric parent before its Hopf restrictions. Its already
generated primitive neutral edge and weak current have identity family
factors. The complete clock/spin, weak source and kinetic factors remain
in the physical word. Thus the marked local family return is $P_a$;
it does not install a projector-weighted terminal mass as a normal action.

\begin{theorem}[Marked mixed pair--current family contraction]
\label{p5c:thm:v2167-mixed-family-coefficient}
Retain the rooted BL2 member $M=M_{BL2}^{\rm geom}$, the charged rays $P_a$,
and the frozen neutral-role projectors
$N_{\rm inv}=|e_2\rangle\langle e_2|$,
$N_{\rm dir}=|e_1\rangle\langle e_1|$,
$N_{\rm Gram}=|e_0\rangle\langle e_0|$. Define
\[
 A_{a;rs}=N_rP_aN_s,\qquad
 Z_{a;rs}=\operatorname{Tr}(M^\dagger A_{a;rs}M A_{a;rs}^{\mathsf T}).
\]
These traces are invariant under simultaneous unitary family
presentation changes and role phases. Complete conjugation of the fixed
labelled packet conjugates $Z$. In particular
\[
 \boxed{Z_{\mu;\rm inv,dir}=-\frac7{288}+\frac i{12}.}
\]
\end{theorem}
\begin{proof}
Under $A\mapsto U^\dagger AU$, $M\mapsto U^\dagger M\bar U$, the
product transforms by $U^{\mathsf T}(\cdot)\bar U$, preserving its trace.
The generated coefficients give
$(P_\mu)_{21}=-1/3+i/4$ and $\bar M_{22}M_{11}=-1/2$.
Thus $Z=-\frac12(-1/3+i/4)^2=-7/288+i/12$.
\end{proof}
This nonzero dimensionless source coefficient is a different invariant
from the ordered current quartet. It is not a charge coefficient,
collision efficiency or final asymmetry. Its physical use retains the
actual support insertion and every spin/mode, metric and propagation
factor. The full role table is in the comparison appendix.

\begin{proposition}[Generated family and port matrix units]
\label{p5c:prop:v2167-generated-coefficient-algebras}
Let $D_i=|e_i\rangle\langle e_i|$ be the three neutral role tests.
All entries of the generated $P_\mu$ are nonzero, and
\[
 E_{ij}=\frac{D_iP_\mu D_j}{(P_\mu)_{ij}}=|e_i\rangle\langle e_j|,
 \qquad \operatorname{alg}^*(I,D_i,P_\mu)=M_3(\mathbb C).
\]
On the local coefficient ports
$R=G_R^\nu$, $L=G_L^\nu$, $E=G_L^e$, put
\[
 V_T=|L\rangle\langle R|\otimes I_G,\quad
 V_W=|E\rangle\langle L|\otimes I_G,\quad
 V_R=P_R,\quad V_L=V_T,\quad V_E=V_WV_T.
\]
Then $V_p(P_R\otimes E_{ij})V_q^\dagger=|p,i\rangle\langle q,j|$,
and the generated joint coefficient algebra is $M_9(\mathbb C)$.
Both simple coefficient blocks have scalar commutants.
\end{proposition}
\begin{proof}
The displayed products give all nine family units and all 81 port/family
units, with their usual multiplication and adjoint rules. A matrix
commuting with all matrix units is scalar. Thus a projection reducing a
simple coefficient block is $0$ or $I$; multiplicity factors can remain.
\end{proof}
Division only identifies an algebra basis. The nine-dimensional port
carrier is not a physical particle count or CAR cutoff. All other action
sectors remain. II-A's least-test result requires already licensed
generators and their application in the same role; this local coefficient
calculation does not select an unresolved K16 generator. Its independent
commutator/star calculation has dimensions $10\to37\to81$.

\begin{corollary}[Actual K16 handoff after the source calculation]
\label{p5c:cor:v2167-cpbl2-handoff-frontier}
The source packet has a nonzero marked mixed family contraction and a
generated complete local coefficient algebra. The remaining action target is
\[
 \boxed{S_{16,\rm normal}=\mathcal R_{16}(S_{F,\rm complete}),}
\]
where $\mathcal R_{16}$ is the actual licensed restriction/reduction on the
pair's inherited field identity. Its complete consumed diagram retains the
pair member, current/source coefficients, spin/mode factors, kinetic metric,
connection, constraints and induced contacts through the required jet order.
On the native component-band domain,
Proposition~\ref{p3f:prop:v2168-component-band-retention} already proves the
identity family extension; this result is reused rather than reopened.

After that actual handoff, evaluate the full mixed response consumed by the
claim, including every physical slot and cross contact. A certified
intertwining/reducing-range guard or the factorized positive-support
criterion is a sufficient shortcut on its stated domain, not an additional
prerequisite for direct evaluation of the complete unfactorized return.
Neither the symbolic $\mathcal R_{16}$ nor this corollary evaluates that map
or its return. A nonzero-incidence claim closes only if its actual evaluated
return is nonzero.
\end{corollary}
\begin{proof}
The pair handoff supplies its inherited pairing coordinate. III-F supplies
the complete normal/source return and the actual component-band family
retention on their printed domains. II-A preterminal relation preservation
and III-A complete-parent reuse retain every mixed block of the licensed
restriction before a noninjective readout is taken. They do not manufacture
an absent physical restriction from an available coefficient. Once the
actual complete word is exhibited, direct differentiation/evaluation gives
its full response; the conditional structural guards are alternative
sufficient evaluations on their admitted subdomains. Co-consumption then
needs no second occurrence or normalization token.
\end{proof}

\begin{corollary}[Claim-scoped direct signed charge-history target]
\label{p5c:cor:v2186-direct-signed-charge-target}
For a claim whose target is the actual signed charge history, a separate
nonzero marked-incidence proxy is not required. On an already licensed
finite physical carrier and regular history interval, compile the actual
complete $H_{16}(t)$ from the inherited action/field restriction, retaining
its pair, normal, interaction, metric/connection and constraint/contact
terms. Evolve the inherited cyclic state $\rho_t$ with this same generator
and use the already owned signed charge $Q(t)$. Use the coordinate derivative
when the moving connection is included in $H_{16}$, or its covariant derivative
with that connection contribution removed from the generator; it is counted
exactly once. Paper~IV's exact complete charge history then gives
\[
 \boxed{q_Q(t_f)-q_Q(t_i)
 =\int_{t_i}^{t_f}\operatorname{Tr}\!\left[
 \rho_t\left(D_tQ+\frac{i}{\hbar}[H_{16},Q]\right)\right]dt.}
\]
Thus direct full-state evaluation, or the existing least complete charge-test
flow on its domain, may discharge this declared history target without
first evaluating a marked mixed-source proxy. For a specifically claimed
CP-reversal comparison, the actual inherited reversal packet and its charge
convention remain consumed data; matrix conjugation alone does not supply
that physical comparison.

An evaluated zero is a completed zero charge result on that domain. It does
not certify a positive asymmetry, the separate nonzero-incidence claim, or a
final baryogenesis yield. No K16 state, coupling, efficiency, collision or
washout law is added. The actual restriction and history remain unevaluated
here; A12's positive claim and A21's final-yield obligations are not promoted.
\end{corollary}
\begin{proof}
The formula is Paper~IV's exact many-body charge balance applied to the
already owned generator, state and charge. Its finite test flow retains all
owned commutator directions and agrees with the full-state evolution, so
it does not need an intermediate source proxy. Retaining a zero output is
ordinary forward evaluation, whereas a positive claim additionally requires
its nonzero certificate. State inheritance and complete-parent reuse forbid
an independent preparation or response rescaling. The physical reversal
clause is the scope of IV's signed complete-packet charge diagnostic.
\end{proof}

Accordingly the actual K16 handoff and its consumed same-history evaluation
remain open. Generic receiver comparisons, including a lossy pair/current
compression or a block that deletes a generated primitive edge, are not
licensed alternative K16 realizations and prove no native non-derivability.
The detailed comparison scope remains in the existing appendices.

% END INLINE SOURCE: Paper5C-CPBL2SourceIncidence-Main-v1_1.tex




% BEGIN INLINE SOURCE: Paper5C-NativeWeakMemory-Main-v1_0.tex
% Framework v2.168: integrated finite charged-source derivations.
\subsection{Even native response tests and odd interaction legs}
\label{p5c:sec:v226-even-odd-tests}
\begin{proposition}[Exact size of the native even and missing odd test spaces]
\label{p5c:prop:v226-even-odd-tests}
For the geometric/protected splitting $\mathbb C^{15}=\mathbb C^{11}
\oplus\mathbb C^4$, the native Hermitian algebra has dimension 137.
The Hermitian off-diagonal space has dimension $2\cdot11\cdot4=88$.
The native 137-effect frame separates every even Hermitian residual but
annihilates this entire odd space. Full $M_{15}$ tomography, if licensed
in an accessible event role, requires at least 225 outcomes.
\end{proposition}
\begin{proof}
Theorem~\ref{p1:thm:v226-coordinate-only} gives the block diagonal algebra.
Its trace pairing with an off-diagonal matrix is zero. An arbitrary complex
$11\times4$ block has 88 real entries and specifies a Hermitian odd matrix
with its adjoint. The full state algebra has affine dimension $15^2-1=224$,
so a normalized POVM needs at least 225 outcomes to separate all its states.
\end{proof}
Thus an actual even Ward, stress or completed same-action return can be
checked by the native dual coefficients. A source current, odd solder or
three-point leg outside that even sort retains its actual larger source
frame and physical incidence. Eliminating an odd source does not justify
omitting its even $BC^{-1}B^*$ return. No new interaction normalization or
physical response value is selected by the operator-algebra closure.

% END FRAMEWORK v2.226: KINEMATIC CLOSURES

% BEGIN FRAMEWORK v2.227: DOWNWARD CLOSURES
\subsection{Relative central energy remains an odd-source coupling datum}
\label{p5c:sec:v227-odd-offset}
\begin{proposition}[Neutral retirement has an explicit source boundary]
\label{p5c:prop:v227-odd-offset}
Suppose the actual native Hamiltonian is block reducing. Under the central
shift $C=\sum_ac_aP_a$, a source component $q_{ab}=P_aqP_b$ has
\[
 \operatorname{ad}_{H+C}(q_{ab})
 =\operatorname{ad}_{H}(q_{ab})+(c_a-c_b)q_{ab}.
\]
Consequently its transition frequencies retain the relative block energy.
They are unchanged by the shift only when that difference vanishes on
every actually consumed nonzero source component.
\end{proposition}
\begin{proof}
Apply Proposition~\ref{p1b:prop:v227-central-boundary} componentwise.
Distinct source blocks are linearly independent, so vanishing of the
extra commutator is equivalent to the printed component condition.
\end{proof}
Actual interaction vertices, source jets and Ward identities keep these
relative coefficients. Neither equality of a neutral response nor the
137-dimensional even native algebra determines an omitted odd source.
If an enlarged total-even action contains odd tensor odd factors, first
verify the actual reducing-algebra hypothesis before applying any central
retirement to that larger system.

% END FRAMEWORK v2.227: DOWNWARD CLOSURES


% BEGIN FRAMEWORK v2.229: ODD NORMAL AND K16 PRODUCTION FIREWALL

\section{Electroweak gauge and electromagnetic source response}
\label{p5c:sec:organized-ewem}
\subsection{Exact defect counts and the gauge-normalization boundary}
\label{p7v2:sec:gauge-normalization}

\begin{lemma}[Exact defect-count identities]
\label{p7v2:lem:defect-counts}
Let
\[
 \wp_8=\frac18,\qquad \wp_9=\frac19,\qquad \wp_7=\frac17.
\]
Then
\begin{equation}
\label{p7v2:eq:defect-counts}
 \delta^2=\wp_8\wp_9=\frac1{72},
 \qquad
 \wp_\star=\wp_8\wp_7=\frac1{56},
\end{equation}
and hence
\begin{equation}
\label{p7v2:eq:defect-sum-difference}
 \delta^{-2}+\wp_\star^{-1}=128,
 \qquad
 \delta^{-2}-\wp_\star^{-1}=16.
\end{equation}
\end{lemma}

\begin{proof}
The first identities are direct products of the protected octet, color-pair,
and reduced-septet normalized counts.  Taking inverses gives
\(\delta^{-2}=8\cdot9=72\) and
\(\wp_\star^{-1}=8\cdot7=56\); their sum and difference are \(128\) and
\(16\).
\end{proof}

\begin{proposition}[Gauge normalization: source-native heat branch and broader calibration chart]
\label{p7v2:prop:strong-seed-status}
The count family
\begin{equation}
\label{p7v2:eq:strong-seed-family}
 \mathfrak A_{\rm gauge}^{\rm norm}
 =\{(d_Ad_P)^{-t}:t>0\}
\end{equation}
remains a legitimate coordinate chart on the one-dimensional absolute gauge
normalization of the broader parity-even kinetic class.  The dimensions and
native symmetries do not select any exponent \(t\).  In particular,
\begin{equation}
\label{p7v2:eq:strong-seed-ansatz}
 \alpha_s^{(1/2)}=(d_Ad_P)^{-1/2}=\delta
\end{equation}
is only the historical \(t=1/2\) calibration member on that broader branch.

On the source-native common-heat branch, however, Papers~III-F and V-A now fix
the least self-adjoint trace carrier and the ordinary trace.  Canonical
normalization gives
\[
 \boxed{k_i g_i^2=4\pi^2,}
\]
and hence at the native common-heat matching scale
\[
 \boxed{\alpha_s^{\rm heat}=\frac\pi2.}
\]
There is no independent \(\CALgauge\) or E16 scalar on this branch.
\end{proposition}

\begin{proof}
The nonselection of \(t\) on the broad branch is the counterfamily in
Theorem~\ref{p7v2:thm:type-separation}.  On the common-heat branch,
Paper~III-F fixes the half-weighted Yang--Mills heat coefficient and Paper~V-A
fixes the rank-48 internal carrier, the least self-adjoint Clifford
completion, the ordinary trace, three family copies, and the representation
indices.  Its common-heat normalization corollary therefore gives
\(k_i g_i^2=4\pi^2\).  With \(k_3=2\) and
\(\alpha_s=g_3^2/(4\pi)\), one obtains \(\alpha_s^{\rm heat}=\pi/2\).
\end{proof}

\begin{remark}[Matching-scale and determinant firewalls]
The number \(\alpha_s^{\rm heat}=\pi/2\) is attached to the native common-heat
matching datum.  Identifying its dimensionful scale with \(M_Z\), a unification
scale, or another observed renormalization scale requires the separately typed
scale/RG matching.  The equality \(\alpha_s=\delta\) is not recovered: the
Paper-III determinant-shell readout has different scheme and variational
ownership.  Thus the old \(t=1/2\) formula is retained only as a broad-branch
calibration/comparison chart, not as a native strong prediction.
\end{remark}

\subsection{Electroweak coupling conventions}
\label{p7v2:sec:ew-conventions}

\begin{theorem}[Native electroweak reduction from the occurring chiral solder]
\label{p7v2:thm:v258-native-ew-solder-reduction}
On the same local clocked uninserted source and gauge-action domain, with a
positive invertible gauge-kinetic metric, retain the occurring rooted parent
of Paper~V-A.  Its rooted incidence
theorem supplies the actual nonzero primitive odd source; no second
primitive-source G0 premise is required.  Paper~V-A retains its primitive Clifford factor \(c(\hat h)\)
and identifies it with the normalized generated scalar/solder direction.  For
the same occurring odd source, Paper~III-F gives the positive gauge Gram
\[
 M_{\rm gauge}^2
 =2C Q_\rho^{-1/2}\mathcal J^\dagger\mathcal JQ_\rho^{-1/2}.
\]
The Paper--V-A source-stabilizer theorem gives
\[
 \ker\mathcal J|_{\mathfrak g_{\rm EW}}=\mathbb R Q_H,
 \qquad Q_H=U_HT_L^3U_H^\dagger+Y.
\]
Therefore
\[
 \boxed{\ker M_{\rm EW}^2=Q_\rho^{1/2}\mathbb R Q_H
 \simeq\mathfrak u(1)_Q,\qquad
 \operatorname{rank}M_{\rm EW}^2=3.}
\]
Including color,
\[
 \boxed{\ker M_{\rm SM}^2
 =Q_\rho^{1/2}(\mathfrak{su}(3)_c\oplus\mathbb R Q_H)
 \simeq\mathfrak{su}(3)_c\oplus\mathfrak u(1)_Q.}
\]
Here \(U_H=\Phi(\widehat H)\) and the isomorphisms explicitly undo
kinetic whitening and use the source-aligned weak frame.  They are not
identities between untransported component vectors.
No \(\Hvac\), scalar magnitude, scalar-potential stationary condition, or
global scalar minimum is consumed.  Absolute pole masses and physical scalar
vertices remain downstream claims.  Stationarity is required for the separate
scalar Hessian Ward-null statement, not for this gauge-Gram kernel.
\end{theorem}

\begin{proof}
Paper~III-F Corollary~\ref{p3f:cor:v111-gauge-gram-no-stationarity} identifies
the kernel of the positive gauge Gram with the stabilizer of the fixed
occurring odd source.  Paper~V-A
Theorem~\ref{p5:thm:v377-occurring-solder-em-stabilizer} computes that
stabilizer as \(\mathbb R Q_H\) on the electroweak factor and leaves color
unbroken.  Positivity gives
\(v^\dagger M^2v=2C\|\mathcal JQ_\rho^{-1/2}v\|^2\),
so \(\ker M^2=Q_\rho^{1/2}\ker\mathcal J\).
Invertible kinetic whitening therefore preserves nullity, with the displayed
transport of its vectors.  Since
\(\dim\mathfrak g_{\rm EW}=4\), the electroweak rank is three.
\end{proof}


% BEGIN INLINE SOURCE: Paper5C-RootedElectroweak-v1_0.tex
% Source: RemainingGapTransportAnalysis v1.0, CLS-T01; audit CC-01, TR-01--02.
\subsubsection{Rooted occurrence and the electroweak gauge-Gram kernel}
\label{p7v2:sec:v2109-rooted-source-transport}
\begin{corollary}[No separate primitive-source occurrence premise on the rooted branch]
\label{p7v2:cor:v2109-rooted-ew-occurrence}
On the same local clocked, uninserted source and gauge-action domain as
Paper~V-A's rooted incidence theorem, retain the occurring parent
\(\mathbb T_*\) and a positive invertible gauge-kinetic metric.  Then
\[
 \ker M_{\rm EW}^2\simeq\mathfrak u(1)_Q,
 \qquad\operatorname{rank}M_{\rm EW}^2=3,
 \qquad
 \ker M_{\rm SM}^2\simeq\mathfrak{su}(3)_c\oplus\mathfrak u(1)_Q.
\]
The identifications use the retained kinetic whitening.  An independently
supplied nonzero primitive odd-source member is not an additional premise
on this rooted specialization.
\end{corollary}
\begin{proof}
Paper~V-A's Theorem~\ref{p5:thm:v392-rooted-srcval-incidence} and the
Paper~III-F rooted member handoff generate the actual nonzero primitive odd
source.  Applying Theorem~\ref{p5:thm:v377-occurring-solder-em-stabilizer},
the stabilizer is \(\mathbb R Q_H\), the weak conjugate of
\(\mathbb R Q\).  In its source-aligned weak frame, invariance under
\(x_1T_L^1+x_2T_L^2+x_3T_L^3+bY\) is exactly
\(x_1=x_2=0\), \(x_3=b\); colour remains in the full stabilizer.
The displayed \(Q\)-identifications use this common transported frame.
For the positive gauge Gram, \(\ker(J^\dagger J)=\ker J\), and invertible
kinetic whitening preserves its nullity.  This proves the three statements.
\end{proof}
This is an occurrence-premise reduction, not a scalar-VEV, stationary-scalar
or fitted-Yukawa assumption.  Gauge pole values and the noncentral family
spectrum remain separate calculations.  Charged-slot, flavour and
baryogenesis consumers inherit this occurring primitive source, but still
require the actual same-action family assembly, terminal-role incidence,
and their own neutral, kinetic and state inputs where consumed.

% END INLINE SOURCE: Paper5C-RootedElectroweak-v1_0.tex


\begin{theorem}[Conventional single-doublet comparison dictionary]
\label{p7v2:thm:ew-mass-dictionary}
On \(\Urole+\Hvac\), use the v1.75 local scalar-action branch of
Paper~III-F and a selected flavor fibre satisfying \(a_Y>0\), so that the
canonical doublet of
Corollary~\ref{p5:cor:v175-yukawa-scalar-action-family} exists.  Rename that
canonically normalized field \(H\).  Use the structurally generated ordinary-hypercharge representation with
\(Y_H=1/2\) and a nonzero background
\(\langle H\rangle=(0,v/\sqrt2)^T\).  Let \(g_Y\) denote ordinary hypercharge
and let
\begin{equation}
\label{p7v2:eq:g1-convention}
 g_1=\sqrt{\frac53}\,g_Y
\end{equation}
denote the GUT-normalized coupling.  Then
\begin{align}
\label{p7v2:eq:ew-mass-dictionary}
 m_W^2&=\frac14g_2^2v^2,
 &
 m_Z^2&=\frac14(g_2^2+g_Y^2)v^2
       =\frac14\left(g_2^2+\frac35g_1^2\right)v^2,\\
 g_2&=\frac{2m_W}{v},
 &
 g_Y&=\frac{2\sqrt{m_Z^2-m_W^2}}{v}.
\end{align}
Moreover,
\begin{align}
\label{p7v2:eq:em-coupling-conventions}
 \sin^2\theta_W
 &=\frac{g_Y^2}{g_2^2+g_Y^2}
  =1-\frac{m_W^2}{m_Z^2},\\
 \alpha_{\rm em}^{-1}
 &=\alpha_2^{-1}+\alpha_Y^{-1}
  =\alpha_2^{-1}+\frac53\alpha_1^{-1},
 \qquad
 \alpha_i=\frac{g_i^2}{4\pi}.
\end{align}
The kinetic normalization is derived only on the \(a_Y>0\) subfamily.
Structural hypercharge and native electroweak reduction are already generated upstream.  This theorem is a conventional coordinate dictionary on a supplied single-doublet stationary chart; it is not the native cause of electroweak reduction and is not consumed by Theorem~\ref{p7v2:thm:v258-native-ew-solder-reduction}.  It does not derive the flavor fibre or a scalar-stationary history.
\end{theorem}

\begin{proof}
Insert the displayed Higgs background into the covariant derivative and
diagonalize the neutral \(W^3\)--\(B\) mass matrix.  The charged eigenvalue is
\(g_2^2v^2/4\), while the nonzero neutral eigenvalue is
\((g_2^2+g_Y^2)v^2/4\).  Equation~\eqref{p7v2:eq:g1-convention} gives the
GUT-normalized forms.  The photon is the orthogonal zero eigenvector, so
\(e^{-2}=g_2^{-2}+g_Y^{-2}\), which is
\eqref{p7v2:eq:em-coupling-conventions}.
\end{proof}

% BEGIN HM v1.79 EW01
\begin{remark}[Geometric response versus canonical mass operator]
\label{p7:rem:v179-hurwitz-action-comparison}
The gauge masses above are the commutator-source Gram of
\eqref{p3f:eq:v179-gauge-mass-commutator} with the scalar and gauge
kinetic metrics retained.  The local radial Higgs direction is the odd
matrix \(X_H\) of Theorem~\ref{p5:thm:v179-canonical-higgs-direction};
its second moment fixes its kinetic term and its fourth moment fixes
the family-valued quartic.  Thus a Hurwitz projector norm is compared
through the same complete action, not by adding that norm as a new
mass coefficient.  The explicit Higgs/muon local matching identities
are \eqref{p6:eq:v179-higgs-muon-action-targets}.
The midpoint orbit conversion is
\eqref{p3c:eq:v179-carrier-orbit-conversion}.  None of these source
identities supplies an additional vacuum or flavor member.
\end{remark}
% END HM v1.79 EW01



\begin{theorem}[Physical \(7\mid9\) owner reduction]
\label{p7:thm:physical-seven-nine-owner-reduction}
Let \(V_{16}=V_7\oplus V_9\) be the orthogonal finite split imported from
Paper~V-B, with projectors \(P_7,P_9\) of ranks seven and nine.  The most
general block-scalar response compatible with this split is
\begin{equation}
 \label{p7:eq:seven-nine-response-family}
 \mathcal J_{a,b}:=aP_7+bP_9,
 \qquad a,b\in\mathbb C.
\end{equation}
Its two Hilbert--Schmidt Gram weights are exactly
\begin{equation}
 \label{p7:eq:seven-nine-gram-weights}
 \|\mathcal J_{a,b}P_7\|_{\rm HS}^2=7|a|^2,
 \qquad
 \|\mathcal J_{a,b}P_9\|_{\rm HS}^2=9|b|^2.
\end{equation}
Suppose, in addition, that one common physical owner identifies these two
blocks with the charged and neutral massive response lines, supplies the same
overall pole/energy normalization, and makes the corresponding quadratic
weights the physical mass-squared readouts.  Then
\begin{equation}
 \label{p7:eq:physical-seven-nine-ratio}
 \boxed{\frac{m_W^2}{m_Z^2}
 =\frac79\rho_{79}^2,
 \qquad \rho_{79}:=\frac{|a|}{|b|}.}
\end{equation}
In particular, the physical \(7/9\) ratio follows if and only if
\(\rho_{79}=1\) on this owner branch.  Rank, block symmetry, and ordinary
trace alone do not impose that equality.

A sufficient equal-amplitude condition, denoted \((\mathrm{EQAMP}_{79})\),
is that both block responses descend as restrictions of one scalar
pre-split operator \(a_*I_{V_{16}}\), up to presentation phases.  Then
\(a=e^{i\phi_7}a_*\), \(b=e^{i\phi_9}a_*\), and \(\rho_{79}=1\).
Papers~I-A and V-B now constitute the canonical common finite parent
\(\mathcal G_{16}=\mathcal G_7\sqcup\mathcal G_9\); restriction of one
scalar pre-split operator on its linear carrier gives equal block amplitudes
mathematically.  The remaining content of
\((\mathrm{EQAMP}_{79})\) is the physical owner theorem that the charged and
neutral massive response lines descend from that common pre-split response.
Counting measure, rank, and ordinary trace do not prove this ancestry.
\end{theorem}

\begin{proof}
Since \(P_7P_9=0\) and \(\Tr P_d=d\),
\[
 \Tr\bigl((\mathcal J_{a,b}P_7)^*
               (\mathcal J_{a,b}P_7)\bigr)=|a|^2\Tr P_7=7|a|^2,
\]
and similarly for the rank-nine block.  The stated physical owner hypotheses
cancel the common overall factor in the ratio, giving
\eqref{p7:eq:physical-seven-nine-ratio}.  The family
\(\mathcal J_{a,b}\) permits arbitrary positive \(|a|/|b|\), proving
nonselection.  Restriction of one scalar pre-split operator fixes both
absolute amplitudes and leaves only irrelevant block phases, proving the
sufficient condition.
\end{proof}

\begin{corollary}[Common-heat/tree-shell same-scale firewall]
\label{p7:cor:common-heat-tree-shell-firewall}
At the source-native common-heat matching scale, Paper~V-A gives
\begin{equation}
 g_1^2=g_2^2=2\pi^2,
 \qquad g_Y^2=\frac35g_1^2=\frac65\pi^2,
\end{equation}
and therefore
\begin{equation}
 \label{p7:eq:common-heat-ew-ratio}
 \sin^2\theta_W=\frac38,
 \qquad
 \frac{m_W^2}{m_Z^2}=\frac58
\end{equation}
whenever the single-doublet dictionary is evaluated at that same scale.
The equal-amplitude shell value \(7/9\) (equivalently
\(\sin^2\theta_W=2/9\)) cannot hold at the same scale and in the same
coupling scheme.  Relating the two is necessarily an RG/threshold matching
problem, not an algebraic identity.
\end{corollary}

\begin{proof}
Insert \(g_Y^2=6\pi^2/5\) and \(g_2^2=2\pi^2\) into
Theorem~\ref{p7v2:thm:ew-mass-dictionary}.  This gives
\(g_Y^2/(g_2^2+g_Y^2)=3/8\) and its complement \(5/8\).
Since \(5/8\ne7/9\), simultaneous same-scale identification is impossible.
\end{proof}


\begin{theorem}[Electroweak characteristic-cone inheritance]
\label{p7v2:thm:electroweak-characteristic-cone-inheritance}
Assume the hypotheses of
Theorem~\ref{p7v2:thm:ew-mass-dictionary} and the Paper~V-A gauge-cone export
domain, including the derived joint connection and upstream same-variable
Lorentzian metric realization.  Then the Weinberg rotation is an invertible
change of basis in internal \((W^3,B)\) space and the Higgs-induced gauge-boson
mass matrix is zeroth differential order.  Hence neither operation changes the
spacetime principal characteristic set.  For physical polarizations,
\[
 \boxed{
 \operatorname{Char}(W^\pm)=\operatorname{Char}(Z)=\operatorname{Char}(\gamma)
 =\{k:g_{\rm op}^{\mu\nu}k_\mu k_\nu=0\}.}
\]
Therefore the front/characteristic speeds obey
\[
 \boxed{
 c_{{\rm front},W}=c_{{\rm front},Z}=c_\gamma=c_{\rm op}.}
\]
The photon is the zero eigenvector already identified in
Theorem~\ref{p7v2:thm:ew-mass-dictionary}; its theoretical vacuum light speed
may consequently be denoted
\[
 \boxed{c_{\rm light}^{\rm th}:=c_\gamma=c_{\rm op}.}
\]
For the massive \(W\) and \(Z\), this is a front-speed statement and does not
assert that every finite-frequency group velocity equals \(c_{\rm op}\).  An
empirical metrological identification of an observed light standard with this
theoretical photon channel is a separate calibration/measurement statement.
\end{theorem}

\begin{proof}
Before symmetry breaking, Paper~V-A's gauge-cone export gives the same metric
principal symbol to \(W^1,W^2,W^3\), and \(B\).  Passing to
\(W^\pm\) and to the neutral photon/\(Z\) basis multiplies the internal field
vector by an invertible matrix, so the spacetime principal factor is unchanged.
The Higgs background adds only the displayed algebraic mass matrix, which is
zeroth order and therefore absent from the principal symbol.  The characteristic
zero set is consequently the inherited metric null cone for every physical
polarization.  Equivalently, the calculation verifies the support/principal
ancestry and no-deformation hypotheses of Paper~II-A,
Theorem~\ref{p2a:thm:undeformed-host-principal-inheritance}.  The speed
identities follow in the calibrated representative
\(g_{\rm op}=-c_{\rm op}^2d\widetilde t^{\,2}+h\).
\end{proof}

\begin{theorem}[Transverse electroweak propagation from the physical Hessian]
\label{p7v2:thm:v244-transverse-ew-wave-operators}
Assume the hypotheses of
Theorems~\ref{p7v2:thm:ew-mass-dictionary} and
\ref{p7v2:thm:electroweak-characteristic-cone-inheritance}, together with the
source-free physical coexact reduction of Paper~III-F,
Corollary~\ref{p3f:cor:v244-transverse-gauge-propagation}.  Then the transverse
physical eigendirections of the single electroweak gauge-field Hessian obey
\begin{align}
\label{p7v2:eq:v244-transverse-ew-wave-operators}
 \boxed{K_{\gamma,T}}
 &=\boxed{\partial_t^2+c_{\rm op}^2\Delta_1^{\rm co}},\\
 \boxed{K_{W,T}}
 &=\boxed{\partial_t^2+c_{\rm op}^2\Delta_1^{\rm co}
 +\frac{m_W^2c_{\rm op}^4}{\hbar^2}},\\
 \boxed{K_{Z,T}}
 &=\boxed{\partial_t^2+c_{\rm op}^2\Delta_1^{\rm co}
 +\frac{m_Z^2c_{\rm op}^4}{\hbar^2}}.
\end{align}
The mass terms are zeroth differential order, so these full operators refine
rather than replace the common characteristic-cone theorem.

If the consuming branch separately constitutes a round closed spatial radius
\(R_x\), the coexact curl spectrum gives, for \(L\ge1\),
\begin{equation}
\label{p7v2:eq:v244-transverse-ew-spectra}
 \boxed{
 \omega_{V,L}^2
 =c_{\rm op}^2\frac{(L+1)^2}{R_x^2}
 +\frac{m_V^2c_{\rm op}^4}{\hbar^2},
 \qquad V\in\{\gamma,W,Z\},\quad m_\gamma=0.}
\end{equation}
Without a dimensionful radius calibration, the theorem retains the exact
spectral form with the eigenvalues of \(\Delta_1^{\rm co}\) and makes no
absolute-frequency claim.

For the photon, writing \(C=*d\) on coexact one-forms gives
\begin{equation}
\label{p7v2:eq:v244-photon-helicity-factorization}
 \boxed{K_{\gamma,T}
 =(\partial_t-i c_{\rm op}C)(\partial_t+i c_{\rm op}C)}.
\end{equation}
Thus the two curl branches give the two first-order transverse helicity
branches.  Since \(H^1(S^3)=0\), there is no harmonic transverse one-form.  On
a finite round slice the first physical coexact level is therefore \(L=1\),
so, when \(R_x\) is physical,
\[
 \boxed{E_{\gamma,\min}=2\hbar c_{\rm op}/R_x}.
\]
This is a compact-volume spectral gap and not a photon rest mass; the same
theorem has \(m_\gamma=0\), and the gap tends to zero as \(R_x\to\infty\).
The massive longitudinal \(W/Z\) polarization is not part of this coexact
calculation; it is derived separately from the same gauge--scalar Hessian in
Theorem~\ref{p7v2:thm:v247-longitudinal-vector-spectrum} below.
\end{theorem}

\begin{proof}
Paper~III-F derives the massless physical transverse differential Hessian
\(\partial_t^2+c_{\rm op}^2\Delta_1^{\rm co}\) from the same Yang--Mills
action after Gauss reduction.  Theorem~\ref{p7v2:thm:ew-mass-dictionary}
computes the charged and neutral Higgs Hessian eigenvalues and the zero photon
eigenvalue.  Adding this algebraic Hessian to the common transverse
differential block gives
\eqref{p7v2:eq:v244-transverse-ew-wave-operators}; because it is zeroth order,
the principal characteristic set is unchanged.  Paper~I-A's coexact spectrum
\(\lambda_L^{\rm co}=(L+1)^2/R_x^2\), its identity
\(C^2=\Delta_1^{\rm co}\), and \(H^1(S^3)=0\) give the remaining statements.
\end{proof}


\begin{theorem}[Gauge Gram kernel and primitive-clock locking]
\label{p7v2:thm:v245-gauge-hessian-kernel-clock-lock}
Let \(\mathfrak g\) be a compact gauge algebra acting unitarily on the scalar
carrier by \(\rho\), and let \(h\) be any fixed occurring scalar/solder
background on a branch with positive gauge and scalar kinetic metrics.  Define
\[
 \mathcal J_h:\mathfrak g\to V_H,
 \qquad
 \mathcal J_h(X)=\rho(X)h.
\]
Write \(Q_\rho>0\) for the raw gauge kinetic metric and
\(B_h>0\) for the scalar kinetic metric in fixed source coordinates.
The canonically normalized gauge-mass Gram is
\[
 \widetilde{\mathcal J}_h:=B_h^{1/2}\mathcal J_h Q_\rho^{-1/2},
 \qquad M_G^2:=\widetilde{\mathcal J}_h^\dagger
                       \widetilde{\mathcal J}_h.
\]
An overall positive action coefficient can be absorbed into \(B_h\).
Define the \emph{raw} stabilizer Lie algebra
\(\mathfrak g_h:=\{X\in\mathfrak g:\rho(X)h=0\}\).
Then the exact, coordinate-typed kernel relation is
\[
 \boxed{\ker(\mathcal J_h^\dagger B_h\mathcal J_h)
   =\ker\mathcal J_h=\mathfrak g_h,
 \qquad \ker M_G^2=Q_\rho^{1/2}\mathfrak g_h.}
\]
Thus the kinetic-coordinate null space is \emph{isomorphic} to the raw
stabilizer, not generally equal to it as an untransported subspace.

On a calibrated constant round branch, a transverse mass eigen-direction
\(X\) has
\[
 E_{n,X}^2=n^2\Eper^2+m_X^2c_{\rm op}^4,
 \qquad n=2,3,\ldots .
\]
Its complete transverse tower lies on primitive scalar-clock characters,
\(E_{n,X}/\Eper\in\mathbb Z\) for every sufficiently large \(n\), if and
only if \(m_X=0\).  Hence
\[
 \boxed{\text{clock-locked raw gauge algebra}
 =Q_\rho^{-1/2}\ker M_G^2=\mathfrak g_h.}
\]
For the Standard--Model occurring solder/source branch of Theorem~\ref{p7v2:thm:v258-native-ew-solder-reduction}, color acts trivially and the electroweak stabilizer is generated by \(Q=T_L^3+Y\).  Therefore
\[
 \boxed{\ker M_{\rm SM}^2
 =Q_\rho^{1/2}(\mathfrak{su}(3)_c\oplus\mathfrak u(1)_{\rm em})
 \simeq\mathfrak{su}(3)_c\oplus\mathfrak u(1)_{\rm em},
 \qquad \dim\ker M_{\rm SM}^2=8+1=9.}
\]
The number nine here is the dimension of an unbroken gauge Lie algebra and is
not the unrelated color-carrier identity \(C\otimes C^*=\mathbf1\oplus\mathbf8\).
\end{theorem}

\begin{proof}
The scalar covariant kinetic term at \(h\) is the positive quadratic
form \(\langle\mathcal J_hX,B_h\mathcal J_hX\rangle\).
For every canonically normalized gauge vector \(Y\),
\[
 \langle Y,M_G^2Y\rangle
 =\|B_h^{1/2}\mathcal J_hQ_\rho^{-1/2}Y\|^2.
\]
Since both kinetic square roots are invertible, this vanishes exactly when
\(Q_\rho^{-1/2}Y\in\ker\mathcal J_h\), proving the displayed
transported kernel identities. The Lie bracket on the canonical null space
is the transported bracket
\([Y_1,Y_2]_{\rm kin}=Q_\rho^{1/2}
[Q_\rho^{-1/2}Y_1,Q_\rho^{-1/2}Y_2]\);
there is no assertion that ordinary component-wise brackets are preserved
by an arbitrary positive whitening.  For the spectral statement put
\(\mu=m_Xc_{\rm op}^2/\Eper\).  If
\(r_n=\sqrt{n^2+\mu^2}\in\mathbb Z\) for all sufficiently large \(n\), then
\[
 r_n-n=\frac{\mu^2}{r_n+n}\longrightarrow0.
\]
The left side is a nonnegative integer, so it is eventually zero and
\(\mu=0\).  The converse is immediate.  The Standard--Model stabilizer follows from the occurring-solder theorem: color is untouched and the chiral source stabilizer is exactly the electromagnetic generator in the electroweak factor.  No scalar-potential stationarity is used in this Gram-kernel statement.
\end{proof}



\begin{theorem}[Goldstone orbit, polar mass map, and longitudinal-degree transfer]
\label{p7v2:thm:v247-goldstone-hessian-transfer}
Under the hypotheses of
Theorem~\ref{p7v2:thm:v245-gauge-hessian-kernel-clock-lock}, perform the same
positive scalar and gauge kinetic whitening used to define the physical gauge
mass operator, and denote the resulting orbit map by
\[
 \widetilde{\mathcal J}_h:\mathfrak g_{\rm kin}\to V_{H,\rm kin}.
\]
Here the canonical gauge coordinate is \(Y=Q_\rho^{1/2}X\),
so \(\ker M_G^2=Q_\rho^{1/2}\mathfrak g_h\) and vectors in the
following polar decomposition are kinetic-coordinate vectors.
Normalize it so that
\[
 \boxed{M_G^2=\widetilde{\mathcal J}_h^\dagger
 \widetilde{\mathcal J}_h.}
\]
Its polar decomposition is
\[
 \boxed{\widetilde{\mathcal J}_h=U\Sigma,
 \qquad \Sigma=(M_G^2)^{1/2}.}
\]
On the broken-generator space
\(\mathfrak b:=(\ker M_G^2)^\perp\), \(U\) is an isometry onto
\(\mathcal G_h:=\operatorname{im}\widetilde{\mathcal J}_h\), the tangent
space to the scalar gauge orbit through \(h\).

Let \(V_{\rm alg}\) denote the gauge-invariant algebraic scalar potential at
fixed background metric. Its Hessian on canonically whitened scalar
coordinates is
\(\operatorname{Hess}_{\rm kin,h}V_{\rm alg}
:=B_h^{-1/2}(\operatorname{Hess}_hV_{\rm alg})B_h^{-1/2}\).
If \(h\) is stationary, \(\nabla V_{\rm alg}(h)=0\), then
\[
 \boxed{\operatorname{Hess}_{\rm kin,h}V_{\rm alg}\,
 \widetilde{\mathcal J}_hY=0\qquad(Y\in\mathfrak g_{\rm kin}).}
\]
Thus every broken gauge-orbit tangent is a scalar Hessian-zero direction;
the theorem does not exclude additional accidental flat directions.

Write the Goldstone fluctuation as \(\varphi_G=U\chi\),
\(\chi\in\mathfrak b\).  To quadratic order the normalized scalar kinetic
term on the broken block is
\[
 \boxed{\|\partial_\mu\chi+\Sigma A_\mu^{\mathfrak b}\|^2.}
\]
Under the linearized broken gauge transformation
\[
 A_\mu^{\mathfrak b}\mapsto A_\mu^{\mathfrak b}-\partial_\mu\epsilon,
 \qquad
 \chi\mapsto\chi+\Sigma\epsilon,
\]
the combination
\[
 \boxed{B_\mu=A_\mu^{\mathfrak b}+\Sigma^{-1}\partial_\mu\chi}
\]
is invariant and the kinetic term becomes \(\|\Sigma B_\mu\|^2\).
Consequently
\[
 \boxed{\operatorname{rank}M_G^2
 =\dim(G\!\cdot\! h)
 =N_{\rm Goldstone}
 =N_{\rm new\ longitudinal}.}
\]
For the Standard-Model Higgs branch,
\(\dim\mathfrak g_{\rm EW}=4\) and
\(\dim\mathfrak u(1)_{\rm em}=1\), so
\[
 \boxed{\operatorname{rank}M_{\rm EW}^2=3.}
\]
Exactly three orbit directions are transferred to the longitudinal
polarizations of \(W^+,W^-,Z\).  Equivalently the local bosonic degree count
is \(4\times2+4=1\times2+3\times3+1=12\).  The Gram rank used here is already fixed without scalar stationarity by Theorem~\ref{p7v2:thm:v245-gauge-hessian-kernel-clock-lock}.  The scalar Hessian-zero/Goldstone and stationary longitudinal interpretation in this theorem consume the explicitly stated stationarity condition; they do not create the native electroweak stabilizer.
\end{theorem}

\begin{proof}
The Gram identity is the normalized form of the scalar covariant-kinetic
mass construction already used in the preceding theorem.  Polar decomposition
then gives the isometry on the orthogonal complement of the kernel.  Gauge
invariance implies
\(\nabla V_{\rm alg}(e^{tX}h)=e^{tX}\nabla V_{\rm alg}(h)\); differentiating
at \(t=0\) and using stationarity gives
\((\operatorname{Hess}_hV_{\rm alg})\mathcal J_hX=0\) for every
raw generator \(X\). On the whitened scalar/gauge coordinates this becomes
\[
 B_h^{-1/2}(\operatorname{Hess}_hV_{\rm alg})B_h^{-1/2}
 \bigl(B_h^{1/2}\mathcal J_hQ_\rho^{-1/2}Y\bigr)=0,
\]
which proves the displayed kinetic-coordinate Hessian-zero identity
without equating raw and whitened generator vectors.
At quadratic order
\(D_\mu(h+\varphi_G)=U(\partial_\mu\chi+\Sigma A_\mu^{\mathfrak b})\),
which proves the Stückelberg combination and rank count.  The Standard-Model
rank is the dimension of the broken electroweak quotient.
\end{proof}

\begin{theorem}[Stationary longitudinal massive-vector spectrum on round \(S^3\)]
\label{p7v2:thm:v247-longitudinal-vector-spectrum}
Let \(m>0\) be one nonzero eigenvalue of the physical gauge mass Hessian in
Theorem~\ref{p7v2:thm:v247-goldstone-hessian-transfer}.  After the Goldstone
reduction, the corresponding gauge-invariant field \(B_\mu\) is a quadratic
Proca block.  On a static round spatial three-sphere let
\[
 \Delta_0Y_{k\alpha}=\lambda_kY_{k\alpha},
 \qquad
 \lambda_k=\frac{k(k+2)}{R_x^2},\qquad k\ge1.
\]
The longitudinal ansatz
\(B_0=u(t)Y_{k\alpha}\),
\(B_i=\phi(t)\nabla_iY_{k\alpha}\) has, in units
\(c_{\rm op}=\hbar=1\), the reduced mode Lagrangian
\[
 L_{\parallel,k}
 =\frac12\left[
 \lambda_k(\dot\phi-u)^2+m^2u^2-m^2\lambda_k\phi^2\right].
\]
The nondynamical temporal component satisfies
\[
 \boxed{u=\frac{\lambda_k}{\lambda_k+m^2}\dot\phi,}
\]
and hence
\[
 \boxed{L_{\parallel,k}^{\rm red}
 =\frac12\frac{\lambda_km^2}{\lambda_k+m^2}\dot\phi^2
 -\frac12m^2\lambda_k\phi^2.}
\]
After canonical normalization and restoration of units,
\[
 \boxed{\omega_{\parallel,k}^2
 =c_{\rm op}^2\frac{k(k+2)}{R_x^2}
 +\frac{m^2c_{\rm op}^4}{\hbar^2}.}
\]
There is no \(k=0\) longitudinal spatial mode because \(dY_0=0\).  The
kinetic coefficient tends to zero as \(m\to0\), so the extra physical
polarization disappears continuously in the massless gauge limit.

Writing \(n=k+1=L+1\ge2\), the already-derived transverse spectrum and the
new longitudinal spectrum are
\[
 \boxed{E_{\perp,n}^2=n^2\Eper^2+m^2c_{\rm op}^4,}
 \qquad
 \boxed{E_{\parallel,n}^2=(n^2-1)\Eper^2+m^2c_{\rm op}^4.}
\]
Therefore finite positive curvature gives the exact polarization splitting
\[
 \boxed{E_{\perp,n}^2-E_{\parallel,n}^2=\Eper^2,}
\]
which vanishes in the flat limit.  Per massive internal direction the
multiplicities are \(2(n^2-1)\) transversely and \(n^2\) longitudinally.
This closes the stationary quadratic longitudinal \(W/Z\) reduction only; it
does not determine interacting pole masses, widths, or nonstationary
Goldstone/gauge mixing.
\end{theorem}

\begin{proof}
For a gradient one-form \(dY_k\), \(d(dY_k)=0\), so the longitudinal magnetic
curvature vanishes and the electric term is
\((\dot\phi-u)dY_k\).  With
\(\int|dY_k|^2=\lambda_k\int|Y_k|^2\), the Proca quadratic action gives the
displayed Lagrangian.  Eliminating \(u\) algebraically yields the reduced
kinetic coefficient.  The ratio of the potential coefficient to that kinetic
coefficient is \(\lambda_k+m^2\) in \(c=\hbar=1\) units, giving the restored
frequency.  The transverse formula is
Theorem~\ref{p7v2:thm:v244-transverse-ew-wave-operators}; using
\(k(k+2)=(k+1)^2-1\) and
\(\hbar c_{\rm op}/R_x=\Eper\) proves the splitting.
\end{proof}

\begin{corollary}[Spectral reconstruction of the gauge mass Hessian]
\label{p7v2:cor:v247-gauge-mass-spectral-reconstruction}
On the calibrated stationary transverse gauge carrier,
\[
 \boxed{H_T^2
 =\Eper^2N_x^2\otimes I_{\mathfrak g}+c_{\rm op}^4M_G^2,}
\]
so
\[
 \boxed{M_G^2=c_{\rm op}^{-4}
 \left(H_T^2-\Eper^2N_x^2\otimes I_{\mathfrak g}\right).}
\]
Consequently, for one transverse mass eigen-direction \(X\),
\[
 \boxed{E_{X,n}^2-E_{X,m}^2=(n^2-m^2)\Eper^2,}
\]
whereas for two eigendirections \(X,Y\) at the same level,
\[
 \boxed{E_{X,n}^2-E_{Y,n}^2
 =(m_X^2-m_Y^2)c_{\rm op}^4.}
\]
The same \(m_X^2\) is reconstructed from the longitudinal tower by
\[
 \boxed{m_X^2c_{\rm op}^4
 =E_{\parallel,X,n}^2-(n^2-1)\Eper^2.}
\]
Thus internal Higgs Hessian data and compact spatial geometry separate into a
level-independent internal residual and a universal harmonic term.
\end{corollary}

\begin{corollary}[Internal tree-shell reconstruction]
\label{p7v2:cor:tree-shell-numerics}
If the common physical owner of
Theorem~\ref{p7:thm:physical-seven-nine-owner-reduction} is supplied,
\((\mathrm{EQAMP}_{79})\) holds so that
\(m_W^2/m_Z^2=7/9\), and the same branch uses
\[
 g_2=\frac{\sqrt{24/25}\,\sqrt7}{4\rho_V},
 \qquad \rho_V=\frac{\pi}{\sqrt{10}},
\]
then the consistently normalized tree-shell values are
\begin{align}
\label{p7v2:eq:tree-shell-numerics}
 g_2&=0.6523411464\ldots,
 &
 g_Y&=0.3486910099\ldots,
 &
 g_1&=0.4501581581\ldots,\\
 \alpha_2^{-1}&=29.52978731\ldots,
 &
 \alpha_1^{-1}&=62.01255336\ldots,
 &
 \alpha_{\rm em}^{-1}&=132.88404292\ldots,\\
 \sin^2\theta_W&=\frac29.
\end{align}
These are a conditional reconstruction from the imported mass dictionary and
equal-amplitude owner premise, not an independent derivation of that premise
or of renormalized couplings at \(M_Z\).
\end{corollary}

\begin{proof}
The mass ratio gives \(g_Y^2/g_2^2=2/7\).  Substitute this relation and
\(\rho_V=\pi/\sqrt{10}\) into
Theorem~\ref{p7v2:thm:ew-mass-dictionary}.
\end{proof}

\begin{remark}[Removed convention conflict]
The previous edition used \(b_1=41/10\), which presupposes the
GUT-normalized \(g_1\), but also inserted \(g_1\) directly into the ordinary
hypercharge mass formula.  Equations
\eqref{p7v2:eq:g1-convention}--\eqref{p7v2:eq:em-coupling-conventions} remove
that conflict.  Any scalar effective potential or thermal mass written with
GUT-normalized \(g_1\) must carry the same factor \(3/5\).
\end{remark}

\subsection{Continuum running and threshold matching}
\label{p7v2:sec:running}

\begin{theorem}[Conditional one-loop Standard-Model transport]
\label{p7v2:thm:conditional-one-loop-running}
Assume \(\Hqft\), three Standard-Model generations, one Higgs doublet, and
GUT-normalized \(g_1\).  In a mass-independent perturbative scheme,
\begin{equation}
\label{p7v2:eq:one-loop-running}
 \alpha_i^{-1}(\mu)
 =
 \alpha_i^{-1}(\mu_0)
 -\frac{b_i}{2\pi}\log\frac{\mu}{\mu_0},
 \qquad
 (b_1,b_2,b_3)
 =
 \left(\frac{41}{10},-\frac{19}{6},-7\right).
\end{equation}
The finite representation carrier fixes the group-theoretic traces entering
the coefficients once \(\Hqft\) is supplied.  It does not derive the
logarithmic loop law or a subtraction scheme.
\end{theorem}

\begin{proof}
Equation~\eqref{p7v2:eq:one-loop-running} is the standard one-loop
renormalization result for the stated field content
\cite{MachacekVaughn1984,PDG2024}.  The coefficients are sums of adjoint,
Weyl-fermion, and scalar Dynkin indices.  The finite carrier determines those
indices, while the loop integral, local action, regulator, counterterms, and
scheme are precisely the content of \(\Hqft\).  A sharp finite spectral-count
flow does not imply the logarithmic loop integral without that interface.
\end{proof}

\begin{corollary}[Conditional common-heat to equal-amplitude shell transport]
\label{p7:cor:conditional-common-heat-shell-transport}
Assume the hypotheses of
Theorem~\ref{p7v2:thm:conditional-one-loop-running}, no intervening threshold,
and the source-native common-heat boundary at \(\mu_h\),
\begin{equation}
 \label{p7:eq:common-heat-running-boundary}
 g_1^2(\mu_h)=g_2^2(\mu_h)=2\pi^2.
\end{equation}
The unique one-loop scale \(\mu_{79}\) at which the running couplings satisfy
the equal-amplitude shell condition
\begin{equation}
 \frac{g_Y^2}{g_2^2}=\frac27,
 \qquad\text{equivalently}\qquad
 \frac{g_1^2}{g_2^2}=\frac{10}{21},
\end{equation}
obeys
\begin{equation}
 \label{p7:eq:common-heat-shell-scale-ratio}
 \boxed{\log\frac{\mu_{79}}{\mu_h}=-\frac{88}{215},
 \qquad
 \frac{\mu_{79}}{\mu_h}=e^{-88/215}.}
\end{equation}
This is a conditional perturbative transport result.  Thresholds, a different
scheme, or failure of the common physical \(7\mid9\) owner change its
interpretation; the finite carrier alone does not identify either absolute
scale.
\end{corollary}

\begin{proof}
Write \(L=\log(\mu/\mu_h)\).  At the boundary
\(\alpha_1^{-1}=\alpha_2^{-1}=2/\pi\), while one-loop running gives
\[
 \alpha_i^{-1}(\mu)=\frac2\pi-\frac{b_i}{2\pi}L.
\]
Because \(g_1^2/g_2^2=\alpha_1/\alpha_2
=\alpha_2^{-1}/\alpha_1^{-1}\), setting this ratio to \(10/21\) and using
\(b_1=41/10\), \(b_2=-19/6\) yields
\[
 11\frac2\pi
 =\left(21b_2-10b_1\right)\frac{L}{2\pi}
 =-\frac{215}{2}\frac{L}{2\pi},
\]
and hence \(L=-88/215\).  The affine inverse couplings make this solution
unique on the threshold-free interval.
\end{proof}

\begin{proposition}[Physical threshold obligation]
\label{p7v2:prop:threshold-obligation}
A heavy threshold correction to
\(\alpha_i^{-1}\) requires the logarithmic determinant of an explicit heavy
fluctuation operator.  At one loop it has the schematic form
\begin{equation}
\label{p7v2:eq:threshold-form}
 \Delta_i(\mu)
 =
 -\frac{\Delta b_i}{2\pi}\log\frac{M}{\mu}
 +c_i^{\rm scheme},
\end{equation}
with spin, multiplicity, representation index, mass \(M\), and matching
constant specified.  A vector made only from geometric edge lengths is not a
threshold theorem.
\end{proposition}

\begin{proof}
Decoupling follows from the heavy-field functional determinant and depends on
the spectrum of the heavy operator
\cite{AppelquistCarazzone1975,Vassilevich2003HeatKernel}.  The data displayed
in \eqref{p7v2:eq:threshold-form} are not recoverable from an untyped edge
length alone.  Therefore the former Born--Pati--Salam edge vector can be used
as a numerical ansatz but not as a derived matching correction.
\end{proof}

\begin{remark}[Near-unification audit]
When the previous edition's one-loop boundary values and edge-vector
corrections are inserted, the arithmetic is reproducible: the pairwise
crossings occur between approximately
\(2.39\times10^{14}\,\mathrm{GeV}\) and
\(3.26\times10^{14}\,\mathrm{GeV}\), and evaluation at the proposed
\(\Lambda_8\) leaves an inverse-coupling spread \(0.22050\).  This is a
status-\(\mathsf B\) near-unification diagnostic, not exact unification.
\end{remark}

\begin{lemma}[Dimensionless data do not fix an absolute cutoff]
\label{p7v2:lem:absolute-scale}
No collection of dimensionless coupling ratios and finite carrier dimensions
can determine a dimensionful cutoff \(\Lambda\) without at least one
dimensionful anchor or a dimensional transmutation prescription.
\end{lemma}

\begin{proof}
Rescaling every dimensionful quantity by a common \(c>0\) preserves all
dimensionless ratios and carrier dimensions while sending
\(\Lambda\mapsto c\Lambda\).  Hence those data alone cannot select one
absolute value.
\end{proof}

\subsection{Electromagnetic finite sum rule}
\label{p7v2:sec:em-sum-rule}

\begin{definition}[Ordinary and normalized determinants]
\label{p7v2:def:normalized-determinant}
For a rank-one projection \(P\) and scalar \(\lambda>-1\),
\begin{equation}
\label{p7v2:eq:ordinary-vs-normalized-det}
 \det(I+\lambda P)=1+\lambda.
\end{equation}
On an eight-dimensional carrier with normalized trace
\(\tau_8=\Tr/8\), define
\begin{equation}
\label{p7v2:eq:normalized-det}
 \det_{\tau_8}(I+\lambda P)
 :=
 \exp\!\left(\tau_8\log(I+\lambda P)\right)
 =(1+\lambda)^{1/8}.
\end{equation}
The exponent \(1/8\) is therefore mathematically legitimate only after the
normalized trace state \(\tau_8\) has been supplied as the endpoint-state
member.  It is not the ordinary
finite-dimensional determinant.
\end{definition}

\begin{definition}[Electromagnetic endpoint realization family]
\label{p7v2:def:em-selector-family}
Before an electromagnetic endpoint is fixed, retain the typed family
\[
 \boxed{
 \mathfrak A_{\rm EM}
 =\{(B_{\rm obs},\tau_{\rm phys},a,\mathcal A_{\rm ang},\mathcal R_{\rm an})\}
 }
\]
of admissible positive endpoint operators, physical trace states, branch
coordinates, angular-return factors, and analytic-remainder prescriptions.
Definition~\ref{p7v2:def:em-candidate} is one declared finite member of this
family.  Neither carrier rank nor the normalized-octet presentation selects the
physical member.
\end{definition}

\begin{definition}[Electromagnetic candidate functional]
\label{p7v2:def:em-candidate}
For \(a>0\), define the finite candidate
\begin{equation}
\label{p7v2:eq:em-transport-candidate}
 \mathcal T_{\rm EM}^{[8]}(a)
 :=
 \exp\!\left(\frac{a^2}{\pi^2}\right)
 \left(1-a^4-a^6+\frac32a^8\right)^{1/8}
\end{equation}
and
\begin{equation}
\label{p7v2:eq:alpha-em-candidate}
 \Aem(a)
 :=
 \left[
   e^{\,4\pi^2-a}\,
   \mathcal T_{\rm EM}^{[8]}(a)
 \right]^{1/8}.
\end{equation}
This is status \(\mathsf A\).  It is not called an electromagnetic
prediction unless it is proved to represent the same-$\Gamma$ electromagnetic
response and the required infrared/analytic remainder is certified.  No
independent endpoint owner is thereby created.
\end{definition}

\begin{proposition}[What the candidate formula proves]
\label{p7v2:prop:em-candidate-conditional}
Assume all of the following:
\begin{enumerate}[label=(\roman*),leftmargin=2.5em]
  \item the physical endpoint trace is \(\tau_8\);
  \item on a domain where it is positive, an explicit boundary operator is
  \[
    B_{\rm obs}(a)
    =
    I+\left(-a^4-a^6+\frac32a^8\right)P;
  \]
  \item its independent angular return is exactly
  \(\exp(a^2/\pi^2)\);
  \item the dual hierarchy branch is \(e^{4\pi^2-a}\);
\end{enumerate}
Then the endpoint inverse coupling equals \(\Aem(a)\).
\end{proposition}

\begin{proof}
Definition~\ref{p7v2:def:normalized-determinant} gives
\[
 \det_{\tau_8}B_{\rm obs}(a)
 =
 \left(1-a^4-a^6+\frac32a^8\right)^{1/8}.
\]
Multiplying by the assumed angular and hierarchy factors and taking the outer
normalized-octet readout gives
\eqref{p7v2:eq:alpha-em-candidate}.  The conclusion is thus exact under the
four displayed hypotheses.  A formal connected logarithm
\(C(a)=(-a^4-a^6+a^8)P\bmod a^{10}\) reproduces the coefficient \(3/2\)
through order \(a^8\), but does not by itself prove that
\(B_{\rm obs}(a)\) is the exact analytic operator.
\end{proof}

\begin{theorem}[Current finite-sum-rule non-closure]
\label{p7v2:thm:em-nonclosure}
The existing GP1--GP3 and Paper-I-A--V-B data do not establish any of the four
load-bearing hypotheses of
Proposition~\ref{p7v2:prop:em-candidate-conditional}.  In particular:
\begin{enumerate}[label=(\alph*),leftmargin=2.5em]
  \item rank-two terminality does not determine the coefficient
  \(1/\pi^2\) in the angular exponential;
  \item Clifford anticommutation does not determine the complete connected
  coefficient sequence \((-1,-1,+1)\);
  \item the ordinary determinant does not produce an octet root;
  \item \(a=\delta\) and
  \(a=\delta-\delta^8/16\) are inequivalent branches;
  \item calculation in \(\mathbb Q[x]/(x^5)\) proves only a formal truncation,
  not an analytic remainder estimate.
\end{enumerate}
Consequently \(\Aem\) remains a finite-sum-rule conjecture.  This theorem
no longer classifies the physical electromagnetic endpoint itself as an
independent owner; the preferred physical route is the same-effective-action
descendant of Corollary~\ref{p7v2:cor:em-same-gamma-limit-route}.
\end{theorem}

\begin{proof}
Items (a)--(c) follow from Definition~\ref{p7v2:def:normalized-determinant}
and from the counterfamily in Theorem~\ref{p7v2:thm:type-separation}.  For
(d), direct evaluation gives different values in
Remark~\ref{p7v2:rem:em-numerics} below.  For (e), equality modulo \(x^5\)
places no bound on a representative's coefficient of \(x^5\) or on the
analytic tail of a lift.  Therefore the stated upstream data admit inequivalent
operators and endpoints.
\end{proof}

\begin{corollary}[Same-effective-action route to electromagnetic endpoint closure]
\label{p7v2:cor:em-same-gamma-limit-route}
Suppose a generated physical effective action $\Gamma_{\rm eff}$ and stable
nonzero electroweak vacuum are owned.  Let $\iota_Q$ denote the generated
unbroken electromagnetic embedding and define the quadratic response
\[
 K_Q:=\iota_Q^*\left.\delta_A^2\Gamma_{\rm eff}\right|_{\Phi_*}\iota_Q.
\]
Then Paper~II-A's generated-constructor chronology, joint-owner saturation and
whole-family reuse make $K_Q$ a same-owner descendant of the complete
effective-action presentation; no independent electromagnetic endpoint owner
is licensed on this route.  If its transverse momentum-space coefficient
$Z_Q(q^2)$ is defined on a punctured neighbourhood of $q^2=0$ and the licensed
infrared routes have the common finite limit
\[
 \lim_{q^2\to0}Z_Q(q^2)=Z_Q(0),
\]
then Paper~VII-C Theorem~\ref{p12:thm:canonical-limit-descendant} makes
$Z_Q(0)$ the unique electromagnetic endpoint descendant.  A finite member of
$\mathfrak A_{\rm EM}$ may then be retained only as a proved representation or
evaluator of that same response.  The present finite candidate formula does
not by itself establish the representation equality or the IR limit.
\end{corollary}

\begin{theorem}[II-A contraction of the electromagnetic endpoint owner]
\label{p7v2:thm:em-owner-contraction}
On any branch satisfying the hypotheses of
Corollary~\ref{p7v2:cor:em-same-gamma-limit-route} up to existence of the
infrared limit, the structural electromagnetic response owner is already
\[
 (\Gamma_{\rm eff},\Phi_*,\iota_Q)\longmapsto K_Q.
\]
Accordingly the physical G17 obligation is not an independent choice of
$B_{\rm obs}$, trace state, branch or angular factor.  Its unresolved content
is upstream ownership of the complete effective action/stable vacuum together
with the analytic and renormalization/matching conditions needed to establish
and identify the unique $q^2\to0$ coefficient.  Any alternative finite
endpoint construction must map to this same response before it can discharge
the physical endpoint claim.
\end{theorem}

\begin{proof}
The Hessian and restriction are licensed descendants of the same complete
action/vacuum/embedding owner.  Paper~II-A therefore retains their generated
joint incidence and forbids reintroduction of an independent endpoint member.
The only remaining selector-free step is analytic endpoint formation; that is
exactly the hypothesis isolated by the canonical-limit theorem.
\end{proof}

\begin{remark}[Numerical audit of the finite sum rule]
\label{p7v2:rem:em-numerics}
With \(\delta=1/\sqrt{72}\), the order-\(a^8\) logarithmic expansion gives
\begin{align}
 \Aem(\delta)
 &=137.035999137806\ldots,\\
 \Aem\!\left(\delta-\frac{\delta^8}{16}\right)
 &=137.035999176726\ldots.
\end{align}
Keeping the compact factor in
\eqref{p7v2:eq:em-transport-candidate} gives
\(137.035999175630\ldots\) on the half-octic branch.  The arithmetic is
correct; the coexistence of two branches proves that the number is not unique.
Moreover,
\[
 \delta^{10}=5.16817865\times10^{-10}.
\]
An omitted term \(C a^{10}\) in the logarithm entering the outer
one-eighth readout shifts the displayed inverse coupling by approximately
\(8.85\times10^{-9}C\).  The printed final digits are therefore uncertified
until \(\Hem\) supplies a bound on \(C\) and the remaining tail.
\end{remark}

\subsection{Cycle-free electromagnetic running}
\label{p7v2:sec:em-running}

\begin{proposition}[Vacuum-polarization interface]
\label{p7v2:prop:em-running-interface}
Assume \(\Hqft\) and a specified renormalization prescription.  Then the
relation between the Thomson-limit coupling and a timelike scale \(q^2\) may
be written
\begin{equation}
\label{p7v2:eq:em-running-interface}
 \alpha(q^2)
 =
 \frac{\alpha(0)}
 {1-\Delta\alpha_{\rm lep}(q^2)
    -\Delta\alpha_{\rm had}^{(5)}(q^2)
    -\Delta\alpha_{\rm top}(q^2)
    -\cdots}.
\end{equation}
The leptonic terms are perturbative QED outputs.  The hadronic term is a
nonperturbative QCD or data-dispersion output and is not computed in
Paper~V-C.
\end{proposition}

\begin{proof}
Equation~\eqref{p7v2:eq:em-running-interface} is the standard
vacuum-polarization definition of the running electromagnetic coupling.  At
low quark momenta the hadronic correlator is not determined by a bare color
carrier or by perturbative threshold counting.  It belongs to the
Paper-VII-A--VII-B execution layer.
\end{proof}

\begin{corollary}[No Paper-V-C value of \(\alpha(M_Z)\)]
\label{p7v2:cor:no-alpha-mz}
Paper~V-C does not import a hadronic response number from Papers VII-A--VII-B and
therefore does not claim a closed numerical \(\alpha(M_Z)\).  A later synthesis
may combine a selected status-\(\mathsf A\) endpoint \(\Aem(a)\) with an
independently computed
\(\Delta\alpha_{\rm had}^{(5)}(M_Z^2)\), but that result is downstream of the
present paper.
\end{corollary}

\label{p7v2:part:qcd-ew}

\subsection{same-source electromagnetic boundary response}
\label{p7v2:sec:v2141-em-boundary-response}

\begin{theorem}[Electromagnetic Hessian is a retained-source descendant]
\label{p7v2:thm:v2141-em-same-source-hessian}
On an occurring branch with the structural $U(1)_Q$ embedding, retain the
background electromagnetic source through the exact finite blocking tower and
define
\begin{equation}
\label{p7v2:eq:v2141-KQ}
 K_{Q,k}:=\iota_Q^*\,\delta_A^2\Gamma_k^{\partial}\,\iota_Q .
\end{equation}
Whenever the Ward/background identity gives the transverse vacuum form,
\begin{equation}
\label{p7v2:eq:v2141-ZQ}
 \boxed{
 K_{Q,k}^{\mu\nu}(q)
 =Z_Q(k,q^2)(q^2g^{\mu\nu}-q^\mu q^\nu)+\cdots .}
\end{equation}
The physical Thomson normalization, when the infrared limit exists, is
\begin{equation}
\label{p7v2:eq:v2141-alpha-thomson}
 \boxed{
 \alpha_{\rm em}^{-1}(0)
 =4\pi\lim_{q^2\to0}Z_Q(0,q^2).}
\end{equation}
The background Hessian and the conserved-current response are two derivatives
of this same retained source: at finite blocking level,
\begin{equation}
\label{p7v2:eq:v2141-current-hessian}
 \boxed{
 K_Q^{\mu\nu}
 =\left\langle\frac{\delta^2S}{\delta A_{Q,\mu}\delta A_{Q,\nu}}\right\rangle
 -\langle J_Q^\mu J_Q^\nu\rangle_c .}
\end{equation}
Thus no independent electromagnetic current owner is required for the
quadratic response once the source/action occurrence is fixed.
\end{theorem}

\begin{corollary}[One-loop electromagnetic boundary tangent]
\label{p7v2:cor:v2141-em-one-loop-tangent}
On the v2.140 local-QFT branch,
\[
 b_{\rm em}=b_2+\frac53b_1
 =-\frac{19}{6}+\frac53\frac{41}{10}
 =\boxed{\frac{11}{3}},
\]
so
\begin{equation}
\label{p7v2:eq:v2141-em-one-loop-tangent}
 \boxed{
 \frac{d\alpha_{\rm em}^{-1}}{d\log k}=-\frac{11}{6\pi}.}
\end{equation}
Equivalently the charged Dirac-fermion, charged-scalar and charged gauge/quotient
pieces are $32/3$, $1/3$ and $-22/3$, respectively.  This is a local Gaussian
tangent of the exact resolution-boundary response, not a long-distance
solution by itself.
\end{corollary}

\begin{proposition}[Strong common-heat obstruction to long one-loop transport]
\label{p7v2:prop:v2141-common-heat-strong-obstruction}
At the native common-heat datum,
\[
 \alpha_1^{-1}=\alpha_2^{-1}=\alpha_3^{-1}=\frac2\pi.
\]
Writing $L=\log(k/\mu_h)$, the one-loop electroweak inverse couplings are
\[
 \alpha_1^{-1}(L)=\frac2\pi-\frac{41}{20\pi}L,
 \qquad
 \alpha_2^{-1}(L)=\frac2\pi+\frac{19}{12\pi}L .
\]
Hence
\begin{equation}
\label{p7v2:eq:v2141-common-heat-no-weak-corridor}
 \boxed{\min(\alpha_1^{-1},\alpha_2^{-1})\le\frac2\pi\quad\text{for all }L.}
\end{equation}
Therefore the common-heat point has no parametrically weak one-loop corridor
that can be extrapolated over a long range to the physical IR.  Exact/nonperturbative
flow, finite matching charts, or a differently typed matching step is required
for absolute transport.
\end{proposition}

\begin{corollary}[Retyping of the historical finite electromagnetic candidate]
\label{p7v2:cor:v2141-em-candidate-retyping}
A Ward-subtracted vacuum-polarization determinant fixes differences/running and
vanishes at its subtraction point; it does not by itself set the absolute
Thomson coefficient.  Consequently the historical finite candidate
$\mathcal A_{\rm em}(a)$ of
Definition~\ref{p7v2:def:em-candidate}, if a same-action derivation is ever
completed, must be typed as an absolute source/matching normalization (or an
equivalent endpoint evaluator), not as ordinary RG running from common heat.
It remains conditional because its operator, normalized trace, branch and
analytic remainder are not derived by the present theorem.
\end{corollary}

\paragraph{G17 contraction.}
For control purposes the electromagnetic endpoint frontier is decomposed as
\[
 \boxed{G17=G17_{\rm source}+G17_{\rm vacuum}+G17_{\rm match}+G17_{\rm IR}.}
\]
The retained-source/Hessian owner is closed on the finite branch above;
$G17_{\rm IR}$ reduces to the charged-gap/uniformity certificate of Paper~VII-C.
The historical symbol $G17_{\rm match}$ denoted the then-unresolved absolute
heat-to-Thomson interface; Release v2.75 below supersedes that coarse typing by
separating the closed no-free-source-normalization theorem from the still-open
endpoint-response evaluation.  Any claim that depends on a particular stable
physical vacuum retains its separately typed vacuum occurrence/stability
requirements.
\paragraph{Finite current measure is a dependent readout.}
Once the physical finite Hamiltonian, state and normalized current are
retained, its current-weighted transition measure is a spectral functional of
that packet, not a second electromagnetic owner.  The same source fixes its
local contacts.  Their execution and the zero-atom, inverse-moment and uniform
IR/UV/refinement conditions of Paper~VII-C remain the physical endpoint tasks.

% END EM-SAME-SOURCE-BOUNDARY-RESPONSE v2.71


% BEGIN BL2-OWNER-CLOSURE-IMPORT v2.72
\subsection{Same-action electromagnetic and source--detector certificates}
\label{p5c:sec:v2210-lit-observable-certificates}
The following conditional results use the existing action and its complete
source derivatives.  They give explicit discharge tests, not substitute
couplings, density matrices, detector effects or scalar vacuum selectors.

\begin{theorem}[A sufficient neutral-response Thomson certificate]
\label{p5c:thm:v2210-lit-thomson-response}
Use the generated unbroken \(U(1)_Q\) restriction of the actual effective
action and a declared physical infrared approach \(t=q^2\downarrow0\).
After the physical transverse and kinetic restriction, take the surviving
massless \(Q\) block to be one-dimensional per physical transverse
polarization, with the same scalar coefficient in each polarization.
Equivalently, its endpoint coefficient is \(\kappa_0\) times the
polarization identity.  Suppose its complete neutral quadratic kernel
has the same-owner block form
\[
 K(t)=\begin{pmatrix}
 t\,a(t)&t\,b(t)\\
 t\,b(t)^\dagger&C+t\,d(t)
 \end{pmatrix},\qquad C>0,
\]
where all blocks are continuous at zero and the massive complement is
invertible nearby.  The source normalization, current vertices and regulator
Ward identity are those of this same action.  Let \(g_Q(t)\) be the
resulting unbroken-current coefficient after eliminating that complement,
and assume \(g_Q(t)\to g_0\) and \(a(0)=\kappa_0>0\).
Then the unique canonically normalized quadratic endpoint is
\[
 \boxed{K_Q^{\rm eff}(t)=t\,\kappa(t),\qquad
 \kappa(t)=a(t)-t\,b(t)(C+t\,d(t))^{-1}b(t)^\dagger,}
\]
\[
 \boxed{\kappa(0)=\kappa_0,\qquad
 e_Q(0)=\frac{g_0}{\sqrt{\kappa_0}},\qquad
 \alpha_Q(0)=\frac{g_0^2}{4\pi\kappa_0}.}
\]
Here \(e_Q\) is dimensionless in the standard canonically normalized
\(-F^2/4\) convention.  Identifying this endpoint with a Thomson
experiment additionally consumes the actual charge-universality, external
source and infrared certificate stated below.
\end{theorem}
\begin{proof}
Exact elimination of the massive block gives
\(K_{QQ}-K_{QX}K_{XX}^{-1}K_{XQ}\), which is the displayed Schur kernel.
Continuity and \(C>0\) bound its inverse near zero; the subtracted term
in \(\kappa(t)\) is therefore \(O(t)\), so \(\kappa(0)=a(0)\).
The action rescaling \(A_{Q,c}=\sqrt{\kappa(t)}A_Q\) sends the complete
current coefficient to \(g_Q(t)/\sqrt{\kappa(t)}\).  Taking the positive
limit proves the formulas.  All coefficients are derivatives/readouts of
the fixed action; none is selected by the desired endpoint value.
\end{proof}

% BEGIN S3 LITERATURE INTEGRATION v3.040 [V-C: ELECTROMAGNETIC READOUT CONTRACT]
\begin{remark}[Same-owner electromagnetic packet and Thomson acceptance]
\label{p5c:rem:v340-s3lit-thomson-packet}
For the actual physical comparison in
Theorem~\ref{p5c:thm:v2210-lit-thomson-response}, require an
explicit finite packet with the selected $H_L$, normalized charged
source state and $U(1)_Q$ background, transverse current operators
$J^i_{{\rm em},L}$, positive current-visible gaps or proved
inverse-moment control, and every second-source contact jet of
that same action.  The packet must include physical transverse
gauge restriction, Ward tests on a gauge-invariant source window,
the spatial momentum response, and uniform refinement/infrared
error bounds before declaring the actual Thomson endpoint.
On a finite temporal reflection-positive chart, the independently
verified current variance $V$, double-commutator $C$ and minimum
visible gap $m_J$ can first be tested against
Corollary~\ref{p12:cor:v340-s3lit-current-envelope}; the current
contribution alone is \emph{not} the complete kinetic/Thomson
coefficient.  Use the source-owned physical form factor, not a
chosen momentum POVM, to test a short-distance Coulomb observable.
No external target value may be used to select a current, weight,
physical pole, source normalization or localized state.  The actual
measured $\alpha_{\rm em}(0)$ and short-distance Coulomb core are
not derived by this execution contract.
\end{remark}
% END S3 LITERATURE INTEGRATION v3.040 [V-C]


\begin{proposition}[Uniform source-limit error bound]
\label{p5c:prop:v2210-lit-thomson-uniform}
For an actual comparison family \(j\), suppose the preceding elimination
is valid uniformly and its completed coefficients satisfy
\[
 \kappa_j(t)\ge\kappa_*>0,\qquad\kappa_0\ge\kappa_*,
\]
\[
 |g_j(t)-g_0|\le r_g(t)+\delta_{g,j},\qquad
 |\kappa_j(t)-\kappa_0|\le r_\kappa(t)+\delta_{\kappa,j},
\]
where \(r_g(t),r_\kappa(t)\to0\) and
\(\delta_{g,j},\delta_{\kappa,j}\to0\) on the declared directed family.
Then
\[
 \boxed{|e_j(t)-e_0|
 \le\frac{r_g(t)+\delta_{g,j}}{\sqrt{\kappa_*}}
 +\frac{|g_0|[r_\kappa(t)+\delta_{\kappa,j}]}{2\kappa_*^{3/2}}.}
\]
Consequently the joint endpoint is unique; either iterated limit, when
defined on this uniform domain, has that same value.
\end{proposition}
\begin{proof}
Add and subtract \(g_0/\sqrt{\kappa_j}\).  The derivative of
\(x^{-1/2}\) has modulus at most \(1/(2\kappa_*^{3/2})\) on the positive
interval, so the mean-value bound gives the inequality.  Its right-hand
side tends to zero on the product-directed approach, which proves the
joint and iterated uniqueness statements.
\end{proof}

\begin{corollary}[Ward/current completion of the Thomson interpretation]
\label{p5c:cor:v2210-lit-thomson-ward}
On a charged-source domain with the complete regulated Ward identity
\[
 q_\mu\Gamma_Q^\mu(p+q,p)
 =e_Q\{Q S^{-1}(p+q)-S^{-1}(p)Q\},
\]
assume a regular zero-transfer vertex limit, the corresponding physical
charged-leg normalization and a compatible Thomson/infrared prescription.
Then contraction with the normalized charged legs fixes the zero-transfer
charge to \(e_Q Q\).  With the endpoint theorem and its actually consumed
uniform limit certificate, \(\Hem\) has no further independent source
normalization or finite-sum factor.
\end{corollary}
\begin{proof}
In a fixed charged block \(Q\) commutes with the inverse propagator.
Divide the Ward identity by the transfer and take the regular limit to get
\(\Gamma_Q^\mu=e_Q Q\partial_{p_\mu}S^{-1}\).  The stated external-leg
normalization contracts this derivative to its physical unit current;
hence its coefficient is \(e_Q Q\).  Charge universality and the
photon/neutral self-energy route are established in their quantum domains
by Dittmaier \cite{LitSMWardThomson2021}; the native application above
exposes the identities and limits it must actually satisfy.
\end{proof}
This does not assume an isolated bare charged-particle pole in an untreated
massless theory.  A charged pole, dressed source or inclusive infrared route
must be the one explicitly certified by the requested observable.  A compact
discrete spectrum alone does not furnish the assumed continuous physical
\(q^2\downarrow0\) approach.  Nor has \(\kappa_0\), the actual Ward
certificate or its uniform errors been numerically evaluated here.  The
finite electromagnetic sum remains only a representation candidate until
identified with this response and its remainder is bounded.  No stationary
scalar equilibrium or fundamental \(v\) is an antecedent of this theorem.

\begin{theorem}[Actual source--propagator--detector contraction]
\label{p5c:thm:v2210-lit-detector-kernel}
Let the actual normalized production density be \(\rho_P\ge0\),
\(\operatorname{Tr}\rho_P=1\), let the generated detector effect obey
\(0\le E_D\le I\), and let \(V\) be the unitary physical propagator on the
consumed finite support.  Then
\[
 \boxed{P_D=\operatorname{Tr}(E_DV\rho_PV^\dagger)\in[0,1].}
\]
On a diagonal phase chart \(V=\operatorname{diag}(e^{-i\Phi_r})\),
\[
 \boxed{P_D=\sum_r(\rho_P)_{rr}(E_D)_{rr}
 +2\operatorname{Re}\sum_{r<s}
 (\rho_P)_{rs}(E_D)_{sr}e^{-i(\Phi_r-\Phi_s)}.}
\]
Common unitary changes of presentation of \(\rho_P,E_D,V\) leave this
probability unchanged.  For equal normalized pure production/detection
boundaries with role weights \(p_r\), if one role \(r\) has phase difference
\(\Delta\) from all remaining phase-coincident roles, the survival
probability reduces exactly to
\[
 \boxed{P_{\rm surv}=1-4p_r(1-p_r)\sin^2(\Delta/2).}
\]
\end{theorem}
\begin{proof}
Positivity and \(E_D\le I\) bound its expectation in the unit-trace evolved
density.  Expand the matrix trace and combine each pair of Hermitian
conjugate entries to obtain the interference formula.  Trace cyclicity
proves common presentation covariance.  In the equal pure-boundary case
the survival amplitude is \(\sum_s p_s e^{-i\Phi_s}\); under phase
coincidence it is, up to a common phase,
\((1-p_r)+p_r e^{-i\Delta}\).  Its squared modulus gives the last identity.
\end{proof}
Wave-packet/QFT production and detection can supply such physically
normalized boundaries \cite{LitSMNeutrinoQFT2010}; reactor EFT matching
\cite{LitSMReactorEFT2019} shows why source/detector coefficients depend on
the actual interaction and process.  These references do not select the
native \(\rho_P,E_D\), or permit fitted replacement angles.

\begin{corollary}[Equal-boundary obstruction and minimum state displacement]
\label{p5c:cor:v2210-lit-reactor-obstruction}
For the recorded rooted current's weight set \((1/4,1/4,1/2)\), the
separated-frequency equal-boundary disappearance amplitude is in
\(\{3/4,1\}\).  Permutations and rephasings cannot alter that set.
If an equal pure-boundary repair has separated-role weight \(q\), the
projective angle from a row with old weight \(p\) is at least
\[
 \arccos\{\sqrt{pq}+\sqrt{(1-p)(1-q)}\},
\]
and its pure-state trace distance is at least
\[
 \boxed{\sqrt{1-[\sqrt{pq}+\sqrt{(1-p)(1-q)}]^2}.}
\]
\end{corollary}
\begin{proof}
Substitute \(p=1/4\) or \(1/2\) into the survival theorem.  Separate two
unit vectors into the one-role component and its orthogonal complement.
Cauchy--Schwarz bounds their overlap by the displayed sum.  For pure states
the projective angle is \(\arccos|\langle\psi,\phi\rangle|\) and the
half-trace-norm distance is
\(\sqrt{1-|\langle\psi,\phi\rangle|^2}\), giving both bounds.
\end{proof}
The Daya Bay final fit reports amplitude \(0.0851\pm0.0024\)
\cite{DayaBayFinal2023}.  Its central small-weight branch is
\(q=(1-\sqrt{1-0.0851})/2=0.02174797439\ldots\); against \(p=1/4\)
the minimum angle is \(0.37559\) radians and trace distance \(0.36682\).
This is a quantitative test of the equal pure-boundary repair, not a bound
on general unequal/mixed production and detector effects or a full event
likelihood.  For those effects the complete interference coefficients above
must be executed.  The new harmonic mass/incidence constructor remains
independent of this historical current test; its current is not selected by
the failed comparison and is not automatically inherited from an obsolete
neutral mass word.

\begin{proposition}[Scalar-coordinate-independent soft response]
\label{p5c:prop:v2210-lit-scalar-coordinate}
On a declared real scalar chart \(h\), let the positive kinetic metric be
\(G(h)\,dh^2\) and let \(K_V(z;h)\) be the actual source-normalized inverse
vector response.  Its unit-direction soft scalar insertion is
\[
 \mathcal V_X(z;h)=G(h)^{-1/2}\partial_h K_V(z;h).
\]
Under a smooth orientation-preserving scalar reparametrization \(y=f(h)\), with
the action, kinetic metric and source transformed together,
\(G_y=G/(f')^2\), \(\mathcal V_X\) is unchanged.  At a stationary
scalar background the normalized scalar Hessian is likewise invariant.
The value of a vector pole mass alone does not determine
\(\mathcal V_X\) or the finite-momentum three-point response.
\end{proposition}
\begin{proof}
The chain rule gives
\(G_y^{-1/2}\partial_y=G^{-1/2}\partial_h\).
For the scalar Hessian, the additional second-derivative coordinate term
is proportional to the first variation and vanishes at stationarity; the
metric rescales by the same two derivative factors.  Finally two response
families can have the same value at \(h_*\) and different first derivative
there, so the pole value does not fix its insertion.  A soft derivative has
zero scalar transfer and is not the whole finite-momentum Higgs vertex.
\end{proof}
This is a native observable comparison route compatible with geometric
scalar/HEFT coordinates \cite{LitSMScalarGeometry2016,LitSMHEFT2014}.
No conventional \(v\) mass generator is imposed.  The physical Higgs
couplings still require the complete two-/three-point response, residues
and external channel normalization; high-energy longitudinal amplitudes
must be tested from that same action.  The finite \(17\) and \(7\mid9\)
supports do not evaluate these derivatives or prove
\(\mathrm{EQAMP}_{79}\).

\paragraph{Updated local open-data ledger.}
\(\Hem\) is narrowed to the actual complete quadratic/current response,
its Ward/source interpretation and the consumed infrared bound of
\ref{p5c:thm:v2210-lit-thomson-response}--\ref{p5c:cor:v2210-lit-thomson-ward}.
Common-heat couplings remain native matching-datum values; the physical
scale, nonperturbative transport and thresholds are still unexecuted.
\(\Hcp\) concerns regulated strong-phase/measure transport compatible with
weak-family CP, while the finite native strong character remains fixed.
The harmonic neutral origin and its actual incidence/propagation/production
and detector map remain separately open.  Complete inherited BL2 charge
history and actual anomalous baryon export, rather than pair number or a CP
quartet alone, remain the signed-charge task.  Optional Majorana/seesaw,
saturation and scalar-equilibrium branches keep their existing scopes.
Physical rates retain the claim-indexed scattering/infrared domain already
stated in \ref{p7v2:prop:claim-scoped-decay-boundary}; finite vertices or
kinematic support do not create an irreversible width.
% END LITERATURE-VERIFIED EM AND OBSERVABLE INTERFACES v2.210



\subsection{no-free electromagnetic source matching and endpoint-response frontier}
\label{p5c:sec:v275-g17-no-free-matching}

\begin{theorem}[No independent heat-to-Thomson source-normalization coefficient]
\label{p5c:thm:v275-no-free-em-matching}
Assume an occurring branch on which the generated electromagnetic embedding
$\iota_Q$ is retained as the same background-source line through the exact
finite blocking tower.  Let the source-native common-heat action be the parent
presentation and write its retained electromagnetic source as $A_Q$.  Suppose
no independently generated owner-changing or source-line deformation acts on
this line.  Then every exact blocked descendant uses the same normalized source
coordinate $A_Q$.  In particular, a putative matching relation
\[
 A_{Q,L}=c_LA_{Q,h}
\]
with an independently chosen $c_L$ is admissible on this branch only for
\[
 \boxed{c_L=1.}
\]
Equivalently, the physical quadratic response is
\[
 \boxed{
 K_{Q,L}=\iota_Q^*D_A^2\,\mathscr R_{hL}\Gamma_h[A_Q] \,\iota_Q ,}
\]
with no additional multiplicative matching coefficient outside the exact
effective-action/blocking response.  Threshold, contact, vacuum-polarization,
and interacting hadronic contributions may change the coefficient of this
Hessian; they are contributions to the same blocked effective action, not a
second normalization of the electromagnetic source.
\end{theorem}

\begin{proof}
Paper~VII-C Theorem~\ref{p12:thm:v2141-blocking-source-naturality} retains
$A_Q$ as a parameter of every fibre integral and gives
$D_A^r\mathscr R_{hL}\Gamma_h=D_A^r\Gamma_L$ for every controlled derivative.
Thus exact blocking transports the action and its electromagnetic Hessian in
one source coordinate; it does not introduce a source-coordinate change.
Paper~II-A Theorem~\ref{p2:thm:gp2-generated-owner-normalization-inheritance}
then applies to this already generated same-owner source line.  A replacement
$A_Q\mapsto c_LA_Q$ with $c_L\neq1$ is a nontrivial endomorphism of its inherited
normalization slot and requires an independently generated deformation.
By hypothesis no such deformation occurs.  The same conclusion follows from
generated co-consumption,
Theorem~\ref{p2a:thm:v3101-generated-coconsumption}: the source incidence and
its normalization cannot be independently rescaled while claiming the same
generated output.  Hence $c_L=1$.  Since the blocked action itself is allowed
to change, the conclusion does not set $K_{Q,L}=K_{Q,h}$; it excludes only an
extra normalization outside the exact response.
\end{proof}

\begin{corollary}[Source-native common-heat electromagnetic reference coefficient]
\label{p5c:cor:v275-em-heat-reference}
On the same $Q$ source line, evaluate the common-heat gauge normalization of
Paper~V-A in the GUT-normalized basis.  Since
\[
 (g_1^{\rm heat})^2=(g_2^{\rm heat})^2=2\pi^2,
 \qquad
 \alpha_i=\frac{g_i^2}{4\pi},
\]
and the electromagnetic coordinate satisfies
$\alpha_{\rm em}^{-1}=\alpha_2^{-1}+\frac53\alpha_1^{-1}$, one obtains
\[
 \boxed{
 \alpha_{Q,h}^{-1}
 =\frac2\pi+\frac53\frac2\pi
 =\frac{16}{3\pi}.}
\]
In the Hessian convention of
Theorem~\ref{p7v2:thm:v2141-em-same-source-hessian},
$\alpha_{\rm em}^{-1}=4\pi Z_Q$, so the corresponding source-native reference
coefficient is
\[
 \boxed{Z_{Q,h}=\frac{4}{3\pi^2}.}
\]
This is the common-heat coefficient on the retained $Q$ source line, not the
physical Thomson endpoint.
\end{corollary}

\begin{corollary}[Exact G17 endpoint-response identity]
\label{p5c:cor:v275-g17-endpoint-identity}
If the physical branch admits the controlled $k\to0$, $q^2\to0$ endpoint of
Theorem~\ref{p7v2:thm:v2141-em-same-source-hessian}, define the same-source
response increment
\[
 \Delta Z_{Q,h\to0}
 :=\lim_{q^2\to0}Z_Q(0,q^2)-Z_{Q,h}.
\]
Then
\[
 \boxed{
 \alpha_{\rm em}^{-1}(0)
 =\frac{16}{3\pi}+4\pi\,\Delta Z_{Q,h\to0}.}
\]
The quantity $\Delta Z_{Q,h\to0}$ is not a new calibration parameter: it is the
net exact same-action response generated by blocking, thresholds, contact
terms and charged current correlations.  Its numerical evaluation still
requires the relevant continuum/IR certificates, including the actual QCD
current-weighted limit where consumed.
\end{corollary}

\begin{corollary}[G17 matching-token retirement]
\label{p5c:cor:v275-g17-match-retired}
The historical aggregate token $\texttt{G17\_MATCH}$ is too coarse and is
retired.  Its two logically distinct contents are
\[
 \boxed{\texttt{G17\_SOURCE\_NORMALIZATION\_MATCH}
 =\texttt{CLOSED\_NO\_FREE\_COEFFICIENT},}
\]
\[
 \boxed{\texttt{G17\_ENDPOINT\_RESPONSE}
 =\texttt{OPEN\_READOUT/ANALYTIC},}
\]
with $\texttt{G17\_IR\_QCD}$ retained as the actual hadronic current-measure
and uniformity subproblem.  Thus the remaining electromagnetic frontier is an
exact-response evaluation problem, not an independently adjustable absolute
matching normalization.
\end{corollary}
% END G17-NO-FREE-MATCHING v2.75

% BEGIN LEPTON-DIRECTION-CP-BRANCH-CORRECTION v2.76
\subsection{response and zero-temperature gauge vertex}
\label{p7:sec:v164-response-gauge-vertex}

\begin{theorem}[Renormalized response owner]
\label{p7:thm:v164-renormalized-response-owner}
On its admitted source, state and physical gauge/constraint quotient, let
\(\Gamma_R\) denote the complete renormalized effective action consumed
by the claim.  Its complete quadratic kernel and proper three-point vertex
are
\[
 K_{R,\rm full}=\delta^2\Gamma_R,
 \qquad V_{R,\rm proper}=\delta^3\Gamma_R,
 \qquad G_R=K_{R,\rm full}^{-1}
\]
on a regular chart, with the stipulated signature/continuation convention.
If a separately defined renormalized core kernel \(K_{R,0}\) is used,
define its complete residual return by
\[
 \boxed{K_{R,\rm full}=K_{R,0}-\Sigma_R,
 \qquad G_R=(K_{R,0}-\Sigma_R)^{-1}.}
\]
The same \(\Sigma_R\) is never subtracted again from
\(\delta^2\Gamma_R\).  For an effective average action with a retained
quadratic regulator, the corresponding identity is instead
\(G_k=(\Gamma_k^{(2)}+R_k)^{-1}\); its regulator is removed only on a
certified limit.
The one-loop \(7/9\) crossing remains a conditional RG diagnostic.  It does
not turn the finite \(7:9\) shell blocks into the charged and neutral
Higgs-Hessian eigenspaces, and it does not identify a finite row with an
interacting pole.
\end{theorem}
\begin{proof}
On a regular source chart the Legendre inverse-Hessian identity gives the
inverse of the complete effective Hessian; with a subtracted quadratic
regulator it gives exactly
Theorem~\ref{p12:thm:v199-rgas-eaa}.  For an exact finite elimination,
Proposition~\ref{p3f:prop:v2167-complete-neutral-return} gives
\(K_{\rm full}=\zeta I-h_R-\Sigma_R\) and
\(K_{\rm full}^{-1}=[(\zeta I-h)^{-1}]_{RR}\).
Thus defining \(K_{R,0}=\zeta I-h_R\), or the corresponding explicitly
chosen pre-self-energy core, gives the displayed decomposition.  Replacing
\(K_{R,0}\) by the already complete Hessian and subtracting the same
return would change this exact inverse.  Proper vertices are derivatives
of the same complete \(\Gamma_R\), including its source contacts.
The source, quotient, domain and any actual regulator/continuation hypotheses
are retained by these identities.
\end{proof}

\begin{corollary}[No independent physical vertex-number selector]
\label{p7:cor:v261-no-independent-vertex-selector}
Once a claim-complete renormalized action/effective action, actual background,
channel inclusion, and gauge quotient are fixed, the proper vertex operator is
generated by the same owner,
\[
 V_{R,\rm proper}=\delta^3\Gamma_R,
\]
just as the complete quadratic kernel is
\(K_{R,\rm full}=\delta^2\Gamma_R\).  Derivatives of a classical or
pre-self-energy action generate its own vertices; they are not renamed
complete proper vertices before their actual quantum/contact transport.
Therefore G19 does
not contain an additional numerical vertex coefficient selector.  Physical
amplitudes and widths are deterministic functionals after the required
external-state normalization, pole/residue, unstable-pole or massless-IR
certificate, and actual state/channel occurrence are supplied.  This
corollary does not assert numerical decay widths or close massless scattering.
\end{corollary}


% BEGIN FRAMEWORK v2.371: same-action rooted Petersen comparison diagnostic.
\begin{proposition}[A nonfitted rooted geometric comparison test]
\label{p5c:prop:v371-petersen-rooted-action-diagnostic}
Paper~I-A, Theorem~\ref{p1:thm:v371-petersen-rooted-leakage},
constructs a rank-six \emph{mathematical} root map
$C_m=P_{\rm rec}R_mP_{\rm loc}$ with
$\|C_m\|_{\rm HS}^2=69/50$.
\emph{Conditionally}, if the actual action/effective action, physical
source and quotient, complete Schur return, pole/kinetic normalization
and state embedding yield a well-typed block
$D_m(z):\mathcal V_{\rm loc}\to\mathcal V_{\rm rec}$ in those same
inner-product coordinates, define
\[
 g_m(z)={50\over69}\operatorname{Tr}
          (C_m^\dagger D_m(z)),\qquad
 \mathcal E_m(z)=D_m(z)-g_m(z)C_m.
\]
Then $\mathcal E_m=0$ if and only if the actual block lies on the
complex line spanned by this particular geometric operator; moreover
\[
 \|\mathcal E_m\|_{\rm HS}^2
 =\|D_m\|_{\rm HS}^2-{50\over69}
       |\operatorname{Tr}(C_m^\dagger D_m)|^2\ge0.
\]
This is a pass/fail residual \emph{after} the action source is
constructed, not a vertex selector or permission to fit decay widths.
If no licensed embedding into $\mathcal V$ exists, the test is not
applicable; the geometry alone does not supply $D_m$.
\end{proposition}
\begin{proof}
Orthogonally project $D_m$ onto the one-dimensional Hilbert--Schmidt
span of $C_m$, using its exact squared norm $69/50$; Pythagoras gives
the identity. The same-owner physical qualifications are those of
the preceding renormalized-response owner and vertex corollary.
\end{proof}
% END FRAMEWORK v2.371: physical insertion remains conditional.

\begin{theorem}[First ordinary scalar-to-gauge source and RG typing]
\label{p7:thm:v164-scalar-gauge-source}
For the declared breathing/RG spurion
\(\tau=\alpha^{-1/2}\phi_c/\bar M_{\rm Pl}\), let \(g_i(\mu)\) denote the
renormalized running coupling in the same scheme as the one-loop coefficient
\(b_i\).  The canonically normalized gauge action gives
\[
 \Delta\mathcal L_{\phi FF}^{(1)}(\mu)
 =-\frac{\alpha^{-1/2}\phi_c}{\bar M_{\rm Pl}}
 \sum_i\frac{b_i g_i^2(\mu)}{32\pi^2}F_{i,c}^{,2}
 +O(g_i^4(\mu)\tau).
\]
Equivalently, in the convention
\(-\lambda_i(\mu)\phi_c F_{i,c}^{,2}/4\),
\[
 \boxed{\lambda_i(\mu)=
 \frac{\alpha^{-1/2}}{\bar M_{\rm Pl}}
 \frac{b_i g_i^2(\mu)}{8\pi^2}.}
\]
At the source-native common-heat boundary \(\mu_h\), write
\(g_i^{\rm heat}:=g_i(\mu_h)\).  In the one-loop GUT-normalized gauge basis
used here the native boundary values give
\(g_1^{\rm heat\,2}=g_2^{\rm heat\,2}=g_3^{\rm heat\,2}=2\pi^2\), and with
\(b=(41/10,-19/6,-7)\) the three weights reduce to
\[
 \boxed{\left(\frac{41}{40},-\frac{19}{24},-\frac7{4}\right).}
\]
These are formal one-loop coefficients evaluated at the given common-heat inputs; the large native boundary coupling does not make a truncated loop expansion nonperturbatively exact.  They are all nonzero and require no thermal occupation factor.  This proves a
zero-temperature gauge-pair vertex on the stated branch.  The native boundary
values are not automatically the physical couplings at the K16 hard or plasma
scale: evaluating \(\lambda_i(\mu)\), a production rate, or a plasma collision
kernel at such a scale consumes the separately typed G16 RG/threshold matching.
\end{theorem}

\begin{remark}[Geometric Weyl coordinate and the distinct breathing source]
\label{p7:rem:v2216-weyl-source-map}
The substitution in Theorem~\ref{p7:thm:v164-scalar-gauge-source} is a
source-coordinate statement.  On the actual matched Legendre chart,
$g_E=Fg_J$ and $\phi_c=M\sqrt{3/2}\log F$, with
$M=\bar M_{\rm Pl}$, the logarithmic metric Weyl coordinate is
\[
 \sigma=\tfrac12\log F=\frac{\phi_c}{\sqrt6M},
 \qquad \tau=\sqrt6\alpha^{-1/2}\sigma.
\]
Thus a running-scale spurion induced solely by this metric transformation
is $\mu\mapsto\mu e^{-\sigma}$, up to the sign of the inverse frame
comparison.  In the same one-loop convention its gauge coefficient is
\[
 \lambda_i^{\rm geom}M=\frac{b_i g_i^2}{8\pi^2\sqrt6},\qquad
 \frac{\lambda_i^\tau}{\lambda_i^{\rm geom}}
       =\sqrt6\alpha^{-1/2}\simeq6.17810\,10^{-9}.
\]
The distinct source $\tau$ is retained as defined.  Identifying its
derivatives with the physical scalaron vertices requires the actual
source/action pullback, including all its jets.  Canonical normalization
and the metric Weyl change alone do not derive its hierarchy factor.
The nonzero vertex on the declared $\tau$ branch remains a source theorem;
its physical production use consumes this incidence and the actual RG,
state and time history.  No multiplicative efficiency is introduced.
\end{remark}

\begin{theorem}[Exact gauge production bound from total source variation]
\label{p7:thm:v2216-gauge-total-variation}
Let a retained transverse coexact $S^3$ mode have the actual quadratic
action
\[
 S_n=\frac12\int d\eta\,Z(\eta)
       \bigl[(A_n')^2-n^2A_n^2\bigr],\qquad n\ge2,
\]
where $Z>0$ is absolutely continuous and
$\mathcal V=\tfrac12\operatorname{TV}_{[\eta_i,\eta_f]}\log Z<\infty$.
For the canonical endpoint annihilators the Bogoliubov coefficient obeys
\[
 \boxed{N_n^{\rm vacuum}=|\beta_n|^2\le\sinh^2\mathcal V.}
\]
No WKB, small rate, thermal state or decay-width hypothesis is needed.
For any inherited state with finite
$N_i=\langle b_i^\dagger b_i\rangle$ and
$m_i=\langle b_i b_i\rangle$, including a non-Gaussian state,
\[
 |N_f-N_i|\le(2N_i+1)\sinh^2\mathcal V
 +2\sqrt{N_i(N_i+1)}\sinh\mathcal V\cosh\mathcal V.
\]
For a finite collection with positive fixed comparison weights $w_j$,
put $E=\sum_jw_jN_{i,j}$ and $W=\sum_jw_j$.
If $\mathcal V_j\le\mathcal V_*$, then
\[
 \Delta E\le(2E+W)\sinh^2\mathcal V_*
 +2\sqrt{E(E+W)}\sinh\mathcal V_*\cosh\mathcal V_*.
\]
These weights compare excitation energies on the declared endpoint
basis; any change of physical endpoint frequencies or volume is retained
separately when forming a work or density budget.
\end{theorem}
\begin{proof}
Set $\Pi_n=ZA_n'$ and
$b_n=\sqrt{nZ/2}A_n+i\Pi_n/\sqrt{2nZ}$.  Hamilton's equations give
\[
 b_n'=-inb_n+\frac{Z'}{2Z}b_n^\dagger.
\]
Writing $b_f=\alpha_nb_i+\beta_nb_i^\dagger$ gives
$|\alpha_n|^2-|\beta_n|^2=1$ and
$d\operatorname{arsinh}|\beta_n|/d\eta\le|Z'|/(2Z)$
almost everywhere.  Integrating from $\beta_n(\eta_i)=0$ proves the
vacuum bound, extending through its zeros by continuity.  The exact
occupation identity is
\[
 N_f-N_i=(2N_i+1)|\beta_n|^2
              +2\operatorname{Re}(\beta_n^*\alpha_n m_i).
\]
State positivity gives $|m_i|^2\le N_i(N_i+1)$, proving the second
bound.  Weighted Cauchy--Schwarz gives
$\sum_jw_j\sqrt{N_{i,j}(N_{i,j}+1)}\le\sqrt{E(E+W)}$;
summing the single-mode bounds proves the last claim.  The proof uses
only quadratic expectations and the canonical map.
\end{proof}

% Framework v2.217 helper integration.
\begin{theorem}[Exact gauge force, coherence and weighted work]
\label{p7:thm:v2217-gauge-force-work}
On the same actual quadratic transverse branch, in proper time, let
\[
 H_n=\frac{P^2/Z+Zn^2Q^2}{2a},\quad [Q,P]=i\hbar,
 \quad n\ge2,\quad Z=Z(\phi)>0.
\]
Here $a,\phi$ are the supplied regular history and $Z$ has no additional
independent time dependence. Let $\omega=n/a$, $k=\dot{\log Z}$,
$L=\partial_\phi\log Z$ and
$b=\sqrt{nZ/(2\hbar)}Q+iP/\sqrt{2\hbar nZ}$.
For any inherited state with finite second moments write
$N=\langle b^\dagger b\rangle$, $M=\langle bb\rangle$ and
$E_n=\langle H_n\rangle$. Exact evolution gives
\begin{align*}
 \dot N&=k\operatorname{Re}M,&
 \dot M&=-2i\omega M+\tfrac{k}{2}(2N+1),\\
 F_{\phi,n}:=\langle\partial_\phi H_n\rangle
   &=\hbar\omega L\operatorname{Re}M,&
 \dot E_n&=-H E_n+F_{\phi,n}\dot\phi,\quad H=\dot a/a.
\end{align*}
One has $|M|^2\le N(N+1)$ and
\[
 \int_{t_i}^{t_f}aF_{\phi,n}\dot\phi\,dt
       =\hbar n(N_f-N_i).
\]
For a supplied initial ground $N_i=0$ this weighted receiver work is
nonnegative. The force on the scalar in the same action is $-F_{\phi,n}$.
\end{theorem}
\begin{proof}
Hamilton's equations and the derivative of the instantaneous $b$ give
$\dot b=-i\omega b+(k/2)b^\dagger$. Differentiating its two quadratic
expectations gives the first line; no factorization is used.
Cauchy--Schwarz applied to $b^\dagger$ and $b$ in the state gives
$|\langle bb\rangle|^2\le\langle bb^\dagger\rangle
\langle b^\dagger b\rangle$. At fixed $Q,P,a$,
$\partial_\phi H_n=L(Zn^2Q^2-P^2/Z)/(2a)$, whose expectation is the
stated force. Since $E_n=\hbar(n/a)(N+1/2)$, differentiation gives the
energy equation and multiplication by $a$ integrates to the work identity.
Variation of the same gauge action with respect to the dynamical scalar
gives the opposite force.
\end{proof}
The classical local gauge metric trace at fixed $\phi$ is zero in four
dimensions, while its scalar variation is generally nonzero. Thus on the
isotropic quadratic radiation branch $p=\rho/3$ does not erase this force.
The complete source additionally retains central, RG/threshold, metric and
preparation contacts. If $Z$ has other time-dependent controls, their work
must be added to the last identity. A nonempty inherited state can give
either sign of $N_f-N_i$; this is not a universal friction or stopping law.


\begin{remark}[K16 common-heat / plasma-coupling firewall]
\label{p7:rem:v262-k16-commonheat-plasma-firewall}
The symbol \(g_i^{\rm heat}\) denotes the source-native boundary datum only.
The physical couplings entering K16 production and subsequent gauge-plasma
kinetics are \(g_i(\mu)\) at specified hard/soft scales.  In particular,
\(\alpha_s^{\rm heat}=\pi/2\) is not a weak-coupling certificate and cannot be
inserted into a weakly-coupled effective kinetic theory without an RG/matching
map.  This retyping changes no source theorem and introduces no new transfer
constant.
\end{remark}

% Framework v2.217 helper integration.
\begin{corollary}[Pair-target nullity requires the actual equivariant incidence]
\label{p7:cor:v2217-pair-target-incidence}
For a symmetric gauge-singlet two-transverse-boson target on complete round
coexact shells, the bivector and spinor target multiplicities are zero by
Paper~I-C, Theorem~\ref{p1b:thm:v2217-complete-shell-pair-targets}.
If the actual reciprocal transport identifies the consumed intrinsic and
spatial $\operatorname{Spin}(4)$ actions and its source vertex is an
intertwiner, that vertex vanishes on the two bivector and two spinor source
summands. On the source $(0,0)\oplus(1/2,0)\oplus(0,1/2)\oplus
(1/2,1/2)\oplus(1,0)\oplus(0,1)$, its rank is at most $1+4=5$.
A lowest fermion--antifermion target instead
permits both bivectors, without establishing their vertex values.
\end{corollary}
\begin{proof}
Schur's lemma kills an intertwiner from a missing irreducible type. Its
remaining rank is bounded by the dimension of the retained source scalar
and vector summands. The mixed spinor target decomposition is the second
formula of the cited theorem.
\end{proof}
On a source with the framework's single scalar $e_0$ and one vector copy,
the bound is $1+4=5$; no second scalar source is created by the target's
two helicity singlets. The parity-even $\phi F^2$ pump exchanges the two
helicities symmetrically. For a parity-balanced inherited quadratic state
their equations and initial data coincide, so their occupations stay equal.
The required actual joint equivariance is not inferred from matching
representation names. Other statistics, gauge tensors or chiral interactions
retain their own targets.






\section{Strong-sector response, CP phase and chiral measure}
\label{p5c:sec:organized-qcd}
\subsection{Structural QCD interface}
\label{p7v2:sec:qcd-structure}

\begin{theorem}[Color representation dictionary]
\label{p7v2:thm:color-dictionary}
On \(\Urole\), import the Paper--V-A color carrier
\(C_{\rm col}\simeq\CC^3\) with the fundamental \(SU(3)\) action and choose
generators normalized by
\(\Tr(T^aT^b)=\tfrac12\delta^{ab}\).  Then
\begin{equation}
\label{p7v2:eq:color-decomposition}
 C_{\rm col}\otimes C_{\rm col}^{*}
 \simeq\mathbf1\oplus\mathbf8,
\end{equation}
and the group-theoretic constants are
\begin{equation}
\label{p7v2:eq:qcd-color-factors}
 C_A=3,\qquad
 C_F=\frac43,\qquad
 T_R=\frac12.
\end{equation}
For a connection \(A_\mu=A_\mu^aT^a\), the curvature is
\begin{equation}
\label{p7v2:eq:qcd-curvature}
 F_{\mu\nu}
 =
 \partial_\mu A_\nu-\partial_\nu A_\mu
 -ig_s[A_\mu,A_\nu].
\end{equation}
These are exact representation-theoretic statements.  They do not fix the
normalization or scale dependence of \(g_s\).
\end{theorem}

\begin{proof}
The decomposition
\(\mathbf3\otimes\overline{\mathbf3}=\mathbf1\oplus\mathbf8\) is the trace plus
traceless decomposition of \(\End(\CC^3)\).  The chosen generator
normalization gives \(T_R=1/2\).  Standard \(SU(N)\) identities give
\(C_A=N\) and \(C_F=(N^2-1)/(2N)\), hence
\eqref{p7v2:eq:qcd-color-factors} at \(N=3\).  Equation
\eqref{p7v2:eq:qcd-curvature} is the curvature of the imported connection.
\end{proof}

\begin{proposition}[Short-distance color channels]
\label{p7v2:prop:color-channels}
For a representation channel
\(R\subset R_1\otimes R_2\),
\begin{equation}
\label{p7v2:eq:two-body-color-factor}
 T_1\!\cdot T_2
 =
 \frac12\left(C_R-C_{R_1}-C_{R_2}\right).
\end{equation}
Consequently,
\begin{align}
\label{p7v2:eq:meson-color-channels}
 \mathbf3\otimes\overline{\mathbf3}
 &=\mathbf1\oplus\mathbf8,
 &
 T_q\!\cdot T_{\bar q}\big|_{\mathbf1}
 &=-\frac43,
 &
 T_q\!\cdot T_{\bar q}\big|_{\mathbf8}
 &=\frac16,\\
\label{p7v2:eq:diquark-color-channels}
 \mathbf3\otimes\mathbf3
 &=\overline{\mathbf3}\oplus\mathbf6,
 &
 T_q\!\cdot T_q\big|_{\overline{\mathbf3}}
 &=-\frac23,
 &
 T_q\!\cdot T_q\big|_{\mathbf6}
 &=\frac13.
\end{align}
These are exact color factors.  Their interpretation as attractive or
repulsive static potentials additionally assumes the leading one-gluon
exchange QFT.
\end{proposition}

\begin{proof}
Expand
\((T_1+T_2)^2=C_R\) on the irreducible channel \(R\) to obtain
\eqref{p7v2:eq:two-body-color-factor}.  Insert
\[
 C_{\mathbf3}=C_{\overline{\mathbf3}}=\frac43,\quad
 C_{\mathbf1}=0,\quad C_{\mathbf8}=3,\quad C_{\mathbf6}=\frac{10}{3}.
\]
The potential interpretation uses the additional perturbative relation
\(V_R(r)\propto\alpha_s(T_1\!\cdot T_2)/r\).
\end{proof}

\begin{proposition}[Perturbative QCD evolution interface]
\label{p7v2:prop:qcd-running}
Assume \(\Hqft\).  Between quark thresholds the one-loop strong coupling obeys
\begin{equation}
\label{p7v2:eq:qcd-one-loop}
 \frac{d\alpha_s}{d\log\mu}
 =
 -\frac{\beta_0}{2\pi}\alpha_s^2+O(\alpha_s^3),
 \qquad
 \beta_0=11-\frac23n_f.
\end{equation}
In the conventional normalization the second universal coefficient is
\begin{equation}
\label{p7v2:eq:qcd-two-loop-coefficient}
 \beta_1=102-\frac{38}{3}n_f.
\end{equation}
Threshold matching must be performed at specified quark masses in a specified
scheme.  If the ansatz
\(\alpha_s(M_Z^{\rm BH})=\delta\) is used, the resulting flow is a
status-\(\mathsf B\) test of that ansatz.
\end{proposition}

\begin{proof}
Substitution of \eqref{p7v2:eq:qcd-color-factors} into the standard Yang--Mills
one-loop coefficient yields
\(\beta_0=(11/3)C_A-(4/3)T_Rn_f=11-2n_f/3\)
\cite{GrossWilczek1973,Politzer1973}.  The starting value is not supplied by
this calculation and has the status established in
Proposition~\ref{p7v2:prop:strong-seed-status}.
\end{proof}

\begin{proposition}[Inclusive electromagnetic quark ratio]
\label{p7v2:prop:inclusive-rgamma}
Assume \(\Hqft\), a scale away from narrow resonances and physical
thresholds, with active quark masses negligible on the selected partonic
branch. Project both external currents to \emph{photon exchange only};
the photon-current cross sections are denoted \(\sigma_\gamma\).
Then the electromagnetic parton ratio is
\begin{equation}
\label{p7v2:eq:inclusive-rgamma}
 \boxed{R_\gamma(s)
 :=\frac{\sigma_\gamma(e^+e^-\to{\rm hadrons})}
          {\sigma_\gamma(e^+e^-\to\mu^+\mu^-)}
 =3\sum_{q=1}^{n_f}Q_q^2
 \left[1+\frac{\alpha_s(s)}{\pi}+O(\alpha_s^2)\right].}
\end{equation}
For the \emph{full} massless Born neutral-current observable, let
\(g_\lambda^f\) be the action-owned chiral \(Z\)-current coefficient
for external fermion \(f\) and chirality \(\lambda=L,R\), and let
\(\zeta_Z(s)\) include the relative neutral-current coupling and
propagator normalization with respect to photon exchange. Define
\[
 \mathcal A_{\lambda\lambda'}^f(s)
 :=Q_eQ_f+\zeta_Z(s)g_\lambda^eg_{\lambda'}^f,
 \qquad
 F_f(s):=\frac14\sum_{\lambda,\lambda'=L,R}
            |\mathcal A_{\lambda\lambda'}^f(s)|^2.
\]
The corresponding unpolarized tree-level inclusive ratio on the same
massless, open-channel domain is the \emph{different} observable
\begin{equation}
\label{p7v2:eq:inclusive-neutral-current-born}
 \boxed{R_{\gamma+Z}^{(0)}(s)
 =3\frac{\sum_{q=1}^{n_f}F_q(s)}{F_\mu(s)}.}
\end{equation}
The exact native neutral-current coefficients and physical pole/width
input must be evaluated to predict this full ratio. Being merely away
from the \(Z\) pole does not eliminate its interference or exchange.
The \emph{photon-only} Born values for \(n_f=3,4,5,6\) are respectively
\begin{equation}
\label{p7v2:eq:inclusive-rgamma-born}
 2,\qquad \frac{10}{3},\qquad \frac{11}{3},\qquad5.
\end{equation}
Numerical corrected values require a separately normalized
\(\alpha_s(s)\).
\end{proposition}

\begin{proof}
The photon-current part follows from the color multiplicity \(3\) and
weights \(Q_q^2\). Summing the quark charges gives
\eqref{p7v2:eq:inclusive-rgamma-born}; the first perturbative QCD
correction is the displayed coefficient on that restricted current.
At massless tree level, the unpolarized helicity sum of each neutral-current
amplitude gives \(F_f\), with the common kinematic prefactor canceling
between hadron and muon final states; color multiplies each quark channel
by three. This proves Eq.~\eqref{p7v2:eq:inclusive-neutral-current-born}.
At \(\zeta_Z=0\), \(F_q=Q_q^2\) and \(F_\mu=1\), reproducing the
photon formula. With nonzero \(Z\)-current couplings the interference
terms and squared \(Z\) terms do not generally cancel, even off pole.
\end{proof}

\begin{remark}[A perturbative handoff is not confinement]
Solving the truncated flow for the scale at which
\(\alpha_s(\mu)=1\) defines a conventional perturbative handoff.  It does not
prove a mass gap, confinement, chiral symmetry breaking, or hadronization:
the truncation is least reliable precisely where the coupling becomes order
one.
\end{remark}

\begin{proposition}[Singlet support does not determine parton measures]
\label{p7v2:prop:singlet-not-pdf}
Color-singlet projection constrains the color tensor of a hadronic state but
does not determine a parton distribution, distribution amplitude, or
fragmentation function.
\end{proposition}

\begin{proof}
Let \(\chi_{\rm sing}\) be any fixed normalized color-singlet tensor and let
\(\psi(x_1,\ldots,x_n)\) be any normalized longitudinal wavefunction on the
parton simplex.  Every state
\[
 \Psi(x_1,\ldots,x_n)
 =
 \chi_{\rm sing}\,\psi(x_1,\ldots,x_n)
\]
has identical color-singlet support, while \(|\psi|^2\) can be varied
arbitrarily subject to normalization and the momentum sum constraint.
Therefore color singletness alone cannot select an equal-momentum delta
measure, a triangular density, or a maximum-entropy distribution.
\end{proof}

\begin{definition}[Admissible QCD boundary ansatz]
\label{p7v2:def:qcd-boundary-ansatz}
Any equal-momentum valence measure, triangular shell-smearing kernel, or
primitive fragmentation endpoint retained for downstream testing is denoted
\(\mathcal H_{\rm QCD}^{\rm ans}\) and has status \(\mathsf A\).  The genuine
nonperturbative boundary
\(\mathcal H_{\rm QCD}(\mu_{\rm int})\) remains a Paper-VII-A--VII-B output.
\end{definition}

\begin{proposition}[Spacelike, timelike, and ERBL evolution handoff]
\label{p7v2:prop:dglap-erbl}
Given \(\Hqft\), a factorization scheme, and the appropriate nonperturbative
spacelike parton, timelike fragmentation, or exclusive distribution-amplitude
boundary at \(\mu_{\rm int}\), the leading-twist spacelike DGLAP, timelike
DGLAP, and ERBL equations evolve that boundary with kernels whose color
factors are fixed by \eqref{p7v2:eq:qcd-color-factors}.  Paper~V-C exports the
representation/color dictionary and perturbative evolution generators, not a
solved hadronic boundary state.
\end{proposition}

\begin{proof}
The leading spacelike, timelike, and exclusive kernels are the standard
factorization-scale evolution generators
\cite{GribovLipatov1972,Dokshitzer1977,AltarelliParisi1977}.  Their numerical
color coefficients follow from
Theorem~\ref{p7v2:thm:color-dictionary}; their initial distributions do not.
\end{proof}

\subsection{Native strong-action character and the strong-CP invariant}
\label{p7v2:sec:strong-cp}

\begin{theorem}[Native strong action has no free topological sector character]
\label{p7v2:thm:bare-theta}
The generated finite strong action and normalized Haar reference measure define
one real positive native sector weight.  Multiplying an index-\(\nu\) sector
by \(e^{i\theta\nu}\) changes the action by
\[
 \Delta S_\theta=-i\theta\nu.
\]
Thus a nontrivial topological character is an additional orientation-odd
action term, not a normalization convention.  Since the current strong-action
constructor generates no such independent port, GP2/no-insertion gives
\begin{equation}
\label{p7v2:eq:bare-theta-zero}
 \boxed{\chi_{\rm top}(\nu)=1,\qquad \theta_{\rm nat}=0.}
\end{equation}
This is stronger than merely setting to zero the coefficient of an operator
absent from a chosen basis: there is no free native \(\theta\)-member to
select on the current action owner.
\end{theorem}
\begin{proof}
The native strong weight is fixed by the generated action.  A sector-dependent
factor \(e^{i\theta\nu}\) is equivalently absorbed into the exponential as
\(S\mapsto S-i\theta\nu\).  Such a term changes relative weights of disconnected
topological components and therefore defines a different action.  No licensed
constructor in the current color owner generates this orientation-odd term.
\end{proof}

\begin{proposition}[Rooted quark determinant orientation]
\label{p7v2:prop:physical-strong-cp}
On the nonzero full-rank quark-source branch, Paper~V-A's rooted orientation
covariance and Paper~V-B's positive-polar separation give
\[
 \boxed{\frac{\det(D_uD_d)}{|\det(D_uD_d)|}=1.}
\]
No anomalous after-the-fact chiral rephasing is used: positivity is a property
of the generated rooted source pair before continuum field redefinitions.
\end{proposition}
\begin{proof}
The two rooted orientation cores are complex conjugates, while every radial
polar factor has positive determinant.  Hence
\(\det(D_uD_d)=(\prod\det P)|\det A_+|^2>0\).
\end{proof}

\begin{corollary}[Native strong-CP invariant]
\label{p7v2:cor:v259-native-strong-cp}
Using the conventional export coordinate
\[
 \zeta_{\rm CP}:=e^{i\theta}\frac{\det(D_uD_d)}{|\det(D_uD_d)|}
 =e^{i\bar\theta},
\]
the native finite theory gives
\[
 \boxed{\zeta_{\rm CP}=1,\qquad \bar\theta=0\pmod{2\pi}.}
\]
Weak-family CP is not removed: on the active direct-mass branch its native
readout is carried by the generated weak-current incidence/current loops.
Relative-eigenframe invariants may be used only as conventional coordinate
dictionaries where their hypotheses are separately established.
\end{corollary}

\begin{remark}[No Standard-Model selector is imported]
The result does not solve a problem by first postulating an arbitrary strong
\(\theta\) and then selecting its value.  It states instead that the present
finite action grammar does not generate such a free coordinate, while the
same rooted source grammar fixes the quark determinant orientation.  A future
new orientation-odd color action constructor would constitute a different
branch and would require re-audit.
\end{remark}


\subsection{Strong-CP export dictionary, continuum chiral line, and the kinetic-state boundary}
\label{p7:sec:cp-kinetic-update}

On the nonzero-determinant branch retain the conventional export coordinate
\begin{equation}
 \label{p7:eq:zeta-cp}
 \zeta_{\rm CP}
 :=e^{i\theta}\frac{\det(D_uD_d)}{|\det(D_uD_d)|}
 =e^{i\bar\theta}.
\end{equation}
For the native action/source branch proved above, \(\zeta_{\rm CP}=1\).
The formula remains useful when comparing with continuum QCD conventions; it
does not create a free native action member.

\begin{proposition}[Topological sector phase is an action descendant, not a G20 selector]
\label{p7:prop:v190-index-determinant-phase}
On any gapped finite overlap component with integer charge \(\nu\), a
hypothetical action character and quark determinant orientation would combine
as \(\zeta_{\rm CP}^{\nu}\).  On the current native branch both factors are
unity, so the physical finite sector phase is
\[
 \boxed{\zeta_{\rm CP}^{\nu}=1.}
\]
Sector probabilities and topological susceptibility remain dynamical outputs
of the strong transfer/action; they are not CP-member selectors.
\end{proposition}

\begin{proposition}[Global chiral determinant line is a continuum/reconstruction criterion]
\label{p7:prop:global-chiral-measure-trivialization-test}
If a continuum or Weyl presentation introduces a gauge-field-dependent
chiral subspace and determinant line \(L_{\det}\), then flatness and trivial
holonomy imply a global parallel measure ray by Paper~IV's
parallel-line theorem.  Failure of those hypotheses is a genuine anomaly or
reconstruction obstruction.

This obligation is attached to the chiral-continuum/regulator certificate
(G39 and ordered chiral/topological limits where consumed).  It is not a
prerequisite for the fixed finite paired vectorlike color measure, whose
Berezin Jacobians cancel on the native carrier.
\end{proposition}

\begin{proposition}[Primitive-sector positivity is not a strong-CP selector]
\label{p7:prop:cp-fixed-positive-sector}
Paper~VII-A supplies open \(\nu=\pm1\) overlap neighborhoods with positive
Haar measure and positive native heat weight.  Hence the native magnitude
measure has populated primitive charge sectors.  Their existence establishes
nontrivial topology and susceptibility capacity; because the native sector
character is already trivial, no additional positive-sector argument is
needed to select \(\bar\theta\).
\end{proposition}

Thus the former G20 wrapper is retired.  Native finite topology, paired
measure and strong determinant orientation are derived; unrestricted
clutching/overlap comparison and global chiral determinant-line integrability
are continuum certification obligations, while weak-family CP remains a G15
descendant.


\begin{corollary}[Finite equilibrium density matrix is not an independent KIN selector]
\label{p7:cor:finite-equilibrium-state-descendant}
On a simple finite physical algebra with supplied Hamiltonian $H$ and inverse
temperature $\beta$, Paper~IV
Theorem~\ref{p4:thm:finite-simple-kms-uniqueness} gives the unique normalized
finite equilibrium state
\[
 \boxed{\rho_\beta=\frac{e^{-\beta H}}{\Tr e^{-\beta H}}.}
\]
Thus G22 does not require an additional equilibrium density-matrix member once
$H$ and $\beta$ are owned.  The result does not generate $\beta$, a
nonequilibrium initial density, collision/washout kernels, spectator maps, or
the cosmological kinetic history.
\end{corollary}

For the kinetic branch, the actual initial density, collision/washout kernels,
spectator map and cosmological history therefore remain process/boundary data,
and the leptogenesis efficiency remains an output functional of those inputs.



\section{Scalar and Higgs branches of the common action}
\label{p5c:sec:organized-scalar}
\subsection{Claim-scoped scalar stationary/effective-action branch}
\label{p7v2:sec:ew-vacuum}

\begin{definition}[Claim-scoped scalar stationary/effective-action family]
\label{p7v2:def:vacuum-selector-family}
Native electroweak reduction is already fixed by the occurring chiral solder.
For claims that additionally require scalar equilibrium, global energetic
preference, equilibrium Hessians, or relaxation/history, retain the
claim-scoped family \(\mathfrak A_{\rm vac}\) of same-owner effective actions
and stationary/history continuations compatible with the complete consumed
word.  Gauge compatibility alone does not license a competing native
potential.  The \(V_\pm\) carrier comparison is confined to
Theorem~\ref{p7v2:thm:type-separation}; the following renormalizable-doublet
calculation is an explicitly supplied comparison target, not a prerequisite
for the native stabilizer or gauge Gram.
\end{definition}

\begin{theorem}[Potential conditional on a supplied broken orbit]
\label{p7v2:thm:conditional-potential}
Fix a compatible \(b_6\in\Acompsix\) whose scalar-normalization coordinate
supplies \(v>0\), and assume \(\Urole+\Hvac\) supplies a physical single
Higgs doublet with a
renormalizable gauge-invariant scalar potential
\begin{equation}
\label{p7v2:eq:higgs-potential}
 V(H)=\mu_H^2H^\dagger H+\lambda_H(H^\dagger H)^2,
 \qquad \lambda_H>0,
\end{equation}
and a supplied nonzero stationary orbit
\(H^\dagger H=v^2/2\) belonging to the chosen \(\mathfrak A_{\rm vac}\)
member.  Then
\begin{equation}
\label{p7v2:eq:conditional-mu}
 \mu_H^2=-\lambda_Hv^2,
 \qquad
 m_h^2=2\lambda_Hv^2
\end{equation}
at tree level.
\end{theorem}

\begin{proof}
Write \(r=H^\dagger H\).  Stationarity at \(r=v^2/2\) gives
\(0=dV/dr=\mu_H^2+2\lambda_Hr\), hence
\(\mu_H^2=-\lambda_Hv^2\).  Expanding
\(H=(0,(v+h)/\sqrt2)^T\) then gives
\(\partial_h^2V|_{h=0}=2\lambda_Hv^2\).
\end{proof}

\begin{remark}[Scalar-stationary diagnostics are not electroweak-constitution premises]
\label{p7v2:rem:v258-scalar-stationary-diagnostics}
The carrier-symmetry counterfamily and the exact positive-heat no-go are
retained in Appendix~\ref{p7:app:exact-heat-vacuum-diagnostic} as diagnostics
for the claim-scoped scalar-equilibrium problem.  They do not obstruct the
already-derived native electroweak solder/stabilizer of
Theorem~\ref{p7v2:thm:v258-native-ew-solder-reduction}.
\end{remark}

\begin{theorem}[Sufficient convex criterion for a unique stable broken orbit]
\label{p7:thm:sufficient-stable-broken-vacuum}
Let a supplied gauge-invariant effective scalar action reduce on constant
orbits to \(V\in C^2([0,\infty))\) with \(r=H^\dagger H\), and assume
\begin{equation}
 \label{p7:eq:stable-broken-vacuum-criterion}
 V'(0)<0,
 \qquad V''(r)>0\ \text{for all }r\ge0,
 \qquad \lim_{r\to\infty}V'(r)>0.
\end{equation}
Then there is a unique \(r_*>0\) with \(V'(r_*)=0\).  It is the unique
global minimum on \([0,\infty)\), is radially stable, and determines the
unique gauge-orbit radius
\begin{equation}
 \label{p7:eq:stable-broken-vacuum-radius}
 \boxed{v_*=\sqrt{2r_*}.}
\end{equation}
Thus \eqref{p7:eq:stable-broken-vacuum-criterion} is a sufficient discharge
criterion for the radial existence, uniqueness, stability, and global
preference parts of \(\Hvac\).  It does not derive the effective action or
its occurrence from the finite carrier.
\end{theorem}

\begin{proof}
Strict positivity of \(V''\) makes \(V'\) strictly increasing.  It starts
negative and has positive limiting value, so continuity gives exactly one
zero \(r_*>0\).  The function decreases before \(r_*\) and increases after
it, proving global minimality; \(V''(r_*)>0\) gives radial stability.  The
standard doublet normalization \(r_*=v_*^2/2\) gives the displayed radius.
\end{proof}

\begin{proposition}[Bare terminal scalar obstruction and actual contact target]
\label{p7:prop:v2211-bare-terminal-scalar-contact-target}
Consider only the displayed bare matched gravity/scalar constituent of
Paper~III-F, on a separately licensed affine terminal restriction \(F=1\),
with no residual scalar force, occupied gauge source or additional
effective-action term.  Assume an admitted Hermitian radial direction
\(\Phi=rT\), \(\operatorname{Tr}T^2=1\), and the positive-gravity domain
\[
 x=\frac{\alpha_\Phi r^2}{3M^2}\in[0,1),\quad
 \Omega=1-x,\quad c_4=\operatorname{Tr}T^4\in(0,1].
\]
Its Einstein radial potential is
\[
 W_0(x)=\frac{M^2C/2-6M^2\Lambda_g^2x
              +9c_4M^4x^2/\alpha_\Phi}{(1-x)^2}.
\]
At the stated bare native coefficients
\(\alpha_\Phi=1/(16\pi^2)\), \(M^2=U/4\),
\(\Lambda_g^2=96\pi^2U\), and \(0<C\le U\), put
\[
 A=M^2C-2304\pi^2M^4<0,
 \qquad B=288\pi^2M^4(c_4-8)<0.
\]
Then
\[
 \boxed{W_0'(x)=\frac{A+Bx}{(1-x)^3}<0,\qquad
 W_0''(x)=\frac{3A+B+2Bx}{(1-x)^4}<0.}
\]
There is no nonzero stationary minimum of this constituent on the
admitted radial domain; its admitted origin direction is tachyonic.
Changing the homogeneous Yukawa family within the same no-offset word
cannot repair this sign: Corollary~\ref{p5:cor:v2211-yukawa-ray-moment}
strengthens its bound to \(1/48\le c_4\le1/2\).

If the actual complete effective action admits on this same physical
radial chart a potential \(W=W_0+\Delta W\), a nonzero minimum with
positive radial Hessian at \(x_*\in(0,1)\) necessarily satisfies
\[
 \boxed{\Delta W'(x_*)=-\frac{A+Bx_*}{(1-x_*)^3}>0,\qquad
 \Delta W''(x_*)>-\frac{3A+B+2Bx_*}{(1-x_*)^4}>0.}
\]
The required \(\Delta W\) and its derivatives must be generated by the
actual source/contact, metric or quantum effective-action transport.
These equations do not insert such a term or a minimum.  A changed
physical action projection, different background, nonaffine retained
scalaron, constant source offset or additional occupied field consumes
its own complete stationary/Hessian calculation.  This bare obstruction
is not an exclusion of that completed interacting theory and does not
reopen the occurring solder or its gauge-Gram rank.
\end{proposition}
\begin{proof}
The matched constituent has Jordan gravitational factor
\(M^2\Omega\), terminal constant potential \(M^2C/2\), scalar
quadratic \(-2\alpha_\Phi\Lambda_g^2r^2\) and quartic
\(\alpha_\Phi c_4r^4\).  Its constant-field Weyl transform divides
the potential by \(\Omega^2\).  Substitute the definition of \(x\)
to obtain \(W_0\).  Differentiation gives
\[
 (1-x)^3W_0'=M^2C-6M^2\Lambda_g^2
 +(18c_4M^4/\alpha_\Phi-6M^2\Lambda_g^2)x=A+Bx.
\]
Since \(C\le U=4M^2\), \(A<0\), while \(c_4\le1\) gives
\(B<0\).  One further derivative proves strict concavity.  At
\(r=0\), \(d^2W_0/dr^2=(2\alpha_\Phi/(3M^2))W_0'(0)<0\).
At \(r>0\), radial stationarity is equivalent to \(W'=0\), and
positive radial Hessian there requires \(W''>0\), since
\(d^2W/dr^2=(dx/dr)^2W''\) at a stationary point.
Apply these identities to \(W_0+\Delta W\).  Stability in any other
physical directions remains a separate necessary part of the complete
claim.  The preceding abstract convex criterion is not applied outside
its supplied action or its own domain \([0,\infty)\).
\end{proof}

\begin{corollary}[Structural stable-vacuum-orbit closure]
\label{p7v2:cor:stable-vacuum-orbit-closure}
Let a generated complete effective action have physically admissible stable
stationary set $\mathcal V$ and generated gauge/symmetry equivalence $G$.  If
\[
 \mathcal V/G=\{[\Phi_*]\},
\]
then Paper~II-A
Theorem~\ref{p2a:thm:compatible-family-observable-closure} removes any further
structural vacuum-orbit selector: the class $[\Phi_*]$ is the unique generated
stable descendant.  This does not create occurrence of the vacuum in a claim
whose ontology separately consumes an occurring state/history.  The radial
criterion of Theorem~\ref{p7:thm:sufficient-stable-broken-vacuum} supplies one
sufficient one-dimensional orbit-radius uniqueness result, but full use of
this corollary requires stability/global admissibility in every physical
direction consumed by the claim.
In particular that abstract criterion does not discharge the actual bare
terminal constituent tested in
Proposition~\ref{p7:prop:v2211-bare-terminal-scalar-contact-target};
any claimed repair consumes the complete generated contact derivatives
and physical action map there.
\end{corollary}

\begin{corollary}[Scope of the scalar stationary branch]
\label{p7v2:cor:ewsb-status}
Native electroweak reduction and the gauge-mass kernel are already generated
by the occurring chiral solder.  \(\Hvac\) is needed only by later claims
which explicitly require a physical stationary scalar equilibrium, global
energetic comparison, equilibrium scalar Hessian, or actual scalar history.
Failure to derive such an equilibrium does not reopen
Theorem~\ref{p7v2:thm:v258-native-ew-solder-reduction}.
\end{corollary}

\begin{theorem}[Derived-later charged physical slot embeddings]
\label{p7v2:thm:charged-slot-physical-later}
Assume the physical-role domain \(\Urole\), fix a compatible
\(b_6\in\Acompsix\) containing \(\mathsf{SLOT}\), \(\mathsf{TMF}\), and
\(\mathsf F\), and use the nonzero primitive odd source generated by the same rooted branch of Paper~V-A.  The actual family member in \(b_6\), its projectors and terminal-row incidence remain the assembly antecedent.  Let \(\widehat H\) be the normalized solder direction supplied by Paper~V-A Theorem~\ref{p5:thm:v377-unit-solder-normalized-scalar} and put
\[
 L_H:=\mathbb C\widehat H,
 \qquad
 L_{\widetilde H}:=\mathbb C\widetilde{\widehat H}.
\]
Let \(H_u,H_d,H_e\) be the charged family operators of the frozen
\(\mathsf F\) member.  On a simple-spectrum branch write
\[
 P_{\alpha i}
 :=\mathbf 1_{\{\lambda_{\alpha i}\}}(H_\alpha),
 \qquad \alpha\in\{u,d,e\},
\]
with the lookup-row labels already bound to these eigenvalues by the same
frozen \(\mathsf F\) member.  Let the family projectors act on the Paper--V-A
family factor and write
\[
 u_{R,i}:=P_{ui}u_R,\qquad
 d_{R,i}:=P_{di}d_R,\qquad
 e_{R,i}:=P_{ei}e_R.
\]
Then the Paper--V-A generated Yukawa supports force, up to unmarked
gauge/presentation unitary equivalence, the charged fermion row embeddings
\[
\begin{array}{rcl}
 u_i&:&
 L_{\widetilde H}\otimes C\otimes\operatorname{Ran}P_{ui}
 \longrightarrow u_{R,i},\\[2pt]
 d_i&:&
 L_H\otimes C\otimes\operatorname{Ran}P_{di}
 \longrightarrow d_{R,i},\\[2pt]
 e_i&:&
 L_H\otimes\ell\otimes\operatorname{Ran}P_{ei}
 \longrightarrow e_{R,i}.
\end{array}
\]
Here the displayed notation denotes the corresponding one-family subspaces of
the Paper--V-A matter blocks; the mass scalar acts as the identity on the
already generated color/lepton multiplicity.  Thus the charged part of the
optional Paper--V-B package satisfies
\[
 \boxed{
 \mathsf{SLOT}_{\rm phys}^{\rm charged}
 \text{ is \textsc{derived-later} from }
 \mathsf F+\text{occurring solder/source}
 }
\]
relative to the frozen lookup and terminal-mass data.  If a charged family
operator is degenerate, only its full generated spectral subspace is fixed;
no artificial individual family line is selected.

The theorem does not derive the high-energy neutral lift \(Y_\nu\) or its
Casas--Ibarra coordinate \(R\).  The heavy spectrum \(D_N\) is already part of
the frozen Paper--V-B base and is not reselected here.  The unresolved lift
remains the independent \(\Hnu\) realization fiber.
\end{theorem}

\begin{proof}
Paper~V-A proves that the only coordinate-linear charged Dirac supports are
\(\bar Q_L\widetilde H u_R\), \(\bar Q_LHd_R\), and
\(\bar L_LHe_R\).  Once the nonzero odd source occurs, the normalized solder theorem generates the lines
\(L_{\widetilde H}\) and \(L_H\) covariantly, so substituting an
opposite weak line or a wrong right-handed block violates the already generated
Yukawa/hypercharge/projector relations.  The frozen \(\mathsf F\) operators
supply their spectral projectors by finite functional calculus; for simple
spectra these projectors bind the row eigenvalues to unique family lines.
Quark mass readouts act as scalars on the color multiplicity, so changing the
color basis is only gauge/presentation equivalence and supplies no physical
selector.  Hence no additional charged lookup-to-matter embedding remains once
\(\mathsf F\) and the occurring solder/source branch are frozen.  Degenerate spectral values
produce only the corresponding higher-rank spectral projector, as claimed.
The neutral high-energy data are not consequences of these charged support
relations and retain their separate \(\Hnu\) neutral-lift-fiber provenance.
\end{proof}

\begin{proposition}[Conditional one-loop stationary diagnostics]
\label{p7v2:prop:one-loop-ew-diagnostic}
Assume \(\Hqft\), a specified gauge, subtraction scheme, scale prescription,
and Goldstone resummation.  Then an RG-improved or
Goldstone--tadpole-resummed effective potential can test the consistency of the
imported orbit.  Such a test does not independently derive that orbit when its
couplings, masses, and stationarity condition are built from the imported
broken phase.
\end{proposition}

\begin{proof}
The effective potential is a functional of the chosen action, scheme, gauge,
and background-dependent masses.  Solving its stationary equation gives a
conditional stationary point.  If those masses and couplings were themselves
reconstructed from the imported \(v\), agreement with \(v\) is a consistency
test rather than an independent selection theorem.  Goldstone resummation does
not change this logical direction \cite{Martin2014Goldstone}.
\end{proof}

\begin{remark}[Reproducible auxiliary-shell audit]
Using the previous edition's unrounded inputs and prescriptions gives
\[
 \mu_{\rm stat}^{\rm GT}=221.0882021\,{\rm GeV},
 \qquad
 \mu_{\rm stat}^{\rm RG}=221.5844320\,{\rm GeV},
\]
a difference of \(-0.4962299\,{\rm GeV}\), and the quoted hard-shift ratio
\(1.019031986119\).  These numbers remain a status-\(\mathsf B\)
one-loop consistency certificate.  They are not a proof of vacuum selection.
\end{remark}

\subsection{Mass-boundary ratios and physical Higgs vertices}
\label{p7v2:sec:higgs-vertices}

\begin{definition}[Mass-boundary diagnostics]
\label{p7v2:def:mass-boundary-diagnostics}
For \(V=W,Z\), define
\begin{equation}
\label{p7v2:eq:mass-boundary-ratio}
 D_V
 :=
 \frac{\widehat m_V^{\,2}}{M_{V,\rm pole}^{\,2}},
\end{equation}
where \(\widehat m_V\) is a named finite-shell mass representative and
\(M_{V,\rm pole}\) the corresponding pole mass.  The quantities \(D_V\) are
mass-boundary diagnostics, not Higgs coupling modifiers.
\end{definition}

\begin{proposition}[A pole-mass ratio is not a vertex modifier]
\label{p7v2:prop:mass-ratio-not-kappa}
Without a renormalized three-point calculation,
\[
 D_V\neq\kappa_{hVV}
\]
in general.  A physical \(hVV\) coupling depends on the background derivative
of the renormalized inverse propagator, its pole residue, and proper vertex
corrections.
\end{proposition}

\begin{proof}
Write a background-dependent transverse inverse propagator schematically as
\[
 \Gamma_{VV}^{(2)}(p^2;v)
 =
 p^2-a v^2-\Sigma_V(p^2;v).
\]
The pole position solves \(\Gamma_{VV}^{(2)}=0\), while the zero-momentum Higgs
insertion contains
\(\partial_v(av^2+\Sigma_V)\), together with wavefunction and proper-vertex
terms.  The value of
\((av^2+\Sigma_V)/M_{V,\rm pole}^2\) does not determine
\(\partial_v\Sigma_V\) or the residue.  Hence the mass ratio cannot be
identified with \(\kappa_{hVV}\).
\end{proof}


\section{Weak currents, complete coherent returns and density receivers}
\label{p5c:sec:organized-weak}
\subsection{Ordered current quartet and the rooted convention}
\label{p5c:sec:v2167-ordered-current-convention}
Freeze the already declared rooted representative
$X_Ge_j=e_{j+1}$ (indices modulo three),
$C_{2,G}=-i(X_G-X_G^\dagger)$, the neutral operation order
$(\mathrm{inv},\mathrm{dir},\mathrm{Gram})$, and
$W_\nu=(e_2,e_1,e_0)$. The charged order is $(e,\mu,\tau)$.
With $C=W_e^\dagger W_\nu$, use Paper~IV's quartet
\[
 J(C)=\operatorname{Im}(C_{11}C_{22}\overline{C_{12}C_{21}}).
\]
This dictionary fixes a signed readout; its magnitude is independent of
row/column orientation reversals.

\begin{lemma}[Correct singlet/plane branch signs]
\label{p5c:lem:v2167-ordered-branch-signs}
For a unit electron vector $g=(x,y,z)^{\mathsf T}$ in $Q_\perp G$, so
$x+y+z=0$, put $\zeta_e=\langle g,C_{2,G}g\rangle/\sqrt3$.
With the fixed neutral order above,
\[
 J=+\frac{\zeta_e}{6\sqrt3}\quad(P_\mu=Q_0),\qquad
 J=-\frac{\zeta_e}{6\sqrt3}\quad(P_\tau=Q_0).
\]
\end{lemma}
\begin{proof}
On the muon-singlet branch its current row is $(1,1,1)/\sqrt3$, up to an
irrelevant role phase. The quartet has imaginary part
$\operatorname{Im}(y\bar z)/3$. Direct use of $x=-y-z$ gives
$\langle g,C_{2,G}g\rangle=6\operatorname{Im}(y\bar z)$,
proving the first formula. Transposing the muon and tau rows changes the
quartet sign and gives the native tau-singlet formula.
\end{proof}

\begin{corollary}[Literal native ordered quartet]
\label{p5c:cor:v2167-literal-current-quartet}
For the generated charged frame of Paper~V-B,
\[
 g_e=\tfrac12(-1-i,i,1)^{\mathsf T},\quad
 g_\mu=\tfrac1{2\sqrt6}(2,-1+3i,-1-3i)^{\mathsf T},\quad
 c=\tfrac1{\sqrt3}(1,1,1)^{\mathsf T},
\]
the declared ordered current satisfies
\[
 C_{11}C_{22}\overline{C_{12}C_{21}}
 =-\frac1{16}-\frac i{12},\qquad
 \boxed{J_\ell^{\rm native,ord}=-\frac1{12},\quad |J_\ell|=\frac1{12}.}
\]
\end{corollary}
\begin{proof}
Insert these rays into $C=W_e^\dagger(e_2,e_1,e_0)$.
Alternatively the preceding lemma and the already generated
$\zeta_e=\sqrt3/2$ give the same signed result.
\end{proof}

% Framework v2.217 helper integration.
\begin{remark}[A current quartet and an inclusive packet test]
\label{p5c:rem:v2217-quartet-inclusive-test}
The nonzero ordered coefficient above is retained. Its inclusive physical
consumption also tests the complete normal, external and inherited-state
packet. The explicit family reality map and complete three-direction
ordered tests are proved in
Theorem~\ref{p5c:thm:v2217-family-reality-ordered-tests} and
Proposition~\ref{p5c:prop:v2217-ordered-trace-reconstruction} below.
They neither erase the quartet nor turn it alone into a produced charge.
\end{remark}



\begin{remark}[Physical mixing comparison]
\label{p5c:rem:v2201-quartet-observation-scope}
This exact ordered quartet is retained. Its direct interpretation as the
physical three-neutrino mixing invariant is excluded by the angle-based
comparison in \ref{p5c:sec:v2201-neutrino-observation-test}. The actual
source/pole/detector calculation remains necessary.
\end{remark}

Simultaneous unitary presentation changes and role phases preserve the
quartet. Complete conjugation of the fixed ordered packet conjugates the
quartet and reverses its imaginary part. A rooted reversal used as a
presentation must transport its ordered role dictionary as well; it cannot
silently be identified with fixed-label conjugation. The sign correction
changes no generated current ray, pair, state lineage or CP magnitude.
The historical branch formulas are corrected with this same dictionary
in the appendix, and the dependent signed readouts use $-1/12$.

% END INLINE SOURCE: Paper5C-OrderedCurrentConvention-Main-v1_0.tex





% BEGIN LEPTON-CP-VALUE-CLOSURE v2.77
\subsection{finite charged-lepton incidence and native CP-value closure}
\label{p5c:sec:v277-lepton-cp-value-closure}
Paper~V-B v5.97 fixes the complete normalized charged left-family frame on the
finite current carrier:
\[
 P_\tau=Q_0,
 \qquad P_e+P_\mu=Q_\perp,
 \qquad
 \boxed{\zeta_e=\frac{\sqrt3}{2}.}
\]
Thus the finite native lepton weak-current incidence no longer contains an
unevaluated charged-plane angle:
\[
 \boxed{\texttt{LEPTON\_WEAK\_CURRENT\_INCIDENCE}
 =\texttt{CLOSED\_NATIVE\_FINITE}.}
\]
Propagation phases, pole/RG transport and optional Majorana dynamics remain
separately typed downstream questions.

\begin{corollary}[Native leptonic CP value]
\label{p5c:cor:v277-native-lepton-cp-value}
On the declared charged order $(e,\mu,\tau)$ and neutral order
$(\mathrm{inv},\mathrm{dir},\mathrm{Gram})$, the explicit ordered reduction
of Lemma~\ref{p5c:lem:v2167-ordered-branch-signs} gives
\[
 J(C_{e\nu})=-\frac{\zeta_e}{6\sqrt3}
\]
and the derived orientation moment give
\[
 \boxed{J_\ell^{\rm native,ord}=J(C_{e\nu})=-\frac1{12},\quad |J_\ell|=\frac1{12}.}
\]
Hence
\[
 \boxed{\texttt{LEPTON\_CP\_VALUE}=\texttt{CLOSED\_NATIVE\_ORDERED}[-1/12;\ |J_\ell|=1/12].}
\]
No measured PMNS entry, oscillation fit or empirical CP phase enters this
finite-current value.
\end{corollary}
\begin{proof}
Substitute the derived $\zeta_e=\sqrt3/2$ into the corrected
tau-singlet formula $J=-\zeta_e/(6\sqrt3)$.
Corollary~\ref{p5c:cor:v2167-literal-current-quartet} independently evaluates
the ordered quartet. Simultaneous unitary presentation changes and role
phases preserve it; fixed-label conjugation reverses its sign. The
rooted presentation and ordered role dictionary must be transported together.
\end{proof}

\begin{remark}[Scope of the closure]
This result closes the finite family-incidence and native current-CP target.
It does not identify terminal geometric energies with low-energy pole masses,
solve neutrino propagation phases, or supply a bulk Majorana interaction.
Those are distinct downstream claims.
\end{remark}
% END LEPTON-CP-VALUE-CLOSURE v2.77

% BEGIN BARYOGENESIS-CYCLIC-STATE-SYNC v2.80
\subsection{Native charged vertices and the complete coherent weak-memory return}
\label{p5c:sec:v2168-native-weak-memory}
The coefficients in this section are generated by the same
connection/current parent and its actual slot pullback. The uninserted
component-band retention antecedent is proved by III-F,
Proposition~\ref{p3f:prop:v2168-component-band-retention}. Later family-dependent
energy or constraint reductions still need their actual transport guard.
The matching datum is the native common-heat datum $\mu_h$; another scale
uses the existing source/RG transport.
\subsubsection{The charged internal gauge coefficients are fixed}
Paper V-A reconstructs $Y=T_R^3+(B-L)/2$ and
$Q=T_L^3+Y$. On the charged lepton slots,
\[
 (T_L^3,Y)_{e_L}=(-1/2,-1/2),\qquad
 (T_L^3,Y)_{e_R}=(0,-1),
\]
so that
\[
 \boxed{Q|_{e_L\oplus e_R}=-I\otimes I_G.}
\]
The gauge representation acts as the identity on $G$. Consequently
$U^\dagger(-I_G)U=-I$ in the existing current frame, independent of the
unknown free source Gram comparison of the free source route.

\begin{proposition}[Native electromagnetic, neutral weak and charged weak coefficients]
\label{p5c:prop:v2168-native-charged-vertices}
On the common-heat matching branch, let $g_2^2=2\pi^2$ and
$g_Y^2=6\pi^2/5$. An orthogonal change of coordinates in the canonically
normalized neutral gauge plane gives
\[
 A_Q=\frac{g_Y W^3+g_2 B}{\sqrt{g_2^2+g_Y^2}},\qquad
 Z_0=\frac{g_2 W^3-g_Y B}{\sqrt{g_2^2+g_Y^2}}.
\]
Put $g_Z=\sqrt{g_2^2+g_Y^2}$, $s^2=g_Y^2/g_Z^2$ and
$e=g_2g_Y/g_Z$. Then the internal coefficient of $A_Q$ is $eQ$ and that of
$Z_0$ is $g_Z(T_L^3-s^2Q)$. In particular
\[
 \boxed{
 e^2=\frac34\pi^2,\qquad g_Z^2=\frac{16}{5}\pi^2,\qquad s^2=\frac38,}
\]
and, choosing the inherited positive coupling convention,
\[
 \boxed{
 t_Q^{e_L}=t_Q^{e_R}=-\frac{\sqrt3\pi}{2},\qquad
 t_Z^{e_L}=-\frac{\pi}{2\sqrt5},\qquad
 t_Z^{e_R}=\frac{3\pi}{2\sqrt5}.}
\]
The charged weak coefficient is $g_2/\sqrt2=\pi$ in the complex $W^\pm$
coordinate, multiplying the generated current map $C_{e\nu}$ or its
adjoint. Equivalently each real $W^1,W^2$ off-diagonal coefficient has
magnitude $g_2/2=\pi/\sqrt2$.
\end{proposition}
\begin{proof}
Substitute the displayed orthogonal gauge coordinates into
$g_2W^3T_L^3+g_YBY$. The $A_Q$ coefficient is
$g_2g_Y(T_L^3+Y)/g_Z=eQ$. The $Z_0$ coefficient is
$(g_2^2T_L^3-g_Y^2Y)/g_Z=g_Z(T_L^3-s^2Q)$.
Using the inherited coupling squares gives $g_Z^2=16\pi^2/5$,
$s^2=3/8$ and $e^2=3\pi^2/4$. Evaluation on the two charged slots gives
the remaining numbers. Finally
$T^1=(T^++T^-)/2$, $T^2=(T^+-T^-)/(2i)$ and
$W^\pm=(W^1\mp iW^2)/\sqrt2$ give the charged weak coefficients.
The weak current pullback is the existing $W_e^\dagger I_GW_\nu$.
\end{proof}

% BEGIN FRAMEWORK v2.267: RIGHT-FIBRE GAUGE AND INTERACTION CONVERSES
\begin{proposition}[Native bare gauge vertices do not split the positive right fibre]
\label{p5c:prop:v267-native-gauge-right-fibre}
The native charged-lepton photon and neutral-weak \emph{right} vertex
coefficients in
Proposition~\ref{p5c:prop:v2168-native-charged-vertices}
are scalar multiples of $I_G$, and charged weak current vertices act
on the left chiral doublet.  They therefore commute with the constant
right-unitary $U_R(\theta)$ of
Corollary~\ref{p6:cor:v267-positive-right-unitary-orbit}.
More precisely, suppose an actual fixed finite, self-adjoint,
minimally coupled quadratic gauge-action chart has
\[
 \mathscr D_\theta=\begin{pmatrix}
 K_L&M_0U_R(\theta)\\
 U_R(\theta)^\dagger M_0^\dagger&K_R
 \end{pmatrix},\qquad [K_R,U_R(\theta)]=0,
\]
with domain-preserving, spacetime-constant $U_R$ and all retained
source and boundary data transported covariantly. Then
\[
 \boxed{\mathscr D_\theta=W_\theta\mathscr D_0W_\theta^\dagger,
 \quad P_Lf(\mathscr D_\theta)P_L=P_Lf(\mathscr D_0)P_L,\quad
 W_\theta=\operatorname{diag}(I_L,U_R(\theta)^\dagger).}
\]
Thus \emph{these generated bare gauge coefficients}, by themselves,
cannot select a positive right-fibre member through left-facing
finite quadratic responses.  The proposition does \emph{not} assert
the physical charged mass source belongs to that geometric fibre or
that added right-sensitive interactions and physical pole reductions
satisfy the commutator premise.
\end{proposition}
\begin{proof}
Each right identity vertex commutes with the rooted constant rotation;
left-only weak vertices are unchanged.  Direct block conjugation gives
the first equality.  Functional calculus and left compression give
the second, on the stated common operator domain.
\end{proof}

\begin{theorem}[Complete Hermitian commutant and finite environment converse]
\label{p5c:thm:v267-finite-right-orbit-converse}
In the ordered rooted $e,\tau,c$ right frame, let the anti-Hermitian
orbit generator be represented, up to its inessential overall sign, by
\[
 J=\begin{pmatrix}0&u&0\\-\bar u&0&0\\0&0&0\end{pmatrix},
 \quad |u|=1,\quad P=\operatorname{diag}(1,1,0),\quad Q=I-P.
\]
For every Hermitian $K$, the complete stabilizer is
\[
 \boxed{[K,J]=0\ \Longleftrightarrow\
 K=aP+ibJ+cQ,\quad a,b,c\in\mathbb R.}
\]
Therefore right family-blindness is sufficient but \emph{not necessary}
for invariance under the positive orbit.
If the actual finite right Schur kernel is
\[
 K_R(z)=zI-H_R-V(zI-H_E)^{-1}V^\dagger,
 \qquad \dim H_E=n<\infty,
\]
its commutation at every regular $z$ is equivalent to the single
nonnegative finite-moment certificate
\[
 \boxed{\mathfrak W=\|[H_R,J]\|_{\rm HS}^2+
 \sum_{k=0}^{n-1}\|[VH_E^kV^\dagger,J]\|_{\rm HS}^2=0.}
\]
Equivalently, $[K_R(z_j),J]=0$ at any $n+1$ \emph{distinct regular}
values $z_0,\ldots,z_n$ is necessary and sufficient for commutation
at every regular $z$.  This finite sampling claim assumes the printed
leading $zI$ term; it does not apply unchanged to an arbitrary
frequency-dependent kinetic coefficient or continuum self-energy.
\end{theorem}
\begin{proof}
The normal $J$ has three distinct eigenvalues $i,-i,0$.  A Hermitian
commuting matrix is diagonal on its eigenspaces and hence has precisely
the stated three-real-dimensional form.  Expansion of the Feshbach
resolvent at infinity yields the moments $VH_E^kV^\dagger$;
Cayley--Hamilton reduces higher moments to $k<n$. The sum of squares
therefore vanishes precisely when the whole right kernel commutes
with $J$.  Finally multiplication of $[K_R(z),J]$ by the monic
$p(z)=\det(zI-H_E)$ produces a matrix polynomial
\[
 -p(z)[H_R,J]-[V\operatorname{adj}(zI-H_E)V^\dagger,J]
\]
of degree at most $n$. Vanishing at $n+1$ distinct regular points
forces that polynomial to vanish identically.
\end{proof}

\paragraph{Actual action and pole scope.}
These finite tests constrain the admissible right environment once
that environment and the rooted mass-source action have been
independently constituted.  A nonzero commutator proves failure of
this \emph{symmetry certificate}; it does not by itself determine a
physical member or produce a physical pole. A vanishing commutator
proves an invariant \emph{fixed quadratic} left Schur response, not
stationary support or Lorentzian analytic continuation.
% END FRAMEWORK v2.267: RIGHT-FIBRE GAUGE AND INTERACTION CONVERSES

These are internal vertex coefficients at the native matching datum
$\mu_h$. They do not identify $e$ with a measured infrared charge or
determine the spatial mode overlaps. The latter are the actual $Z$-slot
pullbacks. At another scale the existing RG/source transport is consumed.
The same-source electromagnetic Hessian generates the current response;
there is no independent electromagnetic current normalization to provide.
The $Z_0$ coordinate above does not assert a physical gauge pole mass.
Charged weak vertices require the neutral receiver and charged bosonic
ports to remain in the complete response calculation.


Write the charged current frame and neutral role order explicitly:
\[
 U=(g_e,g_\mu,c),\quad W_\nu=(e_2,e_1,e_0),\quad C=U^\dagger W_\nu,
\]
\[
 g_e=(-1-i,i,1)^T/2,\quad
 g_\mu=(2,-1+3i,-1-3i)^T/(2\sqrt6),\quad c=(1,1,1)^T/\sqrt3.
\]
The neutral order is inverse, direct, Gram. The inherited ordered quartet
is $-1/12$ and its magnitude is $1/12$. These frames are current tags,
not mass eigenframes.
\subsubsection{The exact native charged/neutral current coefficient}
Direct substitution gives
\[
 \boxed{
 C=\begin{pmatrix}
 \frac12&-\frac i2&\frac{-1+i}{2}\\
 \frac{-1+3i}{2\sqrt6}&\frac{-1-3i}{2\sqrt6}&\frac1{\sqrt6}\\
 \frac1{\sqrt3}&\frac1{\sqrt3}&\frac1{\sqrt3}
 \end{pmatrix},\qquad C^\dagger C=CC^\dagger=I.}
\]
Its modulus-squared coefficient matrix is
\[
 \boxed{\mathcal P=\bigl(|C_{rj}|^2\bigr)=
 \begin{pmatrix}
 1/4&1/4&1/2\\
 5/12&5/12&1/6\\
 1/3&1/3&1/3
 \end{pmatrix}.}
\]
Every row and column sums to one. These are derived current coefficients,
not independent probabilities to install in the complete action.
The native common-heat coefficient in the complex charged weak coordinate
is $g_2/\sqrt2=\pi$. Away from its matching datum $\mu_h$, use the existing
source/RG transport of that coefficient.

\subsubsection{The weak vertex generates its neutral--boson receiver}
Fix an admitted finite stationary source chart of the complete charged
Hamiltonian $H$, with a stationary density $\rho$ and actual normalized
CAR source ports $q_r=a_e^\dagger(j_{\rm ext}\otimes Ue_r)$.
Take their fixed kinetic frame. The complement below is the source
orthogonal complement in IV's stationary anticommutator quotient, not a deletion
of physical gauge, scalar or charge sectors.

Strip only the fixed dimensionless weak coupling from the inherited
vertex; its receiver words retain all clock, spin, spatial, bosonic and
occupation factors and carry the inherited energy units. On the
factor-preserving chart, the weak vertex has canonical terms of the form
\[
 V_W=\pi\sum_{rj\alpha}
   \overline{C_{rj}}\,a_{\nu,j}^\dagger K_\alpha a_{e,r}
                   \otimes X_{\alpha,-}+\mathrm{h.c.}
\]
with $K_\alpha$ the actual external slot coefficient, and
$X_{\alpha,-}$ the retained bosonic port of charge $-1$.
This notation does not assume one bosonic mode, an empty state or regular
infinite CCR. Source-dependent projections and nonlinear inherited terms
are retained in their actual complete word.

\begin{proposition}[Generated weak partner columns]
\label{p5c:prop:v2168-weak-receiver-columns}
Define the actual receiver words
\[
 Y_j=\sum_\alpha a_{\nu,j}^\dagger(K_\alpha j_{\rm ext})
                         \otimes X_{\alpha,-}.
\]
Then exact CAR give
\[
 \boxed{[V_W,q_r]=\pi\sum_j\overline{C_{rj}}Y_j.}
\]
On the displayed CAR chart $\{q_s^\dagger,Y_j\}=0$, so these columns
already lie in the source complement. More generally their inherited
source-complement projection is applied to the complete columns.
\end{proposition}
\begin{proof}
For a normal monomial
$[a_i^\dagger a_k,a_l^\dagger]=\delta_{kl}a_i^\dagger$.
The term annihilating an electron gives the displayed column; its adjoint
annihilates a neutral port and has zero commutator with $q_r$.
The neutral and charged fermionic component subspaces are distinct, giving
the vanishing source anticommutator. Bosonic coefficients commute with
fermionic ports on this actual tensor chart.
\end{proof}

Let $\mathcal Qe_r=[q_r]$, $P=\mathcal Q\mathcal Q^\dagger$ and
$\mathcal L=\operatorname{ad}_H$ on the stationary state quotient. Define
\[
 A=\mathcal Q^\dagger\mathcal L\mathcal Q,\qquad
 B=(I-P)\mathcal L\mathcal Q,\qquad
 C_{\rm mem}=(I-P)\mathcal L(I-P).
\]
The full complement $C_{\rm mem}$ contains the complete interaction word.
Write $F e_j=(I-P)[Y_j]$, and split the generated columns by their actual
action marks:
\[
 B_W=\pi F C^\dagger,\qquad B_R=B-B_W.
\]
$B_R$ includes the free mass/kinetic partners, electromagnetic and neutral
weak vertices, scalar terms, pairing or higher fermion terms where inherited,
and every consumed contact. It is not an adjustable cancellation column.

If a subsequent physical quotient changes the source Gram to $N$, first
transport the actual source and its tag. When the weak-column identity
survives that reduction, $B_W$ in the normalized source frame is
$\pi F C^\dagger N^{-1/2}$ and the following returns acquire the two
outer $N^{-1/2}$ factors. If the reduction produces additional columns,
they belong to the complete $B_R$. No projected CAR identity is assumed.

More explicitly, on the positive active source-Gram support put
$\widehat C=N^{-1/2}C$. In the normalized source frame,
$B_W=\pi F\widehat C^\dagger$. Every marked return below therefore uses
$\widehat C$ in place of $C$. The upstream current matrix $C$ remains
the fixed unitary incidence; $\widehat C$ is its actual source-frame
transport and need not be unitary. If the raw physical tag is $q_r$,
its normalized-frame vector is $t_r=N^{1/2}e_r$: indeed
$\mathcal Q_{\rm raw}=\mathcal QN^{1/2}$. Thus its pole residue tests
$t_r^\dagger u$, rather than changing the tag to $e_r$ after a mixing
normalization. Null source directions are first quotiented. On the stated
CAR chart $N=I$, $\widehat C=C$ and $t_r=e_r$.

\subsubsection{The full marked memory, with all coherent returns}
Put $R(z)=(zI-C_{\rm mem})^{-1}$. Its same-frequency receiver kernels are
\[
\begin{aligned}
 K_W(z)&=F^\dagger R(z)F,&
 L_{WR}(z)&=F^\dagger R(z)B_R,\\
 L_{RW}(z)&=B_R^\dagger R(z)F,&
 K_R(z)&=B_R^\dagger R(z)B_R.
\end{aligned}
\]
These are outputs of $H,\rho,q_r$ and the marked receiver. They are not
additional functions to supply phenomenologically.

\begin{theorem}[Exact native weak-memory and complete mass-root residual]
\label{p5c:thm:v2168-complete-weak-memory}
The marked weak return and the full source inverse are
\[
 \boxed{\Sigma_{WW}(z)=\pi^2 C K_W(z)C^\dagger,}
\]
\[
 \boxed{\begin{aligned}
 S(z)&=zI-A-\pi^2 CK_W(z)C^\dagger\\
 &\quad-\pi C L_{WR}(z)-\pi L_{RW}(z)C^\dagger-K_R(z),
 \end{aligned}}
 \qquad G(z)=S(z)^{-1}.
\]
All four receiver kernels use the full $C_{\rm mem}$ and the same spectral
argument. At a fixed terminal coordinate $\varepsilon_r$ outside
$\spec(C_{\rm mem})$, the remaining finite root statement is
\[
 \boxed{S(\varepsilon_r)u=0\quad\text{for a nonzero }u.}
\]
For a simple kernel with $u^\dagger u=1$, visibility in the actual normalized
source tag $r$ requires $e_r^\dagger u\ne0$, and its residue is
\[
 \boxed{\operatorname*{Res}_{z=\varepsilon_r}G(z)
       =\frac{uu^\dagger}{a},\qquad
 a=1+\|R(\varepsilon_r)Bu\|^2>0.}
\]
\end{theorem}
\begin{proof}
Expand $B^\dagger R(z)B$ with $B=B_W+B_R$ and substitute the generated
weak columns. Exact Schur elimination gives $G=S^{-1}$.
At a real regular point
$S'(z)=I+B^\dagger(zI-C_{\rm mem})^{-2}B$, which gives $a$ and the
simple-kernel residue. This is the fixed-source stationary specialization
of the inherited complete pole/source transport theorem.
\end{proof}

For a general transported physical tag $t_r$, the tagged scalar residue
is $|t_r^\dagger u|^2/a$ and the full marked formula uses $\widehat C$
as just derived. The $C$-only displayed formula and the unitary coherent
coefficient calculations below apply on the declared $N=I$ source chart.

At complex $z$, $L_{RW}(z)$ is not $L_{WR}(z)^\dagger$:
$L_{RW}(z)=L_{WR}(\bar z)^\dagger$. The real-slice adjoint cannot replace
the second same-frequency analytic kernel. At a singular complementary
target retain the offending modes in the chart, or use the charge-sector construction complete
charged matrix $L_{-1}$. Its rank test remains valid there.
Energy-dependent physical sources, moving frames and further reductions
use all of their existing jets rather than the fixed-source specialization.

\begin{corollary}[Positive upper-half-plane marked return]
\label{p5c:cor:v2168-positive-weak-return}
For $z=x+iy$, $y>0$,
\[
 \boxed{-\operatorname{Im}K_W(z)
 =yF^\dagger\bigl[(xI-C_{\rm mem})^2+y^2I\bigr]^{-1}F\succeq0.}
\]
It is strictly positive if $F$ has full column rank on the actual quotient.
The weak marked return has the same positivity and rank because $C$ is
unitary. The complete memory has the corresponding identity with $B$.
\end{corollary}
\begin{proof}
Use the spectral decomposition of the finite Hermitian complement. Each
scalar resolvent has negative imaginary part
$-y/[(x-\omega)^2+y^2]$. The displayed positive middle matrix is invertible;
its Gram vanishes in a direction precisely when that direction is in
$\ker F$.
\end{proof}
Nonzero internal current coefficients do not alone prove $\rank F=3$ after
the actual mode/occupation/constraint return. On a full tensor CAR carrier,
an actually nonzero common external weak receiver is injective in its
three neutral-family coordinates: test it on the fermionic unit and a
bosonic vector on which that receiver is nonzero. A faithful state then
preserves that operator injectivity. A physical nonfaithful state or a
smaller quotient still needs its own overlap test.


% BEGIN FRAMEWORK v2.240: EXACT WEAK MEMORY CERTIFICATES
\begin{theorem}[Finite source-Krylov criterion for diagonal neutral memory]
\label{p5c:thm:v240-memory-krylov-diagonality}
Let $C_{\rm mem}=C_{\rm mem}^\dagger$ be the actual complete finite
complement, with $d=\dim C_{\rm mem}$, and let
$F=(F_1,F_2,F_3)$ be the three generated neutral-role receiver
columns. The exact memory
$K_W(z)=F^\dagger(zI-C_{\rm mem})^{-1}F$ is diagonal for every
regular $z$ if and only if
\begin{equation}\label{p5c:eq:v240-krylov-moment-test}
 \boxed{F_i^\dagger C_{\rm mem}^{\,n}F_j=0
           \quad(i\ne j,\ n=0,\ldots,d-1).}
\end{equation}
Equivalently the source-cyclic Krylov spaces generated by the three
$F_i$ are mutually orthogonal. Terminal neutral energies alone do
not establish this dynamical criterion.
\end{theorem}
\begin{proof}
At large $z$ the resolvent has coefficients
$F^\dagger C_{\rm mem}^nF/z^{n+1}$. If the meromorphic memory
is diagonal, every off-diagonal coefficient is zero. Conversely,
Cayley--Hamilton reduces each power $n\ge d$ to the first $d$
powers. Thus the complete off-diagonal resolvent is identically
zero. The Krylov equivalence follows by writing every inner product
$\langle C_{\rm mem}^aF_i,C_{\rm mem}^bF_j\rangle$
as a moment of order $a+b$.
\end{proof}

\begin{theorem}[Exact electron-line leakage as a neutral-memory variance]
\label{p5c:thm:v240-electron-memory-variance}
On the normalized CAR source chart where the native coherent current
$C$ is unitary, let $e_e$ be its charged electron basis vector and
$v_e=C^\dagger e_e=\tfrac12(1,i,-1-i)^T$. For an arbitrary
possibly nondiagonal, non-Hermitian $3\times3$ return $K_W(z)$,
\begin{equation}\label{p5c:eq:v240-electron-memory-leakage}
\boxed{\begin{aligned}
 \Lambda_e(z)&:=|\Sigma_{\mu e}(z)|^2+|\Sigma_{\tau e}(z)|^2\\
 &=\pi^4\left(\|K_W(z)v_e\|^2
             -|v_e^\dagger K_W(z)v_e|^2\right)\\
 &=\pi^4\|(I-v_ev_e^\dagger)K_W(z)v_e\|^2.
\end{aligned}}
\end{equation}
In particular $\Lambda_e(z)=0$ if and only if
$K_W(z)v_e\in\mathbb Cv_e$. If the actual memory passes the
separate diagonal-return test and
$K_W=\operatorname{diag}(k_1,k_2,k_3)$, then
\begin{equation}\label{p5c:eq:v240-electron-memory-pair-differences}
 \boxed{\frac{\Lambda_e(z)}{\pi^4}
 =\frac{|k_1-k_2|^2}{16}
  +\frac{|k_1-k_3|^2}{8}
  +\frac{|k_2-k_3|^2}{8}.}
\end{equation}
Thus diagonal neutral memory produces no electron-line weak mixing
precisely at the isotropic return $k_1=k_2=k_3$.
The formula concerns the weak part of the memory, not the complete
physical electron pole or oscillation probability.
\end{theorem}
\begin{proof}
As $Cv_e=e_e$, the squared norm of the electron column of
$\Sigma_{WW}=\pi^2CK_WC^\dagger$ is
$\pi^4\|K_Wv_e\|^2$. Subtract its electron-row component
$\pi^4|v_e^\dagger K_Wv_e|^2$. Orthogonal projection proves
\eqref{p5c:eq:v240-electron-memory-leakage}. For diagonal $K_W$,
expand the weighted variance with weights
$|(v_e)_j|^2=(1/4,1/4,1/2)$.
\end{proof}

\begin{corollary}[Coherent current CP-odd return and residual-return test]
\label{p5c:cor:v240-weak-return-cp-and-residual}
For the fixed native $C$ in the stated role order,
\[
 J_C=\operatorname{Im}
  (C_{e1}C_{\mu2}\overline{C_{e2}}\overline{C_{\mu1}})
 =-\frac1{12}.
\]
For real diagonal $K_W(z)=\operatorname{diag}(k_1,k_2,k_3)$ at a
regular real $z$, put
$\Delta(a)=(a_1-a_2)(a_2-a_3)(a_3-a_1)$. With the charged
mass-coordinate matrix $M_{\rm ch}$,
\[
 \boxed{\frac1{6i}\operatorname{tr}
    [M_{\rm ch},\Sigma_{WW}(z)]^3
     =-\frac{\pi^6}{12}\Delta(\varepsilon_e,\varepsilon_\mu,
         \varepsilon_\tau)\Delta(k_1,k_2,k_3).}
\]
This finite matrix invariant is not, by itself, a measured CP
asymmetry. In the full Schur decomposition,
$K_R(z)=B_R^\dagger(zI-C_{\rm mem})^{-1}B_R$ vanishes
identically in $z$ if and only if $B_R=0$; even then the inherited
residual interactions continue to enter the full $C_{\rm mem}$
and hence $K_W$. No free-vacuum normal-ordering argument can delete
those contributions from a compiled member.
\end{corollary}
\begin{proof}
The first value follows from the printed current entries.
The $3\times3$ commutator cubic/Jarlskog identity yields the
second value for the displayed $\Delta$ convention. For the last
claim the leading large-$z$ coefficient of $K_R(z)$ is
$B_R^\dagger B_R/z$; its vanishing forces $B_R=0$, and the
converse is immediate. The complementary operator itself is not
altered by deleting a direct source column.
\end{proof}
% END FRAMEWORK v2.240: EXACT WEAK MEMORY CERTIFICATES

\subsubsection{The coherent family contraction evaluated explicitly}
For arbitrary $K_W$, its exact entry is
\[
 (\Sigma_{WW})_{rs}=\pi^2\sum_{jk}C_{rj}(K_W)_{jk}\overline{C_{sk}}.
\]
If the actual return has already been proved diagonal in the existing
neutral-role tests, write
$K_W=\diag(k_{\rm inv},k_{\rm dir},k_{\rm Gram})$.
This is a dynamical diagonal-return antecedent, not a reconstruction from
the terminal neutral energies. The tagged diagonal contractions are
\[
 \boxed{
 \begin{aligned}
 \Sigma_{ee}/\pi^2&=(k_{\rm inv}+k_{\rm dir})/4+k_{\rm Gram}/2,\\
 \Sigma_{\mu\mu}/\pi^2&=5(k_{\rm inv}+k_{\rm dir})/12+k_{\rm Gram}/6,\\
 \Sigma_{\tau\tau}/\pi^2&=(k_{\rm inv}+k_{\rm dir}+k_{\rm Gram})/3.
 \end{aligned}}
\]
Its three upper mixed entries are
\[
 \boxed{
 \begin{aligned}
 \Sigma_{e\mu}/\pi^2
 &=\frac{(-1-3i)k_{\rm inv}+(3+i)k_{\rm dir}+(-2+2i)k_{\rm Gram}}
          {4\sqrt6},\\
 \Sigma_{e\tau}/\pi^2
 &=\frac{k_{\rm inv}-ik_{\rm dir}+(-1+i)k_{\rm Gram}}{2\sqrt3},\\
 \Sigma_{\mu\tau}/\pi^2
 &=\frac{(-1+3i)k_{\rm inv}+(-1-3i)k_{\rm dir}+2k_{\rm Gram}}
          {2\sqrt{18}}.
 \end{aligned}}
\]
For real $k_j$, the lower entries are their adjoints. For complex $k_j$,
compute the lower entries from the same matrix word; do not conjugate the
spectral argument.

\begin{proposition}[Complete return and its scalar information boundary]
\label{p5c:prop:v2168-coherent-family-return}
The coefficient matrix $\mathcal P$ has rank two and
$\ker\mathcal P=\mathbb C(1,-1,0)^T$. Thus its three diagonal returns alone
discard the inverse/direct difference. The complete contraction is
injective and gives
\[
 K_W=\pi^{-2}C^\dagger\Sigma_{WW}C.
\]
For a diagonal $K_W$, its complete charged return is diagonal exactly when
$k_{\rm inv}=k_{\rm dir}=k_{\rm Gram}$, in which case
$\Sigma_{WW}=\pi^2kI$.
\end{proposition}
\begin{proof}
The first two columns of $\mathcal P$ coincide, while the third is linearly
independent. Its stated kernel and rank follow. Unitarity of $C$ proves
injectivity of the full word. If $C\diag(k)C^\dagger=\diag(d)$, then
$\diag(d)C=C\diag(k)$. Every entry of $C$ is nonzero, so
$(d_r-k_j)C_{rj}=0$ forces every $d_r$ and every $k_j$ to be the same
scalar. The converse is immediate.
\end{proof}

This is a concrete application of Paper II-A's construct-before-terminalize
rule. The full coherent return is generated before any scalar readout.
The $k_j$ are resolvent returns of the complete action and state. Substituting
the terminal N56 coefficients as their propagator denominators would require
the very neutral normal-action/pole dictionary being asserted; no such
substitution is made here. Nor is diagonal weak memory required for a
pole-only mass prediction: the physical null channel can be a mixed source
direction determined downstream by the complete dynamics.


\paragraph{Actual remaining root.}
On the certified same-owner rest-comparison chart, the first finite mass
equality is $S(\varepsilon_r)u=0$ with actual tag visibility. All four receiver
kernels above are outputs of the complete action/state/source word.
Their full cross returns cannot be dropped. At a singular complementary
target use IV's full charge-sector rank test. After a root is certified,
the physical state support, particle interpretation and consumed limits
still follow the A05 handoff of VII-C. The component-band calculation
does not discharge A12's separate K16 normal-action incidence.

% BEGIN C01 V2174 BOSE GRADING WORD

% BEGIN INLINE SOURCE: Paper5C-BosonicGradingWord-Main-v1_0.tex
% v2.174: retain the full remainder in the Bose-grading action word.
\subsubsection{The Bose grading retains the complete remainder}
\begin{lemma}[Exact charged weak grading commutator]
\label{p5c:lem:v2174-bosonic-grading-word}
On an already licensed complete joint representation, retain its actual
Bose grading $N_B$ and write the existing weak word as
$V_W=\pi\sum_{rj\alpha}\bar C_{rj}F_{rj\alpha}\otimes
X_{\alpha,-}+\mathrm{h.c.}$, with
$F_{rj\alpha}=a_{\nu,j}^\dagger K_\alpha a_{e,r}$.
Define
\[
 D_\alpha=[N_B,X_{\alpha,-}],\qquad
 A=\sum_{rj\alpha}\bar C_{rj}F_{rj\alpha}\otimes D_\alpha,
 \qquad H_R=H-V_W.
\]
Then
\begin{equation}
 [N_B,V_W]=\pi(A-A^\dagger),\qquad
 \dot N_B=\frac{i}{\hbar}[H_R,N_B]
                -\frac{i\pi}{\hbar}(A-A^\dagger).
 \label{p5c:eq:v2174-bosonic-grading-word}
\end{equation}
\end{lemma}
\begin{proof}
The fermionic factors commute with this constituted Bose grading.
For self-adjoint $N_B$, $[N_B,X]^\dagger=-[N_B,X^\dagger]$.
Commuting the two adjoint halves therefore gives $\pi(A-A^\dagger)$.
Use $\dot N_B=(i/\hbar)[H,N_B]$ and retain the full $H_R$.
\end{proof}
On a standard charged coordinate constituent,
$D_\alpha=\chi_\alpha(a_{\alpha,-}^\dagger-a_{\alpha,+})$.
No cancellation with the full remainder is asserted. This is a Bose-degree
identity, distinct from the already closed conserved fermion-number law.

% END INLINE SOURCE: Paper5C-BosonicGradingWord-Main-v1_0.tex

% END C01 V2174 BOSE GRADING WORD

% END INLINE SOURCE: Paper5C-NativeWeakMemory-Main-v1_0.tex



% v2.169: generated action/support constructors and consuming scope.

% BEGIN INLINE SOURCE: Paper5C-WeakDensityReceiver-Main-v1_0.tex
% Framework v2.169: Batches 7--13, integrated at the owning source.
\subsection{Native weak density and finite boundary receivers}
\label{p5c:sec:v2169-weak-density}
\subsubsection{The native weak receiver enters the generated moments explicitly}
On the already licensed tensor CAR weak-vertex chart, V-C gives
\begin{equation}
 [V_W,q_r]=\pi\sum_j\overline{C_{rj}}Y_j,\qquad
 Y_j=\sum_\alpha a_{\nu,j}^\dagger(K_\alpha j_{\rm ext})\otimes X_{\alpha,-}.
 \label{p5c:eq:v2169-b7-p11}
\end{equation}
These are actual action columns with the external slots and finite bosonic
occupation factors retained. They are not a new receiver chosen by its
spectrum. The native internal coefficient is at the existing matching datum;
another scale uses the inherited RG/source transport.

\begin{proposition}[Exact trace receiver Gram, including finite occupation]
\label{p5c:v2169:b7:prop:weakgram}
On the full tensor Fock chart with these identical neutral-family external
slots, let $\tau_B$ be the normalized finite bosonic trace. Then
\begin{equation}
 (Y_j,Y_k)_*=\delta_{jk}\kappa,\qquad
 \boxed{\kappa=\sum_{\alpha\beta}
   \langle K_\alpha j_{\rm ext},K_\beta j_{\rm ext}\rangle
                 \tau_B(X_{\alpha,-}^\dagger X_{\beta,-})\geq0.}
 \label{p5c:eq:v2169-b7-p12}
\end{equation}
The charged-source pairing with $Y_j$ vanishes. Thus the trace Gram of the
marked weak first-commutator columns is exactly $\pi^2\kappa I_3$. This
calculation requires neither commuting bosonic ports nor infinite CCR.
The number $\kappa$ is generated by the actual slot and occupation word;
it is positive iff that common receiver is nonzero.
\end{proposition}
\begin{proof}
The finite fermionic trace obeys
$\tau_F(a^\dagger(w)a(v))=\langle v,w\rangle/2$. In IV's coefficient
pairing, the occupation correction has bosonic trace
$\tau_B([X_{\beta,-},X_{\alpha,-}^\dagger])=0$. Neutral family
orthogonality therefore gives \eqref{p5c:eq:v2169-b7-p12}, while distinct charged/neutral
fermionic slots give the zero source pairing. It is also a positive Gram
of the common receiver, proving its sign and faithful nonvanishing test.
Finally the current is unitary, so contraction of the Gram through \eqref{p5c:eq:v2169-b7-p11}
gives $\pi^2 C(\kappa I)C^\dagger=\pi^2\kappa I$.
\end{proof}

This is a calculation of one marked part of the first return. It does not
replace the full second moment by that term, assume a physical trace state,
or make the full memory scalar. Cross returns with the complete remainder
remain in \eqref{p4:eq:v2169-b7-p1}, and all later commutators use the full $H$. A smaller
constraint chart or unequal family slots uses its actual Gram instead of
\eqref{p5c:eq:v2169-b7-p12}. Even exact first-return isotropy does not prove pole isotropy.
Appendix~\ref{p4:v2169:b7:app:weakchart} evaluates the polynomial of an additional
closed receiver chart and exhibits the missing dynamical closure explicitly.

\paragraph{Result.}
The native weak vertex's second source commutator can be evaluated exactly.
On its licensed finite tensor CAR chart it contains a centered
neutral-density receiver in addition to a linear charged column. A nonzero
finite charged bosonic port forces the coefficient of that density receiver
to be nonzero. This is an action-generated input of the next source cyclic
step, not a new interaction to supply. The complete other action marks and
all their mixed commutators remain in the physical source polynomial.
No charged mass equality or physical occupation chart is inferred from this
marked calculation. The terminal coefficients and $m_p=1$ remain unchanged.

\subsubsection{The exact generated weak word}
Use the actual normalized charged and neutral tensor CAR ports on the
factor-preserving chart of V-C, with the already generated unitary current
$C=U^\dagger W_\nu$ in inverse/direct/Gram order. For a common external
one-mode constituent let $e_r^\dagger$ be the charged creation ports,
$c_j^\dagger$ the neutral ports, and $X=X_-$ the actual finite bosonic
operator with $[Q_B,X]=-X$. The inherited common-heat vertex is
\begin{equation}
 V_W=\pi\sum_{rj}
   \left(\overline{C_{rj}}c_j^\dagger e_r\otimes X
       +C_{rj}e_r^\dagger c_j\otimes X^\dagger\right).
 \label{p5c:eq:v2169-b8-w1}
\end{equation}
The coefficient $\pi$ is at the native matching datum; later scales use the
existing RG/source transport. Neither $X$ nor the external slot is selected
by a desired spectrum. The one-mode presentation isolates a constituent of
the actual complete word. The multi-slot expression is evaluated below.

Define the generated neutral presentation
\[
 n_r^\dagger=\sum_j\overline{C_{rj}}c_j^\dagger,
 \qquad n_r=\sum_j C_{rj}c_j.
\]
Unitarity of $C$ proves their CAR. This is a same-owner presentation of the
already occurring neutral fields. Their free action and every other vertex
must transform together; it does not diagonalize a neutral mass operator.
In this presentation $V_W=\pi\sum_r(n_r^\dagger e_rX+
e_r^\dagger n_rX^\dagger)$. Fermionic and bosonic factors commute.

\begin{theorem}[Native weak second source commutator]
\label{p5c:v2169:b8:thm:weaksecond}
Put $q_r=e_r^\dagger\otimes I$, $K=[X^\dagger,X]$ and
$D=(X^\dagger X+XX^\dagger)/2$. Define
\[
 F_{rs}=n_r^\dagger n_s-\tfrac12\delta_{rs}I,
 \qquad Z_r=\sum_s e_s^\dagger F_{rs}.
\]
Then
\begin{equation}
 \boxed{[V_W,q_r]=\pi n_r^\dagger\otimes X,\qquad
 [V_W,[V_W,q_r]]
       =\pi^2(q_rD-Z_r\otimes K).}
 \label{p5c:eq:v2169-b8-w2}
\end{equation}
In the first term $q_rD=e_r^\dagger\otimes D$. No bosonic CCR or
factorization of a physical state is assumed.
\end{theorem}
\begin{proof}
For even fermion bilinears CAR give
$[e_s^\dagger n_k,n_r^\dagger]=\delta_{kr}e_s^\dagger$.
For coefficient operators the exact tensor identity is
\[
 [A\otimes B,O\otimes F]=[A,O]\otimes BF
                            +OA\otimes[B,F].
\]
The forward weak term has zero fermionic commutator with $n_r^\dagger$,
and its bosonic self-commutator vanishes. The adjoint term gives
\[
 [V_W,n_r^\dagger X]
  =\pi e_r^\dagger X^\dagger X
   -\pi\sum_s e_s^\dagger n_r^\dagger n_s[X^\dagger,X].
\]
Insert $n_r^\dagger n_s=F_{rs}+\delta_{rs}/2$ and use
$X^\dagger X-K/2=D$. The first source commutator follows directly from
$[c_j^\dagger e_s,e_r^\dagger]=\delta_{sr}c_j^\dagger$ and \eqref{p5c:eq:v2169-b8-w1}.
This proves \eqref{p5c:eq:v2169-b8-w2}, including the order and sign of the finite coefficient.
\end{proof}

\subsubsection{The next receiver is quantitatively nonzero on the full finite chart}
Let $\tau_F$ and $\tau_B$ be normalized mathematical traces on the full
finite fermion and boson factors. Their tensor trace gives Theorem~\ref{p4:v2169:b7:thm:closure}'s
faithful comparison pairing $(A,B)_*=\tau\{A^\dagger,B\}$.
It is used for operator independence, not as a selected physical state.

\begin{theorem}[Centered density receiver Gram and forced finite charge return]
\label{p5c:v2169:b8:thm:densitygram}
With three neutral-family ports,
\begin{equation}
 \boxed{(Z_r,Z_t)_*=\tfrac34\delta_{rt},\qquad
 (e_u^\dagger,Z_r)_*=0.}
 \label{p5c:eq:v2169-b8-w3}
\end{equation}
The $Z_r$ are orthogonal to every linear fermionic port in the complete
tensor Fock trace, including the weak partner ports. Consequently
\begin{equation}
 \boxed{(Z_r\otimes K,Z_t\otimes K)_*
          =\tfrac34\delta_{rt}\tau_B(K^2).}
 \label{p5c:eq:v2169-b8-w4}
\end{equation}
For every nonzero finite $X$ with $[Q_B,X]=-X$,
\[
 \tau_B(Q_BK)=\tau_B(X^\dagger X)=\kappa>0,\quad \tau_B(K)=0,
\]
and
\begin{equation}
 \boxed{\tau_B(K^2)\geq
   \frac{\kappa^2}{\tau_B((Q_B-\tau_BQ_B)^2)}>0.}
 \label{p5c:eq:v2169-b8-w5}
\end{equation}
Thus the marked weak second commutator generates three independent
higher-fermion receiver directions with a strictly positive trace norm.
\end{theorem}
\begin{proof}
The finite full-Fock trace has $\tau(n_r^\dagger n_s)=\delta_{rs}/2$ and
\[
 \tau(F_{rs}^\dagger F_{tu})=\tfrac14\delta_{rt}\delta_{su}.
\]
This follows either by occupation summation or two CAR exchanges followed by
the trace. Charged-fermion CAR in the pairing of
$e_s^\dagger F_{rs}$ and $e_u^\dagger F_{tu}$ give a $\delta_{su}$ term;
the remaining neutral commutator has zero trace. Summing over $s$ gives
$3\delta_{rt}/4$. Linear charged overlap has zero centered neutral mean.
Linear neutral and conjugate overlaps have mismatched charged/neutral
occupation changes, proving their orthogonality. The tensor trace and
trace cyclicity give \eqref{p5c:eq:v2169-b8-w4}.

Because $Q_BX=X(Q_B-I)$,
$\Tr(Q_BXX^\dagger)=\Tr(X^\dagger Q_BX)
=\Tr((Q_B-I)X^\dagger X)$. Subtracting proves
$\tau_B(Q_BK)=\kappa$; $\tau_B(K)=0$ is cyclicity. Cauchy--Schwarz for
the centered charge and $K$ proves \eqref{p5c:eq:v2169-b8-w5}. Its denominator is positive: zero
charge variance would make $Q_B$ scalar and its nonzero charge commutator
impossible. Faithfulness gives $\kappa>0$.
\end{proof}

The result does not assume a special occupation cutoff shape. It uses the
actual finite charged operator and its actual charge representation.
Equivalently a finite charge-lowering operator is nilpotent, and a nonzero
nilpotent operator cannot be normal. Formula \eqref{p5c:eq:v2169-b8-w5} additionally quantifies
the generated return. It does not imply an arbitrary constraint or
nonfaithful state quotient sees the density column.

\subsubsection{An evaluated marked fourth moment}
For the marked weak word alone define $\mathcal L_W=\operatorname{ad}_{V_W}$ and the
normalized three-source embedding $q_r$. Then
\[
 \mu_1^W=0,\qquad \mu_2^W=\pi^2\kappa I_3,
\]
and the now evaluated second commutator gives
\begin{equation}
 \boxed{\mu_4^W=\pi^4
       \left[\tau_B(D^2)+\tfrac34\tau_B(K^2)\right]I_3,}
 \label{p5c:eq:v2169-b8-w6}
\end{equation}
\begin{equation}
 \boxed{\mu_4^W-(\mu_2^W)^2=\pi^4
       \left[\operatorname{Var}_{\tau_B}(D)
                    +\tfrac34\tau_B(K^2)\right]I_3>0}
 \label{p5c:eq:v2169-b8-w7}
\end{equation}
when $X\ne0$. Orthogonality \eqref{p5c:eq:v2169-b8-w3} removes the two cross terms in the
norm of \eqref{p5c:eq:v2169-b8-w2}, and $\tau_B(D)=\kappa$ gives \eqref{p5c:eq:v2169-b8-w7}.
This evaluates the actual weak mark through order four. It is not the
fourth moment of the complete $H$; every other mark and mixed word remains
part of that calculation. In particular an isotropic first receiver Gram
does not prove an isotropic or two-step closed complete response.

\subsubsection{All external slots and finite coefficient commutators}
To state the actual multi-slot word, retain external mode maps $K_\alpha$,
bosonic charged ports $X_{\alpha,-}$, and a supplied charged source vector
$v$. Write the weak fermion bilinear
$A_{sj\beta}=a_{\nu,j}^\dagger K_\beta a_{e,s}$ and set
$N_{j\alpha}=a_{\nu,j}^\dagger(K_\alpha v)$. Then
\[
 V_W=\pi\sum_{sj\beta}
  \left(\overline{C_{sj}}A_{sj\beta}\otimes X_{\beta,-}
                +C_{sj}A_{sj\beta}^\dagger\otimes X_{\beta,+}\right),
 \quad X_{\beta,+}=X_{\beta,-}^\dagger.
\]
The full second weak partner column, before any state or coefficient readout,
is
\begin{align}
 [V_W,N_{j\alpha}\otimes X_{\alpha,-}]
 =\pi\sum_{sk\beta}\Big(&C_{sk}\delta_{kj}
 a_{e,s}^\dagger(K_\beta^\dagger K_\alpha v)
                   \otimes X_{\beta,+}X_{\alpha,-}\notag\\
 &+C_{sk}N_{j\alpha}A_{sk\beta}^\dagger
                   \otimes[X_{\beta,+},X_{\alpha,-}]\notag\\
 &+\overline{C_{sk}}N_{j\alpha}A_{sk\beta}
                   \otimes[X_{\beta,-},X_{\alpha,-}]\Big).
 \label{p5c:eq:v2169-b8-w8}
\end{align}
Multiply by the first-source coefficient
$\pi\overline{C_{rj}}$ and sum $j,\alpha$ to obtain
$\operatorname{ad}_{V_W}^2q_r$. The last line need not vanish after a joint occupation
compression. The complete $H_B$ also generates
$N_{j\alpha}\otimes[H_B,X_{\alpha,-}]$, and the actual neutral kinetic,
pairing, scalar, electromagnetic and contact words generate their own
commutators. Formula \eqref{p5c:eq:v2169-b8-w8} retains those coefficient orders rather than
replacing them by a scalar memory denominator.

\subsubsection{First unproved application after this calculation}
The density receiver, finite boundary coefficient and marked fourth moment
are now supplied by the existing weak action and licensed finite occupation
operator. II-A requires them to remain in the generated full word. The next
unproved use is evaluation of the complete cyclic closure, including these
columns and all remaining mixed returns, on the named physical chart.
One cannot replace this step by the nine-direction block of Theorem~\ref{p4:v2169:b7:thm:closure} without
proving its full closure or a licensed state-null reduction of the extra
columns. Any cancellation in the complete action must itself be proved;
marked positivity does not prohibit a different mark from sharing and
canceling an operator column. At a physical state the exact effect/support
test is retained. A05's actual terminal root and consumed physical limits
remain unproved by this marked calculation.
On the actual empty-fermion support, the higher-fermion column is instead
state-null equivalent to a linear source: IV,
Corollary~\ref{p4:v2169:b12:cor:onefermion}, and
\eqref{p4:eq:v2169-b12-n5} give its exact reduction. The bosonic
occupation return remains; full-trace independence and this support quotient
are compatible. No empty-support conclusion is transferred to a filled ground.

% END INLINE SOURCE: Paper5C-WeakDensityReceiver-Main-v1_0.tex


% BEGIN INLINE SOURCE: Paper5C-WeakOccupationReturn-Main-v1_0.tex
% Framework v2.169: Batches 7--13, integrated at the owning source.
\subsection{Generated weak coordinate and occupation returns}
\label{p5c:sec:v2169-weak-occupation}
\paragraph{Result and scope.}
Theorem~\ref{p4:v2169:b12:thm:adjacentnumber} closes the extra fermionic-density direction on the empty-fermion
support. This construction calculates the bosonic occupation columns which remain.
On the licensed positive-frequency charged-coordinate constituent, the native
weak action generates an explicit finite block ladder. Its whole occupation
return, including its top boundary, has an exact continued fraction.
The one-boson-cap six-family determinant and each physical tagged numerator
are evaluated using the actual current coefficients. These formulas are
constituent calculations on stated charts; the complete action's other
marks are retained before applying them to a mass claim.

\subsubsection{The coordinate port supplies a nonzero physical receiver}
On an actual common positive-frequency two-real-mode constituent, pass to
the inherited charge modes $a_+,a_-$ with charges $+1,-1$. Retain the actual
complete total-occupation cap $N\ge1$ and define the charged coordinate port
\begin{equation}
 X=\chi(a_-^\dagger+a_+),\qquad
 [Q_B,X]=-X,\qquad Q_B=N_+-N_-,\qquad \kappa=|\chi|^2>0 .
 \label{p5c:eq:v2169-b13-o1}
\end{equation}
The coefficient $\chi$ is the generated external-slot vacuum matrix element.
For a unit coordinate profile of frequency $\omega>0$ in $\hbar=1$ units,
$\kappa=1/(2\omega)$; actual mode/clock overlaps multiply that value.
This is IV's finite free-vacuum covariance, not an adjustable weak coupling.
The weak coefficient remains $\pi$ at the native matching datum.
At another scale use the existing same-source transport.

\begin{theorem}[The first two occupation returns are evaluated]
\label{p5c:v2169:b13:thm:vacuumreturn}
For the bosonic Fock unit,
\begin{equation}
 \boxed{X|0_B\rangle=\chi|-\rangle,\qquad
 X^\dagger X|0_B\rangle
       =\kappa\left(|0_B\rangle+
                    \mathbf1_{N\ge2}|+,-\rangle\right).}
 \label{p5c:eq:v2169-b13-o2}
\end{equation}
Hence the marked weak source moments on the empty fermion--boson support are
\begin{equation}
 \boxed{\mu_2^W=\pi^2\kappa I_3,\qquad
 \mu_4^W=\pi^4\kappa^2(1+\mathbf1_{N\ge2})I_3.} \label{p5c:eq:v2169-b13-o3}
\end{equation}
The fourth-moment defect is zero for $N=1$ and strictly positive for $N\ge2$.
\end{theorem}
\begin{proof}
The annihilation term kills the unit. Apply $a_-+a_+^\dagger$ to the
one-minus vector. The first term returns the unit; the second creates one
plus and one minus, unless excluded by the actual total cap. Both amplitudes
are one and the two outputs are orthogonal. Use Theorem~\ref{p4:v2169:b12:thm:adjacentnumber}'s marked moments.
\end{proof}

A charged coordinate is not a charged annihilation map. Its action on the
free bosonic unit is nonzero even though every annihilator kills that unit.
Thus the dark-port condition $XP_B=0$ cannot be inferred on this coordinate
branch. The constituent's positive covariance proves the first receiver is
physically present. A different actual bosonic state uses its actual $X$,
not this free-unit value.

\subsubsection{The entire finite weak word generates a charge ladder}
Use the common three-family charged and neutral one-fermion slots with
the already generated unitary $C=U^\dagger W_\nu$. On total charge $-1$,
the charged slots have boson occupations $(k,k)$; the neutral slots have
$(k,k+1)$. The retained spaces therefore alternate with total boson degree
$n=0,\ldots,N$, giving exactly $3(N+1)$ dimensions. They are constructed
from the actual one-fermion charge sector, not from an inserted position
lattice or a momentum-rest preparation.

Let $A$ and $B$ be the actual one-body charged and neutral energy matrices
of this constituent, and retain its actual common free boson energy
$\omega(N_++N_-)$. Its charge-sector matrix is block tridiagonal with
\[
 D_{2k}=A+2k\omega I_3,\qquad
 D_{2k+1}=B+(2k+1)\omega I_3,
\]
\begin{equation}
 \boxed{
 W_{2k}=\pi\chi\sqrt{k+1}\,C^\dagger,\qquad
 W_{2k+1}=\pi\overline\chi\sqrt{k+1}\,C.}
 \label{p5c:eq:v2169-b13-o4}
\end{equation}
Here $W_n$ is the matrix from degree $n$ to degree $n+1$ and the reverse
matrix is its adjoint. Retain the actual phase of $\chi$ or transport all
slots together when presenting it as positive.
No terminal neutral coefficients are substituted for $B$.

\begin{theorem}[Exact finite occupation return]
\label{p5c:v2169:b13:thm:ladder}
For the stated complete constituent, let $\mathsf T_N$ be its full
$3(N+1)$ matrix and let $\mathsf E_0$ embed the three bare charged sources.
At regular complex spectral points its matrix source response is
\[
 \boxed{\mathsf G_0(z)=\mathsf E_0^\dagger
                   (z-\mathsf T_N)^{-1}\mathsf E_0=\mathsf R_0(z),}
\]
\begin{equation}
 \boxed{\mathsf R_N=(z-D_N)^{-1},\qquad
 \mathsf R_n=(z-D_n-W_n^\dagger\mathsf R_{n+1}W_n)^{-1}.}
 \label{p5c:eq:v2169-b13-o5}
\end{equation}
The three-source collective cyclic space is the entire ladder when
$\chi\ne0$. Individual tagged polynomials are calculated on that full space;
the determinant alone need not include an individual tag's visibility.
\end{theorem}
\begin{proof}
The two terms of $X$ connect $(k,k)$ with $(k,k+1)$ or $(k-1,k)$,
with amplitudes $\sqrt{k+1}$ and $\sqrt{k}$. Their adjoints give the reverse
links, proving \eqref{p5c:eq:v2169-b13-o4}, including the actual top boundary. Exact tail Schur
elimination gives \eqref{p5c:eq:v2169-b13-o5}. In $\mathsf T_N^n\mathsf E_0$, the highest possible
degree-$n$ block is $W_{n-1}\cdots W_0$, which is invertible.
Induction isolates every three-dimensional block in the collective cyclic
space. This proves its full dimension without selecting an energy root.
\end{proof}
% BEGIN C01 V2174 HIGHEST DEGREE WEAK WORD

% BEGIN INLINE SOURCE: Paper5C-HighestDegreeWeakWord-Main-v1_0.tex
% v2.174: explicit evaluated coefficients of the existing bare ladder.
\begin{corollary}[Exact highest-degree columns of the weak constituent]
\label{p5c:cor:v2174-highest-degree-weak-word}
On exactly the bare constituent of
Theorem~\ref{p5c:v2169:b13:thm:ladder}, let $\mathsf E_n$ embed its
degree-$n$ three-family block. Its actual upward links obey
\[
 W_{2k+1}W_{2k}=\pi^2\kappa(k+1)I_3,
 \qquad \kappa=|\chi|^2,
\]
whenever both links are retained. Consequently
\begin{align}
 \mathsf E_{2k}^\dagger\mathsf T_N^{2k}\mathsf E_0
     &=(\pi^2\kappa)^k k!I_3 &&(2k\le N),\nonumber\\
 \mathsf E_{2k+1}^\dagger\mathsf T_N^{2k+1}\mathsf E_0
     &=\pi\chi\sqrt{k+1}(\pi^2\kappa)^k k!C^\dagger
       &&(2k+1\le N).
 \label{p5c:eq:v2174-highest-degree-weak-word}
\end{align}
For the existing electron current tag,
$C^\dagger e_1=(1,i,-1-i)^T/2$. The corresponding squared column norms
are $\pi^{4k}\kappa^{2k}(k!)^2$ and
$\pi^{4k+2}\kappa^{2k+1}(k+1)(k!)^2$.
\end{corollary}
\begin{proof}
Use $CC^\dagger=I_3$ and the links~\eqref{p5c:eq:v2169-b13-o4}.
In $n$ applications only $n$ upward links can reach degree $n$ from degree
zero. Thus diagonal energies do not enter this column. Pair successive links
and multiply $1\cdot2\cdots k$; for odd degree retain the last
$\pi\chi\sqrt{k+1}C^\dagger$. The displayed current column has norm one.
\end{proof}
These explicit coefficients sharpen the existing constituent calculation;
its collective cyclicity is already proved. All actual $\mathsf V_{\rm rem}$
and additional target blocks remain in the complete action. This result
does not assert noncancellation for full $H^n$ or $\operatorname{ad}_H^n$,
choose the Gibbs degree, or determine a pole. On an already generated complete
sublevel, Proposition~\ref{p4:prop:v2174-graded-port-boundary} recovers the
boundary from the actual graded field port.

% END INLINE SOURCE: Paper5C-HighestDegreeWeakWord-Main-v1_0.tex

% END C01 V2174 HIGHEST DEGREE WEAK WORD


Every other complete-action term restricts by Theorem~\ref{p4:v2169:b12:thm:adjacentnumber} to an actual
$\mathsf V_{\rm rem}$. If it stays on this constituent, the full matrix is
$\mathsf T_N+\mathsf V_{\rm rem}$. If it reaches other modes, retain the
additional target blocks. Formula \eqref{p5c:eq:v2169-b13-o5} for the bare ladder is not asserted
for a dense additional contact or a different boson energy. Such a matrix
uses the complete support-polynomial construction or its actual larger
block elimination. The value of $N$ is the generated occupation boundary,
not chosen by the fourth-moment result.

\subsubsection{The additional diagonal one-cap calculation}
If the actual cap is $N=1$ and the charged and neutral one-body blocks
are diagonal in their already occurring presentations, the coherent native
current yields an explicit sixth-degree determinant $D_6$ and tagged cofactor
$N_{5,r}$. The exact source polynomial is
$D_6/\gcd(D_6,N_{5,r})$; its root includes that tag's visibility.
Appendix~\ref{p5c:app:v2169-occupation-polynomials} evaluates both polynomials
and proves the cancellation. These additional dynamical hypotheses do not
follow from the terminal compiler, and neither $a_r=\varepsilon_r$ nor a
terminal-neutral substitution for the receiver energy is assumed.

\subsubsection{First actual application after the evaluated returns}
The native V-C weak coefficient/current and the actual generated finite
coordinate/occupation chart supply \eqref{p5c:eq:v2169-b13-o1}--\eqref{p5c:eq:v2169-b13-o4}. The remaining application is
to establish which of these constituent hypotheses the claimed state and
complete action satisfy, retain every actual $\mathsf V_{\rm rem}$ and
receiver, and evaluate its tagged polynomial at the \emph{fixed upstream}
$\varepsilon_r$. No equality $a_r=\varepsilon_r$ or physical assignment of
$b_j$ follows from the calculation. Theorem~\ref{p4:v2169:b11:thm:normalstate}'s state equation precedes
Theorem~\ref{p4:v2169:b12:thm:adjacentnumber}'s target reduction when that state is not already certified.
The actual electron, muon and tau root equalities remain unevaluated.

% END INLINE SOURCE: Paper5C-WeakOccupationReturn-Main-v1_0.tex



% v2.171: conversation proofs, with actual history obligations retained.

% BEGIN INLINE SOURCE: Paper5C-NativePairReceiver-Main-v1_0.tex
% Framework v2.171: integrated conversation proof on its stated domain.
\subsection{A positive certificate that does not require numerical mode functions}
Work in one certified, canonically normalized finite chart. The external
carrier $\mathcal M$ contains every consumed mode/spin slot; $G=\mathbb C^3$
is the already generated family carrier. All traces in this section are
ordinary unnormalized traces on the explicitly displayed carriers.
The conversion to an owner's physical readout retains its existing trace
normalization; it introduces no new coefficient.

On the factorized source-return domain let
\[
 \Gamma=B\otimes M,\qquad B^{\mathsf T}=-B,\qquad
 \mathcal A=R\otimes A,\qquad R=R^\dagger\geq0.
\]
Here $\Gamma$ is the physical antisymmetric pair and $\mathcal A$ is the
marked normal-current return. The normal family marker $A$ need not be
Hermitian. For a real regular quadratic Schur source slice, the antecedent
$R\geq0$ is furnished by
\[
 R=F_m^\dagger X_mF_m,\qquad X_m\geq0,
\]
if the actual source return factors $F_E=F_m\otimes I_G$ and the selected
normal source is $X_m\otimes P_\mu$. Moving source/coupling terms must be
included before this particular factorized positive representation is used.
It is not inferred for an arbitrary retarded complex resolvent or a general
interacting reduction.

\begin{theorem}[Exact positive joint-support retention]\label{p5c:v2171:thm:positive}
On the stated actual factorized chart the complete displayed mixed trace is
\begin{equation}\label{p5c:v2171:eq:kappa}
 \Xi=\Tr_{\mathcal M\otimes G}
 (\Gamma^\dagger\mathcal A\Gamma\mathcal A^{\mathsf T})
 =\kappa Z,\qquad
 \kappa=\|R^{1/2}B(R^{1/2})^{\mathsf T}\|_{\mathrm{HS}}^2\geq0.
\end{equation}
Consequently $\operatorname{Im}\Xi=\kappa/12$. It is strictly positive
exactly when
\begin{equation}\label{p5c:v2171:eq:jointsupport}
 P_RBP_R^{\mathsf T}\ne0,\qquad P_R=\operatorname{supp} R.
\end{equation}
If $R\geq r_{\min}I$ on the full external carrier with $r_{\min}>0$, then
\[
 \kappa\geq r_{\min}^2\|B\|_{\mathrm{HS}}^2.
\]
\end{theorem}
\begin{proof}
The tensor trace gives
$\Xi=\Tr_{\mathcal M}(B^\dagger RBR^{\mathsf T})Z$.
Set $C=R^{1/2}B(R^{1/2})^{\mathsf T}$. Hermiticity of $R^{1/2}$ and
cyclicity of trace give
\[
 \Tr C^\dagger C=\Tr B^\dagger RBR^{\mathsf T}.
\]
The square roots are invertible on their supported ranges, so $C=0$ iff
$P_RBP_R^{\mathsf T}=0$. Diagonalize $R$ by a unitary change of external
frame; in the congruently transformed pair coordinates the squared norm is
$\sum_{ij}\lambda_i\lambda_j|B_{ij}|^2$. This also proves the lower bound.
Insert \eqref{p5c:v2171:eq:family} to obtain the stated imaginary part.
\end{proof}

\paragraph{What this discharges.}
On a certified factorized positive-return chart, no numerical external mode
basis or extra efficiency is needed to prove retention. The exact remaining
overlap is \eqref{p5c:v2171:eq:jointsupport}. Nonzero $B$ and nonzero $R$ separately
are insufficient when $R$ is singular. The actual K16 map must establish
these antecedents or be evaluated by its complete, possibly unfactorized,
mixed return. The III-F native component-band theorem establishes a tensor
extension on its declared branch, but expressly does not identify the
additional K16 normal-action handoff.

\subsection{The native pair generates a required density receiver}
Work first on the local rooted coefficient ports $R,L,E$, each with the
two-component spin coefficient and all three family coefficients. This is
an algebra calculation; these labels do not select a physical CAR cutoff,
spatial band, occupation cap, or state. III-F's
\texttt{p3f:prop:v2168-component-band-retention} already retains the internal
family coefficients on the uninserted component-band branch. Additional
physical reductions are still consumed in their actual role.

Use a compatible primitive spin frame for the unitary neutral edge. The
family factor of both primitive edges is the existing $I_G$. Put
\[
 \varepsilon=\begin{pmatrix}0&1\\-1&0\end{pmatrix},\qquad
 \Gamma=M\otimes\varepsilon,\qquad
 P_R(\Gamma)=\tfrac12 a_R^\dagger\Gamma a_R^\dagger.
\]
Family precedes spin in this presentation. Let $s=\mathrm{dir}$ and
$r=\mathrm{inv}$, and suppress identity spin factors in the next formulas.
The oriented primitive coefficient marks and the existing charged-current
insertion are
\begin{align}
 T_s&=a_L^\dagger N_s a_R,&
 T_r^\dagger&=a_R^\dagger N_r a_L,\notag\\
 W_+&=\pi a_E^\dagger a_L\otimes X^\dagger,&
 W_-&=\pi a_L^\dagger a_E\otimes X,\notag\\
 J_\mu^E&=a_E^\dagger P_\mu a_E.&&
 \label{p5c:v2171:eq:receiververtices}
\end{align}
The weak coefficient $\pi$ is the source-owned common-heat matching value;
another scale uses its inherited transport. $X$ is retained as a finite
charged operator, with no infinite CCR substituted. The $N_r,N_s$ and
$J_\mu^E$ identify an already rooted response coefficient. They do not add
a projector-weighted mass or install an independent family potential.
The oriented marks are constituents of a response calculation, not separately
declared Hermitian Hamiltonians.

Define the ordered marked response
\begin{equation}
 \mathcal C=\operatorname{ad}_{T_r^\dagger}
 \operatorname{ad}_{W_-}\operatorname{ad}_{J_\mu^E}
 \operatorname{ad}_{W_+}\operatorname{ad}_{T_s},\qquad
 \operatorname{ad}_A O=[A,O].
 \label{p5c:v2171:eq:markedresponse}
\end{equation}
This is not the complete Liouvillian or a selected time history.
The calculation retains all terms generated by this ordered primitive word;
it does not replace the sum over all physical source words and contacts.
Write
\[
 A=N_rP_\mu N_s,\quad L_A\Gamma=A\Gamma+\Gamma A^{\mathsf T},
 \quad B=P_\mu N_s\Gamma,\quad
 P_{ER}(B)=\sum_{ab}B_{ab}a_{E,a}^\dagger a_{R,b}^\dagger.
\]
Define $Q=P_{ER}(B)(a_R^\dagger N_r a_E)$, and
\begin{equation}
 D_X=\tfrac12(XX^\dagger+X^\dagger X),\quad
 K_X=[X,X^\dagger],\quad
 Y=Q+\tfrac12P_R(L_A\Gamma).
 \label{p5c:v2171:eq:centeredpair}
\end{equation}
The symbol $Q$ here is a receiver monomial, not the conserved charge.

\begin{theorem}[Exact pair and centered density return]
\label{p5c:v2171:thm:pairreceiver}
The complete ordered marked commutator is
\begin{equation}
 \boxed{\mathcal C P_R(\Gamma)
  =\pi^2\bigl(P_R(L_A\Gamma)\otimes D_X+Y\otimes K_X\bigr).}
 \label{p5c:v2171:eq:pairreceiver}
\end{equation}
Let $\tau_F$ be the normalized full-Fock mathematical trace on these
coefficient slots. It is not the native cyclic state. Then
\begin{equation}
 \boxed{\begin{aligned}
 \tau_F(P_R(L_A\Gamma)^\dagger P_R(L_A\Gamma))&=\frac{25}{576},\\
 \tau_F(Y^\dagger Y)&=\frac5{96},\\
 \tau_F(P_R(L_A\Gamma)^\dagger Y)&=0.
 \end{aligned}}
 \label{p5c:v2171:eq:pairreceivergram}
\end{equation}
In particular the marked pair receiver is nonzero for every nonzero $X$.
\end{theorem}
\begin{proof}
The first three commutators give, respectively,
$P_{LR}(N_s\Gamma)$,
$\pi P_{ER}(N_s\Gamma)\otimes X^\dagger$, and
$\pi P_{ER}(P_\mu N_s\Gamma)\otimes X^\dagger$.
For even fermion words use the exact tensor identity
\[
 [H\otimes F,O\otimes G]=[H,O]\otimes FG+OH\otimes[F,G].
\]
The last two commutators therefore give
\[
 \mathcal C P_R(\Gamma)=\pi^2\bigl(
 P_R(L_A\Gamma)\otimes XX^\dagger+Q\otimes K_X\bigr).
\]
The second summand has three creation operators and one annihilator.
The partial full-Fock trace over $E$ gives
$\tau_E(Q)=-P_R(L_A\Gamma)/2$, by CAR and
$\tau_E(a_{E,a}^\dagger a_{E,b})=\delta_{ab}/2$.
Substitution gives \eqref{p5c:v2171:eq:pairreceiver} and $\tau_E(Y)=0$,
which proves orthogonality to the pure $R$ receiver.

For the two nonzero spin components of the pure receiver,
$|(P_\mu)_{21}|^2=25/144$, $|M_{11}|^2=1/2$, and each distinct
two-creation monomial has squared trace norm $1/4$. Thus its norm is
$2(25/144)(1/2)/4=25/576$. For the density term, the finite mathematical
occupation sum gives
\[
 \tau_E(e_i^\dagger e_j e_k^\dagger e_l)
 =\tfrac14(\delta_{ij}\delta_{kl}+\delta_{il}\delta_{jk}).
\]
Contracting the two $R$ pair slots with their factors $1/2$ gives
\[
 \tau_F(Q^\dagger Q)=|M_{11}|^2\left(
 \frac{\sum_\alpha|(P_\mu)_{\alpha1}|^2}{4}
 +\frac{|(P_\mu)_{21}|^2}{8}\right)=\frac{145}{2304}.
\]
The partial-trace identity gives
$\tau_F(P_R(L_A\Gamma)^\dagger Q)=-25/1152$.
Consequently $\tau_F(Y^\dagger Y)=145/2304-25/2304=5/96$.
Here $\sum_\alpha|(P_\mu)_{\alpha1}|^2=(P_\mu)_{11}=5/12$.
These are ordinary finite occupation sums, with no physical Wick-state
assumption.

Finally $D_X\ge0$ is nonzero if $X\ne0$. The partial $E$ trace of
\eqref{p5c:v2171:eq:pairreceiver} is then a nonzero pair coefficient times $D_X$,
so the complete marked operator is nonzero.
\end{proof}
For any finite bosonic carrier and its normalized mathematical trace,
orthogonality also yields the exact squared receiver norm
\begin{equation}
 \tau_{FB}\bigl((\mathcal CP_R)^\dagger\mathcal CP_R\bigr)
 =\pi^4\left(\frac{25}{576}\tau_B(D_X^2)
                 +\frac5{96}\tau_B(K_X^2)\right).
 \label{p5c:v2171:eq:fullreceivernorm}
\end{equation}
The native finite charge relation $[Q_{\rm bos},X]=-X$ forces $K_X\ne0$ when
$X\ne0$: a finite charge-lowering operator is nilpotent, and a nonzero
nilpotent operator is not normal. Thus the density receiver is a required
additional direction on this faithful coefficient chart.

\subsection{The full marked two-return coefficient keeps the CP phase}
\begin{theorem}[Evaluated CP--pair incidence with finite bosonic coefficients]
\label{p5c:v2171:thm:markedmixed}
On the coefficient domain just specified, define
\[
 \Theta_X=\tau_{FB}\bigl(P_R(\Gamma)^\dagger
                              \mathcal C^2P_R(\Gamma)\bigr),\qquad
 B_X=\tau_B((XX^\dagger)^2)+\tau_B((X^2)^\dagger X^2).
\]
Then, retaining all density terms,
\begin{equation}
 \boxed{\Theta_X=\frac{\pi^4}{4}Z_{\mu;\rm inv,dir}B_X,\qquad
             \Im\Theta_X=\frac{\pi^4}{48}B_X>0\quad(X\ne0).}
 \label{p5c:v2171:eq:positivecp}
\end{equation}
No bosonic CCR or physical tracial state is assumed.
\end{theorem}
\begin{proof}
The family-charge grading of this overlap requires two transfers
$\mathrm{dir}\to\mathrm{inv}$. The two current insertions are the only
family-changing operations in the word. Each must therefore use
$(P_\mu)_{21}a_{E,2}^\dagger a_{E,1}$, while the source pair contributes
$M_{11}$ and the adjoint test pair contributes $\overline M_{22}$.
All other family entries remain in the parent calculation but have zero
coefficient in this particular trace. The common family factor is exactly
\[
 \overline M_{22}M_{11}(P_\mu)_{21}^2
   =Z_{\mu;\rm inv,dir}=-7/288+i/12.
\]
Expanding the ten ordered commutators and normal ordering with CAR gives
the following identity before a bosonic trace:
\begin{equation}
 \boxed{\begin{aligned}
 \tau_F(P_R^\dagger\mathcal C^2P_R)
 =\pi^4Z\Bigl[&\frac{
  XX^\dagger XX^\dagger+X(X^\dagger)^2X+X^\dagger XX^\dagger X}{8}\\
 &+\frac{X^2(X^\dagger)^2+(X^\dagger)^2X^2}{16}\Bigr].
 \end{aligned}}
 \label{p5c:v2171:eq:freewordtable}
\end{equation}
Every density monomial is retained in obtaining this table. The trace of
an ordered normal fermion monomial with $q$ creators and the same ascending
annihilator indices is $(-1)^{q(q-1)/2}2^{-q}$; unequal index sets have
zero trace. This rule, the displayed ordered vertices, and the five-word
table give an exact rational calculation independent of any physical state.

Cyclicity groups the two alternating words into
$\tau_B((XX^\dagger)^2)/4$. The remaining three terms group into
$\tau_B((X^\dagger)^2X^2)/4$.
Both are squared Hilbert--Schmidt norms:
$\tau_B((XX^\dagger)^2)>0$ if $X\ne0$, and
$\tau_B((X^2)^\dagger X^2)\ge0$.
This proves \eqref{p5c:v2171:eq:positivecp}. Writing
$b_X=\tau_B(X^\dagger X)>0$ also gives
\[
 B_X\ge b_X^2,\qquad \Im\Theta_X\ge\pi^4b_X^2/48.
\]
Here $b_X$ is a mathematical trace norm, not an assumed mode variance or
physical occupation.
\end{proof}

\begin{corollary}[Exact two-spin overlap factor]
\label{p5c:v2171:cor:spinoverlap}
For a common retained $2\times2$ spin overlap $K$ replace $W_+$ by
$\pi a_E^\dagger K a_L\otimes X^\dagger$ and $W_-$ by its adjoint.
The primitive neutral edge remains unitary. Then
\begin{equation}
 \boxed{\Theta_{X,K}=\frac{\pi^4}{4}\det(K^\dagger K)ZB_X.}
 \label{p5c:v2171:eq:spinoverlap}
\end{equation}
Strict positivity of its imaginary part holds precisely when $X\ne0$
and $K$ has rank two.
\end{corollary}
\begin{proof}
Use singular-value bases on the spin slots, transporting the primitive
neutral edge and both pair tests together. A two-spin antisymmetric pair
changes by the determinant phase, which cancels between the source and
adjoint test. The family current commutes with these spin basis changes.
For diagonal singular values $k_1,k_2$, the surviving family transfer sends
both occupied source spin slots to the corresponding return slots. Its
four weak factors are $k_1^2k_2^2=\det(K^\dagger K)$. The complete free
boson word table is multiplied by this same factor. This proves the formula
for arbitrary $K$, including a singular overlap, by continuity.
\end{proof}
This is an evaluated external factor, not a declaration that an unknown
physical overlap has rank two. Multi-slot weak vertices retain their actual
external matrices and cross-slot bosonic commutators. A single common $K,X$
calculation is not substituted for that complete sum.


Separate field-slot injectivity can retain the pair while erasing this
mixed phase. Appendix~\ref{p5c:v2171:app:fieldslot} gives an exact comparison.
The positive criterion is used only after its factorization and support
antecedents are certified on the actual field chart.

% END INLINE SOURCE: Paper5C-NativePairReceiver-Main-v1_0.tex


% BEGIN INLINE SOURCE: Paper5C-NativeMultiportReturn-Main-v1_0.tex
% Framework v2.171: integrated conversation proof on its stated domain.
\subsection{All weak slots enter the two-spin coefficient before contraction}
The source weak vertex in the supplementary file
\begin{center}
\small\nolinkurl{Paper5C-WeakDensityReceiver-Main-v1_0.tex}
\end{center}
\noindent at \texttt{p5c:eq:v2169-b8-w8} retains adjoint and
same-charge bosonic commutators after joint finite occupation compression. On an admitted
common two-spin component chart, combine the existing slots into
\begin{equation}
 \mathsf V=\sum_\alpha K_\alpha X_{\alpha,-}^\dagger
 =\begin{pmatrix}a^\dagger&b^\dagger\\c^\dagger&d^\dagger\end{pmatrix},
 \quad W_+=\pi\sum_f a_{E,f}^\dagger\mathsf V a_{L,f},
 \quad W_-=W_+^\dagger.
 \label{p5c:v2171:eq:matrixweak}
\end{equation}
For example $a=\sum_\alpha\overline{(K_\alpha)_{00}}X_{\alpha,-}$.
Here $K_\alpha$ denotes the actual oriented $L\to E$ slot matrix, the
adjoint of the $E\to L$ matrix in the source weak-vertex convention.
Here $a,b,c,d$ are bosonic coefficients, not fermion annihilators or new
fields. This compiler retains the complete weak slot sum on the stated
component chart; a further physical mode/constraint reduction remains its
own consumed map. Keep the rooted current, pair, and ordered primitive
neutral-edge marks of \eqref{p5c:v2171:eq:markedresponse}.

\begin{theorem}[Complete operator-valued two-spin marked return]
\label{p5c:v2171:thm:matrixweak}
Let $\|O\|_\tau^2=\tau_B(O^\dagger O)$. The exact two-return trace is
\begin{equation}
 \boxed{\Theta_{\mathsf V}=\pi^4Z\mathcal R(a,b,c,d),\qquad
  \Im\Theta_{\mathsf V}=\frac{\pi^4}{12}\mathcal R(a,b,c,d),}
 \label{p5c:v2171:eq:matrixcp}
\end{equation}
where the real coefficient $\mathcal R$ is evaluated by
\begin{equation}
 \boxed{\begin{aligned}
 8\mathcal R={}&\|[a,b]\|_\tau^2+\|[c,d]\|_\tau^2
 +\|ad-bc\|_\tau^2+\|da-cb\|_\tau^2\\
 &+\|a^\dagger d\|_\tau^2+\|ad^\dagger\|_\tau^2
 +\|bc^\dagger\|_\tau^2+\|b^\dagger c\|_\tau^2\\
 &-2\Re\tau_B(b^\dagger d c^\dagger a+b^\dagger a c^\dagger d).
 \end{aligned}}
 \label{p5c:v2171:eq:matrixcoefficient}
\end{equation}
No bosonic commutator or higher-fermion return is dropped. General positivity
is not asserted without the actual retained coefficient relations.
Appendix~\ref{p5c:v2171:app:matrixscope} gives the complete cyclic table and an
exact scope comparison.
\end{theorem}
\begin{proof}
Family-charge grading still forces both current insertions to be
$(P_\mu)_{21}$; the bosonic spin matrix cannot change this grading.
The factor is therefore $Z=\overline M_{22}M_{11}(P_\mu)_{21}^2$.
Expand all ordered commutators with \eqref{p5c:v2171:eq:matrixweak} and use the
CAR trace rule of \eqref{p5c:v2171:eq:freewordtable}. There are 24 nonzero cyclic
words, each with coefficient $+1/8$ or $-1/8$, displayed in the appendix.
Expanding the norms in \eqref{p5c:v2171:eq:matrixcoefficient} gives exactly this
table. Each word with a complex trace has its adjoint partner, so
$\mathcal R$ is real. The two equations follow.
\end{proof}
For $\mathsf V=KX^\dagger$, this reduces to \eqref{p5c:v2171:eq:spinoverlap}.
Its determinant condition belongs to that common-factor specialization;
it is not imposed on the complete slot sum.

\begin{corollary}[Complementary rank-one spin slots]
\label{p5c:v2171:cor:complementaryslots}
For $b=c=0$, the two diagonal spin slots give
\begin{equation}
 \boxed{\mathcal R_{\rm diag}=\frac18\left(
 \|d^\dagger a\|_\tau^2+\|da\|_\tau^2
 +\|ad\|_\tau^2+\|ad^\dagger\|_\tau^2\right)\ge0.}
 \label{p5c:v2171:eq:complementaryslots}
\end{equation}
It is strictly positive exactly when one of these four actual products is
nonzero. Each individual spin-slot matrix has rank one.
\end{corollary}
\begin{proof}
Set $b=c=0$ in the full formula and use
$\|a^\dagger d\|_\tau=\|d^\dagger a\|_\tau$.
Faithfulness of the mathematical trace gives the strict criterion.
\end{proof}

\begin{corollary}[Density-mediated single-row pair return]
\label{p5c:v2171:cor:rowreturn}
For $c=d=0$, the complete marked return has
\begin{equation}
 \boxed{\mathcal R_{\rm row}=\frac18\|[a,b]\|_\tau^2.}
 \label{p5c:v2171:eq:rowreturn}
\end{equation}
A noncommuting pair of retained same-charge ports supplies a nonzero
CP--pair coefficient through a single charged spin row.
\end{corollary}
\begin{proof}
Substitute into \eqref{p5c:v2171:eq:matrixcoefficient}. This term disappears if
the coefficients are commuted or the density returns are deleted.
The row form isolates a coefficient case; it is not declared to be the
complete physical weak matrix.
\end{proof}


\begin{corollary}[Complete-cap single-row CP coefficient]
\label{p5c:v2171:cor:cap-row-cp}
On the actually occurring two-port complete-occupation chart of
Theorem~\ref{p4:v2171:thm:capcommutator}, the coefficient case $a=X_i$, $b=X_j$,
$c=d=0$ obeys
\begin{equation}
 \Im\Theta_{\rm row}
 =\frac{\pi^4N\kappa_i\kappa_j}{48(d_B+1)}>0.
 \label{p5c:v2171:eq:caprowcp}
\end{equation}
\end{corollary}
\begin{proof}
Insert the complete-cap commutator norm into
\eqref{p5c:v2171:eq:rowreturn} and \eqref{p5c:v2171:eq:matrixcp}.
\end{proof}
The complete physical weak matrix must still be evaluated. This isolated
coefficient does not license removal of other spin slots, words or contacts.

% END INLINE SOURCE: Paper5C-NativeMultiportReturn-Main-v1_0.tex


% v2.189: positive complete source/charge evaluations on their actual domains.

% BEGIN INLINE SOURCE: Paper5C-CompleteNativeChargeJet-Main-v1_0.tex
% Framework v2.189: integrated source-bound complete native charge/current proof chain.
% Source-bound positive results independently checked in next187_cp and scoped in next188_charge.
% Extended relative-direction and spectral diagnostics are retained in the companion appendix.
\subsection{Complete native current returns and the inherited charge jet}
\label{p5c:sec:next187-complete-native-charge-jet}
The rooted current and boundary pair are already generated constituents.
The results below evaluate their complete local coefficient, propagate the
complete source-return constructor, and reduce the initial charge jet on its
stated actual charge-graded domain. They introduce no physical CP parameter,
field, state preparation, kinetic mechanism or efficiency.

\begin{theorem}[Completed local pair--current coefficient]
\label{p5c:thm:next187-completed-local-pair-current}
Retain the rooted member of
Theorem~\ref{p5:thm:v3103-generated-bl2-family-coefficient},
\[
 M=\tfrac12\operatorname{diag}(-1+i,1+i,-1-i),
\]
and the actual charged rays in
Corollary~\ref{p5c:cor:v2167-literal-current-quartet}:
\[
 g_e=\tfrac12(-1-i,i,1)^T,\qquad
 g_\mu=\tfrac1{2\sqrt6}(2,-1+3i,-1-3i)^T,\qquad
 g_\tau=\tfrac1{\sqrt3}(1,1,1)^T.
\]
For $P_a=g_ag_a^\dagger$, the completed same-current contraction is
\begin{equation}
 \boxed{\Xi(P_a):=\operatorname{Tr}(M^\dagger P_aMP_a^{\mathsf T}),
 \qquad (\Xi(P_e),\Xi(P_\mu),\Xi(P_\tau))=(0,2/9,1/18).}
 \label{p5c:eq:next187-completed-local-values}
\end{equation}
More generally, for any complex pair matrix $\Gamma$ and Hermitian $X,Y$,
\begin{equation}
 \operatorname{Tr}(\Gamma^\dagger X\Gamma Y^{\mathsf T})\in\mathbb R,
 \qquad
 \operatorname{Tr}(\Gamma^\dagger X\Gamma X^{\mathsf T})
 =\|X^{1/2}\Gamma(X^{1/2})^{\mathsf T}\|_{\rm HS}^2\ge0
 \quad(X\succeq0).
 \label{p5c:eq:next187-hermitian-pair-return}
\end{equation}
\end{theorem}
\begin{proof}
The rank-one identity gives
$\Xi(P_a)=|g_a^\dagger M\overline{g_a}|^2$; direct substitution gives
the three displayed rational values. In the general statement,
$\Gamma^\dagger X\Gamma$ and $Y^{\mathsf T}$ are Hermitian, and the
trace of their product is real by cyclicity. For positive $X$, expanding
the displayed squared norm and using cyclicity proves its equality.
\end{proof}
These are complete local family values. The directed source coefficient
$Z_{\mu;\mathrm{inv,dir}}=-7/288+i/12$ remains exactly the ordered mark
of Theorem~\ref{p5c:thm:v2167-mixed-family-coefficient}. Reality of the
completed undifferentiated scalar does not set a real CP-sensitive mixed
derivative or a chronological physical response to zero.

\begin{theorem}[Complete coupling jets and common-family covariance]
\label{p5c:thm:next187-complete-coupling-covariance}
Use the actual regular quadratic return of
Proposition~\ref{p3f:prop:v2167-complete-neutral-return},
$\Sigma=\mathsf B G\mathsf B^\dagger$, $G=(\zeta I-h_A)^{-1}$.
At fixed common spectral argument, the full source jet is
\begin{equation}
 \boxed{\Sigma_u=\mathsf B_uG\mathsf B^\dagger+
 \mathsf BG\mathsf B_u^\dagger+
 \mathsf BGh_{A,u}G\mathsf B^\dagger.}
 \label{p5c:eq:next187-complete-coupling-jet}
\end{equation}
For two consumed source directions,
\[
 G_{uv}=Gh_{A,uv}G+Gh_{A,u}Gh_{A,v}G+Gh_{A,v}Gh_{A,u}G,
\]
and
\begin{equation}
\begin{aligned}
 \Sigma_{uv}={}&\mathsf B_{uv}G\mathsf B^\dagger+
 \mathsf B_uG_v\mathsf B^\dagger+\mathsf B_uG\mathsf B_v^\dagger\\
 &+\mathsf B_vG_u\mathsf B^\dagger+\mathsf BG_{uv}\mathsf B^\dagger+
 \mathsf BG_u\mathsf B_v^\dagger\\
 &+\mathsf B_vG\mathsf B_u^\dagger+\mathsf BG_v\mathsf B_u^\dagger+
 \mathsf BG\mathsf B_{uv}^\dagger.
\end{aligned}
\label{p5c:eq:next187-complete-mixed-coupling-jet}
\end{equation}
On a common transported family with
$\mathsf B_\theta=U_R\mathsf BU_A^\dagger$ and
$h_{A,\theta}=U_Ah_AU_A^\dagger$, one has
\begin{equation}
 \boxed{\Sigma_\theta=U_R\Sigma U_R^\dagger,
 \qquad \Sigma_u=i[D_R,\Sigma]}
 \label{p5c:eq:next187-return-covariance}
\end{equation}
at $U_R=e^{i\theta D_R}$, $U_A=e^{i\theta D_A}$, $\theta=0$.
\end{theorem}
\begin{proof}
Differentiate the complete inverse and all three factors. This gives the
first jet and the nine placements of the mixed derivatives. On the common
family, $G_\theta=U_AGU_A^\dagger$; substitution proves covariance.
Equivalently $\mathsf B_u=iD_R\mathsf B-i\mathsf BD_A$ and
$G_u=i[D_A,G]$. Inserting both coupling derivatives in the first jet
cancels every internal $D_A$ term and leaves $i[D_R,\Sigma]$.
\end{proof}
A moving spectral argument retains its actual jets by replacing
$h_{A,u}$ with $h_{A,u}-\zeta_uI$ in inverse differentiation, and
likewise at mixed order. The local $h_R$ jets also remain. The actual
canonical metric, embedding, connection and interaction-slot jets are the
generated descendants of
Theorem~\ref{p3f:thm:v1148-finite-metric-constructor}; neither coupling
derivative is set to zero unless the actual direction has that property.

The complete stationary contact is retained:
\[
 S_4^\partial=S_4(u_0)-\tfrac12\langle j_2,C^{-1}j_2\rangle,
 \qquad j_2=\sum_\alpha j_{2,\alpha}.
\]
For $L=C^{-1}$, $L_u=-LC_uL$ and
$L_{uv}=LC_vLC_uL+LC_uLC_vL-LC_{uv}L$; source differentiation
keeps every placement among $j_2,L,j_2$, including the actual stationary
background and frame jets. All crosscontacts and higher inherited terms
belong in the complete canonical Hamiltonian. The boundary-only pair reuse
of Theorem~\ref{p3f:v2171:thm:boundarypair} does not invent an interior
dependence of that pair or delete the normal contacts.

\begin{theorem}[Matched formal charge-orbit transport]
\label{p5c:thm:next187-matched-charge-orbit}
On the finite domain of
Theorem~\ref{p4:thm:v4102-pair-charge-source}, write
\[
 H_\theta=H_c+e^{i\theta}P+e^{-i\theta}P^\dagger,
 \qquad [Q,H_c]=0,\qquad [Q,P]=rP,\quad r\ne0.
\]
The complete $H_c$ includes its actual charge-preserving interactions,
stationary contacts and central energy. For constant
$R_\theta=e^{i\theta Q/r}$,
\begin{equation}
 \boxed{H_\theta=R_\theta H_0R_\theta^\dagger,
 \qquad q(H_\theta,R_\theta\rho_iR_\theta^\dagger)
 =q(H_0,\rho_i).}
 \label{p5c:eq:next187-matched-orbit-charge}
\end{equation}
If the inherited state commutes with $Q$, the equality holds at that same
fixed state. A time-dependent common presentation additionally retains
$i\hbar\dot R R^\dagger$ in its coordinate Hamiltonian.
\end{theorem}
\begin{proof}
The charge grading gives $R_\theta PR_\theta^\dagger=e^{i\theta}P$.
Every constituent of $H_c$ is unchanged. Uniqueness of finite propagation
gives $U_\theta=R_\theta U_0R_\theta^\dagger$; transporting the state and
using $[R_\theta,Q]=0$ proves the charge equality. Differentiating
the transformed Schr\"odinger equation gives the time-dependent connection
term. The scalar center is retained and its state-evolution phase cancels.
\end{proof}
This is the printed formal charge orbit and a matched whole-packet
identity. It does not identify a physical CP reversal or relative CP
history. Holding an anomalous state fixed while changing the orbit changes
its relative preparation. A relative source direction need not cancel even
when the undifferentiated Hermitian return is real;
Appendix~\ref{p5c:v2189:app:relativecharge} records this distinction.
No diagnostic parameter is installed as a physical K16 CP source.

\begin{theorem}[Complete nontracial initial charge jet]
\label{p5c:thm:next187-complete-nontracial-charge-jet}
On an actual differentiable finite history, use a parallel charge frame and
the complete generator
\[
 H(t)=H_c(t)+P(\Gamma(t))+P(\Gamma(t))^\dagger,
 \quad P(\Gamma)=\tfrac12a^\dagger\Gamma a^\dagger,
 \quad \Gamma^{\mathsf T}=-\Gamma,
\]
with $[Q,H_c]=0$, $[Q,P]=-2P$ and $[\rho_i,Q]=0$ for the actual
inherited initial state on this restriction. This last premise is not inferred
from neutral mean, perturbation birth, or the K16 inheritance identity alone.
Here $\Gamma(t)$ is the actual complex, c-number antisymmetric matrix
on the stated canonical restriction, not a dynamical operator coefficient.
All actual normal quartic, bosonic and stationary contacts are included in
$H_c$ on this stated charge-equivariant domain. Let
$C_{ij}=\omega_i(a_j^\dagger a_i)$ on the pair support. Then
\begin{equation}
 \boxed{\dot q_i=0,\qquad
 \ddot q_i=\frac4{\hbar^2}\omega_i([P,P^\dagger])
 =\frac2{\hbar^2}\operatorname{Tr}[(2C-I)\Gamma\Gamma^\dagger].}
 \label{p5c:eq:next187-complete-initial-charge-jet}
\end{equation}
No Gaussian, Wick or tracial-state hypothesis is required.
\end{theorem}
\begin{proof}
An invariant state annihilates the expectation of every nonzero-charge
word. Since $[H,Q]=2(P-P^\dagger)$, this gives $\dot q_i=0$.
The differentiated expectation contains
$i\omega_i([\dot H,Q])/\hbar-\omega_i([H,[H,Q]])/\hbar^2$.
The first term and the normal--pair commutators have nonzero grade, while
the pair--pair term gives $4\omega_i([P,P^\dagger])/\hbar^2$.
Lemma~\ref{p4:lem:v2167-pair-commutator} gives
$[P,P^\dagger]=a^\dagger\Gamma\Gamma^\dagger a
-\operatorname{Tr}(\Gamma\Gamma^\dagger)I/2$, proving the result.
The complete normal contacts vanish from this jet by their actual grade;
they were not omitted from the generator. The scalar center commutes with
$Q$ and remains in the complete parent for other source targets.
\end{proof}
Thus empty pair support gives a strictly negative second jet for every
nonzero actual pair, $C=I/2$ on that support gives zero, and filled support
gives the opposite sign. These are conditional consequences, not state
selection. Merely half the unweighted total occupation is insufficient for
zero when $\Gamma\Gamma^\dagger$ is nonuniform.
If charged pair coefficients are dynamical operators, their joint
occupation/coefficient correlations remain in the full graded commutator
expectation below; they are not reduced to this $C$-only weighted formula.

\begin{corollary}[Complete charged contacts and direct evaluation]
\label{p5c:cor:next187-complete-graded-charge-evaluation}
If additional charged constituents actually occur, retain the complete
same-grade sums in
\[
 H=H_0+\sum_{r>0}(V_r+V_r^\dagger),\qquad [Q,V_r]=rV_r.
\]
For the same invariant initial-state domain,
\begin{equation}
 \boxed{\ddot q_i=\frac2{\hbar^2}
             \sum_{r>0}r\,\omega_i([V_r^\dagger,V_r]).}
 \label{p5c:eq:next187-complete-graded-charge-jet}
\end{equation}
For the actual signed-history target, irrespective of this initial-jet
specialization, the complete charge flow of
Theorems~\ref{p4:v2171:thm:fullcharge} and
\ref{p4:v2171:thm:flow} evaluates the owned charge directly and needs no
prior nonzero marked-incidence proxy.
\end{corollary}
\begin{proof}
In the double commutator, state charge invariance removes every nonzero
total grade. Each opposite-grade pair leaves
$2r[V_r,V_r^\dagger]$ before the overall minus sign, giving the displayed
formula. Summing constituents within $V_r$ first retains their interference.
The direct-flow statement is the exact full-state charge theorem applied
to that same complete generator and inherited state, as already exported
by Corollary~\ref{p5c:cor:v2186-direct-signed-charge-target}.
\end{proof}

\paragraph{Essential application scope.}
The actual canonical pair, its c-number matrix domain where the weighted
formula is used, complete field/source restriction, inherited
interacting state and consistent metric-connection convention remain the
arguments of the physical calculation. A primordial perturbation birth
state does not certify empty interacting fermion support at K16. The
initial jet is only potentially nontrivial and does not determine a later
endpoint. A genuine CP-sensitive mixed derivative consumes its actual
relative source direction or reversed packet and every induced jet;
neither the directed $i/12$ nor Hermitian reality evaluates it. An evaluated
zero completes a signed-charge calculation but establishes no positive
incidence or final baryon/entropy yield. These conditional evaluations do not assert complete A12/C02 physical
history or a positive baryon yield. The charge-invariant matrix specialization
is distinct from the complete paired-Gibbs cancellation in
Theorem~\ref{p5c:v2188:thm:completegibbsjet}; operator coefficients use the
full graded expectation or Theorem~\ref{p5c:v2188:thm:spectralpairjet}.

% END INLINE SOURCE: Paper5C-CompleteNativeChargeJet-Main-v1_0.tex


% BEGIN INLINE SOURCE: Paper5C-InheritedStateChargeChronology-Main-v1_0.tex
% Framework v2.189: integrated inherited-state and complete chronological charge proof chain.
\subsection{Inherited spectral states and complete charge response}
\label{p5c:v2188:sec:inheritedcharge}
Reuse the actual canonical chart, kinetic metric, connection, complete action
and inherited state of \ref{p4:v2171:eq:fullham} and
\ref{p4:v2171:thm:fullcharge}. No fermion vacuum is substituted.
The K16 inheritance theorem \ref{p9:cor:v187-k16-state-inheritance}
identifies the descendant state, but its expressly free perturbation
interface does not supply the full interacting ordinary-source support.
The fixed owner \(M\) and its pair realization are already generated by
\ref{p5:thm:v3103-generated-bl2-family-coefficient} and
\ref{p5:thm:v3105-k16-bl2-pair-corealization}.

\begin{theorem}[Complete spectral-state stationarity and its charge jet]
\label{p5c:v2188:thm:completegibbsjet}
Let \(H\) be the actual fixed finite self-adjoint Hamiltonian, including
every charged contact and operator-valued coefficient consumed by the
claim. Both \(\rho_\beta=e^{-\beta H}/\operatorname{Tr}e^{-\beta H}\) and
the symmetric actual ground state \(P_0(H)/\operatorname{rank}P_0(H)\)
have constant expectation of every fixed charge \(Q\) under this \(H\).
In particular their complete first and second charge jets vanish:
\[
 \operatorname{Tr}\rho[H,Q]=0,\qquad
 \operatorname{Tr}\rho[H,[H,Q]]=0 .
\]
This does not imply \([\rho,Q]=0\) or that the pair-only contribution
vanishes.
\end{theorem}
\begin{proof}
The IV functional-calculus constructor
\ref{p4:thm:v164-finite-gibbs-vacuum-no-thermalization} gives
\([\rho,H]=0\). Cyclicity proves both identities and the full chronological
stationarity. No covariance closure or factorization is used.
\end{proof}
For a varying actual \(H(t)\), the frozen result does not install an
instantaneous Gibbs reset. For parallel \(Q\),
\[
 q''(t_i)=\frac{i}{\hbar}\operatorname{Tr}\rho_i[\dot H_i,Q]
       -\frac1{\hbar^2}\operatorname{Tr}\rho_i[H_i,[H_i,Q]] .
\]
The first term and the actual inherited transport remain consumed.

\begin{theorem}[Prepair spectral charge jet with the full external word]
\label{p5c:v2188:thm:spectralpairjet}
Let \(H_c=\sum_a E_a\Pi_a\) be the actual complete charge-preserving
prepair Hamiltonian on a licensed finite chart, with \([H_c,Q]=0\).
Let \(P\) denote the sum of all retained charge-\(-2\) insertions, so
\([Q,P]=-2P\), and \(H=H_c+P+P^\dagger\). For its generated prepair
Gibbs state \(\rho_c=\sum_a p_a\Pi_a\), including degeneracy, the initial
charge jets are
\begin{equation}
 q'(t_i)=0,\qquad
 q''(t_i)=\frac4{\hbar^2}\sum_{a,b}(p_b-p_a)
              \|\Pi_bP\Pi_a\|_{\mathrm{HS}}^2 .
 \label{p5c:v2188:eq:spectralpairjet}
\end{equation}
The same formula holds for the normalized ground projection with
\(p_a=1/d_0\) on the ground eigenspace and zero elsewhere. The operator
\(P\) may contain bosonic operator coefficients and all same-grade
contacts; they are assembled before this contraction.
\end{theorem}
\begin{proof}
Functional calculus gives \([\rho_c,Q]=0\). The charge grading eliminates
the first jet and the normal-pair terms in the second, leaving
\(4\operatorname{Tr}\rho_c[P,P^\dagger]/\hbar^2\).
Insert the complete spectral resolution on both sides of \(P\). The two
traces have weights \(p_b\) and \(p_a\), respectively, giving the formula.
No Wick or state-factorization hypothesis occurs.
\end{proof}
For \(\beta>0\), every strictly energy-raising pair transition contributes
negatively. The actual word must establish that energy direction; the
rooted family norm alone does not. Neither this theorem nor state
inheritance says that the actual K16 state is this prepair Gibbs member.

\begin{theorem}[Vanishing complete two-insertion response on the static
chiral spectral branch]
\label{p5c:v2188:thm:chiralchronologicalzero}
On an actually admitted static invertible purely off-diagonal
number-conserving free branch let
\[
 h=\begin{pmatrix}0&T\\T^\dagger&0\end{pmatrix},\quad
 C=F_-(h)\ \hbox{or}\ C=(I+e^{\beta h})^{-1},\quad
 \Gamma(t)=\eta(t)\Gamma_0,\quad
 \Gamma_0^{\mathsf T}=-\Gamma_0 ,
\]
where \(\Gamma_0\) is an actual complex c-number matrix supported on the
right creation ports and \(\eta\) is real and integrable on the consumed
interval. Use the generated filled-negative Slater state or the generated
Gibbs state of this actual free normal Hamiltonian; both have no anomalous
moments. With the complete normal propagator retained, their one- and
two-pair chronological coefficients of
\ref{p4:thm:v2167-inherited-pair-coefficients} both vanish:
\begin{equation}
 q_1=0,\qquad
 q_2=\operatorname{Tr}[(2C-I)\mathcal B\mathcal B^\dagger]=0,\quad
 \mathcal B=\hbar^{-1}\int e^{iht/\hbar}\Gamma(t)e^{ih^{\mathsf T}t/\hbar}\,dt .
 \label{p5c:v2188:eq:chiralchronologicalzero}
\end{equation}
No terminal-mass certificate is required.
\end{theorem}
\begin{proof}
Let \(J=\operatorname{diag}(I,-I)\). Then \(JhJ=-h\),
\(JCJ=I-C\), and \(J\Gamma_0J^{\mathsf T}=\Gamma_0\).
Choose paired positive and negative energy vectors \(u_a^-\!=Ju_a^+\).
Chiral partners of both pair slots have identical coefficients and
opposite sum frequencies. Reality of \(\eta\) makes their squared
integrated moduli equal. The row weights \(2C-I\) are opposite on
those partners, so all contributions to \(q_2\) cancel, including
mixed-energy pairs. Number preservation of either state gives \(q_1=0\).
The already owned half-compression formula
\ref{p4:eq:v2169-b11-l9} is reused; the new conclusion is the integrated
two-insertion response, rather than another state-selection assertion.
\end{proof}
This is not an all-orders assertion for arbitrary complete histories.
Moving kinetic normalization, phase-varying pair coefficients, interacting
normal propagation and operator-valued coefficients use the exact
many-body chronological commutators instead.

\begin{corollary}[Complete primitive chronological cancellation on its
common charge-inversion branch]
\label{p5c:v2188:cor:primitivealltimezero}
On a licensed canonical realization of the compatible primitive neutral
coefficient chart of \ref{p5c:v2171:eq:receiververtices}, put the
unitary spin edge in its identity frame. For each family let
\[
 H_c(t)=m(t)\sum_{s=1}^2(a_{L,s}^\dagger a_{R,s}
                       +a_{R,s}^\dagger a_{L,s}),\quad
 P(t)=\eta(t)g\,a_{R,1}^\dagger a_{R,2}^\dagger ,
\]
with real inherited coefficients \(m,\eta\) and the actual fixed
\(g=(-1+i)/2,(1+i)/2,(-1-i)/2\). If the initial state is the generated
spectral ground or Gibbs state of this actual normal primitive, then the
complete evolution \(H_c+P+P^\dagger\) keeps \(\langle N\rangle=2\) for
each four-mode family at every time. Thus its charge change for \(Q=-N\)
vanishes to all orders.
\end{corollary}
\begin{proof}
Choose \(|t_g|=1\), \(t_g^2=-g/\overline g\). The canonical particle-hole
automorphism
\[
 a_R\mapsto t_g\varepsilon a_R^\dagger,\qquad
 a_L\mapsto-t_g\varepsilon a_L^\dagger,\qquad
 \varepsilon=\begin{pmatrix}0&1\\-1&0\end{pmatrix}
\]
fixes \(H_c(t)\), exchanges \(P(t)\) with \(P(t)^\dagger\), and sends
\(N\) to \(4-N\). It fixes the initial functional-calculus state and
therefore the complete chronological evolution. Taking the expectation
of \(N\) proves the statement.
\end{proof}
The local source chart explicitly does not choose a physical CAR cutoff,
state or kinetic history. This corollary is available when that actual
realization meets its antecedents. It does not assert that the complete
charged-current/contact word has this automorphism, that the filled
occupation baseline is the physical signed birth charge, or that this
particle-hole map is the physical discrete CP map.

\paragraph{Essential application scope.}
The actual direct signed target of
\ref{p5c:cor:v2186-direct-signed-charge-target} can bypass the incidence
proxy by the complete IV charge integral. The results above discharge
its spectral-state and chronology subcases when their actual antecedents
hold; they do not supply a replacement state or cancel a full CP-current
packet without its physical rooted reversal. Appendix~\ref{p5c:v2188:app:chargediagnostics}
exhibits the exact normal/pair cancellation and why a quantum external
coefficient cannot be reduced to half occupation alone. The complete finite
baryon-equivariant owner retains its exact zero from
\ref{p3f:v2171:thm:baryonstationary}. A positive baryon export instead
uses the actually licensed regulated current/history map of
\ref{p12:v2171:prop:baryonexport}. A nonzero pair coefficient alone
does not install that map.

% END INLINE SOURCE: Paper5C-InheritedStateChargeChronology-Main-v1_0.tex


% BEGIN LITERATURE-VERIFIED EM AND OBSERVABLE INTERFACES v2.210
\subsection{Charged source origin in the complete interaction return}
\label{p5c:sec:v228-source-origin}
Every use of the noncentral charged source-value handoff carries the
hierarchy-arrow origin certificate of
Definition~\ref{p5:def:v228-hierarchy-origin}, as scoped by
Proposition~\ref{p5:prop:v228-charged-origin-scope}.
The complete scalar and weak memory constructors remain valid functions
of their supplied actual source. The exact tests of
Theorem~\ref{p1b:thm:v228-character-defect} and
Corollary~\ref{p1b:cor:v228-operator-character-defect} can close its
mass-only marked block once word-class membership is proved.
The physical interacting pole still consumes the full same-action kernel,
the stationary support and the actual source visibility condition.
No mass-only identity deletes a generated interaction return.

% END FRAMEWORK v2.228: CHARGED ORIGIN AND CHARACTER TESTS

\subsection{Coherent-return information and the regular complementary gap}
\label{p5c:app:v2168-weak-memory-diagnostics}
These comparisons isolate information lost by scalar readout and the
gap antecedent of the root bound. They are not alternate native actions.
\paragraph{Information lost by diagonalizing the return too soon.}
For the diagnostic difference $K_W=\diag(t,-t,0)$, all three tagged
diagonal returns vanish. The complete upper mixed return is nevertheless
\[
 \frac1{\pi^2}(\Sigma_{e\mu},\Sigma_{e\tau},\Sigma_{\mu\tau})
 =\left(-\frac{1+i}{\sqrt6}t,
          \frac{1+i}{2\sqrt3}t,\frac{i}{\sqrt2}t\right).
\]
The complete matrix has eigenvalues $t,-t,0$ before the common $\pi^2$
factor. Scalar diagonals therefore cannot substitute for its action return.
This is a relation-retention calculation on the generated coefficient packet,
not a new normal action selected by its terminal values.

\paragraph{The regular-gap bound requires its declared gap.}
The short root bracket cannot cross a pole of the eliminated complement.
At such a crossing one enlarges the retained chart. A small residual does
not establish a physical pole if the gap or source-overlap antecedent fails.
For a degenerate root use its kernel projector and complete residue matrix;
the simple-vector derivative statement is not applied without simplicity.

% END INLINE SOURCE: Paper5C-WeakMemoryDiagnostics-Appendix-v1_0.tex



% v2.169: additional hypotheses and finite comparisons.

% BEGIN INLINE SOURCE: Paper5C-WeakOccupationDiagnostics-Appendix-v1_0.tex
% Framework v2.169: Batches 7--13, integrated at the owning source.

\section{Neutral source, propagation and experimental comparison}
\label{p5c:sec:organized-neutral}
\subsection{current/propagation flavor synchronization}
\label{p7v2:sec:v266-current-propagation-flavor-sync}

\begin{definition}[Lepton weak-current terminal-role map]
\label{p7v2:def:v266-lepton-current-map}
Let $W_e$ and $W_\nu$ be the canonically normalized evaluations of the
charged-lepton terminal support and the neutral N56 role support into the
physical left-handed lepton-doublet current owner.  The framework-native
lepton flavor datum is
\[
 \boxed{C_{e\nu}:=W_e^\dagger\,\mathscr J_W\,W_\nu.}
\]
It is a current/incidence map between already-massed terminal roles.  No
charged-lepton or neutrino mass is defined as an eigenvalue of this map or of
a Yukawa/Majorana matrix.
\end{definition}

\begin{proposition}[Framework-native oscillation/readout kernel]
\label{p7v2:prop:v266-oscillation-current-kernel}
If the downstream propagation theory supplies neutral-role phases $\Phi_r$
then
\[
 \mathcal P_\nu=\sum_r e^{-i\Phi_r}\Pi_r^{\rm role},\qquad
 \mathcal A_\ell=C_{e\nu}\mathcal P_\nu C_{e\nu}^\dagger
\]
is the native charged-current propagation kernel.  A conventional PMNS matrix
is at most a normalized coordinate representative of $C_{e\nu}$; a standard
oscillation phase formula is at most a downstream reduction of the generated
propagation phases.
\end{proposition}

\begin{remark}[Type-I/PMNS formulas are comparison dictionaries]
The Casas--Ibarra, seesaw, Takagi and explicit PMNS formulas retained in this
paper are conditional external Standard--Model/type-I-seesaw comparison
coordinates.  They do not define the framework-native light-neutrino masses,
which are terminal geometric readouts, and they do not close a native mixing
gap.  Any high-energy claim that specifically consumes this external
dictionary retains its stated conditional hypotheses.
\end{remark}
% END CURRENT-PROPAGATION-FLAVOR-SYNC v2.66



% BEGIN LEPTON-CURRENT-INCIDENCE-CLOSURE v2.67
\subsection{lepton current-incidence and propagation closure}
\label{p7v2:sec:v267-lepton-current-incidence-closure}

\begin{theorem}[Generated lepton weak-current incidence]
\label{p7v2:thm:v267-generated-lepton-current-incidence}
On the physical left-handed lepton-doublet current owner, Paper~V-A supplies
the family-blind weak-current coefficient $\mathscr J_W=I_G$.  The charged
branch supplies the already generated, canonically normalized charged-lepton
incidence $W_e$, while Paper~V-B Release v5.93 supplies the canonical neutral
role evaluation $W_\nu$.  Hence
\[
 \boxed{C_{e\nu}=W_e^\dagger\mathscr J_WW_\nu=W_e^\dagger W_\nu}
\]
is a generated current/incidence readout with no remaining continuous family
selector.  Therefore
\[
 \boxed{
 \texttt{LEPTON\_WEAK\_CURRENT\_INCIDENCE}=\texttt{CLOSED}.}
\]
A conventional $U_{\rm PMNS}^{\rm eff}$ is a coordinate readout only
after its physical source, propagation and detector map is established.
Identifying it directly with this current consumes the equal-boundary
and physical pole/phase hypotheses, and fails the comparisons in
\ref{p5c:sec:v2201-neutrino-observation-test}. Structural incidence
closure by itself does not establish observational agreement.
\end{theorem}

\begin{corollary}[Former PMNS structural target is exhausted]
\label{p7v2:cor:v267-former-pmns-target-exhausted}
The former object called a neutral/charged mass-eigenframe or Dirac-PMNS
relative frame carries no independent framework-native obligation.  Its mass
part is already supplied by the terminal charged and N56 role-energy
constructors, while its flavor-changing content is $C_{e\nu}$ above.  Thus
\[
 \boxed{
 \texttt{PMNS\_DIRAC\_RELATIVE\_FRAME}
 =\texttt{RETIRED\_OBSOLETE\_ONTOLOGY}}
\]
with no replacement structural gap.
\end{corollary}

\begin{theorem}[Neutral propagation phase on the existing one-particle domain]
\label{p7v2:thm:v267-neutral-propagation-phase}
Assume the already separate terminal-role-energy to physical rest/pole
identification
\[
 \mu_r c_{\rm op}^2=\mathcal E_r^{\rm term}
\]
on a Paper--IV one-particle block.  Then Paper~IV's positive-frequency law
fixes
\[
 E_r(p)=\sqrt{c_{\rm op}^2\|p\|^2+\mu_r^2c_{\rm op}^4},
 \qquad
 \Phi_r(\gamma)=\hbar^{-1}\int_\gamma E_r\,dt,
\]
so the neutral propagation operator
\[
 \mathcal P_\nu(\gamma)
 =\sum_r e^{-i\Phi_r(\gamma)}\Pi_r^{\rm role}
\]
has no additional flavor parameter.  Hence
\[
 \boxed{\begin{gathered}
 \texttt{NEUTRAL\_PROPAGATION\_PHASE}\\
 =\texttt{CLOSED\_DEPENDENT[ONE\_PARTICLE\_DOMAIN,}\\
 \texttt{POLE\_REST\_BRIDGE]}.
 \end{gathered}}
\]
For equal normalized production/detection boundaries,
\[
 \mathcal A_\ell(\gamma)
 =C_{e\nu}\mathcal P_\nu(\gamma)C_{e\nu}^\dagger,
\]
so
\[
 \boxed{
 \texttt{OSCILLATION\_CURRENT\_READOUT}
 =\texttt{CLOSED\_DEPENDENT[PROPAGATION,BOUNDARY]}.}
\]
Matter effects, source coherence, detector response and experimental
likelihoods remain downstream evaluations.
\end{theorem}

\begin{corollary}[N56 scale-free squared-gap ratio]
\label{p7v2:cor:v267-n56-gap-ratio}
On the same rest/pole bridge, the ordered N56 packet
\(E_{56}\{x^{-1},x,x^2\}\), $x^3=192$, gives
\[
 \boxed{
 \frac{\Delta_{32}}{\Delta_{21}}
 =\frac{x^4-x^2}{x^2-x^{-2}}
 =\frac{x^4}{x^2+1}
 =32.310511556\ldots .}
\]
This ratio is independent of the common $E_{56}$ scale.  It is a propagation
prediction on the stated bridge, not an eigenvalue derivation of the terminal
masses.
\end{corollary}
% END LEPTON-CURRENT-INCIDENCE-CLOSURE v2.67


% BEGIN QUARK-CURRENT-INCIDENCE-IMPORT v2.68
\subsection{quark current-incidence closure and SM-legacy synchronization}
\label{p7v2:sec:v268-quark-current-incidence-closure}

\begin{theorem}[Generated quark charged-current map]
\label{p7v2:thm:v268-generated-quark-current}
Paper~V-A Release v3.101 and Paper~V-B Release v5.94 supply the generated
normalized charged terminal-role incidences $W_u,W_d$ on the common
left-handed quark current owner.  Since the weak raising current has family
coefficient $I_G$,
\[
 \boxed{C_{ud}=W_u^\dagger W_d}
\]
is the complete native quark flavor/current datum.  Thus
\[
\begin{gathered}
 \boxed{\texttt{QUARK\_WEAK\_CURRENT\_INCIDENCE}=\texttt{CLOSED}},\\[3pt]
 \boxed{\texttt{CKM\_RELATIVE\_FRAME}=\texttt{RETIRED\_OBSOLETE\_ONTOLOGY}}.
\end{gathered}
\]
A conventional $V_{\rm CKM}^{\rm eff}$ is only the standard coordinate
section $[C_{ud}]_{\rm std}$.
The statement is on the inherited full-rank current domain.  The actual
quark entry table, its modulus matrix and ordered CP quartet still consume
the explicit same-source output certificate of
Proposition~\ref{p6:prop:v2211-quark-current-output-certificate}.  Closed
ownership supplies no numerical identification with the legacy benchmark.
\end{theorem}

\begin{corollary}[Flavor-sensitive interaction rule]
\label{p7v2:cor:v268-flavor-sensitive-current-rule}
Every active finite/current or downstream QFT formula in this paper that
needs quark or lepton flavor change is to consume $C_{ud}$ or $C_{e\nu}$,
respectively.  CKM/PMNS symbols may appear only inside explicit external
Standard--Model comparison dictionaries and are then defined as coordinate
sections of those generated current maps.  Terminal masses remain direct
geometric readouts and are not regenerated by these matrices.
\end{corollary}

\begin{remark}[Majorana/Takagi and seesaw scope]
\label{p7v2:rem:v268-majorana-seesaw-scope}
The type-I seesaw, Casas--Ibarra, Takagi-column, Majorana-phase,
$m_\beta$ and $m_{\beta\beta}$ formulas retained in Papers~V-B/V-C are
external comparison structures unless a separate native coherent neutral
interaction theorem is explicitly invoked.  They do not own N56 terminal
masses, $C_{e\nu}$, or the quark current map $C_{ud}$.
\end{remark}
% END QUARK-CURRENT-INCIDENCE-IMPORT v2.68



% BEGIN K16-VISIBLE-CONNECTIVITY v2.69
\subsection{tau-singlet current branch and leptonic CP reduction}
\label{p5c:sec:v276-tau-singlet-lepton-cp}
\noindent\textbf{Current scope.} The open $\zeta_e$ status below is superseded by main-text Release v2.77.  Its signed formula is corrected by the explicit ordered dictionary of Release v2.82; this block is a historical conditional branch, not a live orientation residual.
Paper~V-B v5.96 now fixes the discrete charged-lepton placement
\[
 \boxed{P_\tau=Q_0,\qquad P_e+P_\mu=Q_\perp,}
\]
while leaving the oriented decomposition of $Q_\perp$ between $e$ and $\mu$
source-dependent.  Therefore the formal current map
$C_{e\nu}=W_e^\dagger W_\nu$ is physically evaluated in its singlet/plane
placement but not yet in its final in-plane angle:
\[
 \boxed{\begin{gathered}
 \texttt{LEPTON\_WEAK\_CURRENT\_INCIDENCE}\\
 =\texttt{CLOSED\_FORMAL/OPEN\_INPLANE}\\
 \texttt{\_PHYSICAL\_EVALUATION}.
 \end{gathered}}
\]

\begin{corollary}[One-scalar CP reduction on the derived tau-singlet branch]
\label{p5c:cor:v276-tau-singlet-zeta-jarlskog}
Keep the declared neutral role order $(\mathrm{inv},\mathrm{dir},\mathrm{Gram})$
and charged role order $(e,\mu,\tau)$.  Define
\[
 \zeta_e=\frac1{\sqrt3}\operatorname{Tr}(P_eC_{2,G}).
\]
Then on the derived branch above,
\[
 \boxed{J(C_{e\nu})=-\frac{\zeta_e}{6\sqrt3}.}
\]
Hence
\[
 \boxed{
 \texttt{LEPTON\_CP\_CONSTRUCTOR}=\texttt{CLOSED},\qquad
 \texttt{LEPTON\_CP\_VALUE}=\texttt{OPEN\_ON\_ZETA\_E}.}
\]
No measured PMNS entry or CP phase is used.
\end{corollary}
\begin{proof}
Corollary~\ref{p5c:cor:v274-zeta-jarlskog-reduction} gives
$J=+\zeta_e/(6\sqrt3)$ on the historical branch with $P_\mu=Q_0$ and
$P_e+P_\tau=Q_\perp$.  Passing to the derived branch exchanges the charged
$\mu$ and $\tau$ rows while leaving $P_e$, the neutral ordering and
$\zeta_e$ fixed.  A Jarlskog invariant changes sign under one row
transposition, giving the displayed minus sign.  Paper~V-B v5.96 leaves
$\zeta_e$ as the sole in-plane source readout.
\end{proof}

\begin{remark}[No premature $R_G$ closure]
The exact structural relation $R^-\mapsto R_G^\dagger$ remains useful
geometric data, but its physical charged-current consumption is not yet a
frozen action word.  It therefore cannot be used to set a numerical
$\zeta_e$ in this release.
\end{remark}
% END LEPTON-DIRECTION-CP-BRANCH-CORRECTION v2.76


% BEGIN INLINE SOURCE: Paper5C-CPBL2Comparisons-Appendix-v1_0.tex
% Framework v2.167: supplementary coordinate table and licensed-branch comparisons.
\subsection{Exact family-reality obstruction and the neutral-role odd signature}
\label{p5c:sec:v229-complete-family-reality}

\begin{corollary}[Complete labelled family-reality obstruction]
\label{p5c:cor:v229-complete-family-reality-obstruction}
Fix the already generated symmetric pair $M$, the three individually
labelled charged ray projectors $P_e,P_\mu,P_\tau$ and a Hermitian normal
family matrix $h$.  Then
\[
 \boxed{\begin{aligned}
 &\exists\,Y\in U(3):\quad
   Y\bar M Y^{\mathsf T}=M,\quad
   Y\bar P_rY^\dagger=P_r\ (r=e,\mu,\tau),\\
 &\hspace{37mm}Y\bar hY^\dagger=h
 \end{aligned}
 \quad\Longleftrightarrow\quad\Xi(h)=0.}
\]
Here $\Xi=2(u^2+v^2+w^2)$ is the exact quadratic norm of
Proposition~\ref{p5c:prop:v2217-ordered-trace-reconstruction}.
The iff concerns the declared \emph{family coefficient chart}; it does
not identify a full dynamical spin/state/charge conjugation and does not
assert a nonzero cosmological $B-L$ yield.
\end{corollary}
\begin{proof}
By Theorem~\ref{p5c:thm:v2217-family-reality-ordered-tests},
the first pair-plus-three-ray conditions force $Y=\pm X$.
The last condition is then exactly $h=h^\sharp$, equivalently to
$h_-=0$.  In the common-real charged frame $h_-=iK$, hence this
is equivalent to $\Xi=2(u^2+v^2+w^2)=0$.
\end{proof}

\begin{corollary}[One-dimensional signature of the neutral direct/inverse gap]
\label{p5c:cor:v229-neutral-gap-odd-trace-signature}
For the restricted neutral-diagonal normal difference
$h_-=\delta(N_{\rm dir}-N_{\rm inv})$, the common-real charged frame
coordinates and its three independent ordered traces are
\[
 \boxed{(u,v,w)=\delta\left(\frac1{\sqrt3},
 -\frac1{\sqrt6},-\frac1{\sqrt2}\right),}
\]
\[
 \boxed{F=-\frac89\delta,\quad H=\frac49\delta,\quad
 G=0,\quad\Xi=2\delta^2.}
\]
Thus $F+2H=G=0$ singles out this one-dimensional odd-normal image;
it does not identify all three odd-normal directions with a single
neutral occupation difference.  For the stated charged rephasing $V$,
$X\bar V=-V$ (not $-iV$), though the induced conjugation on normal
coefficients remains ordinary complex conjugation.
\end{corollary}
\begin{proof}
Transform the diagonal neutral difference into the existing common-real
charged frame and compare with the skew matrix $K$ of
Proposition~\ref{p5c:prop:v2217-ordered-trace-reconstruction}.
Substitution of the resulting $(u,v,w)$ in its invertible linear map
for $(F,H,G)$ gives the displayed values and norm.
\end{proof}

\begin{remark}[An odd coefficient is not an automatically produced charge]
\label{p5c:rem:v229-k16-charge-firewall}
The family criterion $\Xi>0$ is equivalent to \emph{failure of the
specified common family reality}, not to a nonzero physical
$Y_{B-L}$.  Indeed the existing finite chronological comparison has
\[
 \delta K(s)=-\frac29\sin(2\delta s)
 (y_\tau-y_\mu)(y_\tau+2y_\mu-3y_e),
\]
which vanishes for $y_\tau=y_\mu$ even when $\delta\ne0$.
The actual $K16$ restriction, inherited cyclic state and complete
ordered signed-charge evolution must be evaluated.  Conversely a
vanishing $\Xi$ on this family chart does not eliminate hypothetical
CP-odd owners outside that chart.
\end{remark}
% END FRAMEWORK v2.229: ODD NORMAL AND K16 PRODUCTION FIREWALL


\paragraph{Production--propagation--detection requirement.}
A geometric current is not itself an experimentally observed
flavor or oscillation probability. A physical claim must specify
\[P_D(t)=\operatorname{Tr}\bigl(E_D V(t)\rho_P V(t)^\dagger\bigr)\]
in the occurring source, kinetic evolution, and detector domain.
The optional external neutrino comparisons and their unresolved
matching conditions are documented in Paper~V-D.

\section{Signed charge, kinetic normalization and baryon history}
\label{p5c:sec:organized-baryo}
\subsection{Finite traces, thermal states, and kinetic efficiency}
\label{p7v2:sec:thermal-traces}

\begin{theorem}[Block traces and full-carrier invariant states]
\label{p7v2:thm:thermal-trace-nonunique}
Let
\[
 \mathcal V_{\rm th}
 =
 \mathcal V_7\oplus\mathcal V_9
\]
with native block symmetry \(U(7)\times U(9)\), and let \(P_7,P_9\) be the
central block projections.  On the block observable algebra
\[
 \mathcal A_{\rm blk}
 =\End(\mathcal V_7)\oplus\End(\mathcal V_9),
\]
the complete family of normalized tracial states is
\begin{equation}
\label{p7v2:eq:thermal-trace-family}
 \tau_w(A)
 =
 w\frac{\Tr A_7}{7}
 +(1-w)\frac{\Tr A_9}{9},
 \qquad 0\leq w\leq1.
\end{equation}
On the full carrier algebra
\(\mathcal A_{\rm full}=\End(\mathcal V_{\rm th})\simeq M_{16}(\CC)\), define
the same family by
\begin{equation}
\label{p7v2:eq:thermal-full-state-family}
 \tau_w(A)
 =w\frac{\Tr(P_7AP_7)}7
 +(1-w)\frac{\Tr(P_9AP_9)}9.
\end{equation}
These are exactly the \(U(7)\times U(9)\)-invariant states on
\(\mathcal A_{\rm full}\), but they are tracial on the full algebra if and only
if
\[
 \boxed{w=\frac7{16}},
\]
in which case \(\tau_w=\Tr/16\).  Thus the observable algebra is part of the
type of every thermal trace claim.
\end{theorem}

\begin{remark}[Thermal-state selector atlas]
Equations~\eqref{p7v2:eq:thermal-trace-family} and
\eqref{p7v2:eq:thermal-full-state-family} encode different algebra types.  The
native symmetry alone leaves the invariant-state family
\[
 \mathfrak A_{\rm th}^{\rm state}=\{\tau_w:0\leq w\leq1\}.
\]
If the physical observable algebra is independently shown to be
\(\mathcal A_{\rm full}\) and the state is required to be tracial on it, then
\(w=7/16\) and \(I_{16}/16\) are derived rather than selected.  A general
thermal state need not be tracial, however, so this mathematical reduction
does not select the physical preparation.  The larger \(\Hkin\) family also
contains the Hamiltonian, initial condition, collision/washout kernels,
spectator map, and cosmological history.
\end{remark}

\begin{proof}
Every normalized trace on \(\mathcal A_{\rm blk}\) assigns weights
\(w=\tau(P_7)\) and \(1-w=\tau(P_9)\) to the two central simple summands.  Its
restriction to each nonzero summand is the unique normalized matrix trace,
which proves \eqref{p7v2:eq:thermal-trace-family} and its completeness.

On \(\mathcal A_{\rm full}\), a state invariant under block-unitary
conjugation has a density matrix in the commutant of
\(U(7)\times U(9)\).  That commutant is
\(\mathbb C P_7\oplus\mathbb C P_9\), so positivity and normalization give
exactly \eqref{p7v2:eq:thermal-full-state-family}.  To test cyclicity, choose
unit vectors \(u\in\mathcal V_7\), \(v\in\mathcal V_9\), let
\(A=|u\rangle\langle v|\), and let \(B=|v\rangle\langle u|\).  Then
\[
 \tau_w(AB)=\frac w7,
 \qquad
 \tau_w(BA)=\frac{1-w}{9}.
\]
These are equal exactly when \(w=7/16\).  Substitution gives
\(\tau_{7/16}(A)=\Tr(A)/16\) for every full-algebra element, completing the
proof.
\end{proof}

\begin{corollary}[Rank-one overlap identity]
\label{p7v2:cor:rank-one-overlap}
On the full-algebra tracial branch of
Theorem~\ref{p7v2:thm:thermal-trace-nonunique}, the state is necessarily
\(\rho_{16}=I_{16}/16\), and every rank-one projection \(P\) satisfies
\begin{equation}
\label{p7v2:eq:rank-one-overlap}
 \Tr(\rho_{16}P)=\frac1{16}.
\end{equation}
Thus the \(1/16\) coefficient needs no additional state-weight selector on
that branch.  It is still an initial overlap probability, not a late-time
leptogenesis efficiency.
\end{corollary}

\begin{proof}
\(\Tr P=1\), so
\(\Tr(I_{16}P/16)=1/16\).  No Hamiltonian, production rate, collision term,
or expansion history appears in this identity.
\end{proof}

\begin{corollary}[Finite-shell origin of the tracial rank-one atom]
\label{p7v2:cor:rank-one-overlap-shell-origin}
On the full-algebra tracial branch,
\begin{equation}
\label{p7v2:eq:rank-one-overlap-shell-origin}
 \boxed{\Tr\!\left(\frac{I_{16}}{16}P\right)
 =\frac1{16}
 =\omega_{2T}
 =\frac1{|\mathcal G_{16}|}.}
\end{equation}
Thus the tracial rank-one coefficient equals the normalized counting atom of
the upstream Hurwitz shell complement.  The identity is mathematical.  A
general thermal state remains in the larger state family and need not be
tracial.
\end{corollary}

\begin{remark}[The overlap is not an equilibrium abundance]
For one relativistic Majorana state, a Maxwell--Boltzmann convention gives
\[
 Y_N^{\rm eq}(0)=\frac{45}{\pi^4g_*}
 \simeq\frac{0.46197}{g_*},
\]
while Fermi--Dirac statistics gives
\[
 Y_N^{\rm eq}(0)
 =\frac{135\zeta(3)}{4\pi^4g_*}
 \simeq\frac{0.41648}{g_*}.
\]
Both differ from \(1/(16g_*)=0.0625/g_*\).  Therefore the finite projector
overlap cannot be reinterpreted as the standard equilibrium yield.
\end{remark}

\begin{definition}[Minimal kinetic boundary problem]
\label{p7v2:def:kinetic-system}
A full two-state resonant calculation requires, at minimum, equations of the
form
\begin{align}
\label{p7v2:eq:kinetic-system}
 \frac{d\rho_N}{dz}
 &=
 -i[H_N(z),\rho_N]
 -\frac12\{\Gamma_N(z),\rho_N-\rho_N^{\rm eq}(z)\}
 +\mathcal C_N[\rho_N,Y_\Delta],\\
 \frac{dY_{\Delta_\alpha}}{dz}
 &=
 S_\alpha[\rho_N](z)
 -\sum_\beta W_{\alpha\beta}(z)Y_{\Delta_\beta},
 \qquad \alpha=e,\mu,\tau,
\end{align}
together with initial data, the \(N_1\) channel, spectator relations, and the
cosmological relation between \(z\), time, and temperature.  Their complete
specification is \(\Hkin\).
\end{definition}

\begin{theorem}[Efficiency is an output functional]
\label{p7v2:thm:efficiency-output}
For fixed microscopic source data, the late-time efficiency
\(\kappa_{\rm kin}\) is a functional of the initial state and of every kernel
in \eqref{p7v2:eq:kinetic-system}.  It cannot be derived from the rank of a
finite source projector alone.
\end{theorem}

\begin{proof}
Choose identical source projectors but three admissible washout kernels:
\(W=0\), \(W=W_0\), and \(W=cW_0\) with \(c\to\infty\).  The corresponding
solutions have, in general, no washout, finite suppression, and vanishing
late-time asymmetry.  The projector rank is unchanged.  Thus the efficiency
depends on the kinetic data and is not the trace identity
\eqref{p7v2:eq:rank-one-overlap}.
\end{proof}

\subsection{Thermal coefficients and baryon benchmark}
\label{p7v2:sec:baryon-benchmark}

\begin{proposition}[Standard hot Standard-Model count]
\label{p7v2:prop:gstar}
Assume all Standard-Model particles are relativistic and in a common thermal
bath with the equilibrium statistics used below and negligible chemical
potentials.  In this statement the three right-neutral slots of the full
Paper--V-A carrier are absent from the occupied relativistic bath (for
example by an actually established heavy or unequilibrated regime).
No other light species are counted.  Then
\begin{equation}
\label{p7v2:eq:gstar}
 g_*=
 g_{\rm bos}+\frac78g_{\rm ferm}
 =
 28+\frac78(90)
 =
 \frac{427}{4}
 =
 106.75.
\end{equation}
This is a status-\(\mathsf Q\) equilibrium field-content count.
\end{proposition}

\begin{proof}
The unbroken-phase gauge bosons contribute \(12\times2=24\) helicity states and
the complex Higgs doublet contributes \(4\), giving \(g_{\rm bos}=28\).
The Standard-Model Weyl fermions and antifermions contribute
\(g_{\rm ferm}=90\).  Fermionic energy density carries the factor \(7/8\).
\end{proof}

\begin{corollary}[Occupied right-neutral bath count]
\label{p7:cor:v2211-right-neutral-bath-count}
If, in addition to that Standard-Model bath, exactly \(n_R\) of the
constituted right-neutral Weyl slots and their antiparticle modes are
relativistic and equilibrated at the same temperature and negligible
chemical potential, then
\[
 \boxed{g_*(n_R)=\frac{427}{4}+\frac74n_R,
 \qquad n_R\in\{0,1,2,3\}.}
\]
In particular the fully occupied 48-Weyl carrier has \(g_*=112\).
The carrier count does not choose \(n_R\), a distribution or a temperature.
Massive, unequilibrated or differently heated slots consume their actual
energy-density integral instead of this common-temperature count.
\end{corollary}
\begin{proof}
One additional Weyl slot and its antiparticle contribute two relativistic
fermion states, hence \((7/8)2=7/4\).  The full carrier has 48 rather
than 45 Weyl slots.  Add their contributions only under the occupied-bath
hypotheses.
\end{proof}

\begin{proposition}[Sphaleron and photon conversions]
\label{p7v2:prop:standard-conversions}
Assume the standard high-temperature chemical-equilibrium conditions for three
families and one Higgs doublet, with the Standard-Model charged Yukawa and
sphaleron constraints active, and with no right-neutral Yukawa equilibrium
constraint modifying that system.  Here \(B-L\) denotes the asymmetry of
this active Standard-Model sector.  Then
\begin{equation}
\label{p7v2:eq:sphaleron}
 Y_B=\frac{28}{79}Y_{B-L}.
\end{equation}
Under the standard late entropy-to-photon conversion,
\begin{equation}
\label{p7v2:eq:entropy-photon}
 \eta_B\simeq7.04\,Y_B.
\end{equation}
These are continuum thermodynamic interfaces, not new finite-count
derivations.
\end{proposition}

\begin{proof}
Equation~\eqref{p7v2:eq:sphaleron} follows by solving the Standard-Model
chemical-potential constraints with sphalerons, hypercharge neutrality, and
equilibrated Yukawa interactions.  Equation
\eqref{p7v2:eq:entropy-photon} is the standard late cosmological conversion
between entropy-normalized and photon-normalized abundances
\cite{KapustaGale2006,PDG2024}.
\end{proof}

\begin{proposition}[Right-neutral spectator and chemical-equilibrium domains]
\label{p7:prop:v2211-right-neutral-spectator-domain}
Keep the charged-sector equilibrium domain of
Proposition~\ref{p7v2:prop:standard-conversions} and one common entropy
normalization.  If unequilibrated right-neutral modes merely store a
lepton asymmetry \(Y_{L_R}\), with no additional active exchange
constraint, then
\[
 \boxed{Y_B=\frac{28}{79}
 (Y_{B-L_{\rm tot}}+Y_{L_R}).}
\]
On a different, generation-equilibrated Dirac branch where all three
right-neutral Yukawa reactions are in chemical equilibrium, all relevant
species are relativistic, and no lepton-number-violating equilibrium
constraint is added, the corresponding result is
\[
 \boxed{Y_B=\frac14Y_{B-L_{\rm tot}}.}
\]
These are domain-specific consequences of the constituted charges.  They
do not establish either reaction equilibrium, a sphaleron history, a
right-neutral population or a baryon asymmetry.  Partial equilibration,
Majorana violation and additional spectators consume their actual charge
constraint system.
\end{proposition}
\begin{proof}
For stored unequilibrated charge,
\(B-L_{\rm active}=B-L_{\rm tot}+L_R\); substitute this identity in
the preceding conversion.  For the fully Dirac-equilibrated comparison,
write the generation-independent chemical potentials as
\(q,u,d,\ell,e,n,h\).  The constituted hypercharges, the charged and
neutral Yukawa slots, and the sphaleron constraint give
\[
 u=q+h,\quad d=q-h,\quad e=\ell-h,\quad n=\ell+h,
 \qquad 3q+\ell=0,
\]
\[
 3(q+2u-d-\ell-e)+2h=0.
\]
Right-neutral slots have zero hypercharge.  The factor two in the scalar
term is its Bose charge susceptibility relative to a Weyl state.  The
solution has \(\ell=-3q\), \(h=-12q/7\).  In the common fermionic
susceptibility unit,
\[
 B=3(2q+u+d)=12q,\qquad
 L_{\rm tot}=3(2\ell+e+n)=12\ell=-36q.
\]
Thus \(B-L_{\rm tot}=48q\), giving the ratio \(1/4\).
The same equations without \(n\) in the lepton readout give
\(B/(B-L_{\rm active})=28/79\).  No coefficient is chosen to replace
an unevaluated reaction domain.
\end{proof}

\begin{proposition}[Arithmetic status of the former baryon result]
\label{p7v2:prop:baryon-benchmark}
If one imposes
\begin{equation}
\label{p7v2:eq:baryon-benchmark-inputs}
 \varepsilon_{23}^{\rm cand}=4.1945421\times10^{-7},
 \quad
 \kappa_{\rm proj}=\frac1{16},
 \quad
 g_*=106.75,
 \quad
 c_{\rm sph}=\frac{28}{79},
 \quad
 s/n_\gamma=7.04,
\end{equation}
then the arithmetic chain is
\begin{align}
\label{p7v2:eq:baryon-benchmark-chain}
 Y_{B-L}^{\rm bench}
 &=
 \frac{\varepsilon_{23}^{\rm cand}}{16g_*}
 =2.45582090\times10^{-10},\\
 Y_B^{\rm bench}
 &=
 \frac{28}{79}Y_{B-L}^{\rm bench}
 =8.70417535\times10^{-11},\\
 \eta_B^{\rm bench}
 &=
 7.04\,Y_B^{\rm bench}
 =6.12773944\times10^{-10}.
\end{align}
The equalities are exact consequences of the five benchmark inputs.  They do
not constitute a baryogenesis prediction.
\end{proposition}

\begin{proof}
Substitute \eqref{p7v2:eq:baryon-benchmark-inputs} successively into the three
definitions in \eqref{p7v2:eq:baryon-benchmark-chain}.  The quantity
\(\varepsilon_{23}^{\rm cand}\) is an independently declared benchmark source;
no electromagnetic or Hopf formula is used to derive it.  The coefficient
\(1/16\) is the projector overlap of
Corollary~\ref{p7v2:cor:rank-one-overlap}, not the solution of
\eqref{p7v2:eq:kinetic-system}.  Therefore the numerical agreement is a
status-\(\mathsf B\) arithmetic benchmark.
\end{proof}

\begin{proposition}[Dark-radiation boundary]
\label{p7v2:prop:dark-radiation-boundary}
On \(\Urole\) and a compatible light-spectrum member
\(b_6\in\Acompsix\), the imported finite Standard-Model carrier contains
exactly the three named active light-neutrino species and no additional light
sterile carrier in that member.  This is a conditional finite-role statement,
not a thermal abundance theorem.  A prediction for \(N_{\rm eff}\) or
\(\Delta N_{\rm eff}\) additionally requires decoupling, reheating, and
entropy-transfer dynamics and therefore belongs to \(\Hkin\).
\end{proposition}

\begin{proof}
Carrier counting determines which light species are available.  Their
cosmological energy densities depend on distribution functions and
temperatures, which are solutions of a thermal history rather than carrier
dimensions.
\end{proof}

\subsection{native BL2 owner closure and downstream CP frontier}
\label{p5c:sec:v272-bl2-owner-closure}

\begin{proposition}[Native charge-two boundary owner is co-realized]
\label{p5c:prop:v272-bl2-owner-corealized}
Paper~V-A,
Theorem~\ref{p5:thm:v3105-k16-bl2-pair-corealization}, closes the generated
neutral-right K16 pairing member on the already occurring particle/internal
boundary owner.  Hence the historical v2.70 status is superseded by
\[
 \boxed{
 \texttt{BL2\_BOUNDARY\_ACTION\_OCCURRENCE}
 =\texttt{CLOSED\_COREALIZED}.}
\]
The old Dirac-only zero-projection statement remains a valid comparison
subbranch, not the active pair-owner branch.
\end{proposition}

\begin{corollary}[First live native baryogenesis frontier after occurrence]
\label{p5c:cor:v272-native-baryogenesis-phase-frontier}
Occurrence and pair normalization introduce no additional efficiency or
family matrix.  The first live native baryogenesis obligation is now
\[
 \boxed{
 \texttt{BL2\_PHASE\_LOCK}=\texttt{OPEN\_OWNER\_BINDING},}
\]
meaning the physical relative connection of the protected source section and
the neutral pair line.  It is followed by the complete CP/state-history and
washout/entropy readout.  The finite charge-propagation functional is already
determined once its owned data are fixed, as shown in
Theorem~\ref{p5c:thm:v273-bl2-charge-response}.  A constant complex coordinate
of the generated pair is not by itself a CP-odd history.
\end{corollary}

\begin{remark}[Heavy-neutral branch remains conditional]
\label{p5c:rem:v272-heavy-neutral-still-conditional}
The type-I/seesaw/heavy-Majorana decay machinery remains available on its own
independently occurring high-energy neutral owner, but it is not required to
supply the now-closed native K16 charge-two boundary pair.  Its CP/washout
history remains a conditional downstream branch.
\end{remark}
% END BL2-OWNER-CLOSURE-IMPORT v2.72


% BEGIN BL2-FINITE-CHARGE-RESPONSE v2.73
\subsection{finite BL2 charge response and relative phase connection}
\label{p5c:sec:v273-bl2-charge-response}

\begin{theorem}[Dependent finite BL2 charge-response closure]
\label{p5c:thm:v273-bl2-charge-response}
Retain the occurring K16 neutral pair owner of Paper~V-A,
Theorem~\ref{p5:thm:v3105-k16-bl2-pair-corealization}, its support insertion,
kinetic metric and connection.  On a declared clocked history set
\[
 A_{BL2}(t)=\mathfrak a^\dagger\Delta_{BL2}(t)\mathfrak a^\dagger,
 \qquad H(t)=H_0(t)+\frac12(A_{BL2}+A_{BL2}^\dagger),
\]
and suppose $[Q_{B-L},H_0]=0$.  On the neutral-right convention in which each
created neutrino has $B-L=-1$,
\[
 [Q_{B-L},A_{BL2}]=-2A_{BL2};
\]
the coefficient's compensating source weight is $+2$.  Paper~IV,
Theorem~\ref{p4:thm:v4102-pair-charge-source}, therefore gives
\begin{equation}
\label{p5c:eq:v273-bl2-charge-source}
 \boxed{
 \frac d{dt}\langle Q_{B-L}\rangle
 =-\frac2\hbar\operatorname{Im}\langle A_{BL2}\rangle
 =\frac2\hbar\langle\partial_\theta H_\theta\rangle.}
\end{equation}
The initial state, Hamiltonian history and retained connection determine the
exact finite charge functional.  On the quadratic CAR branch its covariance
propagation and readout are precisely
Eqs.~\eqref{p4:eq:v4102-covariance-propagation}--\eqref{p4:eq:v4102-covariance-charge-readout}.
Consequently
\[
 \boxed{\texttt{BL2\_FINITE\_CHARGE\_RESPONSE}
 =\texttt{CLOSED\_DEPENDENT[HISTORY,STATE,CONNECTION]}.}
\]
There is no independent efficiency multiplying this same-owner charge source.
\end{theorem}

\begin{proof}
The pair and its conjugate are already required constituents of the occurring
owner.  Their charge commutators fix the source sign and coefficient in the
Paper--IV Ward identity.  Finite propagation is unique after its history,
connection and initial state are supplied.  Paper~II-A generated
co-consumption preserves these same-owner descendants and excludes an
independent rescaling.  This is a dependent response closure, not selection
of the history, state or relative phase connection.
\end{proof}

\begin{corollary}[Covariant phase target on the fixed-owner pair branch]
\label{p5c:cor:v273-bl2-relative-phase-target}
Paper~III-F's fixed-owner successor-cell transport supplies the radial pair
response $\Delta_{BL2}=\eta\Delta_*$ on its declared fixed-datum domain.
For the purely radial, parallel fixed-member specialization
$D_K\Delta_*=0$, one obtains
\[
 \operatorname{Im}\langle\Delta_{BL2},D_K\Delta_{BL2}\rangle_{\rm HS}=0.
\]
The scalar response law alone does not determine the retained connection
along a physical history.
For a separately specified physical phase $\varphi(K)$ and connection,
one may write $\Delta_{BL2}=a e^{i\varphi}\Delta_*$ and
$D_K\Delta_*=-i\mathcal A_K\Delta_*$.  The general relative diagnostic then is
\[
 \boxed{J_K=\|\Delta_*\|_{\rm HS}^2a^2
                    (\varphi'-\mathcal A_K),}
\]
by Paper~IV, Corollary~\ref{p4:cor:v498-phase-locked-pairing-invariant}.
This is a conditional covariant identity, not an assignment of $\varphi$
from Reeb transport or the octet clock.  On the fixed-owner radial/parallel
branch no additional pair phase is generated, as the subsequent active
retirement corollary makes explicit.
\end{corollary}

\begin{remark}[Scope of the positive contraction]
\label{p5c:rem:v273-bl2-response-scope}
Vanishing radial $J_K$ excludes only the pair's self-trajectory phase
diagnostic, not all relative-sector CP effects or every possible asymmetry.
The anomalous state expectation in Eq.~\eqref{p5c:eq:v273-bl2-charge-source}
must still be propagated.  With quartic interactions use the exact finite
many-body evolution, not a closed quadratic covariance equation.
Initial-state/history occurrence, collision or washout reduction, spectator
maps and final entropy normalization remain where separately required; no
numerical $Y_{B-L}$ follows here.  In particular
$\|\Delta_*\|_{\rm HS}^2$ retains the actual support insertion and kinetic
metric and is not silently replaced by the family norm
$\|M_{BL2}^{\rm geom}\|_{\rm HS}^2=3/2$.
\end{remark}
% END BL2-FINITE-CHARGE-RESPONSE v2.73




% BEGIN INLINE SOURCE: Paper5C-OrderedCurrentConvention-Main-v1_0.tex
% Framework v2.167: explicit quartet dictionary; corrects both branch signs.
\subsection{native baryogenesis frontier and cyclic-state synchronization}
\label{p5c:sec:v280-baryogenesis-cyclic-state-sync}

\begin{corollary}[Current native baryogenesis frontier on the cyclic lineage]
\label{p5c:cor:v280-native-baryogenesis-frontier}
Retain the co-realized K16 neutral-right pair owner of
Paper~V-A, Theorem~\ref{p5:thm:v3105-k16-bl2-pair-corealization}, the exact
finite charge Ward law of Theorem~\ref{p5c:thm:v273-bl2-charge-response}, and
the native charged-current invariant of
Corollary~\ref{p5c:cor:v277-native-lepton-cp-value}.  The active finite inputs
therefore include
\[
\begin{gathered}
 \texttt{BL2\_BOUNDARY\_ACTION\_OCCURRENCE}
 =\texttt{CLOSED\_COREALIZED},\\
 J_\ell^{\rm native,ord}=-\frac1{12},\qquad |J_\ell|=\frac1{12}.
\end{gathered}
\]
The historical fixed-pair self-phase target is retired as stated in
Corollary~\ref{p5c:cor:v274-bl2-phase-lock-retired}.

On the native cyclic lineage, Paper~VI-B,
Corollary~\ref{p9:cor:v187-k16-state-inheritance}, gives
\[
 \omega_{16}=\mathcal U(t_{16},t_0)_*\omega_0,
 \qquad
 \mathsf{State}^{\rm G0}_{16}
 =\texttt{RETIRED\_AS\_INDEPENDENT\_ON\_CYCLIC\_LINEAGE}.
\]
Hence the state argument in the dependent BL2 charge functional is supplied
upstream on that lineage rather than constituting a second K16 preparation or
a new baryogenesis selector.  This synchronization does not remove the state
premise on generic/noncyclic branches.

The first native incidence test belongs to the earlier complete common
interaction word, before kinetic reduction: does that physical action consume
the already closed relative lepton-current CP incidence and the already
occurring BL2 charge-two pair in one carrier/history, retaining every nonzero
mixed block?  If it does, Paper~III-A,
Corollary~\ref{p3a:cor:complete-parent-descendant-reuse}, and claim-complete
action co-consumption close that binding without a second occurrence, phase
selector or efficiency.  The remaining task is evaluation of the CP-sensitive
charge response and any actually consumed kinetic reduction.
The current charge theorem's general \(H_0\) and the visible interaction
graph do not certify this particular neutral-right joint word.  The actual
K16 restriction must be exhibited on the inherited complete field/source
diagram; the already proved component-band family retention is reused on
its domain.  A full unfactorized response can be evaluated directly without
an additional intertwining or factorization prerequisite.  For a claimed
signed charge history, Corollary~\ref{p5c:cor:v2186-direct-signed-charge-target}
also permits direct complete-state evaluation instead of a separate marked
incidence proxy.  A zero evaluated charge is retained; it does not establish
a positive asymmetry.  No actual K16 restriction or charge history is
evaluated by these statements.  Hence
\[
 \boxed{
 \texttt{CP\_BL2\_KINETIC\_INCIDENCE}
 =\texttt{OPEN[COMPLETE\ PARENT\ ACTION\ INCIDENCE]}.}
\]
Collision/washout or spectator reductions are consumed only by claims using
those reduced descriptions.  A final charge/entropy endpoint additionally
consumes its actual material normalization and continuation.  No perturbative
Dyson order, mixed CP invariant, washout coefficient or numerical $Y_{B-L}$ is
fixed by this corollary.
\end{corollary}

\begin{remark}[Native finite charge and final-yield scope]
\label{p5c:rem:v2187-native-charge-yield-scope}
For an actual complete finite interaction word, Paper~IV,
Theorem~\ref{p4:v2171:thm:fullcharge}, already generates
\(\mathcal J_Q=D_tQ+i[H,Q]/\hbar\) and its exact signed expectation history.
Represented redistribution and relaxation are part of that same evolution;
no independent kinetic law or efficiency is inserted at the charge consumer.
A finite-event charge target does not require a thermal-equilibrium or
freeze-out construction.  A positive baryon yield additionally consumes an
actually generated baryon-export law and its same-history expectation:
Theorem~\ref{p3f:v2171:thm:baryonstationary} gives zero baryon charge on its
complete baryon-equivariant split-plus-BL2 branch.  Any regulated export must
retain its actual Jacobians, contacts, residues and boundary flux under
Paper~VII-C, Proposition~\ref{p12:v2171:prop:baryonexport}; its generator and
current representations are not two additive sources.  The final material
entropy and continued endpoint remain the same-history readouts of
Paper~VI-B.  These are actual construction/execution obligations, not a
phenomenological replacement mechanism, and no nonzero yield is asserted.
\end{remark}

\begin{remark}[Optional heavy-neutral branch and generic-state scope]
\label{p5c:rem:v280-heavy-neutral-generic-scope}
The type-I/heavy-neutral decay construction remains an optional conditional
branch when its own high-energy neutral owner occurs.  The native BL2 branch
does not require that branch to constitute its charge-two pair.  On branches
not covered by the cyclic-lineage root-state theorem, the initial state and
history remain explicit inputs to Theorem~\ref{p5c:thm:v273-bl2-charge-response}.
\end{remark}
% END BARYOGENESIS-CYCLIC-STATE-SYNC v2.80

\subsection{Kinetic normalization precedes a charge-history claim}
In fixed source field coordinates let a raw normal coefficient be $h_b$,
a raw pair coefficient $\Delta_b$, and the actual positive kinetic metric
$G$. After the existing metric construction, $\psi=Z\chi$ with
$Z=G^{-1/2}$ gives
\begin{equation}
 h_c=Z^\dagger h_bZ\quad\hbox{plus the inherited clock/connection terms},
 \qquad
 \Gamma_c=Z^\dagger\Delta_b\overline Z.
 \label{p5c:v2171:eq:normalization}
\end{equation}
An actual embedding is included in $Z$ before this formula is applied.
It must not be replaced by a comparison embedding. Pair coefficients use
creation-slot congruence, while normal coefficients use ordinary Hermitian
pullback. Source variations differentiate the whole expressions.

Therefore \eqref{p8:v2171:eq:chargeblock} is not by itself a theorem that
$\Gamma_c(t)=\eta(t)\Gamma_c(t_{16})$. That stronger equality follows on
a certified fixed normalized field frame, or from the actual metric formula;
it is not an additional independent profile to choose. For raw fixed forms
that both scale by $\eta$, their four-slot mixed contraction scales by
$\eta^4$. This raw identity is not promoted to a CP-sensitive charge
coefficient after arbitrary kinetic normalization.

The local rooted family source invariant remains
\begin{equation}
 Z_{\mu;\mathrm{inv,dir}}
 =\Tr(M^\dagger N_{\rm inv}P_\mu N_{\rm dir}
       M(N_{\rm inv}P_\mu N_{\rm dir})^{\mathsf T})
 =-\frac7{288}+\frac{i}{12}.
 \label{p5c:v2171:eq:family}
\end{equation}
The even stationary reduction above does not introduce a new family projection
on the retained local fermion fields. It identifies their quadratic source
coefficient at the actual stationary background. Its complete physical use
still retains the actual field trace, spin/mode coefficients, metric,
connection and every consumed source-background derivative. An unknown
physical contraction is not replaced by the number in \eqref{p5c:v2171:eq:family}.

\subsection{visible high-temperature interaction connectivity}
\label{p5c:sec:v269-k16-visible-connectivity}

\begin{theorem}[Connected visible interaction graph]
\label{p5c:thm:v269-visible-interaction-graph}
On the generated Standard--Model interaction domain consider the block nodes
\[
 \mathcal V_{\rm vis}
 =\{g,W,B,H,Q_L,u_R,d_R,L_L,e_R\}.
\]
The rooted gauge action supplies nonzero QCD edges from $g$ to
$Q_L,u_R,d_R$, weak edges from $W$ to $Q_L,L_L,H$, and hypercharge edges
from $B$ to every hypercharged matter/Higgs block.  The charged Yukawa
supports additionally contain
\[
 \bar Q_L\widetilde H u_R,\qquad
 \bar Q_LH d_R,\qquad
 \bar L_LH e_R.
\]
Therefore the visible interaction graph is connected.
\end{theorem}

\begin{corollary}[Microscopic criterion for visible primitivity modulo exact charges]
\label{p5c:cor:v269-visible-primitivity}
The connected block graph of Theorem~\ref{p5c:thm:v269-visible-interaction-graph}
is a structural connectivity statement, \emph{not} by itself a scalar-commutant
certificate on the occurring microscopic Hilbert space. For an exact
conserved-charge component \(\mathcal H_q\), let \(\mathcal M_q\) be the
unital \(*\)-algebra generated by the \emph{actual} same-owner Hamiltonian
and microscopic positive-rate transition operators and their adjoints.
If the following operator-level condition is verified,
\[
 \boxed{\mathcal M_q'=\mathbb C\mathbf1_{\mathcal H_q}},
\]
then its microscopic interaction algebra is irreducible on \(\mathcal H_q\).
A sufficient finite certificate is a resolved orthonormal basis
\(\{|i\rangle\}_{i=1}^{d_q}\) whose rank-one diagonal projectors
\(|i\rangle\langle i|\) and rank-one edge matrix units
\(|i\rangle\langle j|\) on a connected \emph{microscopic} transition graph
belong to \(\mathcal M_q\). If, moreover, the occurring finite
detailed-balanced quantum Markov semigroup has a faithful stationary state,
is reversible on this block, and has fixed-point algebra precisely
\(\mathcal M_q'\), then it is primitive/ergodic on this exact-charge block.
These \emph{microscopic} conditions are actual operator/source obligations;
coarse species connectivity alone does not establish them.
\end{corollary}
\begin{proof}
Products of adjacent rank-one matrix units along a path give
\(|i\rangle\langle j|\) for arbitrary \(i,j\). The connected microscopic
unit graph and the diagonal projections therefore generate
\(\operatorname{End}(\mathcal H_q)\), whose commutant is scalar. On the
stated reversible finite semigroup branch the generator has real
nonpositive spectrum in the stationary inner product; a one-dimensional
fixed-point space then implies convergence to the unique faithful
stationary state. In contrast, on
\(\mathbb C^2_{\rm species}\otimes\mathbb C^2_{\rm spectator}\), the jumps
\(|b\rangle\langle a|\otimes I\) and their adjoints connect the species
nodes while every \(I\otimes B\) commutes with the jump algebra. Thus
species-graph connectivity alone cannot prove the original irreducibility
or primitivity assertion.
\end{proof}

\begin{remark}[Gauge-singlet heavy neutral]
\label{p5c:rem:v269-heavy-neutral-firewall}
A gauge-singlet heavy neutral $N_R$ is not needed for visible-bath
connectivity.  It joins the visible graph only through its separately declared
neutral/Yukawa lift, e.g.
\[
 \bar L_LY_\nu\widetilde HN_R.
\]
Consequently visible thermalization does not imply heavy-neutral thermalization;
that sector may remain a nonthermal input to the later K24 $B-L$ kinetics.
\end{remark}
% END K16-VISIBLE-CONNECTIVITY v2.69

% BEGIN BOUNDARY-BARYOGENESIS-ANCESTRY-RETYPING v2.70
\subsection{boundary-baryogenesis ancestry retyping}
\label{p5c:sec:v270-boundary-baryogenesis-ancestry}

\begin{proposition}[Native baryogenesis first missing owner]
\label{p5c:prop:v270-native-baryogenesis-first-owner}
Paper~V-A supplies the generated structural coefficient
$M_{BL2}^{\rm geom}$ on the unique neutral-right $|\Delta(B-L)|=2$ support,
while Paper~III-F proves that the active Dirac-only internal neutral response
has vanishing charge-two projection.  Hence the first missing native
baryogenesis arrow is
\[
 \boxed{
 \texttt{BL2\_BOUNDARY\_ACTION\_OCCURRENCE}=\texttt{OPEN\_ACTION\_OWNER}.}
\]
Until this owner occurs, no choice of history profile, density matrix,
Bogoliubov basis, CP coordinate or washout kernel can generate a nonzero
$B-L$ yield from the charge-preserving branch.
\end{proposition}

\begin{remark}[Heavy-neutral/type-I branch is conditional, not the native owner]
\label{p5c:rem:v270-heavy-neutral-conditional-branch}
The type-I/seesaw/heavy-Majorana decay formulas retained in this paper remain
valid conditional comparison/dynamical machinery when a corresponding
high-energy neutral interaction owner is independently generated or supplied.
They do not constitute the missing native charge-two boundary action owner.
Accordingly a K16--K24 heavy-neutral CP/washout history is retained as an
optional downstream branch rather than the unique or prerequisite origin of
baryogenesis.
\end{remark}
% END BOUNDARY-BARYOGENESIS-ANCESTRY-RETYPING v2.70

% BEGIN BARYOGENESIS-INCIDENCE-CORRECTION v2.74
\subsection{baryogenesis incidence correction and native lepton CP target}
\label{p5c:sec:v274-baryogenesis-incidence-correction}
\noindent\textbf{Current scope.} The current physical singlet/plane branch is the tau-singlet branch derived in main-text Release v2.76.  The v2.74 mu-singlet Jarlskog formula is retained as a conditional algebraic branch and is not the active physical placement.

\begin{corollary}[Lepton weak-current existence is not its physical evaluation]
\label{p5c:cor:v274-lepton-current-formal-open}
The identity
\[
 C_{e\nu}=W_e^\dagger\mathscr J_WW_\nu=W_e^\dagger W_\nu
\]
remains the framework-native weak-current readout.  Paper~V-B Release v5.95
corrects the scope of the charged factor: the printed source chain defines
and normalizes $W_e$ for a supplied role-resolved evaluation but does not yet
evaluate that physical map.  Therefore the stronger v2.67 status is
superseded by
\[
 \boxed{
 \texttt{LEPTON\_WEAK\_CURRENT\_INCIDENCE}
 =\texttt{CLOSED\_FORMAL/OPEN\_PHYSICAL\_EVALUATION}.}
\]
No PMNS matrix or new family selector is introduced by this retyping.
\end{corollary}

\begin{theorem}[Current-based native leptonic CP invariant]
\label{p5c:thm:v274-native-lepton-cp-invariant}
Let
\[
 Q_e=\sum_{a=e,\mu,\tau}q_{e,a}P_a,
 \qquad
 Q_\nu=\sum_{r}q_{\nu,r}\Pi_r^{\rm role}
\]
be the nondegenerate charged and neutral terminal role-weight operators, and
let $C_{e\nu}$ be the physical current incidence above.  Then
\[
 \boxed{
 \chi_\ell^{\rm cur}:=\frac1{6i}\operatorname{Tr}
 [Q_e,C_{e\nu}Q_\nu C_{e\nu}^\dagger]^3}
\]
is invariant under simultaneous family-presentation changes and under
independent role-line rephasings.  For simple spectra,
\[
 \boxed{
 \chi_\ell^{\rm cur}=\Delta_e\Delta_\nu J(C_{e\nu}),}
\]
with the Vandermonde factors taken in the declared role ordering.  Thus the
CP constructor is fixed even though its physical value still depends on the
unevaluated charged-lepton incidence.
\end{theorem}

\begin{proof}
The commutator transforms by simultaneous conjugation, so its cubic trace is
presentation invariant.  Paper~IV's rank-three normalized Jarlskog identity
gives the second formula.
\end{proof}

\begin{corollary}[One-scalar CP reduction on the singlet/plane branch]
\label{p5c:cor:v274-zeta-jarlskog-reduction}
Use the declared neutral role order $(\mathrm{inv},\mathrm{dir},\mathrm{Gram})$
and the rooted presentation in which $W_\nu$ is the real odd permutation
$(g_2,g_1,g_0)$.  On the partial charged placement
\[
 P_\mu=Q_0,\qquad P_e+P_\tau=Q_\perp,
\]
let
\[
 \zeta_e=\frac1{\sqrt3}\operatorname{Tr}(P_eC_{2,G}).
\]
Then, with the fixed role-order convention,
\[
 \boxed{J(C_{e\nu})=+\frac{\zeta_e}{6\sqrt3}.}
\]
Equivalently the full continuous counterfamily of Paper~V-B has
$J=+(6\sqrt3)^{-1}\cos2\theta$.  Hence
\[
 \boxed{
 \texttt{LEPTON\_CP\_CONSTRUCTOR}=\texttt{CLOSED},\qquad
 \texttt{LEPTON\_CP\_VALUE}=\texttt{OPEN\_ON\_ZETA\_E}.}
\]
\end{corollary}

\begin{proof}
Choose the normalized $C_3$ character basis of $Q_\perp$ used in Paper~V-B
v5.95, with the ordered dictionary now made explicit in
Lemma~\ref{p5c:lem:v2167-ordered-branch-signs}. Substitution of its $U(2)$ counterfamily into
$C_{e\nu}=W_e^\dagger W_\nu$ gives
$J(C_{e\nu})=+\sqrt3\cos2\theta/18$, while
$\zeta_e=\cos2\theta$.
\end{proof}

\begin{corollary}[BL2 phase-lock target retired]
\label{p5c:cor:v274-bl2-phase-lock-retired}
The fixed-owner BL2 response already has the real form
$\Delta_{BL2}=\eta\Delta_*$ on its declared radial domain.  By Paper~IV
Release v4.103 this response adds no independent covariant pair phase, and on
the evaluated parallel branch its self-trajectory diagnostic vanishes.
Appending a protected scalar phase to this same frozen pair would define a
new deformation rather than a consequence of the radial transport.  Therefore
\[
 \boxed{
 \texttt{BL2\_SELF\_PHASE}=\texttt{CLOSED\_ZERO\_ON\_FIXED\_PARALLEL\_BRANCH},}
\]
\[
 \boxed{
 \texttt{BL2\_PHASE\_LOCK}=\texttt{RETIRED\_WRONG\_TARGET}.}
\]
This supersedes Corollary~\ref{p5c:cor:v273-bl2-relative-phase-target} as the
active frontier statement.  The exact BL2 charge Ward law and finite CAR
propagation remain unchanged.
\end{corollary}

\begin{corollary}[Correct native baryogenesis frontier]
\label{p5c:cor:v274-native-baryogenesis-frontier}
The source architecture is now typed as
\[
 |\Delta(B-L)|=2\quad\text{from the BL2 pair},
\]
independently of
\[
 \chi_\ell^{\rm cur}\quad\text{from relative charged/neutral current incidence}.
\]
The next source-level CP obligation is the owner evaluation of $\zeta_e$, not
a self-phase of the BL2 pair.  The subsequent joint CP--BL2 kinetic incidence,
initial state, washout and final charge/entropy history remain downstream.
\end{corollary}
% END BARYOGENESIS-INCIDENCE-CORRECTION v2.74


% BEGIN G17-NO-FREE-MATCHING v2.75

\section{UV scalaron sources and cycle-free cosmological exports}
\label{p5c:sec:organized-cosmic}
\subsection{Cycle-free cosmology handoff}
\label{p7v2:sec:cosmology-handoff}

\begin{definition}[Paper--V-C thermal export]
\label{p7v2:def:thermal-export}
The cycle-free thermal export is
\begin{equation}
\label{p7v2:eq:thermal-export}
 \mathfrak S_{5C}^{\rm th}
 =
 \left(
  \mathcal H_N,\,
  \mathcal M_N^{\rm aff},\,
  D_\nu,\,
  C_{e\nu},\,
  g_*^{\rm SM},\,
  c_{\rm sph}^{\rm SM},\,
  \text{obligations }\Hnu,\Hkin
 \right).
\end{equation}
The carrier \(\mathcal H_N\) is theorem-level.  The displayed
\(\mathcal M_N^{\rm aff}\) and \(D_\nu\) retain their stated conditional Paper--V-B status,
while \(C_{e\nu}\) is the generated native lepton current map.  A conventional
\([U_{\rm PMNS}]\) may be introduced only inside an explicitly external comparison
dictionary.  The export does not identify \(E_N\) with a reheating temperature,
select a thermal state or efficiency, or determine \(\eta_B\).
Its \(g_*^{\rm SM}\) slot consumes precisely the occupied-bath domain of
Proposition~\ref{p7v2:prop:gstar}; additional occupied right-neutral modes
instead use Corollary~\ref{p7:cor:v2211-right-neutral-bath-count} or their
actual energy-density integrals.  A total rather than active-sector
\(B-L\) readout also consumes
Proposition~\ref{p7:prop:v2211-right-neutral-spectator-domain} on its
actual equilibrium/exchange domain.
\end{definition}

\begin{proposition}[No reheating backflow]
\label{p7v2:prop:no-reheating-backflow}
Paper~VI-A may compute a reheating history and temperature using the
source-side gravity packet and the Standard-Model count
\eqref{p7v2:eq:gstar}.  That downstream temperature may then be tested against
or matched to \(E_N\) in a later synthesis.  It cannot be imported into
Paper~V-C to prove that the heavy anchor was already fixed upstream.
\end{proposition}

\begin{proof}
Paper~VI-A consumes \(g_*^{\rm SM}\) from the Paper--V-C boundary.  Using its
output to establish a premise of that same boundary would create a directed
cycle.  The acyclic order \eqref{p7v2:eq:dependency-order} permits only a later
matching theorem that names both independently obtained quantities.
\end{proof}

\begin{remark}[Chronological Paper--VI-A background discharge]
At the Paper--V-C stage, the background geometry is part of the kinetic/
cosmology interface.  On the current later branch, Paper~VI-A independently
derives the closed-FLRW geometric background.  That later result discharges
only the background-geometry component of \(\Hkin\); it does not select the
thermal-state member, decay law, collision or washout kernels, spectator map,
or baryogenesis efficiency.  The Paper--V-C ledger therefore remains
historically unresolved in those components.
\end{remark}

\label{p7v2:part:integrated}

\subsection{Coherent UV scalaron sources on the closed carrier (framework v2.219)}
\label{p5c:sec:v2219-coherent-source}
On the conditional two-sided UV quadratic branch the conformal Weyl
coordinate is $\sigma=\phi/(\sqrt6\bar M_{\rm Pl})$.  For
minimally coupled Jordan matter the leading Einstein-frame coupling is
$-\phi T/(\sqrt6\bar M_{\rm Pl})$ in the stated metric convention.
It differs from the previously constituted breathing spurion
$\tau=\alpha^{-1/2}\phi/\bar M_{\rm Pl}$, whose declared gauge-pair
coefficient has the ratio
\[
 \boxed{\lambda_i^\tau/\lambda_i^{\rm geom}
 =\sqrt6\,\alpha^{-1/2}\simeq6.18\times10^{-9}.}
\]
The nonconformal photon/Higgs/other particle channels therefore require
the actual same-action source map $\partial_\phi K_P$ and the
already constitued $\partial_\eta K_P$, not a guessed equality of
breathing and Weyl coordinates.  Closed $S^3$ expansion by itself
creates no photons in a constant-normalization free conformal Maxwell
sector.  Creation of charge-conjugate pairs alone supplies neither CP
violation nor a baryon-number asymmetry; the baryogenesis obligation
remains separate from energy production.


% BEGIN FRAMEWORK v2.228: CHARGED ORIGIN AND CHARACTER TESTS

\section{Observable ledger, action hypotheses and open measurements}
\label{p5c:sec:organized-results}
\subsection{Complete result ledger}
\label{p7v2:sec:result-ledger}

\begin{longtable}{@{}L{0.22\textwidth}L{0.17\textwidth}L{0.53\textwidth}@{}}
\caption[Complete Paper--V-C result ledger]{Status of the interaction claims
after the rigorous revision.}\label{p7v2:tab:result-ledger}\\
\toprule
Sector & Status & Result\\
\midrule
\endfirsthead
\toprule
Sector & Status & Result\\
\midrule
\endhead
Active realization
& \(\mathsf T\)
& The canonical support chain is
\(2T\mapsto K=32\), \(C_6\mapsto K=40\), and scalar reclosure
\(\mapsto K=56\); the exact heavy/light split grammars are imported from
Papers~I-A, III, and V-B.  Scalar reclosure is
\(\mathbf1_{\{7\}}(\mathcal Q_{\rm act})=P_0=\mathbb E_{\rm IR}\), so the
infrared tag is a derived certification.\\
Defect counts
& \(\mathsf T\)
& Equations~\eqref{p7v2:eq:defect-counts} and
\eqref{p7v2:eq:defect-sum-difference}.\\
Strong seed
& \(\mathsf T\) on the common-heat branch; \(\mathsf A\) on the broader kinetic branch
& Common heat gives \(k_i g_i^2=4\pi^2\) and
\(\alpha_s^{\rm heat}=\pi/2\) at its native matching datum.  The old
\(72^{-t}\) family is only a broad-branch calibration chart; no count theorem
selects \(t=1/2\).\\
Conventional one-doublet comparison dictionary
& \(\mathsf Q\) on a supplied stationary comparison chart
& Exact standard-coordinate mass--coupling reconstruction, with \(g_Y\) and
GUT-normalized \(g_1\) distinguished.  This row is not the native cause of
electroweak reduction.\\
One-loop gauge flow
& \(\mathsf Q\)
& Standard coefficients from the imported representation traces after a local
QFT and scheme are supplied.\\
Born thresholds
& \(\mathsf B/\mathsf O\)
& Reproducible near-unification audit; physical heavy determinant absent.\\
Electromagnetic endpoint
& \(\mathsf A\)
& Exact finite candidate \(\Aem(a)\); operator, trace selection, branch, and
analytic tail unresolved.\\
Electromagnetic running
& \(\mathsf Q/\mathsf O\)
& Perturbative vacuum-polarization interface retained; hadronic response is a
downstream QCD output.\\
QCD representation
& \(\mathsf T\) on \(\Urole\)
& The abstract \(3\otimes\bar3\) algebra is structural; its named physical
\(SU(3)\) color role and color factors use the role bridge.\\
QCD flow and kernels
& \(\mathsf Q\)
& Perturbative beta and factorization equations conditional on the continuum
QFT and nonperturbative boundary state.\\
Primitive PDFs
& \(\mathsf A\)
& Color singletness does not select momentum measures.\\
Strong CP
& \(\mathsf T\) native finite; \(\mathsf Q/\mathsf O\) chiral completion
& Native complete action and rooted determinant orientation give
\(\zeta_{\rm CP}=1\); a separately claimed chiral-continuum completion
retains its own regulated measure certificate.\\
Electroweak potential
& \(\mathsf Q/\mathsf O\), claim-scoped scalar target
& The occurring-solder stabilizer and gauge Gram are already native.
Potential reconstruction uses the same-action coefficients in its declared
scalar-normalization coordinate; actual extrema/history do not follow from
the gauge-Gram kernel.\\
Charged physical slots
& \(\mathsf T\) on
  \(\Urole+(b_6\in\Acompsix)+\mathsf F+\mathsf{G0}_{\rm odd}\)
& The Paper--V-A Yukawa supports, the physical Higgs weak lines, and charged
family spectral projectors derive the charged lookup-to-matter embeddings;
no independent charged \(\mathsf{SLOT}_{\rm phys}\) remains downstream.
On the current rooted charged branch, source and assembly are reused from
Paper~V-A, Corollary~\ref{p5:cor:v397-charged-preterminal-assembly-closure};
this generic conditional row adds no new source selector or pole-mass claim.
Degenerate family spectra retain only their spectral subspaces.\\
Higgs modifiers
& \(\mathsf B/\mathsf O\)
& Mass-boundary ratios retained under new names; physical vertices require
renormalized correlators.\\
Rare processes
& \(\mathsf Q\)
& Standard consequences conditional on the named Standard-Model or parent
effective theory.\\
Seesaw
& \(\mathsf T\)
& Exact realization identity and a bijection of the positive nonsingular fiber
with \(O(3,\CC)\); orthogonality is not a second premise.\\
Resonant source
& \(\mathsf T/\mathsf A_{\rm opt}\)
& The punctured-domain bound is exact and the regulator is a readout once
masses and widths are physical; saturation remains an optional ansatz.\\
Hopf factor
& \(\mathsf T/\mathsf A_{\rm opt}\)
& Former coisometry is dimensionally impossible; label and Hilbert traces are
distinct candidate measures, and no Hopf premise is needed by the standard
baryogenesis target.\\
Heavy-cell solve
& \(\mathsf T/\mathsf B\)
& The angle is closed form and exact saturation reduces to a monotone scalar
cubic with a unique certified root; source and saturation remain optional
benchmark inputs.\\
Thermal \(1/16\)
& \(\mathsf T\) on the full-algebra tracial branch
& Block traces and full-algebra invariant states are separated.  Full-algebra
traciality derives \(w=7/16\) and the rank-one overlap \(1/16\), but no
late-time efficiency.\\
Baryon number
& \(\mathsf B/\mathsf O\)
& \(6.12774\times10^{-10}\) is the arithmetic benchmark; source and kinetic
evolution remain open.\\
\bottomrule
\end{longtable}

\subsection{Selector and interface ledger}
\label{p7v2:sec:selector-ledger}

\begin{longtable}{@{}L{0.20\textwidth}L{0.25\textwidth}L{0.47\textwidth}@{}}
\caption[Paper--V-C selector ledger]{Selector-specific partition after
synchronization with Papers~V-A--V-B.}\label{p7v2:tab:selector-ledger}\\
\toprule
Class & Family/status & Paper--V-C meaning\\
\midrule
\endfirsthead
\toprule
Class & Family/status & Paper--V-C meaning\\
\midrule
\endhead
Exact finite core
& \(N_{\rm structural}=0\)
& \(\Istr\), active anchors, defect counts, structural color incidence, and
exact algebraic inequalities/trace distinctions require no member selector.
Named physical roles use \(\Urole=\Ustr\); no representation-role residual is added.\\
Imported Paper--V-B branch
& \(b_6\in\Acompsix\): selector remains upstream
& \(\mathsf{SLOT},\mathsf{TMF},\mathsf{CLK}_{\rm rem},\CMPexp,
\mathsf H_{\rm norm},\mathsf F\).  Numerical V-B spectrum/flavor data are
conditional on a compatible member.  \(\REGsig\), \(\mathrm{ANC}_E\), and
\(\mathrm{CAL}_{E_0}\) are typed separately for physical energy readouts.  The separate charged
\(\mathsf{SLOT}_{\rm phys}\) embedding is \textsc{derived-later} on the
\(\Urole+(b_6\in\Acompsix)+\mathsf F+\mathsf{G0}_{\rm odd}\) simple-spectrum branch by
Theorem~\ref{p7v2:thm:charged-slot-physical-later}.\\
Imported gauge certificate
& \(\Rconn=(\JOINTfiveocc)\): derived theorem output
& A canonical nonflat joint-bundle connection and the full unmarked Dirac
family are generated.  Factor same-datum relations are derived;
\(\Rssqft\) is an analytic theorem domain.\\
V-C-owned minimal families
& Five package-level families
& EM, VAC, CP, NU, KIN.  Here NU is the high-energy neutral lift fiber over
the frozen Paper--V-B base; its Casas--Ibarra \(R\) is a coordinate on a fibre
member, not a second selector beside \(Y_\nu\).  The derived gauge certificate
is already recorded in the preceding row.  The five families encode the new
physical operators, states, action/measure structures, or dynamics owned by
Paper~V-C.\\
Optional extensions
& Not in the minimal closure package
& RES saturation and HOPF expectation are retained as optional model branches.
The standard regulator is derived from physical masses and widths, and no
typed Hopf-to-source map is used.\\
Analytic/QFT domains
& Not selectors
& \(\Rssqft\), named subtraction/factorization/equilibrium domains and other
proof hypotheses restrict theorem applicability but do not choose a finite
realization member.\\
Calibrations
& Not structural selectors
& Overall gauge action scale, dimensionful shell/energy units and later
matching relations are explicit physical calibrations.\\
Downstream residues
& Moved downstream
& Nonperturbative QCD/hadronic boundary, continuum vertices, full interacting
reconstruction and later cosmological matching.\\
\bottomrule
\end{longtable}

\subsection{Typed physical-interface atlas}
\label{p7:sec:v236-interface-atlas}
\begin{theorem}[Retyping the five Paper--V-C interface families]
\label{p7:thm:v236-five-interface-types}
The retained symbols \(\mathrm{EM},\mathrm{VAC},\mathrm{CP},\mathrm{NU},\mathrm{KIN}\)
name five interfaces, but they are not five equal kinds of independent
physical root.
\begin{enumerate}[label=(\roman*),leftmargin=2.2em]
\item \(\mathrm{EM}\): on the actually occurring chiral-solder branch the
unbroken generator is already \(Q=T_L^3+Y\), and the native relation
\(e^{-2}=g_2^{-2}+g_Y^{-2}\) is the conventional low-energy coupling dictionary.  The remaining precision
low-energy claim is the claim-specific endpoint/readout fibre
\(\mathfrak A_{\rm EM}\), and only the quotient visible to the requested
observable must be fixed.
\item \(\mathrm{VAC}\): a claim-complete effective action generates its
stationary equations and Hessian.  VAC is retained only as a claim-scoped
scalar-equilibrium/history interface; it is not a headline electroweak-constitution
owner and does not sit upstream of the \(U(1)_Q\) stabilizer or gauge Gram.
\item \(\mathrm{CP}\): ordinary finite complex Dirac duality/CAR/determinant
structure is generated.  On the native complete-action/source domain,
Corollary~\ref{p7v2:cor:v259-native-strong-cp} and
Proposition~\ref{p7:prop:v190-index-determinant-phase} fix the finite
sector phase.  A chiral-continuum action/measure/index certificate is consumed
only by claims that explicitly use that completion.  Optional Majorana
structure remains a separate conditional branch.
\item \(\mathrm{NU}\): the high-energy neutral sector is a compatible lift
fibre downstream of the physical low-energy flavor/neutral data and, where
used, the heavy/Majorana branch.  Claims consume only their
observable-sensitive quotient.
\item \(\mathrm{KIN}\): for fixed \(H\) and \(\beta\) the finite Gibbs/KMS
state is generated.  Native nonequilibrium charge claims consume the actual
initial member and complete evolution; kinetic kernels are certified
reductions only where the requested target uses them.  A continued or
asymptotic target additionally consumes its flow/limit certificate.
\end{enumerate}
Consequently these five interface labels remain visible for source compatibility,
but VAC is claim-scoped and EM is a derived-dependent endpoint.  CP, NU and
KIN name the conditional comparison/interface classes of
\eqref{p7v2:eq:minimal-physical-package}, not three mandatory native gaps.
\end{theorem}

\begin{corollary}[Minimal current baryogenesis ancestry]
\label{p7:cor:v236-baryogenesis-ancestry}
On the high-energy neutral branch the minimal member ancestry is
\[
 \mathrm{SrcVal}_{\rm rooted}^{\rm der}\longrightarrow\mathrm{SVD}_{\rm same\ action}\longrightarrow G21\longrightarrow G22,
\]
with action-derived microscopic vertices/kernels inserted where their owner
has been proved.  The precision EM endpoint is not a premise of this chain,
and the optional resonant/Hopf benchmark package is not promoted into the
minimal ancestry.  A finite rank overlap such as \(1/16\) is an initial
probability on the stated tracial branch and is not, by itself, a kinetic
efficiency.
\end{corollary}
% END v2.36 RETYPED PHYSICAL INTERFACES


\subsection{Economy of the remaining assumptions}
\label{p7v2:sec:economy}

The active Clifford realization is not among the remaining assumptions.  Its
three canonical stages are generated by the Born-even, contact, and finite
Hurwitz response operators, while G0 supplies the occurring conjugate
representative.  The open ledger therefore begins only after finite support has been
realized.  It separately records compatible
Paper--V-B members, the energy ancestry/scheme/calibration domains, and the
derived joint-connection certificate \(\Rconn\).  Paper~V-C's minimal physical
owner families are exactly the five on the premise-minimal common-heat branch in
\eqref{p7v2:eq:minimal-physical-package}.  Exact resonance saturation and a
Hopf expectation remain available only as optional benchmark/model branches.

The smallest coherent extension is not a separate selector for each desired
number.  One functorially constructed finite action could, in principle,
simultaneously determine the gauge kinetic trace, electromagnetic boundary
operator, scalar Hessian, and Majorana/Yukawa maps.  Its admissibility would
have to be proved from the upstream carrier rather than tuned sector by sector.
Even such an action would not determine a cosmological initial state or solve
\eqref{p7v2:eq:kinetic-system}; the thermal history remains a separate
boundary-value problem.

\begin{remark}[Chronological selector firewall]
No successful Paper--V-C benchmark and no later Paper--VI-A/VII-A/VII-B/VII-C output is
allowed to reselect
\[
 \mathcal H_p,\quad m_p,\quad K=32,40,56,\quad
 \mathsf{TMF},\quad\mathsf F,\quad\REGsig,
\]
or any other earlier unresolved member.  A later paper may change a later
status only when it first constitutes independent data that genuinely
distinguish the earlier family.  Such a discharge never rewrites the earlier
paper as historically unique.
\end{remark}

\subsection{Finite interaction boundary, not premature closure}
\label{p7v2:sec:conclusion}

Paper~V-C now has a rigorous positive result: it supplies a cycle-free,
fully typed boundary between the structural finite carrier, the separately
named physical-role and energy domains, the compatible Paper--V-B
spectrum/flavor branch, and the continuum
interaction, QCD, and cosmological calculations.  The exact algebraic
content is retained, the standard continuum formulas are retained with their
actual hypotheses, and every previously hidden normalization, state, or
dynamical input is exposed.  The canonical active support itself is now
closed upstream as \(2T\to C_6\to\mathbb R\mathbf1\) at
\(K=32,40,56\).  Its terminal member is the scalar Haar projector, so neither
active realization nor infrared extraction is an entry in the open-data
ledger.

The revision deliberately does not preserve the former unconditional closure
claim.  It imports the derived \(\JOINTfiveocc\) connection/family certificate
and its factor gauge relations, the complete positive nonsingular
Casas--Ibarra fiber coordinate, the natural resonant domain, the exact
saturated benchmark root, and the full block/full thermal state
classification.  On the common-heat branch the gauge matching coefficient is
derived with \(\JOINTfiveocc\); only the broader kinetic-class branch retains
\(\CALgauge\).  The electromagnetic operator/state, scalar-action/vacuum
member, measure symmetry, high-energy
neutral lift, and thermal evolution cannot be recovered by rephrasing
GP1--GP3.  They must be constructed.  Resonant saturation and Hopf expectation
are optional rather than irreducible.  This premise contraction
is what makes the surviving Paper--V-C theorem publication-defensible and
gives the later papers a non-circular target.



% BEGIN INLINE SOURCE: Paper5C-CPBL2SourceIncidence-Main-v1_1.tex
% Framework v2.186: preserves exact source coefficients and clarifies the actual K16 handoff and claim-scoped direct signed charge target.
\subsection{Source-integrity contract}
\label{p7v2:app:source-integrity}

This revision:
\begin{enumerate}[label=(\roman*),leftmargin=2.5em]
  \item uses the canonical \texttt{Macros-v2\_7.tex};
  \item records the v2.10 proof-repair revision in
  \texttt{Paper7-SM\_Interaction-v2\_10.tex};
  \item contains no backward auxiliary import from Papers VI-A, VII-A, VII-B, or VII-C;
  \item uses \(g_Y\) for ordinary hypercharge and \(g_1=\sqrt{5/3}\,g_Y\)
  throughout;
  \item does not identify an ordinary determinant with a normalized-trace
  determinant;
  \item does not call a projector trace a kinetic efficiency;
  \item does not call a mass ratio a Higgs vertex modifier;
  \item does not call a bare \(\theta\) boundary a proof of physical
  \(\bar\theta=0\).
\end{enumerate}

\subsection{Geometry-first consequences and physical scope (framework v2.201)}
\label{p5c:sec:v2201-geometry-first-consequences}
\paragraph{Geometry ancestry.}
Rooted charged and neutral family incidences $\to$ generated weak current $\to$ inherited propagation and occurring production/detection $\to$ flavor and signed-charge readouts.
\paragraph{Status.}
These are consequences on the domains stated below.
A mathematical certificate or an executed restriction retains every
physical source, state, clock and readout hypothesis of its inputs.

\subsubsection{D07: All nine vacuum current probabilities and the CP envelope}
\label{p5c:sec:v2201-derived-d07}
\noindent\textit{Status: proved result on stated domain.}

Use the fixed generated current $C=W_e^\dagger W_\nu$ and define
\[
 s_{rs}=\sin^2\frac{\Phi_r-\Phi_s}{2},\qquad
 A=\frac{\sin\Delta_{12}-\sin\Delta_{13}+\sin\Delta_{23}}6,
 \quad \Delta_{rs}=\Phi_r-\Phi_s.
\]
For equal normalized coherent boundaries the complete probability matrix
(rows: initial $e,\mu,\tau$; columns: detected $e,\mu,\tau$) is
\[
\boxed{\begin{pmatrix}
1-\frac14s_{12}-\frac12(s_{13}+s_{23})&
\frac14s_{12}+\frac16(s_{13}+s_{23})+A&
\frac13(s_{13}+s_{23})-A\\[2pt]
\frac14s_{12}+\frac16(s_{13}+s_{23})-A&
1-\frac{25}{36}s_{12}-\frac5{18}(s_{13}+s_{23})&
\frac49s_{12}+\frac19(s_{13}+s_{23})+A\\[2pt]
\frac13(s_{13}+s_{23})+A&
\frac49s_{12}+\frac19(s_{13}+s_{23})-A&
1-\frac49(s_{12}+s_{13}+s_{23})
\end{pmatrix}.}
\]
\begin{proof}
Expand $|\sum_r C_{\beta r}C_{\alpha r}^*e^{-i\Phi_r}|^2$.
The coefficient of $s_{rs}$ is
$-4\operatorname{Re}(C_{\beta r}C_{\alpha r}^*C_{\beta s}^*C_{\alpha s})$;
the sine coefficient is twice its imaginary part. Substitution of the
literal current gives the displayed rational coefficients. Row and column
sums are one by unitarity. Direct matrix evaluation checks every entry.
\end{proof}

The CP-conjugate $e\to\mu$ difference is
\[
 \boxed{P_{e\mu}(C)-P_{e\mu}(\bar C)
 =\frac43\sin\frac{\Delta_{12}}2
             \sin\frac{\Delta_{23}}2
             \sin\frac{\Delta_{13}}2.}
\]
Its supremum over independent phases is $\sqrt3/2$. The maximum of the
three-sine product is $3\sqrt3/8$, obtained when the two adjacent
half-phases are $\pi/3$ modulo the corresponding symmetries. In the
admitted leading quadratic $L/E$ phase specialization, the N56 adjacent
gap ratio is irrational: it is $36863/(192x-1)-1$, $x^3=192$.
The resulting phase winding is dense on the phase torus and has the same
supremum. This statement consumes that phase specialization; it is not a
finite-baseline experimental likelihood.
\par\noindent\textit{Inherited inputs:} \ref{p5c:cor:v2167-literal-current-quartet}, \ref{p7v2:thm:v267-generated-lepton-current-incidence}, \ref{p7v2:thm:v267-neutral-propagation-phase}.

\subsubsection{D25: Closed entropy and produced occupation entropy are different readouts}
\label{p5c:sec:v2201-derived-d25}
\noindent\textit{Status: proved result on stated domain.}

For exact finite unitary evolution $\rho_t=V_t\rho_0V_t^\dagger$
with $i\hbar\,\partial_tV_t=H_{\rm phys}(t)V_t$ and a Schr\"odinger-picture
observable $Q(t)$ with explicit derivative $\partial_tQ$,
\[
 S(\rho_t)=S(\rho_0),\qquad
 \frac{d}{dt}\langle Q(t)\rangle
 =\operatorname{Tr}\rho_t
 \left\{\partial_t Q+\frac{i}{\hbar}[H_{\rm phys},Q]\right\}.
\]
If the complete root is pure, its total von Neumann entropy remains zero.
For a specified reduced fermionic subsystem with one-body occupation
eigenvalues $n_i$, however,
\[
 S(\rho_{\rm red})\le\sum_i[-n_i\log n_i-(1-n_i)\log(1-n_i)],
\]
with equality for its quasifree maximum-entropy state.
\begin{proof}
Unitarity preserves the density eigenvalues. Trace cyclicity
gives the charge derivative. The entropy bound follows from positivity of
relative entropy against the quasifree Gibbs density with the specified
one-body expectations; its logarithm is a number-operator linear combination.
\end{proof}
Thus the finite charge, coarse entropy and late baryon-export denominator
must be calculated on the same specified reduction/history. Nonzero pair
support or a CP invariant alone does not evaluate $\langle B-L\rangle$.
The actual complete normal/pair/interaction word is the remaining input for
the BL2 signed history. No efficiency, replacement root state or retired
Reeb phase pulse is inserted. The exact finite identity can return a
completed zero as well as a nonzero value.
\par\noindent\textit{Inherited inputs:} \ref{p5c:cor:v2186-direct-signed-charge-target}, \ref{p3f:cor:v1139-bl2-fixed-owner-transport}, \ref{p8:thm:v2114-fresh-root-neutral-qroot}.

\subsubsection{A1: exact survival and phase-averaged flavor transfer}
\label{p5c:sec:v2201-derived-a1}
\noindent\textit{Status: proved corollary on stated domain.}

The recorded current has squared weights
\[
\mathcal P=(|C_{ar}|^2)=
\begin{pmatrix}
1/4&1/4&1/2\\5/12&5/12&1/6\\1/3&1/3&1/3
\end{pmatrix}.
\]
On the paper's equal normalized production/detection boundary, with neutral
propagation diagonal in the declared roles,
\[
\mathcal A=C\operatorname{diag}(e^{-i\Phi_1},e^{-i\Phi_2},e^{-i\Phi_3})C^\dagger.
\]
Expanding its electron diagonal gives the exact identity
\[
\boxed{P_{ee}=1-\tfrac14\sin^2\tfrac{\Phi_1-\Phi_2}{2}
 -\tfrac12\sin^2\tfrac{\Phi_1-\Phi_3}{2}
 -\tfrac12\sin^2\tfrac{\Phi_2-\Phi_3}{2}.}
\]
If the phase differences dephase so the cross terms average to zero, then
\[
\boxed{\overline{P}_{ab}=\sum_r\mathcal P_{ar}\mathcal P_{br}
=\begin{pmatrix}
3/8&7/24&1/3\\7/24&3/8&1/3\\1/3&1/3&1/3
\end{pmatrix}.}
\]
The transfer has eigenvalues \(1,1/12,0\), with eigenvectors proportional to
\((1,1,1)\), \((1,-1,0)\) and \((1,1,-2)\). Consequently any normalized
initial charged-flavor mixture gives a final tau fraction of \(1/3\) on
this complete-dephasing boundary. The full coherent current is unitary;
this rank loss is in the averaged probability transfer, not in \(C\).

\textbf{Proof.} For the survival amplitude, write
\(\mathcal A_{ee}=\sum_r p_r e^{-i\Phi_r}\), \(p=(1/4,1/4,1/2)\), and
expand its modulus squared using \(\sum_rp_r=1\). Averaging the general
amplitude deletes only unequal-role cross terms and yields
\(\mathcal P\mathcal P^T\). Direct multiplication gives the matrix and
the three eigenvectors. No particular standard neutrino mechanism is used.
\par\noindent\textit{Inherited inputs:} \ref{p7v2:prop:v266-oscillation-current-kernel}, \ref{p7v2:thm:v267-generated-lepton-current-incidence}, \ref{p7v2:thm:v267-neutral-propagation-phase}.

\textbf{Physical test.} If a resolved oscillation frequency separates one
role from two nearly phase-coincident roles, its electron amplitude is
\(4p_r(1-p_r)\in\{3/4,1\}\), unchanged as a set by role permutations.
The Daya Bay final kilometer-baseline analysis reports the much smaller
three-flavor fitted amplitude \(\sin^2(2\theta_{13})=0.0851\pm0.0024\)
\cite{DayaBayFinal2023}. This is a serious conditional compatibility test of the
\emph{unmodified} pole/phase and equal-boundary identification. It is not an
unconditional exclusion of the whole framework: the actual generated
source/detector response, matter effects, coherence and phase domains must
be matched to that experiment. A repair must derive the relevant effective
weights/phases from the same action; phases or basis rephasings alone cannot
change the recorded weight set. The full spectral/readout bridge therefore
deserves priority over another scalar CP comparison.

\subsubsection{Physical neutrino comparison and the remaining readout map}
\label{p5c:sec:v2201-neutrino-observation-test}
The fixed current is a generated algebraic object. Its equality to an
experimentally inferred three-neutrino mixing matrix is a separate,
testable physical identification. On the equal coherent boundaries of
\ref{p5c:sec:v2201-derived-d07}, its electron weights are
\[
 (|C_{e1}|^2,|C_{e2}|^2,|C_{e3}|^2)=(1/4,1/4,1/2).
\]
When a resolved frequency separates one state from two nearly coincident
phases, its disappearance amplitude is
\[
 4p_r(1-p_r)\in\{3/4,1\}.
\]
The final Daya Bay result is $\sin^2(2\theta_{13})=0.0851\pm0.0024$
\cite{DayaBayFinal2023}. Therefore the unchanged equal-boundary current,
with the conventional physical three-neutrino pole/phase identification,
fails this comparison. This conclusion is conditional on that identification;
it does not discard the generated geometric current theorem.

An independent constraint concerns the invariant magnitude. The current has
$|J|=1/12\simeq0.08333$, whereas the NuFIT~6.0 fitted mixing angles give
$J_{\rm CP}^{\max}=0.0333\pm0.0007$ at one standard deviation, with a quoted
three-standard-deviation variation $\pm0.0017$, allowing the CP phase to
vary \cite{NuFIT2024}. These are constraints from the angles, not a claim
that the CP phase has been settled. The same direct identification fails
this test too. NuFIT~6.0 is a named published benchmark here, not a claim
about which global fit is newest.

Permutations preserve the electron weight set and $|J|$; rephasings
preserve both. Family-equivariant canonical normalization preserves the
family current by \ref{p1b:sec:v2201-derived-d08}. Agreement therefore
requires execution of the actual production density, detector effect,
physical propagator and any nonreducing mixed return from the existing
action. The relevant observable is
\[
 P_D(t)=\operatorname{Tr}\{E_DV(t)\rho_PV(t)^\dagger\}.
\]
For diagonal propagation its interference coefficients are
$(\rho_P)_{rs}(E_D)_{sr}$. Their evaluation must account for the observed
weights and energy/baseline dependence together. No fitted replacement
mixing angle or freely chosen normalization is introduced. An unexecuted
bridge supplies no evidence of experimental agreement.

% END GEOMETRY-FIRST CONSEQUENCES v2.201

% BEGIN FRAMEWORK v2.226: KINEMATIC CLOSURES

% BEGIN FRAMEWORK v2.321: PETERSEN MATTER INTERACTION CONTACT AND NO-GO
\section[Petersen source incidence and the strong-interaction interface]{Petersen source incidence and\\ the strong-interaction interface}
\label{p5c:sec:v321-petersen-interaction}

Paper~I-A's Petersen finite packet and Paper~V-A's rooted comparison now
supply a local $A_4$-equivariant $3+12$ representation, but that is not
yet an evaluation of the actual gauge, Dirac, scalar, or color-singlet
interaction action. The constitutive geometry is not licensed to insert
extra physical vertices solely because a finite intertwiner exists.
In particular the extra Petersen packet is available only on the
icosahedral-support branch and does not retroactively select that branch.

\begin{proposition}[Three permitted color-preserving comparison arrows]
\label{p5c:prop:v321-petersen-color-arrows}
Let $C\cong\mathbb C^3$ be the already constituted fundamental color
representation, $\ell\cong\mathbb C$ its singlet companion, and
$G\cong\mathbb C^3$ the gauge-trivial family replication carrier. For the
established color action on $Q=\ell\oplus C$,
\[
 \boxed{\operatorname{Hom}_{SU(3)}
        \bigl(C,(\ell\oplus C)\otimes G\bigr)
       \cong G,\qquad \dim_{\mathbb C}=3.}
\]
These three complex intertwiner coefficients are the full set of
color-equivariant \emph{comparison possibilities} from an auxiliary
color triplet into the existing $Q\otimes G$ carrier. They are not
three additional gauge fields, fermion generations, Yukawa eigenvalues,
or generated vertices.
\end{proposition}
\begin{proof}
With $SU(3)$ acting trivially on $G$, a map from the fundamental $C$ to
$\ell\otimes G$ vanishes by inequivalence of the irreducibles. Every map
$C\to C\otimes G$ is $v\mapsto v\otimes g$ for a unique $g\in G$ by
Schur's lemma. The displayed result follows. A physical application
additionally requires the marked geometric $C$ comparison to be transported
into the actual gauged color support and all weak, chirality, source and
kinetic data to be consumed on the same action owner.
\end{proof}

\begin{proposition}[No generated mirror or strong-CP insertion]
\label{p5c:prop:v321-petersen-no-automatic-insertion}
The Petersen mirror representation $\mathbf3'$ is an independent
$A_5$-geometric sector, not a newly constituted gauge multiplet of the
current Paper--V action. The original coordinate Higgs operator remains
$\mathcal V_H\otimes 1_Q\otimes1_G$, and its representation-level identity
extension cannot select a preferred Petersen mirror coefficient or a new
family mass ratio. The central $\mathbb Z_2$ signed-cycle holonomy acts on
a quaternionic geometry/transport comparison; it is not the $\mathbb Z_3$
center of color $SU(3)$, nor an additional native strong topological
character. Hence the existing strong-CP character result is unmodified.
\end{proposition}
\begin{proof}
The first assertion follows from the source-type distinction between the
signed $A_5$ edge representation and the actually generated matter/gauge
representation. The Higgs identity extension is the existing Paper--V-A
coordinate theorem. The groups $\{\pm1\}$ and the color center have
different orders, and no term coupling the geometric Wilson sign to the
physical topological or matter determinant weight has been generated.
Adding such a term would change the action grammar and requires a separate
same-owner construction and chiral-measure audit.
\end{proof}

\begin{remark}[Concrete next action check]
The decisive computation is a source-inserted matrix element of the
\emph{existing} full interaction action, with one marked tetrahedral
root, the proposed $C\oplus(Q\otimes G)$ intertwiner, and a complete
color-singlet receiver. One must explicitly evaluate the first/second
source jets, the regulated kinetic return, and any pole residue. Until
that check succeeds, the Petersen construction changes no physical
coupling, Yukawa singular value, one-loop coefficient, strong-CP phase,
hadronic mass, or measured rate.
\end{remark}
% END FRAMEWORK v2.321




% BEGIN FRAMEWORK v2.338: SOURCE-TYPING BOUNDARY FOR THE NATIVE PAIR BRIDGE
\subsection{Two-source Petersen bridge and the physical gauge boundary}
\label{p5c:sec:v338-two-source-gauge-boundary}
The geometric map of Proposition~\ref{p5:prop:v338-bilinear-edge-map}
provides an equivariant alternating route from two oriented pair
amplitudes into the established Petersen Clifford-edge carrier.
Its rooted twelve-edge complement is representation-equivalent to
the comparison module $Q\otimes G$ of Paper~V-A.  It does not
thereby become the physical color-singlet two-constituent state.

\begin{proposition}[Symmetry does not select the matter intertwiner]
\label{p5c:prop:v338-intertwiner-ambiguity}
Let $U=\Lambda^2\mathbb C^5$, restricted to the rooted tetrahedral
symmetry $A_4$, and $M=Q\otimes G\cong\operatorname{Reg}(A_4)$.
Then
\[
 U|_{A_4}=\mathbf1\oplus3\mathbf3,\qquad
 M=\mathbf1\oplus\mathbf1'\oplus\mathbf1''\oplus3\mathbf3,
 \qquad \dim_{\mathbb C}\operatorname{Hom}_{A_4}(U,M)=10.
\]
Therefore the geometric equivalence and even an isometric
intertwiner do not fix a unique physical insertion: the multiplicity
three triplet sector has a $U(3)$ family of possible isometries,
plus a scalar phase.  The quadratic cyclic source of
Theorem~\ref{p5:thm:v338-quadratic-cyclic-source} provides the
missing representation \emph{types}, not a canonical value for the
physical two-body coupling.
\end{proposition}
\begin{proof}
Schur's lemma gives one map on the common trivial line, nine maps
between the three identical triplet copies, and no maps from the
missing nontrivial one-dimensional targets.  The physical same-action,
color, charge and kinetic conditions are not statements of $A_4$
representation equivalence and must be checked separately.
\end{proof}
\begin{remark}[No source-free coefficient or hadron mass readout]
The projector $\mathsf J^\dagger\mathsf J$ has only zero and unit
eigenvalues.  Multiplying it by an arbitrary rest-energy scale, or
identifying the norm ratio $25/2$ of two differently typed geometric
arrows with a pion/kaon mass-squared ratio, is not licensed by this
construction.  The actual finite source-action Hessian and its
color-singlet image must be evaluated before any rest-energy claim.
\end{remark}
% END FRAMEWORK v2.338


% BEGIN FRAMEWORK v2.348: CANONICAL GRAM EIGENVALUES AND COHERENT MEMORY NORM
\clearpage
\section{Exact electroweak Gram eigenvalues and neutral-memory contrasts}
\label{p5c:sec:v348-gram-memory}
\begin{theorem}[Electroweak orbit Gram in a canonically whitened gauge frame]
\label{p5c:thm:v348-ew-gram-eigenvalues}
Use the occurring unit solder $\widehat H=(0,1)^T$, ordinary
hypercharge $Y_H=1/2$, and positive common orbit coefficient
$c_G>0$.  Assume the positive gauge-kinetic metric is canonically
whitened and that its coupling coordinates are
$(g_2W^1,g_2W^2,g_2W^3,g_YB)$ with $g_2,g_Y>0$.
The linear orbit map in these coordinates is
\[
 J_{\widehat H}(W,B)=\frac12
 \begin{pmatrix}g_2(W^1-iW^2)\\-g_2W^3+g_YB\end{pmatrix}.
\]
The corresponding real Gram quadratic form is represented by
\[
 \boxed{M^2_{\rm EW}=\frac{c_G}{2}
 \begin{pmatrix}
 g_2^2&0&0&0\\
 0&g_2^2&0&0\\
 0&0&g_2^2&-g_2g_Y\\
 0&0&-g_2g_Y&g_Y^2
 \end{pmatrix},}
\]
with complete eigenvalue multiset
\[
 \boxed{\operatorname{spec}M^2_{\rm EW}=
 \bigl\{0,\tfrac{c_G}2g_2^2,\tfrac{c_G}2g_2^2,
              \tfrac{c_G}2(g_2^2+g_Y^2)\bigr\}.}
\]
In this \emph{quadratic response} (not yet an interacting pole)
$M_{Z,G}^2/M_{W,G}^2=1+g_Y^2/g_2^2$ and
$\rho_G:=M_{W,G}^2/(M_{Z,G}^2\cos^2\theta_W)=1$, where
$\sin^2\theta_W=g_Y^2/(g_2^2+g_Y^2)$.  On the
GUT-normalized common-heat branch $g_1^2=g_2^2=2\pi^2$ and
$g_Y^2=(3/5)g_1^2$, the Gram ratio is $8/5$ and
$\sin^2\theta_W=3/8$.
\end{theorem}
\begin{proof}
The positive Gram is $2c_G J_{\widehat H}^\dagger
J_{\widehat H}$.  Expanding the squared Euclidean norm gives the
matrix displayed.  The neutral $2\times2$ block is the outer
product $(g_2,-g_Y)^T(g_2,-g_Y)$, so has eigenvalues zero and
$g_2^2+g_Y^2$, while the transverse weak entries are $g_2^2$.
Division gives $\rho_G=1$ and the heat ratio.  More general
kinetic metrics must be whitened with their actual positive square
roots before comparing eigenvalues; rank and nullity alone are
metric-independent.
\end{proof}

\begin{theorem}[Exact Frobenius law for diagonal neutral-role memory]
\label{p5c:thm:v348-weak-memory-coherence-norm}
On the normalized source chart $N=I$, take the native unitary
current $C$ and its modulus matrix $\mathcal P=|C|^2$ of
Proposition~\ref{p5c:prop:v2168-coherent-family-return}.  If the
\emph{evaluated} neutral-return matrix is diagonal in the declared
inverse/direct/Gram role basis,
$K_W(z)=\operatorname{diag}(k_1,k_2,k_3)$, with arbitrary complex
$k_j=k_j(z)$, and $\Sigma_{WW}=\pi^2 CK_WC^\dagger$, then
\[
 \boxed{\begin{aligned}
 \frac{\|\operatorname{OffDiag}\Sigma_{WW}\|_F^2}{\pi^4}
 &=\|k\|_2^2-\|\mathcal P k\|_2^2\\
 &=\frac12|k_1-k_2|^2
   +\frac{11}{72}|k_1+k_2-2k_3|^2.
 \end{aligned}}
\]
With $\bar k=(k_1+k_2+k_3)/3$, the sharp bounds are
\[
 \boxed{\frac{11}{12}\pi^4\sum_j|k_j-\bar k|^2
 \le \|\operatorname{OffDiag}\Sigma_{WW}\|_F^2
 \le \pi^4\sum_j|k_j-\bar k|^2.}
\]
In particular $\operatorname{OffDiag}\Sigma_{WW}=0$
if and only if $k_1=k_2=k_3$.  No Hermiticity or reality of the
same-frequency resolvent returns is assumed.
\end{theorem}
\begin{proof}
Unitary similarity gives
$\|\Sigma_{WW}\|_F^2=\pi^4\sum_j|k_j|^2$;
its diagonal is $\pi^2\mathcal Pk$.  Direct multiplication gives
\[
 \mathcal P^T\mathcal P=
 \begin{pmatrix}
 25/72&25/72&11/36\\
 25/72&25/72&11/36\\
 11/36&11/36&7/18
 \end{pmatrix},\quad
 \operatorname{spec}(\mathcal P^T\mathcal P)=\{1,1/12,0\}.
\]
The orthogonal contrast vectors $(1,-1,0)$ and $(1,1,-2)$
have eigenvalues $0$ and $1/12$, respectively.  Subtracting
the diagonal Frobenius norm from the full one proves the first
identity.  On the orthogonal complement to $(1,1,1)$, the two
eigenvalues of $I-\mathcal P^T\mathcal P$ are $1$ and $11/12$;
this proves both sharp bounds and the vanishing criterion, for
complex $k$ as well.
\end{proof}

\begin{corollary}[Two different information-loss directions]
\label{p5c:cor:v348-coherent-diagonal-blindness}
The same native matrix has the exact left and right null lines
\[
 \boxed{\ker(\mathcal P\mathcal P^T)=\mathbb C(1,1,-2)^T,
 \qquad \ker(\mathcal P^T\mathcal P)=\mathbb C(1,-1,0)^T.}
\]
Therefore phase-averaged charged-flavor readout discards a different
contrast from the inverse/direct contrast erased by diagonal
neutral-memory readout.  The full coherent matrix $CK_WC^\dagger$
retains both where its actual source/receiver reduction preserves
this chart.  This is a finite exact information statement, not a
physical PMNS or oscillation-probability claim.
\end{corollary}
\begin{proof}
The columns of $\mathcal P$ give its right kernel, and its third
row is the arithmetic mean of its first two, giving the left
kernel.  Since $\mathcal P$ has rank two, these are the complete
null lines.  Unitary inversion of the full matrix is the one in
Proposition~\ref{p5c:prop:v2168-coherent-family-return}.
\end{proof}

\begin{remark}[Source/observable guard and corrected one-loop weights]
\label{p5c:rem:v348-weak-memory-physical-guard}
The coherence identity applies only to the diagonal $K_W$ antecedent
and $N=I$ source chart; in the actual system $K_W(z)$ may be
nondiagonal and the complete Schur return includes cross terms with
$B_R$, as in Theorem~\ref{p5c:thm:v2168-complete-weak-memory}.
Even diagonal normalized current coefficients are not automatically
physical neutrino oscillation probabilities.  The source-native
one-loop scalar-to-gauge weights in
Theorem~\ref{p7:thm:v164-scalar-gauge-source} are
$(41/40,-19/24,-7/4)$; this repairs the erroneous extra factor
of four in the previous display.  Their exact arithmetic at the
chosen common-heat inputs must not be confused with a
nonperturbatively exact vertex or an experimentally matched running
coupling.
\end{remark}
% END FRAMEWORK v2.348

\printbibliography

% BEGIN FRAMEWORK v2.365: CONDITIONAL COMPLETE-DAUGHTER WEAK INTERFACE
\section[Proper-action weak current and complete daughter-channel typing]{Proper-action weak current and\\ complete daughter-channel typing}
\label{p5c:sec:v365-decay-hyperedge}
A particle decay channel is a directed hyperedge from an actual
parent pole state to a \emph{multiset} of outgoing particle and
antiparticle states.  It is not the unoriented adjacency of the
Hurwitz shell and is not the sum of terminal mass coefficients.
The finite source calculus of Paper IV-B supplies particle--antiparticle
Fock sectors, while this paper supplies the native current maps
$C_{e\nu}$ and $C_{ud}$ on their declared domains.

\begin{proposition}[Unitarity of inclusive family factors, not rates]
\label{p5c:prop:v365-family-decay-factor}
On a normalized three-family current chart in which $C=C_{e\nu}$
is unitary, for any charged-family indices $r,s$,
\[
\boxed{\sum_{j,k=1}^{3}|\overline{C_{rj}}C_{sk}|^2=1,
\qquad\sum_{j=1}^{3}C_{rj}\overline{C_{sj}}=\delta_{rs}.}
\]
The first formula is an incoherently summed two-neutrino source
weight; the second is the coherent contraction that gives an
orthogonality cancellation for $r\ne s$.  Neither equation alone
provides an exclusive decay width, lifetime, neutrino spectrum or
physical state selection.
\end{proposition}
\begin{proof}
The first sum factorizes as
$(\sum_j|C_{rj}|^2)(\sum_k|C_{sk}|^2)=1$.  The second is
the $(r,s)$ entry of $CC^\dagger=I$.
\end{proof}

\begin{definition}[Same-action complete-channel contraction]
\label{p5c:def:v365-decay-source-contract}
For a physical parent preparation projector $P_i$, a final
multiset projector $Q_{D,\alpha}$ in the correctly charged,
spin-coupled and antisymmetrized Fock carrier, and a full
proper-action weak transition $\mathcal T_W(t)$, define the
finite-time inclusive channel functional
\[
\mathcal P_{i\to\alpha}(t)
 :=\operatorname{Tr}\!\left(
 Q_{D,\alpha}\mathcal T_W(t)P_i\rho_iP_i
 \mathcal T_W(t)^\dagger Q_{D,\alpha}\right).
\]
Only channels with physical energy/pole and source support, global
charges, outgoing-state normalization and detector prescription
may be compared with measured partial rates.  A one-fermion
number-preserving family ladder is \emph{not} already the
number-changing process $\mu^-\to e^-+\bar\nu_e+\nu_\mu$;
the current must be lifted to the appropriate field/Fock
sector and its action-owned two-vertex contraction executed.
For quarks other than the top, such elementary arrows are
partonic current transitions, not free on-shell quark decays.
\end{definition}

\begin{remark}[No physical branching promotion]
\label{p5c:rem:v365-channel-physical-firewall}
\leavevmode\par\noindent
The source-defined quark current
\(C_{ud}=Z_u^{-1/2}E_u^\dagger E_dZ_d^{-1/2}\)
remains the native quark current; numerical CKM or PMNS comparison
sections are not substituted as independent inputs.  A finite
Hermitian carrier by itself does not produce an irreversible
exponential lifetime.  Scalar $25/108$ neutral-current fits and
25-cell operator norms do not become proper gauge-current vertices
by reinterpretation.
\end{remark}
% END FRAMEWORK v2.365


% BEGIN FRAMEWORK v3.027: GEOMETRIC-TO-FOCK DECAY CERTIFICATE
\section{Necessary certificates for importing geometric transition restrictions}
\label{p5c:sec:v3027-geometric-decay-certificate}
\begin{definition}[Typed geometric-to-physical decay interface]
\label{p5c:def:v3027-geometric-decay-certificate}
The complete daughter-multiset channel functional in
Definition~\ref{p5c:def:v365-decay-source-contract} may inherit a
Hurwitz/$600$-cell comparison selection rule only after source
certificates establish all of the following, with matching carriers:
\begin{enumerate}
\item \textbf{SOURCE\_600\_EDGE:} an action-owned gauge/Clifford
$600$-cell edge operator, or a demonstrated equivalent already-licensed
nonlocal operator;
\item \textbf{SOURCE\_RESOLVENT:} its actual spin/color/clock
intermediate resolvent, including higher-order return shells and
Clifford-grade remainders rather than an assumed scalar denominator;
\item \textbf{SOURCE\_PATH\_ASSEMBLY:} the actual coherent
path weights and any nonlinear unitary polar operation, or an
independent same-action substitute with identical observable output;
\item \textbf{STATE\_MAP:} physical parent pole preparation and
complete, properly charged, spin-coupled, antisymmetrized daughter
multiset projectors in the licensed Fock carrier;
\item \textbf{PHYSICAL\_WIDTH:} source-owned matrix elements,
spectral/continuum or open-channel interface, residues, normalization,
and detector prescription for a measured partial width.
\end{enumerate}
These are verification obligations, not newly assumed independent
postulates or imported standard-model decay mechanisms.
\end{definition}

% BEGIN FULL-CARRIER GEOMETRIC DECAY CERTIFICATE IMPORT v1.1.
\begin{remark}[Scoped core-pair flavor imports into the decay contract]
\label{p5c:rem:fullcarrier-decay-source-imports}
The mathematical current $\mathscr K(F_A)$ in
Theorem~\ref{p1b:thm:fullcarrier-protected-local-flavor} is a
\texttt{GEOM\_COREPAIR\_CURRENT} candidate.  The minimum
$28$-dimensional $SU(3)$ action and the larger exact-current
$U(3)$ dilation are two \emph{alternative}
\texttt{GEOM\_FULL\_WARD\_BRANCH} constructions:
the former changes off-flavor graph current elements, whereas the
latter has no certified full $A_5$-equivariant embedding.  Neither
supplies \texttt{GAUSS\_SAME\_ACTION\_SOURCE}.
These imports may be consumed by the existing
\texttt{SOURCE\_600\_EDGE}, \texttt{SOURCE\_RESOLVENT},
\texttt{SOURCE\_PATH\_ASSEMBLY}, \texttt{STATE\_MAP} and
\texttt{PHYSICAL\_WIDTH} checks only after the actual operator,
Gauss, Clifford, charge, spin, Fock and physical-source carriers have
been matched.  The source must still be a complete daughter-multiset
current of one actual action; the finite geometry's representation
strengths do not become rate inputs.
\end{remark}
% END FULL-CARRIER GEOMETRIC DECAY CERTIFICATE IMPORT v1.1.

\begin{remark}[Three noninterchangeable levels of selection]
\label{p5c:rem:v3027-transition-promotion-boundary}
The exact compact $Q_8/C_3$ character zero modes of the scalar
$600$-cell adjacency, the nullities of an extended Cartan link
comparison, and the nonzero character leakage in the native
$864$-dimensional QCD clock operator live on distinct carriers.
None alone selects a physical parent or complete daughter channel;
therefore no hadron lifetime, branching ratio, or all-time physical
forbidden transition is promoted by these finite tests.
The relevant observable is the properly contracted complete-channel
probability, not resemblance to a conventional decay mechanism.
\end{remark}
% END FRAMEWORK v3.027 V-C MODULE


% BEGIN HOMOGENEOUS MAXWELL SHAPE RESPONSE v1.17
\section{Isolated Hopf-photon occupation and generated shear source}
\label{bh:berger:v117:p5c:photon-shear}
\begin{proposition}[Three homogeneous Maxwell polarizations on Milnor space]
\label{bh:berger:v117:p5c:maxwell-triplet}
On the position-space volume-preserving metric
$g_3=a^2\eta^{\mathsf T}e^{2B}\eta$ with
$\eta_i=\sigma_i/2$, $B=B^{\mathsf T}$, $\operatorname{tr}B=0$,
retain the isolated source-normalized homogeneous Maxwell modes
$A=f_i(t)\eta_i$. Their exact reduced Lagrangian is
\[
 L_\gamma=\pi^2a\,\dot f^{\mathsf T}e^{-2B}\dot f
              -\frac{4\pi^2}{a}f^{\mathsf T}e^{2B}f.
\]
In a fixed canonical coframe, the corresponding quantum normal-mode
Hamiltonian, retaining the displayed central oscillator contribution, is
\[
 H_\gamma(B)=\frac{2\hbar}{a}
 \left[\mathbf a^\dagger e^{2B}\mathbf a
        +\frac12\operatorname{tr}(e^{2B})\right],
 \qquad [N_\gamma,H_\gamma(B)]=0.
\]
At a fixed occupation covariance $C_{ij}=\langle a_j^\dagger a_i\rangle$,
the leading traceless metric force is
\[
 Q_\gamma=\operatorname{Re}C-\tfrac13\operatorname{tr}(C)I_3,
 \qquad \delta E_\gamma=\frac{4\hbar}{a}
                       \operatorname{tr}(Q_\gamma B).
\]
For $A_C=\operatorname{Re}C+I_3/2>0$, the isolated matter-energy
minimum at fixed volume is
$e^{2B_*}=(\det A_C)^{1/3}A_C^{-1}$.
This is not a fully constrained gravitational equilibrium and does
not supply physical initial occupation or other action contacts.
\end{proposition}
\begin{proof}
The electric and magnetic quadratic forms in the fixed Hopf coframe
carry inverse and direct metric factors. Their canonical Legendre
transform has the same $e^{2B}$ matrix in the momentum and coordinate
quadratic forms, so its number-conserving oscillator representation
has the stated Hamiltonian. The shape force is its first metric
derivative. For the matter-only minimum apply matrix AM--GM to
$A_C^{1/2}e^{2B}A_C^{1/2}$ of fixed determinant $\det A_C$.
\end{proof}
\begin{remark}[Source production and constraints]
\label{bh:berger:v117:p5c:production}
Photon number is conserved only for the isolated subsystem: an actual
charged current changes it. The leading emitted covariance is a
positive matrix of the actual current transition residues, not an
assigned occupancy. Require the emitted mode's transverse Maxwell
constraint, the gravitational momentum constraint, allowed energy
support and any requisite source/quantum contacts. The free static
round one-body Dirac sector does not have an on-shell one-photon decay
merely because a current matrix element is algebraically nonzero.
\end{remark}
% END HOMOGENEOUS MAXWELL SHAPE RESPONSE v1.17




% BEGIN FRAMEWORK v3.068 GEOMETRY BRIDGE SOURCE INTEGRATION
\section{Higgs-dressed local chiral sources and Gauss-compatible weak-current contact}
\noindent\emph{Scope.} The statements below retain their declared finite geometric, source-chart, and action-relative hypotheses. They do not assert physical particle decay or unconditional source occupation.

\begin{proposition}[Complete three-quark chiral current and light-source leakage]
\label{p5c:prop:v368-chiral-full-current-and-leakage}
On the one-orbital left-handed local color-antisymmetric three-quark
source carrier
$\mathcal B_L=\operatorname{Sym}^3(\mathbb C^2_L\otimes
(\mathbb C_u^3\oplus\mathbb C_d^3))$ of dimension $364$,
let $J_W^+=d\Gamma(C_{ud})$ denote the established unitary-family
charged raising current between the $d$ and $u$ role summands.
Then
\[
\operatorname{spec}(J_W^-J_W^+)=
\{4^{(56)},3^{(112)},1^{(70)},0^{(126)}\},\qquad
\operatorname{rank}J_W^+=238.
\]
The compression to the light $(u,d,s)$ source
$\operatorname{Sym}^3(\mathbb C^2\otimes\mathbb C^3)$
has rank $32$ for $x=|(C_{ud})_{ub}|^2<1$, and
its nonzero Gram spectrum equals
$(1-x)\{4^{(4)},3^{(8)},2^{(12)},1^{(8)}\}$.
The complementary flavor leakage $\Lambda$ satisfies
\[
 \operatorname{tr}(\Lambda^\dagger\Lambda)=56(1+x)\ge56.
\]
Thus the light source does not form an invariant subspace for the
full unitary current. The exterior role component is not by itself
an on-shell physical heavy daughter.
\end{proposition}
\begin{proof}
A unitary change of up-family basis reduces $C_{ud}$ to the
identity on the full six-Weyl-doublet carrier.  Schur--Weyl
then gives $56V_{3/2}\oplus70V_{1/2}$; the raising-current
squares have weights $4,3,1,0$ with the shown multiplicities.
On the light sector rotate the two down directions to a weak-active
combination of squared up-light amplitude $1-x$ and its orthogonal
combination. The one-body light compression is $\sqrt{1-x}$ times
the standard flavor raising matrix. On the symmetric cube,
$\Lambda^\dagger\Lambda$ is the complementary number operator
with an exact independent trace $56(1+x)$.
\end{proof}
\subsection{Conditional local Higgs-dressed weak-source and joint adjoint test}
\begin{theorem}[Local weak-singlet source multiplicity with the generated scalar]
\label{p5c:thm:v368-higgs-dressed-quark-source}
On the generated left-handed three-quark local CAR module with one orbital,
three families and two Weyl components, put
\(M=\mathbb C^2_{\rm Weyl}\otimes\mathbb C^3_{\rm family}\),
\(\mathcal B_Q=\operatorname{Sym}^3(\mathbb C^2_{\rm weak}\otimes M)\).
The color-singlet projection is the existing antisymmetric three-color tensor.
The scalar and weak roles are the previously generated carriers of
Theorems~\ref{p5:thm:generated-coordinate-higgs} and
\ref{p5:thm:v3101-generated-quark-current-incidence}.
On the full-rank generated weak current domain,
\[
\mathcal B_Q\cong
56V_{3/2}\oplus70V_{1/2},\quad \dim\mathcal B_Q=364.
\]
After tensoring with \emph{one} copy of the generated pseudoreal Higgs
conjugate \(\widetilde H\) of hypercharge \(-1/2\),
\[
\dim\bigl(\mathcal B_Q\otimes \mathbb C^2_{\widetilde H}\bigr)^{SU(2)_L}=70.
\]
The resulting correctly contracted local \(QQQ\widetilde H\) operator has
zero total hypercharge and is color and weak singlet.  It remains
a fermionic spinor-valued interpolating source with an uncontracted
Lorentz spinor index, not a Lorentz scalar observable. A three-quark local source
without a weak-fundamental or other odd-center dressing cannot be weak singlet,
even if multiplied by any number of weak-adjoint gauge bosons.
This theorem concerns local CAR/representation source types and does not prove
Higgs-vacuum occurrence or physical baryon poles.
\end{theorem}
\begin{proof}
Schur--Weyl gives
\(
\operatorname{Sym}^3(\mathbb C^2\otimes M)=
V_{3/2}\otimes \operatorname{Sym}^3M\oplus
V_{1/2}\otimes S_{(2,1)}M
\), with dimensions \(56\) and \(70\) of the multiplicity spaces.
The Clebsch--Gordan product of \(V_{1/2}\) with one scalar \(V_{1/2}\)
contains exactly one singlet; \(V_{3/2}\otimes V_{1/2}\) contains none.
The weak center acts by minus one on three quark doublets and plus one on
any adjoint tensor word. Hypercharge cancels because \(3/6-1/2=0\).
\end{proof}
\begin{proposition}[Exact nonzero spatial joint-current contact on the minimal dressed sector]
\label{p5c:prop:v368-joint-spatial-current}
On the one-family one-orbital \(20\)-dimensional quark carrier and the
same generated scalar doublet, let \(\mathcal S\) be the \(2\)-dimensional
local weak-singlet subspace. Using the already generated Weyl and weak
matrices, write \(J_{a,k}=d\Gamma((\sigma_a/2)\otimes(-\sigma_k))\)
for the three real weak indices and a fixed spatial current slot \(k\).
Then
\[
P_{\mathcal S}\sum_{a=1}^3 J_{a,k}^{\dagger}J_{a,k}P_{\mathcal S}
=\frac{19}{4}P_{\mathcal S}\qquad(k=1,2,3).
\]
On a \emph{specified} common positive-frequency adjoint tangent-boson chart
with orthonormal one-mode vectors \(|1_a\rangle\) and
\(X_a|0_B\rangle=\chi|1_a\rangle\), the jointly transformed source
\(g_2\sum_a J_{a,k}\otimes X_a\) has nonzero norm squared
\(19g_2^2|\chi|^2/4\) on every normalized source \(s\in\mathcal S\).
This is an exact local joint source matrix element, not a pole or width.
\end{proposition}
\begin{proof}
The source states transform as total weak isospin zero; the spatial current
transforms as the adjoint. Its contraction with the co-transforming adjoint
bosonic tangent mode is invariant. In the symmetric-cube four-mode carrier,
\(C_2^{SU(4)}=63/8\), whence
\(
\sum_{a,k}J_{a,k}^2=63I/4-J_{\rm weak}^2-S_{\rm Weyl}^2
\).
On the \((1/2,1/2)\) source component the right side is \(57I/4\);
rotational symmetry makes the fixed-\(k\) contribution \(19I/4\).
The orthogonality of one-boson vectors gives the displayed squared norm.
The \(SU(4)\) calculation is an identity on the ambient matrix algebra,
not a new physical gauge symmetry.
\end{proof}
\begin{lemma}[Exact native 24-Hurwitz joint adjoint projection]
\label{p5c:lem:v368-Hurwitz-joint-Reynolds}
In the defining weak \(SU(2)\) representation and its adjoint \(R(g)\),
\[
\frac1{24}\sum_{g\in2T}R(g)_{ai}R(g)_{bj}
=\frac13\delta_{ab}\delta_{ij}.
\]
Thus a naked weak-adjoint vector averages to zero but the joint
\(\mathrm{adjoint}\otimes\mathrm{adjoint}\) contraction has a nonzero
invariant component. This is the same second-adjoint moment as full
\(SU(2)\) Haar average; it is not the same as a \(Q_8\)-only average.
\end{lemma}
\begin{remark}[Physical action and Gauss scope]
The connection tangent must transform jointly with its reference background;
its creation as a kinetically normalized proper weak boson is not assumed.
No Higgs vacuum, physical parent pole, absolute boson clock/mass, complete
Gauss/BRST reduction, or Fock daughter multiset has been derived here.
The geometric Petersen source remains a candidate conditional pullback
and must be checked against the complete Paper V-C Schur word and
proper-channel certificate before physical use.
\end{remark}

\begin{remark}[No physical geometric-to-decay certificate discharge]
\label{p5c:rem:v368-geometry-source-certificate-frontier}
The local $70$-source algebra, Higgs dressing, reference boson-port
variable $\kappa=|\chi|^2$, and finite current and Petersen Grams
are input/source compatibility statements only. They do not prove a
physical Higgs or weak-boson occupied state, the gauge-invariant
interacting parent/daughter pole projectors, the full proper-time
propagator, the necessary energy support, or the complete nonzero
channel residues. Keep the physical conditions in
Definitions~\ref{p5c:def:v365-decay-source-contract} and
\ref{p5c:def:v3027-geometric-decay-certificate} open unless a
same-action complete spectral certificate is separately supplied.
Do not identify the 70 source directions with baryon species,
$19/4$ with a decay width, or the rank-two geometric comparison
with a normalized weak interaction vertex.
\end{remark}
% END FRAMEWORK v3.068 GEOMETRY BRIDGE SOURCE INTEGRATION

\end{document}
